%
% File: "BIBANA.26"       Version #3  for arXiv.org        6 July 2008  
% ====                    
%                 
%
 \documentclass [12pt,a4paper      ]{article}
%\documentclass [12pt,a4paper,draft]{article}
%
\usepackage{graphics}
%
\usepackage{times}
\renewcommand{\familydefault}{times}
\DeclareFontFamily{OT1}{times}{}
\DeclareFontShape {OT1}{times}{m }{n }{ <-> ptmr }{}
\DeclareFontShape {OT1}{times}{bx}{n }{ <-> ptmb }{}
\DeclareFontShape {OT1}{times}{m }{it}{ <-> ptmri}{}
\DeclareFontShape {OT1}{times}{bx}{it}{ <-> ptmbi}{}
\usepackage{amsmath}
\usepackage{amsfonts}
\usepackage{amssymb}
\usepackage{latexsym}
%
%
%-------------------------------------------------------------------
% ISRI      special quaternion definitions (version 40)        START
%-------------------------------------------------------------------
\newcommand{\cl}{C \kern -0.1em \ell} % Clifford algebra
\newcommand{\bbR}{\mathbb{R}}         % Real numbers
\newcommand{\bbC}{\mathbb{C}}         % Complex numbers (bireals)
\newcommand{\bbH}{\mathbb{H}}         % Quaternions
\newcommand{\bbB}{\mathbb{B}}         % Biquaternions
\newcommand{\bbO}{\mathbb{O}}         % Octonions

\newcommand{\DEF}{:=}                 % definition

\newcommand{\CON}{\overline}          % quaternion (vector)conjugate

\newcommand{\Scal}{\mathbb{S}}        % scalar part unary  S[AB] = s
\newcommand{\Vect}{\mathbb{V}}        % vector part unary  V[AB] = v
\newcommand{\scal}{\circ}             % scalar part binary   AoB = s
\newcommand{\vect}{\wedge}            % vector part binary   A^B = v

\newcommand{\SCA}{\langle}            % SCA-lar left  bracket   <... 
\newcommand{\LAR}{\rangle}            % sca-LAR right bracket   ...> 
\newcommand{\VEC}{\vec{\kern +.1em[}} % VEC-tor left  bracket   |>.. 
\newcommand{\TOR}{\vec{\kern +.2em]}} % vec-TOR right bracket   ..>| 
\newcommand{\BRA}{\langle\kern -.2em\langle} % Dirac BRA      <<...| 
\newcommand{\KET}{\rangle\kern -.2em\rangle} % Dirac KET      |...>> 

\newcommand{\REA}{\operatorname{Re}}  % real      part
\newcommand{\IMA}{\operatorname{Im}}  % imaginary part
\newcommand{\ARG}{\operatorname{arg}} % argument or phase
\newcommand{\e  }{\operatorname{e}}   %  e  = 2.718281
\newcommand{\Oh}{\tfrac{1}{2}}        % 1/2 = 0.5
\newcommand{\Th}{\tfrac{3}{2}}        % 3/2 = 1.5
\newcommand{\tr}{\pmb{\perp}}         % transverse/perpendicular 
\newcommand{\lo}{\pmb{\parallel}}     % longitudinal/parallel

\newcommand{\Q}{[\hspace{1.mm}]} % quaternion operand A{\Q}B -> A[]B
\newcommand{\A}{(\hspace{.5mm})} % argument of function F{\A} -> F()

\newcommand{\REV}{\sim}               % order reversal: Q^\REV
\newcommand{\DUA}{\widetilde}         % dual: \DUA{Q} = dual of Q

\newcommand{\ASS}{\thickapprox}       % associate function  
\newcommand{\TRA}{{\pmb{=}}}          % transpose function
\newcommand{\ADJ}{\dagger}            % adjoint function  
\newcommand{\INV}{{-1}}               % inverse function  
%-------------------------------------------------------------------
% ISRI           special quaternion definitions                 STOP
%-------------------------------------------------------------------
%
% Options
%--------
\setlength{\parskip}{2.5mm}
%
%%\numberwithin{equation}{section}
%

\begin{document}

\title{\bf\vspace{-2.5cm} Quaternions in mathematical physics (2):\\
                             \emph{Analytical bibliography} }
%                         ========================================

\author{Andre Gsponer and Jean-Pierre Hurni\\ ~\\ 
\emph{Independent Scientific Research Institute}\\ 
\emph{Oxford, OX4 4YS, England}}

\renewcommand{\today}{6 July 2008}
\date{ISRI-05-05.26 ~~ \today}

\maketitle



\begin{abstract}

This is part two of a series of four methodological papers on (bi)quaternions and their use in theoretical and mathematical physics: 1~- Alphabetical bibliography, 2~- Analytical bibliography, 3~- Notations and definitions, and 4~- Formulas and methods.  

This quaternion bibliography will be further updated and corrected if necessary by the authors, who welcome any comment and reference that is not contained within the list.

\end{abstract}

~\\

\begin{center}
\resizebox{3cm}{!}{ \includegraphics{logo.eps}}
%--------------------------------------------------------
\end{center}

~\\

\noindent Living report, to be updated and corrected by the authors, first published on the occasion of the bicentenary of the birth of William Rowan Hamilton (1805--2005).
\newpage

%\tableofcontents

\section{\Huge Table of contents}
%================================
\label{Table of contents}

The table of content is essentially the list of the main KEYWORDs (always written in the singular) used for classifying the items in the bibliography.


\renewcommand{\labelitemi}{}   %Companion p.59
\renewcommand{\labelitemii}{}
\begin{itemize}

\item \ref{Table of contents}.         Table of contents 

\item \ref{Introduction}.              Introduction 

\item \ref{MATHEMATICS}.               MATHEMATICS
   \begin{itemize}
   \item \ref{MATH-VARIA}.             MATH-VARIA
   \item \ref{ALGEBRA}.                ALGEBRA
   \item \ref{INTEGRAL-QUATERNION}.    INTEGRAL-QUATERNION
   \item \ref{EQUATION}.               EQUATION
   \item \ref{LINEAR-FUNCTION}.        LINEAR-FUNCTION
   \item \ref{MATRIX}.                 MATRIX
   \item \ref{DETERMINANT}.            DETERMINANT
   \item \ref{GROUP-THEORY}.           GROUP-THEORY
   \item \ref{ANALYSIS}.               ANALYSIS
   \item \ref{ANALYTICITY-VARIA}.      ANALYTICITY-VARIA
   \item \ref{ANALYTICITY-H}.          ANALYTICITY-H
   \item \ref{ANALYTICITY-B}.          ANALYTICITY-B
   \item \ref{ANALYTICITY-CLIFFORD-R}. ANALYTICITY-CLIFFORD-R
   \item \ref{ANALYTICITY-CLIFFORD-C}. ANALYTICITY-CLIFFORD-C
   \item \ref{MANIFOLD}.               MANIFOLD
   \end{itemize}

\item \ref{RELATIVITICS}.          RELATIVISTICS
   \begin{itemize}
   \item \ref{SPECIAL-RELATIVITY}. SPECIAL-RELATIVITY
   \item \ref{CONFORMALITY}.       CONFORMALITY
   \item \ref{TENSOR}.             TENSOR
   \item \ref{SPINOR}.             SPINOR
   \item \ref{TWISTOR}.            TWISTOR
   \item \ref{GENERAL-RELATIVITY}. GENERAL-RELATIVITY
   \end{itemize}

\item \ref{FIELDS}.                        FIELDS
   \begin{itemize}
   \item \ref{SPIN-1}.                     SPIN-1 (MAXWELL, PROCA)
   \item \ref{SPIN-1/2}.                  SPIN-1/2 (DIRAC, LANCZOS, PAULI, WEYL)
   \item \ref{SPIN-3/2}.                   SPIN-3/2
   \item \ref{ANALYTICITY-MAXWELL}.        ANALYTICITY-MAXWELL
   \item \ref{ANALYTICITY-DIRAC}.          ANALYTICITY-DIRAC
   \end{itemize}

\item \ref{PHYSICS}.                        PHYSICS
   \begin{itemize}
   \item \ref{PHYSICS-VARIA}.               PHYSICS-VARIA
   \item \ref{MECHANICS}.                   MECHANICS
   \item \ref{HYDRODYNAMICS}.               HYDRODYNAMICS
   \item \ref{ELECTRODYNAMICS}.             ELECTRODYNAMICS
   \item \ref{LEPTODYNAMICS}.               LEPTODYNAMICS
   \item \ref{HADRODYNAMICS}.               HADRODYNAMICS
   \item \ref{PARTICLE-PHYSICS}.            PARTICLE-PHYSICS
   \end{itemize}

\item \ref{QUANTICS}.                        QUANTICS
   \begin{itemize}
   \item \ref{QUANTUM-PHYSICS}.              QUANTUM-PHYSICS
   \item \ref{QUANTUM-ELECTRODYNAMICS}.      QUANTUM-ELECTRODYNAMICS
   \item \ref{QUATERNIONIC-QUANTUM-PHYSICS}. QUATERNIONIC-QUANTUM-PHYSICS
   \end{itemize}

\item \ref{ALLIED FORMALISMS}.   ALLIED FORMALISMS
   \begin{itemize}
   \item \ref{OCTONION}.         OCTONION
   \item \ref{GRASSMANN}.        GRASSMANN
   \item \ref{CLIFFORD}.         CLIFFORD
   \item \ref{EDDINGTON}.        EDDINGTON
   \item \ref{SEMIVECTOR}.       SEMIVECTOR
   \item \ref{HESTENES}.         HESTENES
   \end{itemize}

\item \ref{MISCELLANEA}.     MISCELLANEA
   \begin{itemize}
   \item \ref{HISTORY}.      HISTORY and APPRECIATION
   \item \ref{BIBLIOGRAPHY}. BIBLIOGRAPHY
   \end{itemize}

\item \ref{Conventions}.     Conventions used in the bibliography
 
\end{itemize}
\renewcommand{\labelitemi}{$\bullet$}    %Companion p.59
\renewcommand{\labelitemii}{\bfseries --}


\newpage


\section{\Huge Introduction}
%===========================
\label{Introduction}


\emph{The purpose of the present analytical bibliography is to present a selection, but as comprehensive as possible, of the use of biquaternions in theoretical and mathematical physics, with an emphasis on their applications to fundamental rather than applied topics.}\footnote{The only exceptions to this rule are papers or books of general interest, and papers included for completeness when they are important to understand other papers.}   

This bibliography is already available as an alphabetic list \cite{GSPON2005D}. But, in that form, it is not much more useful than the internet.
What is needed is a logically sorted bibliography  --- an analytical bibliography that can be used by physicists to solve a problem, and by mathematicians to see what is of interest to physicists.


As is well known to anybody who has tried to classify a large set of scientific papers, sorting a bibliography is a very difficult and time-consuming task: It took us over ten years to bring this bibliography into its present form.  

Of course, there is a certain amount of subjectivity and arbitrariness in designing any classification scheme.  For this reason there will be a few sentences of introduction at the beginning of every subset of papers, explaining what is being collected in every chapter and section of the bibliography.  

Similarly, in this general introduction, we would like to explain what we mean by \emph{``mathematical physics,''} the concept which appears in the titles of this series of methodological papers \cite{GSPON2005D, GSPON2005E, GSPON2005F, GSPON2005G}, and which is our main thread in compiling and sorting our alphabetical and analytical bibliographies.

According to Ludvig Faddeev, \emph{``the main goal of mathematical physics is the use of mathematical intuition for the derivation of really new results in fundamental physics''} \cite{FADDE2000-}.  In the present case, the mathematical tool is complex quaternion algebra and analysis, which is so well suited to physics in our four-dimensional world that there is no important fundamental result which cannot be elegantly formulated and concisely derived using biquaternions, and only few quaternionic generalizations of fundamental theories which do not correspond to physical reality.  This is, in a forceful way, the confirmation of the validity of ``Hamilton's conjecture,'' the intuition that motivated Hamilton's dedication to quaternions, and their applications to physics, for most of the second half of his life (see Ref.~\cite{GSPON1993B} in Sec.~\ref{PHYSICS-VARIA}).


  However, this bibliography is not restricted to just papers in which quaternions or biquaternions are used explicitly: it also covers papers in which a hypercomplex coordinate-free whole-symbol system allied to quaternions is used (e.g., Clifford-numbers, Pauli vector-matrices, Eddington-numbers, semivectors, two-component spinors, twistors, etc.), and papers in which a quaternion or biquaternion structure plays a central role.  The kind of papers that are not included are those in which a strictly conventional matrix-type formalism is used (e.g., the Pauli- or Dirac-matrix formalisms), and papers which have not been published (or would not qualify to appear) in peer-reviewed journals.

  An important criterion used in compiling our bibliography is that it includes only papers which we have read, so that we were able to attach a few keywords to each entry in the reference list.  These keywords have the format \%\%KEYWORD, where ``\%'' is the symbol used for comments in \TeX~so that they  do not appear in the compiled bibliography, and where KEYWORD is always written in capital letters and in the singular.  However, the keywords are visible and can be searched for in the  \TeX-source of the bibliography.  As a matter of fact, this is how the present ``analytical'' bibliography was created starting from the ``alphabetical'' one.   For this reason the titles of the following chapters, Chaps.~\ref{MATHEMATICS} to \ref{MISCELLANEA}, and their sections, are nothing but the main keywords attached to every reference listed in them.

Finally, in Chap.~\ref{Conventions}, at the end of the bibliography, we detail the conventions used for the labels and styles of all types of references, of which typical examples are given in Sec.~\ref{Format}.  


\begin{thebibliography}{9999}
%============================

\hyphenation{Alphabetical}

\bibitem{GSPON2005D} A. Gsponer and J.-P. Hurni, \emph{Quaternion in mathematical physics (1): Alphabetical bibliography}, Report ISRI-05-04 (4 March 2006) 100~pp.;  e-print \underline{ arXiv:math-ph/0510059 }.

\bibitem{GSPON2005E} A. Gsponer and J.-P. Hurni, \emph{Quaternion in mathematical physics (2): Analytical bibliography}, Report ISRI-05-05 (4 March 2006) 113~pp.;  e-print \underline{ arXiv:math-ph/0511092 }.

\bibitem{GSPON2005F} A. Gsponer and J.-P. Hurni, \emph{Quaternion in mathematical physics (3): Notations and definitions}, Report ISRI-05-06.

\bibitem{GSPON2005G} A. Gsponer and J.-P. Hurni, \emph{Quaternion in mathematical physics (4): Formulas and methods}, Report ISRI-05-07.

\bibitem{FADDE2000-} L.D. Faddeev, \emph{Modern mathematical physics: What it should be}, in: A. Fokas et al., Mathematical Physics 2000 (Imperial College Press, 2000) 1--8.

\end{thebibliography}


%\newpage


\section{\Huge MATHEMATICS}
%==========================
\label{MATHEMATICS}

This chapter contains a selection of mathematical quaternion-papers which are of direct interest to mathematical physics. It not does contain however numerous papers or books in which quaternions, biquaternions, or quaternion structures are primarily studied or used in the context of ``pure mathematics.''

While the concept of quaternion as defined by Hamilton has a universal acceptance, it should be stressed that there are several definitions for related concepts such as ``biquaternions'' and ``pseudo-'' or ``generalized-'' quaternions.  As is explained in our paper on notations and terminology, Ref.~\cite{GSPON2005F} of Chap.~\ref{Introduction}, we remain as much possible consistent with Hamilton's original definitions.  For example, just like quaternions will always be elements $x \in \bbH$, the term biquaternion will always refer to Hamilton's complexified quaternions (i.e., $x\in \bbB$), and not to Cayley's (which are now called ``octonions,'' $x\in \bbO$), or to Clifford's (which are in a way an anticipation of Penrose's ``twistors'').\footnote{There are four Clifford algebras of dimension 8 over the reals: $\cl_{3,0}$ and $\cl_{1,2}$ are isomorphic to $\mathbb{B}$, whereas $\cl_{0,3}$ and $\cl_{2,1}$ are Clifford's misnamed ``biquaternions.''}


\subsection{MATH-VARIA}
%----------------------
\label{MATH-VARIA}

The books by W.R.\ Hamilton, P.G.\ Tait, and C.J.\ Joly listed in this section also contain chapters or sections on the applications of quaternions to classical physics.

\begin{enumerate}

\bibitem{CAYLE1845-} A. Cayley, \emph{On certain results relating to quaternions}, Phil. Mag. and J. of Science {\bf 26} (1845) 141--145.  %%QUATERNION, %%MATH-VARIA, %%OLINDE-RODRIGUES, %%DETERMINANTS, %%QUADRIPERIODIC-FUNCTION, %%ULTRA-ELLIPTIC-FUNCTION, %$D14032002.

\bibitem{CAYLE1848-} A. Cayley, \emph{On the application of quaternions to the theory of rotations}, Phil. Mag. {\bf 33} (1848) 196--200.  %%MATH-VARIA, %%OLINDE-RODRIGUES, %$D24062002.

\bibitem{TAIT-1867-} P.G. Tait, An elementary Treatise on Quaternions (Clarendon, Oxford, 1867) 321~pp. %%BOOK, %%QUATERNION, %%MATH-VARIA, %$D24062002.

\bibitem{LAISA1881-} C.A. Laisant, Introduction \`a la M\'ethode des Quaternions (Gauthier-Villars, Paris, 1881) 242~pp. %%BOOK, %%QUATERNION, %%MATH-VARIA, %$D24062002.

\bibitem{TAIT-1886-} P.G. Tait, \emph{Quaternions}, Encyclopedia Britannica {\bf } (1886)  SP-2:445--455.  %%QUATERNION, %%MATH-VARIA, %$D24062002.

\bibitem{HAMIL1891-} W.R. Hamilton, Elements of Quaternions, Vol.~ I et II (First edition 1866; second edition edited and expanded by C.J. Joly 1899-1901; reprinted by Chelsea Publishing, New York, 1969) 1185~pp.  %%BOOK, %%QUATERNION, %%MATH-VARIA, %%PHYSICS-VARIA, %%HAMILTON, %$D20032002.

\bibitem{COMBE1898-} G. Combebiac, \emph{Sur l'application du calcul des biquaternions \`a la g\'eom\'etrie plane}, Bull. Soc. Math. France {\bf 26} (1898) 259--263. %%QUATERNION, %%MATH-VARIA, %$D06072008.

\bibitem{JOLY-1905-} C.J. Joly, A Manual of Quaternions (MacMillan, London, 1905) 320~pp.  %%QUATERNION, %%MATH-VARIA, %$D24062002.

\bibitem{MACFA1910B} A. Macfarlane, \emph{Unification and development of the principle of the algebra of space}, in: A. Macfarlane, ed., Bull. of the Inter. Assoc. for promoting the study of quaternions and allied systems of mathematics (New Era Printing Company, Lancaster PA, 1910) 41--92. %%BOOK, %%QUATERNION, %%MATH-VARIA, %%EXPONENTIAL, %%LOGARITHM, %%GRASSMANN, %%FEYNMAN, %$D20032002, %$C Page 63: Ordering in exponentation as a means for univocally defining the logarithm.

\bibitem{MADEL1925-} E. Madelung, Die Mathematischen Hilfsmittel des Physikers (Springer Verlag, Berlin, 1925, 6th edition 1957) 535~pp.  %%QUATERNION, %%MATH-VARIA, %%LANCZOS, %$D06072008. 
 
\bibitem{QUADL1979-} D. Quadling, \emph{Q for quaternions}, Math. Gazette {\bf 63} (1979) 98--110. %%QUATERNION, %%MATH-VARIA, %$D14032002.

\bibitem{GERAR1985-} P. Gerardin and W.C.W. Li, \emph{Fourier transforms of representations of quaternions}, J. Reine. Angew. Math. {\bf 359} (1985) 121--173. %%QUATERNION, %%MATH-VARIA, %%GROUP-THEORY, %$D24062002.

\bibitem{GLUCK1986-} H. Gluck, \emph{The geometry of Hopf fibrations}, L'Enseignement Math\'ematique {\bf 32} (1986) 173--198.   %%QUATERNION, %%MATH-VARIA, %$D06072008.

\bibitem{GROGE1992-} D. Gr\"oger, \emph{Homomorphe Kopplungen auf des K\"orper der reellen Quaternionen}, Arch. Math. {\bf 58} (1992) 354--359. %%QUATERNION, %%MATH-VARIA, %$D14032002.

\bibitem{HORAD1993-} A.F. Horadam, \emph{Quaternion recurrence relations}, Ulam Quarterly {\bf 2} (1993) 23--33.   %%QUATERNION, %%MATH-VARIA, %$D06072008.

\bibitem{GENTI1994-} G. Gentili, S. Marchiafava, and M. Pontecorvo, eds., Proc. of the Meeting on Quaternionic Structures in Mathematics and Physics (SISSA, Trieste, Italy, September 5-9, 1994) 270 pp. Available at\\
\underline{ http://www.math.unam.mx/EMIS/proceedings/QSMP94/contents.html }.   %%QUATERNION, %%MATH-VARIA, %%MANIFOLD, %%KAHLER, %%ANALYSIS, %$D04032006.

\bibitem{WARD-1997-} J.P. Ward, \emph{Quaternions and Cayley numbers} (Kluwer, Dordrecht, 1997) 237~pp. %%BOOK, %%QUATERNION, %%MATH-VARIA, %%OCTONION, %$D20032002.

\bibitem{STROP1998-} M. Stroppel, \emph{A characterization of quaternion planes, revisited}, Geometriae Dedicata {\bf 72} (1998) 179--187.   %%QUATERNION, %%MATH-VARIA, %$D06072008.

\bibitem{MARCH2001-} S. Marchiafava, P. Piccinni and M. Pontecorvo, eds., Proceedings of the 2nd meeting on ``Quaternionic structures in mathematics and physics,'' Rome, 6--10 September 1999 (World
Scientific, Singapore, 2001) 469~pp. %%BOOK, %%QUATERNION, %%MATH-VARIA, %%PHYSICS-VARIA, %$D04032006.

\bibitem{TRAVE2001-} L. Traversoni, \emph{Image analysis using quaternion wavelets}, in: E.B Corrochano and G. Sobczyk, eds., Geometric Algebra with Applications in Science and Engineering (Birk\"auser, Basel, 2001) 326--345.  %%QUATERNION, %%MATH-VARIA, %$D04032006.

\bibitem{JANOV2003-} D. Janovska and G. Opfer, \emph{Given's transformation applied to quaternion valued vectors}, BIT Numer. Math. {\bf 43} (2003) 991--1002.   %%QUATERNION, %%MATH-VARIA, %$D06072008.

\bibitem{JIANG2003A} T. Jiang and M. Wei, \emph{Equality constrained least squares problem over quaternion field}, Appl. Math. Lett. {\bf 16} (2003) 883--888.  %%QUATERNION, %%MATH-VARIA, %$D06072008.

\bibitem{LAM--2003-} T.Y. Lam, \emph{Hamilton's quaternions}, in: M. Hazewinkel, ed., Handbook of Algebra {\bf 3} (Elsevier, Amsterdam, 2003) 429--454. %%QUATERNION, %%MATH-VARIA, %$D04032006.

\bibitem{LYONS2003-} D.W. Lyons, \emph{An elementary introduction to the Hopf   fibration}, Mathematics Mag. {\bf 76} (2003) 87--98.  %%QUATERNION, %%MATH-VARIA, %$D06072008.

\bibitem{DANCE2004-} A.S. Dancer, H.R. Jorgensen and A.F. Swann, \emph{Metric geometries over the split quaternions} (2004) 23~pp.; e-print \underline{ arXiv:math/0412215 }. %%QUATERNION, %%MATH-VARIA, %$D06072008.

\bibitem{GOLDM2004-} W.M. Goldman, \emph{An exposition of results of Fricke} (2004) 18~pp.; e-print \underline{ arXiv:math/0402103 }.   %%QUATERNION, %%MATH-VARIA, %$D06072008.

\bibitem{JIN--2005-} D.Jin and G Jin, \emph{Matrix maps for substitution sequences in the biquaternion representation}, Phys. Rev. {\bf B 71} (2005) 014212.  %%QUATERNION, %%MATH-VARIA, %$D06072008.

\bibitem{ELKIE2006-} N. Elkies, L.M. Pretorius and K.J. Swanepoel, \emph{Sylvester-Gallai theorems for complex numbers and quaternions}, Discrete \& Computational Geom. {\bf 35} (2006) 361--373.; e-print \underline{ arXiv:math/0403023 }.  %%QUATERNION, %%MATH-VARIA, %$D06072008.

\bibitem{IVANO2006-} S. Ivanov, I. Minchev and D. Vassilev, \emph{Quaternionic contact Einstein structures and the quaternionic contact Yamabe problem} (2006) 51~pp.; e-print \underline{ arXiv:math/0611658 }.   %%QUATERNION, %%MATH-VARIA, %$D06072008.

\bibitem{ELL--2005-} T.A. Ell and S.J. Sangwine, \emph{Quaternion involutions}, Computers and Math. Appl. {\bf 53} (2007) 137--143; e-print \underline{ arXiv:math/0506034 }.  %%QUATERNION, %%MATH-VARIA, %$D06072008.

\bibitem{WANG-2007-} D. Wang, \emph{The largest sample eigenvalue distribution in the rank 1 quaternionic spiked model of Wishart ensemble}, 47~pp.; e-print \underline{ arXiv:0711.2722 }.   %%QUATERNION, %%MATH-VARIA, %$D06072008.


\end{enumerate}


\subsection{ALGEBRA}
%-------------------
\label{ALGEBRA}

Emphasis on algebraic operations and basic properties of algebras.  

See also GROUP-THEORY, Sec.~\ref{GROUP-THEORY}.

\begin{enumerate}

\bibitem{HAMIL1844-} W.R. Hamilton, \emph{On a new species of imaginary quantites connected with a theory of Quaternions}, Proc. Roy. Irish Acad {\bf 2} (1843) 424--434. %%QUATERNION, %%ALGEBRA, %$D24062002, %$C The seminal paper on quaternions, read on 13 Novembre 1843.

\bibitem{HAMIL1853-} W.R. Hamilton, \emph{On the geometrical interpretation of some results obtained by calculation with biquaternions}, Proc. Roy. Irish Acad. {\bf 5} (1853) 388--390. %%QUATERNION, %%ALGEBRA, %%BIVECTOR, %%BIQUATERNION, %$D14032002.

\bibitem{PEIRC1881-} B. Peirce, \emph{Linear associative algebras}, Amer. J. Math. {\bf 4} (1881) 97--229.  %%QUATERNION, %%ALGEBRA, %$D24062002, %$C Origin of "Peirce decomposition".

\bibitem{BUCHH1885-} A. Buchheim, \emph{A memoir on biquaternions}, Am. J. Math. {\bf 7} (1885) 293--326. %%QUATERNION, %%ALGEBRA, %$D06072008. %$C Various kinds of biquaternions...

\bibitem{MACFA1899-} A. Macfarlane, \emph{Hyperbolic quaternions}, Proc. Roy. Soc. Edinburgh {\bf 23} (1899/1900) 169--181. %%QUATERNION, %%ALGEBRA, %%PSEUDOQUATERNION, %$D14032002.

\bibitem{COMBE1902-} G. Combebiac, Calcul des Triquaternions (Gauthier-Villars, Paris, 1902) 122~pp. Reviewed by C.J. Joly, The Mathematical Gazette {\bf 2} No 35 (1902) 202--204. %%QUATERNION, %%ALGEBRA, %$D06072008.

\bibitem{CARTA1908-} E. Cartan, \emph{Nombre complexes: Expos\'e d'apr\`es l'article allemand de E. Study}, Encyclop. Sc. Math. {\bf 15} (1908); reprinted in: E. Cartan, Oeuvres Compl\`etes, Partie II (Editions du CNRS, Paris, 1984) 107--467. %%QUATERNION, %%ALGEBRA, %$D04032006.

\bibitem{DICKS1918-} L.E. Dickson, \emph{On quaternion and their generalization and the history of the eight square theorem}, Ann. of Math. {\bf 20} (1918) 155--171.  %%QUATERNION, %%ALGEBRA, %%HISTORY, %$D30112005.

\bibitem{MOORE1922B} C.L.E. Moore, \emph{Hyperquaternions}, J. of Math. and Phys. {\bf 1} (1922) 63--77. %%QUATERNION, %%ALGEBRA, %$D04032006.

\bibitem{DUPAS1928-} L.G. DuPasquier, \emph{Sur une th\'eorie nouvelle des id\'eaux de quaternion complexes}, Atti Congr. Int. Mat. Bologna (3--10 Settembre, 1928) Vol.~2, p.135--143. %%QUATERNION, %%ALGEBRA, %$D24062002.

\bibitem{GRIZE1932-} J. Grize, \emph{Sur les corps alg\'ebriques dont les nombres s'expriment rationnellement \`a l'aide de racines carr\'ees et sur les quaternions complexes}, Th\`ese (Universit\'e de Neuch\^atel, 1932) 95~pp.  %%$C Complex here means here that the quaternion coordinates are not rational numbers.

\bibitem{FUETE1933-} R. Fueter, \emph{Quaternionringe}, Comm. Math. Helv. {\bf 6} (1933/1934) 199--222. %%QUATERNION, %%ALGEBRA, %$D24062002.

\bibitem{ORE--1933-} O. Ore, \emph{Theory of non-commutative polynomials}, Ann. of Math. {\bf 34} (1933) 480--508. %%QUATERNION, %%ALGEBRA, %$D14032002.

\bibitem{FUETE1935B} R. Fueter, \emph{Zur Theorie der Brandtschen Quaternionenalgebren}, Math. Annalen. {\bf 110} (1935) 650-661. %%QUATERNION, %%ALGEBRA, %$D24062002.

\bibitem{BROWN1940-} D.M. Brown, \emph{Arithmetics of rational generalized quaternion algebras}, Bull. Amer. Math. Soc. {\bf 46} (1940) 899--908. %%QUATERNION, %%ALGEBRA, %$D04032006.

\bibitem{BRAND1942-} L. Brand, \emph{The roots of a quaternion}, Amer. Math. Monthly {\bf 49} (1942) 1519--1520. %%QUATERNION, %%ALGEBRA, %$D14032002.

\bibitem{NIVEN1942-} I. Niven, \emph{The roots of a quaternion}, Amer. Math. Monthly {\bf 49} (1942) 386--388. %%QUATERNION, %%ALGEBRA, %$D14032002.

\bibitem{WILLI1945-} C.S. Williams and G. Pall, \emph{The thirty-nine systems of quaternions with a positive norm-form and satisfactory factorability}, Duke Math. J. {\bf 12} (1945) 527--539. %%QUATERNION, %%ALGEBRA, %$D04032006.

\bibitem{SALZE1952-} H.E. Salzer, \emph{An elementary note on powers of quaternions}, Amer. Math. Monthly {\bf 59} (1952) 298--300. %%QUATERNION, %%ALGEBRA, %$D14032002.

\bibitem{LEUM-1953-} M. Leum and M.F. Smiley, \emph{A matric proof of the fundamental theorem of algebra for real quaternions}, Amer. Math. Monthly {\bf 60} (1953) 99--100. %%QUATERNION, %%ALGEBRA, %$D04032006.

\bibitem{PALL-1957-} G. Pall and O. Taussky, \emph{Application of quaternions to the representations of a binary quadratic form as a sum of four squares}, Proc. Roy. Irish Acad. {\bf A 58} (1957) 23--28. %%QUATERNION, %%ALGEBRA, %$D14032002.

\bibitem{AEBER1959-} G. Aeberli, \emph{Der Zusammenhang zwischen quatern\"aren quadratischen Formen und Idealen in Quaternionenring}, Comm. Math. Helv. {\bf 33} (1959) 212--239. %%QUATERNION, %%ALGEBRA, %$D14032002.
 
\bibitem{SMITH1959-} A.C. Smith, \emph{Hamiltonian algebras}, Rev. Fac. Sci. Istanbul Univ. {\bf 24} (1959) 69--79. %%QUATERNION, %%ALGEBRA, %$D04032006.

\bibitem{GROSS1960-} H. Gross, \emph{Darstellungsanzahlen von quatern\"aren quadratischen Stammformen mit quadratischer Diskriminante}, Comm. Math. Helv. {\bf 34} (1960) 198--221. %%QUATERNION, %%ALGEBRA, %$D14032002.

\bibitem{EBERL1963-} W.F. Eberlein, \emph{The geometric theory of quaternions}, Amer. Math. Monthly {\bf 70} (1963) 952--954. %%QUATERNION, %%ALGEBRA, %$D04032006.

\bibitem{HORAD1963-} A.F. Horadam, \emph{Complex Fibonacci numbers and Fibonacci quaternions}, Amer. Math. Monthly {\bf 70} (1963) 289--291. %%QUATERNION, %%ALGEBRA, %$D04032006.

\bibitem{KAPLA1969-} I. Kaplanski, \emph{Submodules of quaternion algebras}, Proc. London Math. Soc. {\bf 19} (1969) 219--232.  %%QUATERNION, %%ALGEBRA, %$D06072008. %$C Kaplanski est l'un des pionniers des superalg�bres de Lie

\bibitem{OKUBO1978-} S. Okubo, \emph{Pseudo-quaternion and pseudo-octonion algebras}, Hadronic J. {\bf 1} (1978) 1250--1278. %%QUATERNION, %%ALGEBRA, %$D04032006, %$C These algebra circumvent the Hurwitz theorem. "Pseudo-" has here a very specific meaning.

\bibitem{LEWIS1980-} D.W. Lewis, \emph{A product of Hermitian forms over quaternion division-algebras}, J. London Math. Soc. {\bf 22} (1980) 215--220.  %%QUATERNION, %%ALGEBRA, %$D30112005.

\bibitem{ESTES1981-} D.R.Estes and O. Taussky, \emph{Remarks concerning sums of three squares and quaternion commutators identities}, Linear Alg. Appl. {\bf 35} (1981) 279--285. %%QUATERNION, %%ALGEBRA, %$D14032002.

\bibitem{WALKE1981-} G. Walker, \emph{Estimates for the complex and quaternionic James numbers}, Quart. J. Math. {\bf 32} (1981) 467--489. %%ABSTRACT-ONLY %%QUATERNION, %%ALGEBRA, %$D30112005.

\bibitem{JANTZ1982-} R.T. Jantzen, \emph{Generalized quaternions and spacetime symmetries}, J. Math. Phys. {\bf 23} (1982) 1741--1746.  %%QUATERNION, %%ALGEBRA, %%PSEUDOQUATERNION, %%GROUP-THEORY, %%GENERAL-RELATIVITY, %$D14032002.

\bibitem{FARN1984-} R. Farnsteiner, \emph{Quaternionic Lie algebra}, Linear Alg. Appl. {\bf 61} (1984) 1225--231.  %%QUATERNION, %%ALGEBRA, %%GROUP-THEORY, %$D14032002.

\bibitem{PLEBA1988-} J.F. Plebanski and M. Przanowski, \emph{Generalizations of the quaternion algebra and Lie-algebras}, J. Math. Phys. {\bf 29} (1988) 529--535. %%QUATERNION, %%ALGEBRA, %$D30112005, %$C "quaternion-like algebras".

\bibitem{MINAC1989A} J. Minac, \emph{Quaternion fields inside Pythagorean closure}, J. of Pure and Appl. Algebra {\bf 57} (1989) 79--82.  %%QUATERNION, %%ALGEBRA, %$D04032006.

\bibitem{MINAC1989B} J. Minac, \emph{Classes of quaternion algebras in the Brauer group}, Rocky Mountains J. of Math. {\bf 19} (1989) 819--831.  %%QUATERNION, %%ALGEBRA, %%GROUP-THEORY, %$D04032006.

\bibitem{SHAW-1989-} R. Shaw, \emph{A two-dimensional quaternionic construction of an 8-dimensional ternary composition algebra}, Nuovo Cim. {\bf B 104} (1989) 163--176.  %%QUATERNION, %%ALGEBRA, %%OCTONION, %%GURSEY, %$D30112005.

\bibitem{SERGE1991-} V.V. Sergeichuk, \emph{Classification of sesquilinear forms, pairs of Hermitian-forms, self-conjugated and isometric operators over the division ring of quaternions}, Math. Notes {\bf 49} (1991) 409--414.  %%QUATERNION, %%ALGEBRA, %$D30112005.

\bibitem{LOUNE1993B}  P. Lounesto, \emph{On invertibility in lower dimensional Clifford algebras}, Adv. Appl. Clifford Alg. {\bf 3} (1993) 133--137. %%QUATERNION, %%ALGEBRA, %%CLIFFORD, %$D24062002.

\bibitem{THOMA1995-} O. Thomas, \emph{A local-global theorem for skew-Hermitian forms over quaternion algebras}, Commun. in Algebra {\bf 23} (1995) 1679--1704.  %%QUATERNION, %%ALGEBRA, %$D04032006.

\bibitem{MCCON1998-} J.C. McConnell, \emph{Division algebras --- Beyond the quaternions}, Amer. Math. Monthly {\bf 105} (1998) 154--162.  %%QUATERNION, %%ALGEBRA, %$D30112005.

\bibitem{ARENS1999-} R. Arens, M. Goldberg, and W.A.J. Luxemburg, \emph{Stable norms on complex numbers and quaternions}, J. of Algebra {\bf 219} (1999) 1--15.  %%QUATERNION, %%ALGEBRA, %$D14032002. 

\bibitem{ERIKS1999-} S.L. Eriksson-Bique, \emph{The binomial theorem for hypercomplex numbers}, Ann. Acad. Scient. Fennicae Math. {\bf 24} (1999) 225--229. %%QUATERNION, %%ALGEBRA, %$D14032002.
 
\bibitem{HOLM-1999-} T. Holm, \emph{Derived equivalence classification of algebras of dihedral, semidihedral, and quaternion type}, J. Algebra, {\bf 211} (1999) 159--205.  %%QUATERNION, %%ALGEBRA, %$D04032006. 

\bibitem{SEROD2001-} R. Serodio, E. Pereira and J. Vitoria, \emph{Computing the zeros of quaternion polynomials}, Computers and Mathematics with Applications {\bf 42} (2001) 1229--1237.   %%QUATERNION, %%ALGEBRA, %$D06072008.

\bibitem{GSPON2002D} A. Gsponer, \emph{Explicit closed-form parametrization of $SU(3)$ and $SU(4)$ in terms of complex quaternions and elementary functions}, Report ISRI-02-05 (22 November 2002) 17~pp.;  e-print \underline{ arXiv:math-ph/0211056 }.  %%QUATERNION, %%ALGEBRA, %%HADRODYNAMICS, %$10092005.

\bibitem{DEUTS2005-} J.I. Deutsch, \emph{A quaternionic proof of the representation formula of a quaternary quadratic form}, J. Number Th. {\bf 113} (2005) 149--174; e-print \underline{ arXiv:math/0406434 }.  %%QUATERNION, %%ALGEBRA, %$D06072008.

\bibitem{BOGUS2006-} A.A. Bogush and V.M. Red'kov, \emph{On unique parametrization of the linear group $GL(4,\mathbb{C})$ and its subgroups by using the Dirac matrix algebra basis} (2006) 23~pp.; e-print \underline{ arXiv:hep-th/0607054 }. %%QUATERNION, %%GROUP-THEORY, %%DIRAC, %%PARITY, %%SEMIVECTOR, %$D06072008. 

\bibitem{DELEO2006-} S. De Leo, G. Ducati and V. Leonardis, \emph{Zeros of unilateral quaternionic polynomials}, Electronic J. Linear Algebra {\bf 15} (2006) 297--313. %%QUATERNION, %%ALGEBRA, %$D06072008.

\bibitem{SANGW2006-} S.J. Sangwine and N. Le Bihan, \emph{Quaternion singular value decomposition based on bidiagonalization to a real matrix using quaternion Householder transformations}, Applied Math. Computation {\bf 182} (2006) 727--738; e-print \underline{ arXiv:math/0603251 }.  %%QUATERNION, %%ALGEBRA, %$D06072008.

\bibitem{JANOV2007-} D. Janovska and G. Opfer, \emph{Computing quaternionic roots by Newton's method}, Electronic Trans. Numer. Anal. {\bf 26} (2007) 82--102.   %%QUATERNION, %%ALGEBRA, %$D06072008.

\bibitem{JURIA2007-} S.O. Juriaans, I.B.S. Passi, and A.C. Souza-Filho, \emph{Hyperbolic unit groups and quaternion algebras} (2007) 15~pp.; e-print \underline{ arXiv:0709.2161 }.  %%QUATERNION, %%ALGEBRA, %$D06072008.

\bibitem{SANGW2008B} S.J. Sangwine and N. Le Bihan, \emph{Quaternion polar representation with a complex modulus and complex argument inspired by the Cayley-Dickson form} (2008) 4~pp.; e-print \underline{ arXiv:0802.0852 }.  %%QUATERNION, %%ALGEBRA, %$D06072008.

\end{enumerate}


\subsection{INTEGRAL-QUATERNION}
%-------------------------------
\label{INTEGRAL-QUATERNION}

See also GROUP-THEORY, Sec.~\ref{GROUP-THEORY}.

\begin{enumerate}


\bibitem{LIPSC1886-} R. Lipschitz, \emph{Recherches sur la transformation, par des substitutions r\'eelles, d'une somme de deux ou trois carr\'es en elle-m\^eme}, J. de Math\'ematiques {\bf 2} (1886) 373--439.  %%QUATERNION, %%INTEGRAL-QUATERNION, %%PRIME-QUATERNION, %$D04032006, %$C Lipschitz's definition of integral quaternions.

\bibitem{HURWI1896-} A. Hurwitz,  \emph{Uber die Zahlentheorie der Quaternionen}, in: Mathematische Werke von Adolf Hurwitz, Vol.~2 (Birkh\"auser, Basel, 1963) 303--330.  %%QUATERNION, %%INTEGRAL-QUATERNION, %$D14032002.

\bibitem{DUPAS1916-} L.-G. DuPasquier, \emph{Sur l'arithm\'etique des nombres hypercomplexes}, L'Ensei\-gnement Math\'ematique {\bf 18} (1916) 201--259.  %%QUATERNION, %%INTEGRAL-QUATERNION, %$D06072008.

\bibitem{HURWI1919-} A. Hurwitz, \emph{Vorlesung \"uber die Zahlentheorie der Quaternionen} (Springer, Berlin, 1919) 74~pp. %%BOOK, %%QUATERNION, %%INTEGRAL-QUATERNION, %$D14032002.

\bibitem{DICKS1921-} L.E. Dickson, \emph{Arithmetic of quaternions}, Proc. London Math. Soc. {\bf 20} (1921) 225--232. %%QUATERNION, %%INTEGRAL-QUATERNION, %$D14032002.

\bibitem{DICKS1924-} L.E. Dickson, \emph{On the theory of numbers and generalized quaternions}, Amer. J. Math. {\bf 46} (1924) 1--16. %%QUATERNION, %%INTEGRAL-QUATERNION, %$D04032006.

\bibitem{OLSON1930-} H.L. Olson, \emph{Doubly divisible quaternions}, Ann. of Math. {\bf 31} (1930) 371--374. %%QUATERNION, %%INTEGRAL-QUATERNION,  %$D04032006.

\bibitem{OCONN1939-} R.E. O'Connor and G. Pall, \emph{The quaternion congruence $\overline{t}at = b (\operatorname{mod} ~ g)$}, Am. J. Math. {\bf 61} (1939) 487--508. %%QUATERNION, %%INTEGRAL-QUATERNION, %$D24062002.

\bibitem{PALL-1940-} G. Pall, \emph{On the arithmetic of quaternions}, Trans. Amer. Math. Soc. {\bf 47} (1940) 487--500. %%QUATERNION, %%INTEGRAL-QUATERNION, %$D04032006.

\bibitem{PALL-1942-} G. Pall, \emph{Quaternions and sums of three squares}, Amer. J. Math. {\bf 64} (1942) 503--513. %%QUATERNION, %%INTEGRAL-QUATERNION, %$D04032006.

\bibitem{BENNE1943-} G. Benneton, \emph{Sur l'arithm\'etique des quaternions et des biquaternions (octonions)}, Ann. Sci. Ec. Norm. Sup. {\bf 60} (1943) 173--214. %%QUATERNION, %%INTEGRAL-QUATERNION, %$D14032002.

\bibitem{NIVEN1946-} I. Niven, \emph{A note on the number theory of quaternions}, Duke Math. J. {\bf 13} (1946) 397--400. %%QUATERNION, %%INTEGRAL-QUATERNION, %$D04032006.

\bibitem{RUTLE1952-} W.A. Rutledge, \emph{Quaternions and Hadamard matrices}, Proc. Amer. Math. Soc. {\bf 3} (1952) 625--630. %%QUATERNION, %%MATRIX, %%INTEGRAL-QUATERNION, %$D04032006.

\bibitem{PIZER1976A} A. Pizer, \emph{On the arithmetic of quaternion algebras I}, Acta Arith. {\bf 31} (1976) 61--89.  %%QUATERNION, %%INTEGRAL-QUATERNION, %$D04032006.

\bibitem{PIZER1976B} A. Pizer, \emph{On the arithmetic of quaternion algebras II}, J. Math. Soc. Japan {\bf 28} (1976) 676--688.  %%QUATERNION, %%INTEGRAL-QUATERNION, %$D04032006.

\bibitem{MUSES1980-} C. Mus\`es, \emph{Hypernumbers and quantum field theory with a summary of physically applicable hypernumber arithmetics and their geometries}, Appl. Math. and Comput. {\bf 6} (1980) 63--94. %%QUATERNION, %%INTEGRAL-QUATERNION, %$D19032002.

\bibitem{THOMP1988-} R.C. Thompson, \emph{Integral quaternion matrices}, Lin. Algebra Appl. {\bf 104} (1988) 183-185. %%QUATERNION, %%INTEGRAL-QUATERNION, %%MATRIX, %$D14032002.

\bibitem{KOCA1989A} M. Koca and N. Osdes, \emph{Division algebras with integral elements}, J. Phys. {\bf A22} (1989) 1469--1493. %%QUATERNION, %%INTEGRAL-QUATERNIONS, %%OCTONION, %$D24062002.

\bibitem{KOCA1989B} M. Koca, \emph{Icosian versus octonions as descriptions of the $E_8$ lattice}, J. Phys. {\bf A22} (1989) 1949--1952. %%QUATERNION, %%INTEGRAL-QUATERNION, %%OCTONION, %$D24062002.

\bibitem{KOCA1989C} M. Koca, \emph{$E_8$ lattice with icosians and $Z_5$ symmetry}, J. Phys. {\bf A22} (1989) 4125--4134. %%QUATERNION, %%INTEGRAL-QUATERNION, %%OCTONION, %$D24062002.

\bibitem{KOCA1992-} M. Koca, \emph{Symmetries of the octonionic root system of $E_8$}, J. Math. Phys. {\bf 33} (1992) 497--510. %%QUATERNION, %%INTEGRAL-QUATERNION, %%OCTONION, %%STATISTICAL-MECHANICS, %$D19032002.

\bibitem{CONRA2002-} E. Conrad, \emph{Jacobi's Four Square Theorem} (13 March 2002) 11~pp.  Avail. at\\ \underline{ http://www.math.ohio-state.edu/~econrad/Jacobi/sumofsq/sumofsq.html }. %%QUATERNION, %%INTEGRAL-QUATERNION, %$D06072008. 

\bibitem{GOREN2005-} Y. Goren, \emph{Quaternions and Arithmetic}, Colloquium, UCSD (October 27, 2005) 14~pp.  Available at\\ \underline{ http://www.math.mcgill.ca/goren/PAPERSpublic/quaternions.pdf }.   %%QUATERNION, %%INTEGRAL-QUATERNIONS, %$D06072008. 


\end{enumerate}



\subsection{EQUATION}
%--------------------
\label{EQUATION}

See also DETERMINANT, Sec.~\ref{DETERMINANT}.

\begin{enumerate}

\bibitem{CAYLE1885-} A. Cayley, \emph{On the quaternion equation $qQ-Qq'=0$}, Mess. of Math. {\bf 14} (1885) 108--112.  %%QUATERNION, %%EQUATION, %$D14032002.

\bibitem{ORE--1931-} O. Ore, \emph{Linear equations in non-commutative fields}, Ann. of Math. {\bf 32} (1931) 463--477. %%QUATERNION, %%EQUATION, %$D24062002.

\bibitem{NIVEN1941-} I. Niven, \emph{Equations in quaternions}, Amer. Math. Monthly {\bf 48} (1941) 654--661. %%QUATERNION, %%EQUATION, %$D14032002.

\bibitem{EILEN1944-} S. Eilenberg and I. Niven, \emph{The ``fundamental theorem of algebra'' for quaternions}, Bull. Am. Math. Soc. {\bf 509} (1944) 246--248. %%QUATERNION, %%EQUATION, %$D14032002.

\bibitem{JOHNS1944-} R.E. Johnson, \emph{On the equation $x\alpha = x+\beta$ over an algebraic division ring}, Bull. Am. Math. Soc. {\bf 50} (1944) 202--207. %%QUATERNION, %%EQUATION, %$D14032002.

\bibitem{POLLA1960-} B. Pollak, \emph{The equation $\bar{t}at=b$ in a quaternion algebra},  Duke Math. J. {\bf 27} (1960) 261--271. %%QUATERNION, %%EQUATION, %$D14032002.

\bibitem{UTZ--1979-} W.R. Utz, \emph{The matrix equation $X^2=A$}, Amer. Math. Month. {\bf 86} (1979) 855--856. %%QUATERNION, %%EQUATION, %$D14032002.

\bibitem{BECK-1979-} B. Beck, \emph{Sur les \'equations polynomiales dans les quaternions}, L'enseignement math\'ematique {\bf 25} (1979) 193--201. %%QUATERNION, %%EQUATION, %$D14032002.

\bibitem{KOBAL1991-} D. Kobal and P. Semrl, \emph{A result concerning additive maps on the set of quaternions and an application}, Bull. Austral. Math. Soc. {\bf 44} (1991) 477--482. %%QUATERNION, %%EQUATION, %$D14032002.

\bibitem{LIPIN1996-} Huang Liping, \emph{The matrix equation $AXB-GXD=E$ over the quaternion field}, Linear Alg. and its Appl. {\bf 234} (1996) 197--208. %%QUATERNION, %%EQUATION, %$D14032002.

\bibitem{LIPIN1998-} Huang Liping, \emph{The quaternion matrix equation $\sum AXB=E$}, Acta. Math. Sinica {\bf 14} (1998) 91--98. %%QUATERNION, %%EQUATION, %$D14032002.

\bibitem{SCHWA2000-} A. Schwarz, \emph{Noncommutative algebraic equations and the noncommutative eigenvalue problem}, Lett. Math. Phys. {\bf 52} (2000) 177--184. %%QUATERNION, %%EQUATION, %$D14032002.

\bibitem{JIANG2003-} T. Jiang and M. Wei, \emph{On solutions of the matrix equations $X-AXB=C$ and $X-A\overline{X} B=C$}, Lin. Alg. Appl. {\bf 367} (2003) 225--233.  %%QUATERNION, %%EQUATION, %$D06072008.

\bibitem{JIANG2005A} T. Jiang and M. Wei, \emph{On a solution of the quaternion matrix equation $X-A{\overline X} B=C$ and its application},  Acta Math. Sinica {\bf 21} (2005) 483--490.  %%QUATERNION, %%EQUATION, %$D06072008.

\bibitem{JANOV2006-} D. Janovska and G. Opfer, \emph{Linear equations in quaternions}, in: Proceedings of ENUMATH 2005, the 6th European Conference on Numerical Mathematics and Advanced Applications  Santiago de Compostela, Spain, July 2005, Numerical Mathematics and Advanced Applications (Springer, Berlin, 2006) 945--953.   %%QUATERNION, %%EQUATION, %$D06072008.

\end{enumerate}



\subsection{LINEAR-FUNCTION}
%---------------------------
\label{LINEAR-FUNCTION}

See also MATRIX, Sec.~\ref{MATRIX}.

\begin{enumerate}


\bibitem{JOLY-1894-} C.J. Joly, \emph{The theory of linear vector functions}, Trans. Roy. Irish Acad. {\bf 30} (1892/1896) 597--647.  %%QUATERNION, %%LINEAR-FUNCTION, %$D14032002.

\bibitem{JOLY-1895-} C.J. Joly, \emph{Scalar invariants of two linear vector functions}, Trans. Roy. Irish Acad. {\bf 30} (1892/1896) 707--728. %%QUATERNION, %%LINEAR-FUNCTION, %$D14032002.

\bibitem{JOLY-1896-} C.J. Joly, \emph{Quaternion invariants of linear vector functions and quaternions determinants}, Proc. Roy. Irish Acad. {\bf 4} (1896) 1--15. %%QUATERNION, %%LINEAR-FUNCTION, %%DETERMINANT, %$D14032002.

\bibitem{TAIT-1896-} P.G. Tait, \emph{On the linear and vector function}, Proc. Roy. Soc. Edinburgh {\bf } (18 May and 1 June 1896) SP-2:406--409.  %%QUATERNION, %%LINEAR-FUNCTION, %$D24062002.

\bibitem{TAIT-1897A} P.G. Tait, \emph{On the linear and vector function}, Proc. Roy. Soc. Edinburgh {\bf } (1 March 1897) SP-2:410--412.  %%QUATERNION, %%LINEAR-FUNCTION, %$D24062002.

\bibitem{TAIT-1897B} P.G. Tait, \emph{Note on the solution of equations in linear and vector functions}, Proc. Roy. Soc. Edinburgh {\bf } (7 June 1897) SP-2:413--419.  %%QUATERNION, %%LINEAR-FUNCTION, %$D24062002.

\bibitem{TAIT-1899A} P.G. Tait, \emph{On the linear and vector function}, Proc. Roy. Soc. Edinburgh {\bf 23} (1 May 1899) SP-2:424--426.  %%QUATERNION, %%LINEAR-FUNCTION, %$D24062002.

\bibitem{JOLY-1902D} C.J. Joly, \emph{The multi-linear quaternion function}, Proc. Roy. Irish Acad. {\bf A 8} (1902) 47--52. %%QUATERNION, %%LINEAR-FUNCTION, %$D14032002.

\bibitem{HITCH1920-} F.L. Hitchcock, \emph{A study of the vector product $V \phi \alpha \theta \beta$}, Proc. Roy. Irish Acad. {\bf A 35} (1920) 30--37. %%QUATERNION, %%LINEAR-FUNCTION, %$D14032002.

\bibitem{ELL--2007-} T.A. Ell, \emph{On systems of linear quaternion functions} (2007) 11~pp.; e-print \underline{ arXiv:math/0702084 }.  %%QUATERNION, %%LINEAR-FUNCTION, %$D06072008.

\bibitem{SANGW2008A} S.J. Sangwine, \emph{Canonic form of linear quaternion functions} (2008) 4~pp.; e-print \underline{ arXiv:0801.2887 }.  %%QUATERNION, %%LINEAR-FUNCTION, %$D06072008.


\end{enumerate}



\subsection{MATRIX}
%------------------
\label{MATRIX}

See also DETERMINANT, Sec.~\ref{DETERMINANT}, EQUATION, Sec.~\ref{EQUATION}, and INTEGRAL-QUATERNION, Sec.~\ref{INTEGRAL-QUATERNION}.

\begin{enumerate}

\bibitem{JOLY-1902B} C.J. Joly, \emph{Quaternion arrays}, Trans. Roy. Irish Acad. {\bf A 32} (1902-1904) 17--30. %%QUATERNION, %%MATRIX, %%LINEAR-FUNCTION, %$D14032002.
  
\bibitem{WOLF-1936-} L.A. Wolf, \emph{Similarity of matrices in which the elements are real quaternions}, Bull. Amer. Math. Soc. {\bf 42} (1936) 737--743. %%QUATERNION, %%MATRIX, %$D24062002.

\bibitem{LEE--1949-} H.C. Lee, \emph{Eigenvalues and canonical forms of matrices with quaternion coefficients}, Proc. Roy. Irish Acad. {\bf A 52} (1949) 253--260. %%QUATERNION, %%MATRIX, %%EIGENVALUES, %$D14032002.

\bibitem{BRENN1951-} J.L. Brenner, \emph{Matrices of quaternions}, Pacific J. Math. {\bf 1} (1951) 329--335. %%QUATERNION, %%MATRIX, %%EIGENVALUES, %$D14032002.

\bibitem{WIEGM1955-} N.A. Wiegmann, \emph{Some theorems on matrices with real quaternion elements}, Can. J. Math. {\bf 7} (1955) 191--201. %%QUATERNION, %%MATRIX, %$D14032002.

\bibitem{ILAME1965-} Y. Ilamed, \emph{Hamilton-Cayley theorem for matrices with non-commutative elements}, in: W.E. Brittin, A.O. Barut, eds., Lect. in Th. Phys. {\bf 7A}, {Lorentz Group} (University of Colorado, Boulder, 1965) 295--296.   %%QUATERNION, %%MATRIX, %$D14032002.

\bibitem{DEPIL1971-} J. DePillis and J. Brenner, \emph{Generalized elementary symmetric functions and quaternion matrices}, Linear Algebra Appl. {\bf 4} (1971) 55--69. %%QUATERNION, %%MATRIX, %$D24062002.

\bibitem{WOOD-1985-} R.M.W. Wood, \emph{Quaternionic eigenvalues}, Bull. London Math. Soc. {\bf 17} (1985) 137--138. %%QUATERNION, %%MATRIX, %$D30112005.

\bibitem{BUNSE1989-} A. Bunse-Gerstner, R. Byers, and V. Mehrmann, \emph{A quaternion QR algorithm}, Numerische Mathematik {\bf 55} (1989) 83--95. %%QUATERNION, %%MATRIX, %$D30112005.

\bibitem{CHATT1992-} A.W. Chatters, \emph{Matrices, idealizers, and integer quaternions}, J. Algebra {\bf 150} (1992) 45--56. %%QUATERNION, %%MATRIX, %%INTEGRAL-QUATERNION, %$D30112005.

\bibitem{MACKE1995A} N. Mackey, \emph{Hamilton and Jacobi meet again: quaternions and the eigenvalue problem}, SIAM J. Matrix Anal. Appl. {\bf 16} (1995) 421--435. %%QUATERNION, %%MATRIX, %%EIGENVALUES, %$D14032002.

\bibitem{MACKE1995B} N. Mackey, \emph{Hamilton and Jacobi meet again: quaternions and the eigenvalue problem}, Proc. Roy. Irish Acad. {\bf 95A}. Suppl. (1995) 59--66. %%QUATERNION, %%MATRIX, %%EIGENVALUES, %$D03062002.

\bibitem{CONNE1997-} A. Connes and A. Schwarz, \emph{Matrix Vieta theorem revisited}, Lett. Math. Phys. {\bf 39} (1997) 349--353. %%QUATERNION, %%MATRIX, %$D30112005.

\bibitem{THOMP1997-} R.C. Thompson, \emph{The upper numerical range of a quaternionic matrix is not a complex numerical range}, Lin. Algebra Appl. {\bf 254} (1997) 19--28. %%QUATERNION, %%MATRIX, %$D14032002.

\bibitem{ZANGH1997-} F. Zhang, \emph{Quaternions and matrices of quaternions}, Lin. Algebra Appl. {\bf 251} (1997) 21--57. %%QUATERNION, %%MATRICES, %$D30112005.

\bibitem{BAKER1999-} A. Baker, \emph{Right eigenvalues for quaternionic matrices}, Lin. Alg. Appl. {\bf 286} (1999) 303--309.  %%QUATERNION, %%MATRIX, %$D30112005.

\bibitem{JIANG1999-} T.S. Jiang and C. Li, \emph{Generalized diagonalization of matrices over quaternion field}, Appl. Math. Mech. (Engl.) {\bf 20} (1999) 1297--1304.  %%QUATERNION, %%MATRIX, %$D30112005.

\bibitem{JOHNS1999-} N.W. Johnson and A.I. Weiss, \emph{Quaternionic modular groups}, Lin. Alg. Appl. {\bf 295} (1999) 159--189.  %%QUATERNION, %%MATRIX, %%INTEGRAL-QUATERNION, %%GROUP-THEORY, %$D30112005.

\bibitem{MERIN1999-} D.I. Merino and V.V. Sergeichuk, \emph{Littlewood's algorithm and quaternion matrices}, Linear Algebra Appl. {\bf 298} (1999) 193--208; e-print \underline{ arXiv:0709.2466 }.  %%QUATERNION, %%MATRIX, %$D06072008.

\bibitem{JIANG2004A} T.S. Jiang, \emph{An algorithm for eigenvalues and eigenvectors of quaternion matrices in quaternionic quantum mechanics}, J. Math. Phys. {\bf 45} (2004) 3334--3338. %%QUATERNION, %%MATRIX, %%QQPH, %$D04032006.

\bibitem{JANOV2005-} D. Janovska and G. Opfer, \emph{Fast Givens transformation for quaternion valued matrices applied to Hessenberg reductions}, Electronic Trans. Numer. Anal. {\bf 20} (2005) 1--26.   %%QUATERNION, %%MATRIX, %$D06072008.

\bibitem{OPFER2005-} G. Opfer, \emph{The conjugate gradient algorithm applied to quaternion-valued matrices}, Z. Angew. Math. Mech. {\bf 85} (2005) 660--672. %%QUATERNION, %%MATRIX, %$D06072008.

\bibitem{TIAN-2005-} Y. Tian and G.P.H. Styan, \emph{Some inequalities for sums of nonnegative definite matrices in quaternions}, J. Inequalities and Appl. {\bf 5} (2005) 449--458. %%QUATERNION, %%MATRIX, %$D06072008.

\bibitem{DJOKO2007-} D.Z. Djokovic and B.H. Smith, \emph{Quaternionic matrices: Unitary similarity, simultaneous triangularization and some trace identities} (2007) 26~pp.; e-print \underline{ arXiv:0709.0513 }. %%QUATERNION, %%MATRIX, %$D06072008.


\end{enumerate}



\subsection{DETERMINANT}
%-----------------------
\label{DETERMINANT}


See also MATRIX, Sec.~\ref{MATRIX}.


\begin{enumerate}

\bibitem{STUDY1920-} E. Study, \emph{Zur Theorie der linearen Gleichungen}, Acta Math. {\bf 42} (1920) 1--61. %%QUATERNION, %%DETERMINANT, %%EQUATION, %$D14032002.

\bibitem{MOORE1922-} E.H. Moore, \emph{On the determinant of an hermitian matrix of quaternionic elements}, Bull. Am. Math. Soc. {\bf 28} (1922) 161--162. %%QUATERNION, %%DETERMINANT, %$D14032002.

\bibitem{DIEUD1943-} J. Dieudonn\'e, \emph{Les d\'eterminants sur un corps non commutatif}, Bull. Soc. Math. France {\bf 71} (1943) 27--45. %%QUATERNION, %%DETERMINANT, %$D14032002.

\bibitem{DYSON1972A} F. Dyson, \emph{Quaternion determinants}, Helv. Phys. Acta {\bf 45} (1972) 289--302. %%QUATERNION, %%DETERMINANT, %$D14032002, %$C On page 301 Dyson gives a short appreciation of Hamilton.

\bibitem{VANPR1989-} P.V. vanPraag, \emph{Sur les d\'eterminants des matrices quaternioniennes}, Helv. Phys. Acta {\bf 62} (1989) 42--46. %%QUATERNION, %%DETERMINANT, %$D14032002.

\bibitem{GELFA1991-} I. Gelfand and V. Retakh, \emph{Determinants of matrices over non-commutative rings}, Funct. Annal. Appl. {\bf 25} (1991) 91--102. %%QUATERNION, %%DETERMINANT, %$D24062002.

\bibitem{GELFA1992-} I. Gelfand and V. Retakh, \emph{A theory of non-commutative determinants and characteristic functions of graphs}, Funct. Annal. Appl. {\bf 26} (1992) 231--246. %%QUATERNION, %%DETERMINANT, %$D24062002.

\bibitem{ASLAK1996-} H. Aslaken, \emph{Quaternionic determinants}, The Mathematical Intelligencer {\bf 18} (1996) 57--65. %%QUATERNION, %%DETERMINANT, %$D14032002.

\bibitem{CLEVE1998-} J. Cleven, \emph{Norms and determinants of quaternionic line bundles}, Arch. Math. {\bf 71} (1998) 17--21.  %%QUATERNION, %%DETERMINANT, %$D30112005.

\bibitem{NAGAO1999-} T. Nagao and P.J. Forrester, \emph{Quaternion determinant expressions for multilevel dynamical correlation functions of parametric random matrices}, Nucl. Phys. {\bf B 563} (1999) 547--572. %%QUATERNION, %%DETERMINANT, %$D14032002.

\bibitem{COHEN2000-} N. Cohen and S. DeLeo, \emph{The quaternionic determinant}, The Electronic J. of Lin. Algebra {\bf 7} (2000) 100--111. %%QUATERNION, %%DETERMINANT, %$D04032006.

\bibitem{GARIB2004-} R.S. Garibaldi, \emph{The characteristic polynomial and determinant are not ad hoc constructions}, Amer. Math. Monthly {\bf  
111} (2004) 761--778. %%QUATERNION, %%DETERMINANT, %$D24062002.

\end{enumerate}




\subsection{GROUP-THEORY}
%------------------------
\label{GROUP-THEORY}

See also ALGEBRA, Sec.~\ref{ALGEBRA}.


\begin{enumerate}

\bibitem{STRIN1881-} W.I. Stringham, \emph{Determination of finite quaternion groups}, Am. J. Math. {\bf 4} (1881) 345--357. %%QUATERNION, %%GROUP-THEORY, %%INTEGRAL-QUATERNION, %$D30112005, %$C Cited by Girard, but not by Du Val.

\bibitem{WEYL-1931-} H. Weyl, The Theory of Groups and Quantum Mechanics, Transl. by H.P. Robertson (Methuen, London, 1931; Dover Publs., NY, 1950) 422 pp. %%QUATERNION, %%GROUP-THEORY, %$D06072008.  %$C Quaternions on pages 138 and 182.

\bibitem{CARTA1940-} E. Cartan, \emph{Sur les groupes lin\'eaires quaternioniens}, Beiblatt zur Vierteljahrsschrift des Naturforschenden Gesellschaft in Zurich {\bf 32} (1940) 191--203. %%QUATERNION, %%GROUP-THEORY, %$D14032002.

\bibitem{MILLE1947-} G.A. Miller, \emph{Abstract group generated by the quaternion units}, Proc. Nat. Acad. Sci. {\bf 33} (1947) 236--237.   %%QUATERNION, %%GROUP-THEORY, %$D06072008.

\bibitem{DYSON1962B} F. Dyson, \emph{The threefold way. Algebraic structure of symmetric groups and ensembles in quantum mechanics}, J. Math. Phys. {\bf 3} (1962) 1199--1215. %%QUATERNION, %%GROUP-THEORY, %%QUANTUM-PHYSICS, %$D14032002.

\bibitem{FINKE1963A} D. Finkelstein, J.M. Jauch, S. Schiminovich, and D. Speiser, \emph{Quaternionic representations of compact groups}, J. Math. Phys. {\bf 4} (1963) 136--140. %%QUATERNION, %%GROUP-THEORY, %%QQPH, %$D14032002.

\bibitem{SUDBE1984-} A. Sudbery, \emph{Division algebras, (pseudo)orthogonal groups and spinors}, J. Phys. {\bf A17} (1984) 939--955. %%QUATERNION, %%GROUP-THEORY, %$D14032002.

\bibitem{ALTMA1986-} S.L. Altmann, Rotations, Quaternions, and Double Groups (Clarendon, Oxford, 1986) 303~pp. %%BOOK, %%QUATERNION, %%GROUP-THEORY, %$D20032002.

\bibitem{DUVAL1964-} P. DuVal, \emph{Homographies, quaternions and rotations} (Clarendon, Oxford, 1964) 116~pp.  %%BOOK, %%QUATERNION, %%GROUP-THEORY, %$D20032002. %$C Finite groupes of rotations.

\bibitem{GILMO1974-} R. Gilmore,  Lie Groups, Lie Algebras, and Some of Their Applications (Wiley, New York, 1974) 587~pp. %%BOOK, %%QUATERNION, %%GROUP-THEORY, %$D20032002, %$C This is the first modern text book in which real numbers, complex numbers, and quaternions are used in a systematical and comprehensive way to study and build Lie groups and Lie algebras.

\bibitem{MIRKA1978-} M.R. Mir-Kasimov and I.P. Volobujev, \emph{Complex quaternions and spinor representations of de Sitter groups SO(4,1) and SO(3,2)}, Acta Phys. Polonica {\bf B9} (1978) 91--105. %%QUATERNION, %%GROUP-THEORY, %$D14032002.

\bibitem{ILAME1981-} Y. Ilamed and N. Salingaros, \emph{Algebras with three anticommuting elements. I. Spinors and quaternions}, J. Math. Phys. {\bf 22} (1981) 2091--2095.  %%QUATERNION, %%GROUP-THEORY, %$D14032002.

\bibitem{KIM--1981-} S.K. Kim, \emph{A unified theory of point groups and their general irreducible representations}, J. Math. Phys. {\bf 22} (1981) 2101--2107.  %%QUATERNION, %%GROUP-THEORY, %$D30112005.

\bibitem{SALIN1981A} N. Salingaros, \emph{Realization, extension, and classification of certain physically important groups and algebras}, J. Math. Phys. {\bf 22} (1981) 226--232.  %%QUATERNION, %%GROUP-THEORY, %$D14032002.

\bibitem{SALIN1981B} N. Salingaros, \emph{Algebras with three anticommuting elements. II. Two algebras over a singular field}, J. Math. Phys. {\bf 22} (1981) 2096--2100.  %%QUATERNION, %%GROUP-THEORY, %$D14032002.

\bibitem{KOCA1989D} M. Koca, \emph{Quaternionic and octonionic orbifolds}, Phys. Lett. {\bf B 104} (1989) 163--176.  %%QUATERNION, %%GROUP-THEORY, %%OCTONION, %$D30112005.

\bibitem{NOWIC1988-} A. Nowicki, \emph{Quaternionic strange superalgebras and the description of nonrelativistic spin}, Mod. Phys. Lett. A {\bf 3} (1988) 179--185. %%QUATERNION, %%GROUP-THEORY, %%PARTICLE-PHYSICS, %%QUANTUM-PHYSICS, %$D24062002.

\bibitem{WARE1990-} R. Ware, \emph{A note on the quaternion group as Galois group}, Proc. Am. Math. Soc. {\bf 108} (1990) 621--625. %%QUATERNION, %%GROUP-THEORY, %$D04032006.

\bibitem{ALTHO1992-} S.C. Althoen, K.D. Hansen, and L.D. Kugler, \emph{Rotational scaled quaternion division algebras}, J. of Alg. {\bf 146} (1992) 124--143.  %%QUATERNION, %%GROUP-THEORY, %%ALGEBRA, %$D14032002.

\bibitem{FOOT-1992-} R. Foot and G.C. Joshi, \emph{An application of the division algebras, Jordan algebras and split composition algebras}, Int. J. Mod. Phys. {\bf A 7} (1992) 4395--4413.  %%QUATERNION, %%GROUP-THEORY, %$D14032002.

\bibitem{DORAN1993A} C. Doran, D. Hestenes, F. Sommen, and N. vanAcker, \emph{Lie Groups as spin groups}, J. Math. Phys. {\bf 34}  (1993) 3642--3669.  %%QUATERNION, %%GROUP-THEORY, %%SPINOR, %$D10022002.

\bibitem{DING1994-} H.M. Ding, \emph{Bessel-functions on quaternionic Siegel domains} J. Funct. Anal. {\bf 120} (1994) 1--47.  %%QUATERNION, %%GROUP-THEORY, %%QUANTUM-PHYSICS, %$D30112005.

\bibitem{DELEO1995A} S. DeLeo and P. Rotelli, \emph{Representations of U(1,q) and constructive quaternion tensor products}, Nuovo Cim. {\bf B 110} (1995) 33--51. %%QUATERNION, %%GROUP-THEORY, %$D15032002.

\bibitem{SCOLA1995-} G. Scolarici and L. Solombrino, \emph{Notes on quaternionic group representations}, Int. J. Theor. Phys. {\bf 34} (1995) 2491--2500. %%QUATERNION, %%GROUP-THEORY, %%QQPH, %$D04022002.

\bibitem{BARBE1997-} M.L. Barberis, \emph{Hypercomplex structures on four-dimensional Lie groups}, Proc. Am. Math. Soc. {\bf 125} (1997) 1043--1054. %%QUATERNION, %%GROUP-THEORY, %$D04032006. 

\bibitem{SCOLA1997-} G. Scolarici and L. Solombrino, \emph{Quaternionic representations of magnetic groups}, J. Math. Phys. {\bf 38} (1997) 1147--1160. %%QUATERNION, %%GROUP-THEORY, %%QQPH, %$D04022002.

\bibitem{NEBE-1998-} G. Nebe, \emph{Finite quaternionic matrix groups}, Represent. Theory {\bf 2} (1998) 106--223.  %%QUATERNION, %%GROUP-THEORY, %$D04032006.

\bibitem{DELEO1999A} S. DeLeo and G. Ducati, \emph{Quaternionic groups in physics}, Int. J. Th. Phys. {\bf 38} (1999) 2197--2220.  %%QUATERNION, %%GROUP-THEORY, %$D14032002.

%%\bibitem{SCOLA1999-} G. Scolarici and L. Solombrino, \emph{Quaternionic group representations and their classifications}, Proc. of Second Meeting on Quaternionic Structures in Mathematics and Physics (6--10 September 1999) 365--375. %%QUATERNION, %%GROUP-THEORY, %%QQPH, %$D24062002.

\bibitem{LOKE-2000-} H.Y. Loke, \emph{Restrictions of quaternionic representations}, J. Funct. Anal. {\bf 172} (2000) 377-403.  %%QUATERNION, %%GROUP-THEORY, %$D30112005.

\bibitem{SCOLA2000C} G. Scolarici and L. Solombrino, \emph{Central projective quaternionic representations}, J. Math. Phys. {\bf 41} (2000) 4950--4603. %%QUATERNION, %%GROUP-THEORY, %%QQPH, %$D04022002.

\bibitem{PUTA-2001-} M. Puta, \emph{Optimal control problems on the Lie group $SP(1)$}, in: S. Marchiafava et al., eds., Proceedings of the 2nd Meeting on Quaternionic Structures in Mathematics and Physics (World Scientific, Singapore, 2001) 339--347. %%QUATERNION, %%GROUP-THEORY, %$D04032006.

\bibitem{SCOLA2001-} G. Scolarici and L. Solombrino, \emph{Quaternionic group representations and their classifications}, in: S. Marchiafava et al., eds., Proceedings of the 2nd Meeting on Quaternionic Structures in Mathematics and Physics (World Scientific, Singapore, 2001) 365--375. %%QUATERNION, %%GROUP-THEORY, %%QQPH, %$D04032006.

\bibitem{HESTE2002A} D. Hestenes, \emph{Point groups and space groups in geometric algebra}, in: L. Dorst et al., eds, Applications of Geometric Algebra in Computer Science and Engineering (Birkh\"auser, Boston, 2002) 3--34. %%QUATERNION, %%GROUP-THEORY, %%HESTENES, %$D10022002.

\bibitem{JURIA2006-} S.O. Juriaans and A.C. Souza-Filho, \emph{Hyperbolicity of orders of quaternions algebras and group rings} (2006) 6~pp.; e-print \underline{ arXiv:math/0610699 }. %%QUATERNION, %%GROUP-THEORY, %$D06072008.

\bibitem{PAVSI2008-} M. Pavsic, \emph{A novel view on the physical origin of $E_8$} (26 Jun 2008) 14~pp.; e-print \underline{ arXiv:0806.4365 }.  %%QUATERNION, %%GROUP-THEORY, %PARTICLE-PHYSICS, %$D06072008. 

\end{enumerate}




\subsection{ANALYSIS}
%--------------------
\label{ANALYSIS}

Analysis, differentiation, integration, differential equations, differential forms, etc. 

\begin{enumerate}

\bibitem{TAIT-1870A} P.G. Tait, \emph{Note on linear partial differential equations}, Proc. Roy. Soc. Edinburgh {\bf } (6 June 1870) SP-1:151--152. %%QUATERNION, %%ANALYSIS, %$D24062002.

\bibitem{TAIT-1870B} P.G. Tait, \emph{Note on linear partial differential equations in quaternions}, Proc. Roy. Soc. Edinburgh {\bf } (20 December 1870) SP-1:153--158.  %%QUATERNION, %%ANALYSIS, %$D24062002.

\bibitem{TAIT-1870C} P.G. Tait, \emph{On some quaternion integrals}, Proc. Roy. Soc. Edinburgh {\bf } (10 December 1870) SP-1:159--163.  %%QUATERNION, %%ANALYSIS, %$D24062002.

\bibitem{BRILL1887-} J. Brill, \emph{A new geometrical interpretation of the quaternion analysis}, Proc. Cambridge Phil. Soc. {\bf 6} (1887) 156--169. %%QUATERNION, %%ANALYSIS, %$D04032006.

\bibitem{BRILL1890-} J. Brill, \emph{Note on the application of quaternions to the discussion of Laplace's equation}, Proc. Cambridge Phil. Soc. {\bf 7} (1890) 120--125. %%QUATERNION, %%ANALYSIS, %$D04032006.

\bibitem{CARVA1890-} E. Carvallo, \emph{Formules de quaternions pour la r\'eduction des int\'egrales multiples les unes des autres}, Bull. Soc. Math. France {\bf 18} (1890) 80--90. %%QUATERNION, %%ANALYSIS, %$D04032006.

\bibitem{BRILL1891A} J. Brill, \emph{Note on the application of quaternions to the discussion of Laplace's equation}, Camb. Phil. Soc. {\bf 7} (1891) 120--125.  %%QUATERNION, %%ANALYSIS, %$D24062002.

\bibitem{BRILL1891B} J. Brill, \emph{On quaternion functions, with especial reference to the discussion of Laplace's equation}, Camb. Phil. Soc. {\bf 7} (1891) 151--156.  %%QUATERNION, %%ANALYSIS, %$D24062002.

\bibitem{BRILL1896-} J. Brill, \emph{On the generalization of certain properties of the tetrahedron}, Proc. Cambridge Phil. Soc. {\bf 9} (1896) 98--108. %%QUATERNION, %%ANALYSIS, %$D04032006, %$C The method bears a close analogy to that used in Ref. \cite{BRILL1887-}.

\bibitem{LAROS1896-} M. Larose, \emph{D\'emonstration du th\'eor\`eme de M. Vaschy sur une distribution quelconque de vecteur}, Bull. Soc. Math. France {\bf 24} (1896) 177--180. %%QUATERNION, %%ANALYSIS, %$D04032006, %$C Refers to \cite{CARVA1890-}.

\bibitem{JOLY-1899B} C.J. Joly, \emph{Some applications of Hamilton's operator $\nabla$ in the calculus of variations}, Proc. Roy. Irish Acad. {\bf 5} (1899) 666. %%QUATERNION, %%ANALYSIS, %$D130302002.


\bibitem{MACFA1900-} A. Macfarlane, \emph{Differentiation in the quaternion analysis}, Proc. Roy. Irish Acad. {\bf 6} (1900) 199--215.  %%QUATERNION, %%ANALYSIS, %$D17052001.

\bibitem{JOLY-1902C} C.J. Joly, \emph{Integrals depending on a single quaternion variable}, Proc. Roy. Irish Acad. {\bf A 8} (1902) 6--20. %%QUATERNION, %%ANALYSIS, %$D17052001.

\bibitem{WANG-1911-} K.T. Wang, \emph{The differentiation of quaternion functions}, Proc. Roy. Irish Acad. {\bf 29 A} (1911) 73--80. %%QUATERNION, %%ANALYSIS, %$D17052001.

\bibitem{FULLE1954-} F.B. Fuller, \emph{Harmonic mappings}, Proc. Nat. Acad. Sci. {\bf 40} (1954) 987--991.   %%QUATERNION, %%ANALYSIS, %$D06072008.

\bibitem{HESTE1968A} D. Hestenes, \emph{Multivector calculus}, Journal of Mathematical Analysis and Applications {\bf 24} (1968) 313--325. %%QUATERNION, %%ANALYSIS, %%HESTENES, %$D10022002.

\bibitem{HESTE1968B} D. Hestenes, \emph{Multivector functions}, J. Math. Anal. Appl. {\bf 24} (1968) 467--473. %%QUATERNION, %%ANALYSIS, %%CLIFFORD-R, %%HESTENES, %$D09022002.

\bibitem{SALIN1979-} N. Salingaros and M. Dresden, \emph{Properties of an associative algebra of tensor fields. Duality and Dirac identities}, Phys. Rev. Lett. {\bf 43} (1979) 1--4. %%QUATERNION, %%ANALYSIS, %%DIFFERENTIAL-FORM, %%DIRAC, %%MAXWELL, %$D06022002.

\bibitem{PICCI1981-} P. Piccini, \emph{Quaternionic differential forms and symplectic Pontrjagin  classes}, Ann. Mat. Pur. Appl. {\bf 129} (1981) 57--68. %%QUATERNION, %%ANALYSIS, %%DIFFERENTIAL-FORMS, %$D30112005.

\bibitem{STOYA1986-} D.T. Stoyanov, \emph{On the classical solutions of the Liouville equation in a four-dimensional space}, Lett. Math. Phys. Lett. {\bf 12} (1986) 93--96. %%QUATERNION, %%ANALYSIS, %%H, %%LIOUVILLE-EQUATION, %$D17052001.

\bibitem{BREDI1987-} A. Bredimas, \emph{N-dimensional general solutions of the Liouville type equation $A u = e^u$ with applications}, Phys. Lett. A {\bf 121} (1986) 283-286. %%QUATERNION, %%ANALYSIS, %%H, %%LIOUVILLE-EQUATION, %$D17052001.

\bibitem{DIMAK1991-} A. Dimakis and F. M\"uller-Hoissen, \emph{Clifform calculus with applications to classical field theories}, Class. Quantum Grav. {\bf 8} (1991) 2093--2132. %%QUATERNION, %%ANALYSIS, %$D13032002.

\bibitem{LUGOJ1991-} S. Lugojan, \emph{Quaternionic derivability}, Anal. Univ. Timisoara {\bf 29} (1991) 175--190.  %%QUATERNION, %%ANALYSIS, %$D01062001.  %$C Continuation of this paper is \cite{LUGOJ1994-}.

\bibitem{CHUNG1992-} W.S. Chung, J.J. Lee, and J.H. Cho, \emph{Quaternion solutions of four-dimensional Liouville and Sine-Gordon equations}, Mod. Phys. Lett. {\bf 7} (1992) 2527--2533. %%QUATERNION, %%ANALYSIS, %%H, %%LIOUVILLE-EQUATION, %%SINE-GORDON-EQUATION, %%QQPH, %$D17052001.

\bibitem{SOBCZ1992-} G. Sobczyk, \emph{Simplicial calculus with geometric algebra}, in: A. Micali et al., eds., Clifford Algebras and their Applications in Mathematical Physics (Kluwer Academic Publishers, Dordrecht, 1992) 279--292.  %%QUATERNION, %%ANALYSIS, %$D10022002.

\bibitem{VONIN1992-} M. vonIns, \emph{An approach to quaternionic elliptic integrals},  Ph.D. thesis (Universit\"at Bern, 1992) 43~pp. %%BOOK, %%QUATERNION, %%ANALYSIS, %%QUADRIPERIODIC-FUNCTIONS, %$D14032002.

\bibitem{HESTE1993A} D. Hestenes, \emph{Differential forms in geometric calculus}, in:  F. Brackx et al, ed., Clifford Algebras and their Applications in Mathematical Physics (Kluwer Academic Publishers, Dordrecht, 1993) 269-285. %%QUATERNION, %%ANALYSIS, %%DIFFERENTIAL-FORM, %%HESTENES, %$D10022002.

\bibitem{LUGOJ1994-} S. Lugojan, \emph{Quaternionic derivability II}, in: G. Gentili et al., Proc. of the Meeting on Quaternionic Structures in Mathematics and Physics (SISSA, Trieste, 1994) 189--196. %%QUATERNION, %%ANALYSIS, %$D04032006.

\bibitem{KAMBE1997-} G.I. Kamberov, \emph{Quadratic differentials, quaternionic forms, and surfaces} (1997) 11~pp.; e-print \underline{ arXiv:dg-ga/9712011 }.  %%QUATERNION, %%ANALYSIS, %$D06072008.

\bibitem{VAZ--1997B}  J. Vaz, Jr., and W.A. Rodrigues, Jr., \emph{On the equation $\vec{\nabla} \times \vec{a} \kappa \vec{a}$}, Bol. Soc. Paran. Mat. {\bf 17} (1997) 19--24. %%QUATERNION, %%ANALYSIS, %%MAXWELL, %$D17032002.

\bibitem{HEMPF2000-} T. Hempfling, \emph{On the radial part of the Cauchy-Riemann operator}, in: J. Ryan and W. Spr\"ossig, eds., Clifford Algebra and their Applications in Mathematical Physics, Vol.~2: \emph{Clifford Analysis} (Birkh\"auser, Boston, 2000) 261--273. %%QUATERNION, %%ANALYSIS, %%CLIFFORD-R, %$D12052001.

\bibitem{DELEO2001B} S. DeLeo and G. Ducati, \emph{Quaternionic differential operators}, J. Math. Phys. {\bf 42} (2001) 2236--2265.   %%QUATERNION, %%ANALYSIS, %%EQUATION, %%QQPH, %$D14032002.

\bibitem{KRAVC2001C} V. Kravchenko, V.V. Kravchenko, and B. Williams, \emph{A  quaternionic generalization of the Riccati differential equation},  in: Clifford Analysis and its Applications, NATO Sci. Ser. II, Math. Phys. Chem. {\bf 25} (Kluwer, Dordrecht, 2001) 143--154. %%QUATERNION, %%ANALYSIS, %$D130302002. 

\bibitem{BAS--2003-} P. Bas, N. Le Bihan and J. Chassery, \emph{Utilisation de la transform\'ee de fourier quaternionique en tatouage d'images couleur}, in: Actes du 19�me Colloque sur le traitement du signal et des images (GRETSI'03, Paris, 2003) 191--194. %%QUATERNION, %%ANALYSIS, %$D06072008.

\bibitem{DELEO2003-} S. De Leo and G. Ducati, \emph{Solving simple quaternionic differential equations}, J. Math. Phys. {\bf 44} (2003) 2224--2233. %%QUATERNION, %%ANALYSIS, %$D06072008.

\bibitem{KRAVC2003B} V.G. Kravchenko and V.V. Kravchenko, \emph{Quaternionic factorization of the Schroedinger operator and its applications to some first order systems in mathematical physics}, J. Phys. A: Math. Gen. {\bf 36} (2003) 1285--1287. %%QUATERNION, %%ANALYSIS, %%MAXWELL, %%DIRAC, %$D30112005.

\bibitem{KRAVC2003C} V.V. Kravchenko, \emph{On Beltrami fields with nonconstant proportionality factor}, J. Phys. A: Math. Gen. {\bf 36} (2003) 1515--1522. %%QUATERNION, %%ANALYSIS, %$D30112005.

\bibitem{KUNZE2004-} K. Kunze and H. Schaeben, \emph{The Bingham distribution of quaternions and its spherical radon transform in texture analysis}, Math. Geology {\bf 36} (2004) 917--943.  %%QUATERNION, %%ANALYSIS, %$D06072008.

\bibitem{SAID-2006-} S. Said, N. Le Bihan and S.J. Sangwine, \emph{Fast complexified quaternion Fourier transform} (2006) 6~pp.; e-print \underline{ arXiv:math/0603578 }.  %%QUATERNION, %%ANALYSIS, %$D06072008.

\bibitem{HITZE2007-} E.M.S. Hitzer, \emph{Quaternion Fourier transform on quaternion fields and generalizations}, Adv. Appl. Cliff. Alg. {\bf 17} (2007) 497--517.  %%QUATERNION, %%ANALYSIS, %$D06072008.

\bibitem{LAKEW2008-} D.A. Lakew, \emph{Mollifiers in Clifford analysis} (21 April 2008) 13~pp.; e-print \underline{ arXiv:0802.1539 }. %%QUATERNION, %%ANALYSIS, %%CLIFFORD, %$D06072008.  


\end{enumerate}


\subsection{ANALYTICITY-VARIA}
%-----------------------------
\label{ANALYTICITY-VARIA}

Various extensions of complex analyticity to spaces of higher dimensions.

\begin{enumerate}

\bibitem{CARVA1896-} E. Carvallo, \emph{G\'en\'eralisation et extension \`a l'espace du th\'eor\`eme des r\'esidus de Cauchy}, Bull. Soc. Math. France {\bf 24} (1896) 180--184. %%QUATERNION, %%ANALYTICITY-VARIA, %$D04032006, %$C Refers to \cite{LAROS1896-}.

\bibitem{DIXON1904-}A.C. Dixon, \emph{On the Newtonian potential}, Quarterly J. of Math. {\bf 35} (1904) 283--296. %%QUATERNION, %%ANALYTICITY-VARIA, %%TRIPERIODIC-FUNCTIONS, %$D04032006. %$C Cited by \cite{RYAN-1994B}.

\bibitem{TAYLO1938-} A.E. Taylor, \emph{Biharmonic functions in abstract space}, Am. J. Math. {\bf 60} (1938) 416--422. %%QUATERNION, %%ANALYTICITY-VARIA, %$D01062001.

\bibitem{WARD-1940-} J.A. Ward, \emph{A theory of analytic functions in linear associative algebras}, Duke Math. J. {\bf 7} (1940) 233--248. %%QUATERNION, %%ANALYTICITY-VARIA, %$D01062001.

\bibitem{WEISS1946-} P. Weiss, \emph{An extension of Cauchy's integral formula by means of a Maxwell's stress tensor}, J. London Math. Soc. {\bf 21} (1946) 210--218. %%QUATERNION, %%ANALYTICITY-VARIA, %$D01062001.

\bibitem{TRAMP1960-} A. Trampus, \emph{Differentiability and analyticity of functions in linear algebras}, Duke Math. J. {\bf 27} (1960) 431--441. %%QUATERNION, %%ANALYTICITY-VARIA, %$D01062001.

\bibitem{STEIN1968-} E.M. Stein and G. Weiss, \emph{Generalization of the Cauchy-Riemann equations and representations of the rotation group}, Am. J. Math. {\bf 90} (1968) 163--196. %%QUATERNION, %%ANALYTICITY-VARIA, %$D01062001.

\bibitem{RYAN-1982A} J. Ryan, \emph{Topics in hypercomplex analysis}, Ph.D. thesis  (University of York, 1982) 241~pp.  %%BOOK, %%QUATERNION, %%ANALYTICITY-VARIA, %%H, %%B, %%CLIFFORD-R, %%CLIFFORD-C, %$D01062001.

\bibitem{IMAED1986A} M. Imaeda, \emph{On regular functions of a power-associative hypercomplex variable}, in: J.S.R. Chisholm, A.K. Common, eds.,  \emph{Clifford algebras and their applications in mathematical physics} (Reidel, Dordrecht, 1986) 565--572. %%QUATERNION, %%ANALYTICITY-VARIA, %$D01062001.

\bibitem{LEUTW1995-} H. Leutwiler, \emph{More on modified quaternionic analysis in R(3)}, Forum Math. {\bf 7} (1995) 279--305. %%QUATERNION, %%ANALYTICITY-VARIA, %$D30112005, %$C "reduced quaternions".

\bibitem{ERIKS1998-} S.L. Eriksson-Bique and H. Leutwiler, \emph{On modified quaternionic analysis in $R_3$}, Arch. Math. {\bf 70} (1998) 228--234.   %%QUATERNION, %%ANALYTICITY-VARIA, %$D30112005, %$C "reduced quaternions". 

\bibitem{SOMME1999B} F. Sommen, \emph{An extension of Clifford  Analysis  towards  super-symmetry}, in J. Ryan, W. Spr\"ossig. eds., Clifford Algebras and their Applications in Mathemetical Physics, Vol. 2 (Birkh\"auser, Boston, 1999) 199--224.  %%QUATERNION, %%ANALYTICITY-VARIA, %$D04022002.

\bibitem{MALON2000-} H. Malonek, \emph{Hypercomplex derivability --- The characterization of monogenic functions in $R^{n+1}$ by their derivative}, in: J. Ryan and W. Spr\"ossig, eds., Clifford Algebra and their Applications in Mathematical Physics, Vol.~2: \emph{Clifford Analysis} (Birkh\"auser, Boston, 2000) 273--285. %%QUATERNION, %%ANALYTICITY-VARIA, %%DERIVATION, %%CLIFFORD-R, %$D12052001.

\bibitem{SOMME2001B} F. Sommen, \emph{Clifford  Analysis  on  super-space}, Adv. in Appl. Clifford Alg. {\bf 11 (S1)} (2002) 291--304. %%QUATERNION, %%ANALYTICITY-VARIA, %$D04022002.

\bibitem{SOMME2002-} F. Sommen, \emph{Analysis using abstract vector variables}, in: L. Dorst  
et al., eds., Applications of Geometric Algebra in Computer Science and Engineering (Birkh\"auser, Boston, 2002) 119--128. %%QUATERNION, %%ANALYTICITY-VARIA, %$D04022002.

\bibitem{KRAVC2003D} V.V. Kravchenko, Applied Quaternionic Analysis, Research and Exposition in Mathematics Series {\bf 28} (Heldermann-Verlag, Lemgo, 2003) 127~pp. %%BOOK, %%QUATERNION, %%ANALYTICITY-VARIA, %$D30112005.

\bibitem{EASTW2007-} M. Eastwood and J. Ryan, \emph{Monogenic functions in conformal geometry}, SIGMA (Symmetry, Integrability and Geometry: Methods and Applications) {\bf 3} (2007) 084;; e-print \underline{ arXiv:0708.4172 }.  %%QUATERNION, %%ANALYTICITY-VARIA, %$D06072008.

\bibitem{FOKAS2007-} A.S. Fokas and D.A. Pinotsis, \emph{Quaternions, evaluation of integrals and boundary value problems}, Computational Methods and Function Theory {\bf 7} (2007) 443--476. %%QUATERNION, %%ANALYTICITY-VARIA, %$D06072008. 

\bibitem{LIBIN2007-} M. Libine, \emph{Hyperbolic Cauchy integral formula for the split complex numbers} (2007) 6~pp.; e-print \underline{ arXiv:0712.0375 }.  %%QUATERNION, %%ANALYTICITY-VARIA, %$D06072008.


\end{enumerate}




\subsection{ANALYTICITY-H}
%-------------------------
\label{ANALYTICITY-H}


Real quaternions analyticity only --- $\mathbb{H}$ alone

\begin{enumerate}

\bibitem{FUETE1928-} R. Fueter, \emph{Uber Funktionen einer Quaternionenvariablen}, Atti Congr. Int. Mat. Bologna (3--10 Settembre, 1928) Vol.~2, p.145. %%QUATERNION, %%ANALYTICITY-H, %$D Seminal paper on the subject, %$D24062002.

\bibitem{MOISI1931-} G.C. Moisil, \emph{Sur les quaternions monog\`enes}, Bull. Sci. Math. {\bf 55} (1931) 168--174. %%QUATERNION, %%ANALYTICITY-H, %$D29052001.

\bibitem{FUETE1932A} R. Fueter, \emph{Analytische Funktionen einer Quaternionenvariablen}, Comm. Math. Helv. {\bf 4} (1932) 9--20. %%QUATERNION, %%ANALYTICITY-H, %$D17052001.

\bibitem{FUETE1934-} R. Fueter, \emph{Die Funktionentheorie der Differentialgleichungen $\triangle u=0$ und $\triangle \triangle u=0$ mit vier reellen Variablen}, Comm. Math. Helv. {\bf 7} (1934/1935) 307--330. %%QUATERNION, %%ANALYTICITY-H, %$D17052001.

\bibitem{FUETE1935A} R. Fueter, \emph{Uber die analytische Darstellung der regul\"aren Funktionen einer Quaternionen variablen}, Comm. Math. Helv. {\bf 8} (1935/1936) 371--378. %%QUATERNION, %%ANALYTICITY-H, %$D17052001.

\bibitem{FUETE1936A} R. Fueter, \emph{Zur Theorie der regul\"aren Funktionen einer Quaternionvariablen}, Monatsch. f\"ur Math. und Phys. {\bf 43} (1936) 69--74.  %%QUATERNION, %%ANALYTICITY-H, %$D17052001.

\bibitem{FUETE1936B} R. Fueter, \emph{Die Singularit\"aten der eindeutigen regul\"aren Funktionen einer Quaternionvariablen. I.}, Comm. Math. Helv. {\bf 9} (1936/1937) 320--334.  %%QUATERNION, %%ANALYTICITY-H, %$D17052001.

\bibitem{FUETE1937A} R. Fueter, \emph{Integralsatze f\"ur regul\"are Funktionen einer Quaternionen-Variablen}, Comm. Math. Helv. {\bf 10} (1937/1938) 306--315. %%QUATERNION, %%ANALYTICITY-H, %$D17052001.

\bibitem{FUETE1937C} R. Fueter, \emph{Die Theorie der regul\"aren Funktionen einer Quaternionvariablen}, in: Compte Rendus Congr. Intern. des Math. Oslo 1936, Tome 1 (1937) 75--91. %%QUATERNION, %%ANALYTICITY-H, %$D17052001.

\bibitem{SCHUL1937-} B. Schuler, \emph{Zur Theorie der regul\"aren Funktionen einer Quaternionen-Variablen}, Comm. Math. Helv. {\bf 10} (1937/1938) 327-342. %%QUATERNION, %%ANALYTICITY-H, %$D17052001. 

\bibitem{FERRA1938-} V.C.A. Ferraro, \emph{On functions of quaternions}, Proc. Roy. Irish Acad. {\bf A 44} (1938) 101--108. %%QUATERNION, %%ANALYTICITY-H, %$D01062001. %$C Shows that q'=a+bq is analytic.

\bibitem{FUETE1939A} R. Fueter, \emph{Uber ein Hartogs'schen Satz}, Comm. Math. Helv. {\bf 12} (1939/1940) 75--80. %%QUATERNION, %%ANALYTICITY-H, %%FUETER, %%HARTOGS, %$D17052001.

\bibitem{FUETE1939B} R. Fueter, \emph{Uber vierfachperiodische Funktionen}, Monatsch. f\"ur Math. und Phys. {\bf 48} (1939) 161--169. %%QUATERNION, %%ANALYTICITY-H, %%QUADRIPERIODIC-FUNCTION, %$D17052001.

\bibitem{NEF--1943B} W. Nef, \emph{Die unwesentlichen Singularitaten der regul\"aren Funktionen einer Quaternionenvariabeln},  Comm. Math. Helv. {\bf 16} (1943/1944) 284--304. %%QUATERNION, %%ANALYTICITY-H, %%QUADRIPERIODIC-FUNCTION, %$D17052001.

\bibitem{FUETE1945B} R. Fueter, \emph{Uber die Quaternionenmultiplikation regul\"arer vierfachperiodischer Funktionen}, Experientia {\bf 1} (1945) 51.  %%QUATERNION, %%ANALYTICITY-H, %%QUADRIPERIODIC-FUNCTION, %$D24062002.

\bibitem{FUETE1948A} R. Fueter, \emph{Uber die Funktionentheorie in einer hyperkomplexen Algebra}, Element der Math. {\bf 3} (1948) 89--94.   %%QUATERNION, %%ANALYTICITY-H, %$D24062002.

\bibitem{FUETE1948B} R. Fueter and E. Bareiss, \emph{Funktionen Theorie im Hyperkomplexen} (Mathematisches Institut der Universit\"at, Zurich, 1948--1949) 318~pp. %%BOOK, %%QUATERNION, %%ANALYTICITY-H, %%B, %%CLIFFORD-R, %%CLIFFORD-C, %%DIRAC, %%HARTOGS, %$D13062001.

\bibitem{FUETE1949-} R. Fueter, \emph{Uber Abelsche Funktionen von zwei Komplexen Variablen}, Ann. Math. Pura. Appl. {\bf 4} (1949) 211--215.  %%QUATERNION, %%ANALYTICITY-H %%QUADRIPERIODIC-FUNCTION, %$D17052001.

\bibitem{RINEA1960-} R.F. Rinehart, \emph{Elements of a theory of intrinsic functions on algebras}, Duke Math. J. {\bf 27} (1960) 1--19.  %%QUATERNION, %%ANALYTICITY-H, %$D01062001.

\bibitem{HOSHI1962-} S. Hoshi, \emph{On some theories of quaternion functions}, Memoirs Fac. Engineering Miyazaki Univ. {\bf 3} (1962) 70~pp. %%BOOK, %%QUATERNION, %%ANALYTICITY-H, %$D01062001.

\bibitem{CULLE1965-} C.G. Cullen, \emph{An integral theorem for analytic intrinsic functions on quaternions}, Duke. Math. J. {\bf 32} (1965) 139--148. 
%%QUATERNION, %%ANALYTICITY-H, %$D01062001.

\bibitem{HABET1976-} K. Habetha, \emph{Eine Bemerkung zur Funktionentheorie in Algebren}, Lect. Notes in Math. {\bf 561} (Springer, Berlin, 1976) 502--509. %%QUATERNION, %%ANALYTICITY-H,  %$D17052001.

\bibitem{DEAVO1973-} C.A. Deavours, \emph{The quaternion calculus}, Amer. Math. Monthly {\bf 80} (1973) 995--1008.  %%QUATERNION, %%ANALYTICITY-H  %$D17052001.

\bibitem{SUDBE1979-} A. Sudbery, \emph{Quaternionic analysis}, Math. Proc. Cambridge Phil. Soc. {\bf 85} (1979) 199--225. %%QUATERNION, %%ANALYTICITY-H,  %$D17052001.

\bibitem{NONO-1982-} K. Nono, \emph{Hyperholomorphic functions of a quaternion variable}, Bull. of Fukuoka University of Education {\bf 32} (1982) 21--37.  %%QUATERNION, %%ANALYTICITY-H, %$D04032006.

\bibitem{SOUCE1984-} V. Soucek, \emph{$\mathbb{H}$-valued differential forms on $\mathbb{H}$}, Rendiconti Circ. Mat. Palermo--Suppl. {\bf 3} (1984) 293--294. %%QUATERNION, %%ANALYTICITY-H, %%DIFFERENTIAL-FORMS, %$D01062001.

\bibitem{GURLE1986-} K. G\"urlebeck, \emph{Hypercomplex factorization of the Helmholz equation}, Zeitschr. f\"ur Anal. und ihre Anwendung {\bf 5} (1986) 125--131. %%QUATERNION, %%ANALYTICITY-H, %$D29052001.

\bibitem{KOCHE1986-} R.R. Kocherlakota, \emph{Functions of a quaternion variable which are gradients of real-valued functions}, Aeq. Math. {\bf 31} (1986) 109--117. %%QUATERNION, %%ANALYTICITY-H, %$D30112005.

\bibitem{DATTA1987-} B. Datta and S. Nag, \emph{Zero-sets of quaternionic and octonionic analytic functions with central coefficients}, Bull. London Soc. Math. {\bf 19} (1987) 329--336.  %%QUATERNION, %%ANALYTICITY-H, %%O  %$D17052001.

\bibitem{NONO-1987-} K. Nono, \emph{Runge's theorems for complex valued harmonic and quaternion valued hyperholomorphic functions}, Rev. Roumaine Math. Pures Appl. {\bf 32} (1987) 155--158.  %%QUATERNION, %%ANALYTICITY-H, %$D04032006.

\bibitem{STRUP1998-} D.C. Struppa, \emph{Gr\"obner bases in partial differential equations}, London Math. Soc. Lect. Note Ser. {\bf 251} (1998) 235--245. %%QUATERNION, %%ANALYTICITY-H, %$D04022002.

\bibitem{GURLE1990-} K. G\"urlebeck and W. Spr\"ossig, Quaternionic Analysis and Elliptic Boundary Value Problems (Birkh\"auser, Basel, 1990) 253~pp.  %%BOOK, %%QUATERNION, %%ANALYTICITY-H, %$D20032002.

\bibitem{PERTI1990-} D. Pertici, \emph{Traces de fonctions r\'eguli\`eres de plusieurs variables quaternioniennes}, C.R. Acad. Sci. Paris. {\bf 311} (1990) 37--40.  %%QUATERNION, %%ANALYTICITY-H, %$D24062002.

\bibitem{PERTI1991-} D. Pertici, \emph{Trace theorems for regular functions of several quaternion variables}, Forum Math. {\bf 3} (1991) 461--478.  %%QUATERNION, %%ANALYTICITY-H, %$D17052001.

\bibitem{KROLI1992-} W. Krolikowski and E. Ramirez de Arellano, \emph{Fueter-H\"urwitz regular mappings and an integral representation}, in: A. Micali et al., eds., Clifford algebras and their Applications in Mathematical Physics (Kluwer, Dordrecht, 1992) 221--237.  %%QUATERNION, %%ANALYTICITY-H, HURWITZ-PAIR %$D17052001.

\bibitem{AUBER1993-} G. Auberson, \emph{Monogenic continuation for vector fields}, J. Math. Phys. {\bf 34} (1993) 3151--3161. %%QUATERNION, %%ANALYTICITY-H, %%WEISS, %$D01062001.

\bibitem{PERTI1993-} D. Pertici, \emph{Quaternion regular functions and domains of regularity}, Boll. Union Matem. Italiana {\bf B 7} (1993) 973--988. %%QUATERNION, %%ANALYTICITY-H,  %$D17052001.

\bibitem{KRAVC1994B} V.G. Kravchenko and V.V. Kravchenko, \emph{On some nonlinear equations generated by Fueter type operators}, Zeitschr. f\"ur Anal. und ihre Anwend. {\bf 4} (1994) 599--602.  %%QUATERNION, %%ANALYTICITY-H, %$D17052001.

\bibitem{PERNA1994-} L. Pernas, \emph{About some operators in quaternionic analysis}, in: G. Gentili et al., Proc. of the Meeting on Quaternionic Structures in Mathematics and Physics (SISSA, Trieste, 1994) 237--246.  %%QUATERNION, %%ANALYTICITY-H, %$D04032006.

\bibitem{BEDDI1995A} S. Bedding and K. Briggs, \emph{Iteration of quaternion maps}, Int. J. Bifurc. Chaos {\bf 3} (1995) 877--881. %%QUATERNION, %%ANALYTICITY-H, %%ITERATION, %$D24062002.

\bibitem{BEDDI1995B} S. Bedding and K. Briggs, \emph{Regularly iterable linear quaternion maps}, La Trobe University research report No. 95-2, Submitted to J. Austr. Math Soc. (31 May, 1995) 13~pp. %%QUATERNION, %%ANALYTICITY-H, %%ITERATION, %$D24062002.

\bibitem{ABDEL1996-} K. Abdel-Kahlek, \emph{Quaternion analysis} (18 November 1996) 8~pp.;  e-print \underline{ arXiv:hep-th/9607152 }.  %%QUATERNION, %%ANALYTICITY-H, %$D30112005.

\bibitem{ADAMS1996-} W.W. Adams, C.A. Berenstein, P. Loutaunau, I. Sabadini, and D.C. Struppa, \emph{On compact singularities for regular functions of one quaternionic variable}, Complex Variables {\bf 31} (1996) 259--270.  %%QUATERNION, %%ANALYTICITY-H, %%HARTOGS, %$D17052001.

\bibitem{BEDDI1996-} S. Bedding and K. Briggs, \emph{Iteration of quaternion functions}, Amer. Math. Monthly {\bf 103} (1996) 654--664.  %%QUATERNION, %%ANALYTICITY-H, %%ITERATION, %$D24062002, %$C SPIN 3/2 page 660.

\bibitem{BEREN1996-} C.A. Berenstein, I. Sabadini, and D.C. Struppa, \emph{Boundary values of regular functions of quaternionic variables},
Pitman Res. Notes Math. Ser. {\bf 347} (1996) 220--232. %%QUATERNION, %%ANALYTICITY-H, %$D04022002.

\bibitem{ADAMS1997-} W.W. Adams,  P. Loutaunau, V.P. Palamodov, and D.C. Struppa, \emph{Hartogs's phenomenon for polyregular functions and projective dimensions of related modules over a polynomial ring}, Ann. Inst. Fourier {\bf 47} (1997) 623--640.  %%QUATERNION, %%ANALYTICITY-H, %%HARTOGS, %$D17052001.

\bibitem{GURLE1997-} K. G\"urlebeck and W. Spr\"ossig, Quaternionic and Clifford Calculus for Physicists and Engineers (John Wiley, New York, 1997) 371~pp.  %%BOOK, %%QUATERNION, %%ANALYTICITY-H, %$D20032002.

\bibitem{NAPOL1997-} D. Napoletani and D.C. Struppa, \emph{On a large class of supports for quaternionic hyperfunctions in one variable}, Pitman  
Res. Notes Math. Ser. {\bf 394} (1997) 170--175. %%QUATERNION, %%ANALYTICITY-H, %$D04022002.

\bibitem{SABADI1997-} I. Sabadini and D.C. Struppa, \emph{Some open problems on the analysis of the Cauchy-Fueter system in several variables}, Surikaisekiken Kyusho Kokyuroku {\bf 1001} (1997) 1--21. %%QUATERNION, %%ANALYTICITY-H, %$D04022002.

\bibitem{ADAMS1999-} W.W. Adams, C.A. Berenstein, P. Loutaunau, I. Sabadini, and D.C. Struppa, \emph{Regular functions of several quaternionic variables and the Cauchy-Fueter complex}, J. Geom. Analysis {\bf 9} (1999) 1--15.  %%QUATERNION, %%ANALYTICITY-H, %%HARTOGS, %$D17052001.

\bibitem{BERNS1999-} S. Bernstein, \emph{The quaternionic Riemann problem}, Contemporary Mathematics {\bf 232} (1999) 69--83. %%QUATERNION, %%ANALYTICITY-H, %$D08082001.

\bibitem{ADAMS2000-} W.W. Adams and P. Loutaunau, \emph{Analysis of the module determining the properties of regular functions of several quaternionic variables}, Pacific J. of Math. {\bf 196} (2000) 1--15.  %%QUATERNION, %%ANALYTICITY-H, %%HARTOGS, %$D17052001.

\bibitem{SABAD2000A} I. Sabadini, M.V. Shapiro, and D.C. Struppa, \emph{Algebraic analysis of the Moisil-Theodorescu system}, Complex Variables {\bf 40} (2000) 333--357.  %%QUATERNION, %%ANALYTICITY-H, %%HARTOGS, %$D17052001.

\bibitem{SABADI2000B} I. Sabadini, F. Sommen, and D.C. Struppa, \emph{Computational algebra and its promise for analysis}, Quaderni  
di Matematica {\bf 7} (2000) 293--320. %%QUATERNION, %%ANALYTICITY-H, %%HARTOGS, %$D04022002.

\bibitem{ERIKS2001-} S.-L. Eriksson-Bique, \emph{Hyperholomorphic functions in $\mathbb{R}^4$}, in: S. Marchiafava et al., eds., Proceedings of the 2nd Meeting on Quaternionic Structures in Mathematics and Physics (World Scientific, Singapore, 2001) 125--135. %%QUATERNION, %%ANALYTICITY-H, %$D04032006.

\bibitem{LUDKO2003-} S.V. L\"udkovsky and F. vanOystaeyen, \emph{Differentiable functions of quaternion variables}, Bull. Sci. Math. {\bf 127} (2003) 755--796.  %%QUATERNION, %%ANALYTICITY-H, %%NONCOMMUTATIVE-GEOMETRY, %$D30112005.

\bibitem{ZIEGL2003-} M. Ziegler, \emph{Quasi-optimal arithmetic for quaternion polynomials}, Proc. 14th ISAAC, Springer LNCS {\bf 2906} (2003) 705--715;  e-print \underline{ arXiv:cs.SC/0304004 }.  %%QUATERNION, %%ANALYTICITY-H, %%NONCOMMUTATIVE-GEOMETRY, %$D30112005.

\bibitem{ALAYO2005-} D. Alay\'on-Solarz, \emph{On some modifications of the Fueter operator} (11 November 2005) 11~pp.;  e-print \underline{ arXiv:math.AP/0412125 }.  %%QUATERNION, %%ANALYTICITY-H, %$D30112005.

\bibitem{PEROT2006-} A. Perotti, \emph{Quaternionic regularity and the $\overline{\partial}$-Neumann problem in $C^2$}, to appear on Complex Variables and Elliptic Equations (2006) 13~pp.; e-print \underline{ arXiv:math/0612092 }.  %%QUATERNION, %%ANALYTICITY-H, %$D06072008.

\bibitem{PINOT2006-} D.A. Pinotsis, The Dbar formalism, Quaternions and Applications, PhD Thesis (University of Cambridge, 2006) 124~pp. %%QUATERNION, %%ANALYTICITY-H, %$D06072008. 

\bibitem{COLOM2007A} F. Colombo, G. Gentili, I. Sabadini and D.C. Struppa, \emph{Non commutative functional calculus: bounded operators} (2007) 18~pp.; e-print \underline{ arXiv:0708.3591 }.  %%QUATERNION, %%ANALYTICITY-H, %$D06072008.

\bibitem{COLOM2007B} F. Colombo, G. Gentili, I. Sabadini and D.C. Struppa, \emph{Non commutative functional calculus: unbounded operators} (2007) 13~pp.; e-print \underline{ arXiv:0708.3592 }. %%QUATERNION, %%ANALYTICITY-H, %$D06072008.

\bibitem{COLOM2007C} F. Colombo, I. Sabadini and D.C. Struppa, \emph{A new functional calculus for non-commuting operators} (2007) 18~pp.; e-print \underline{ arXiv:0708.3594 }.  %%QUATERNION, %%ANALYTICITY-H, %$D06072008.

\end{enumerate}

\subsection{ANALYTICITY-B}
%-------------------------
\label{ANALYTICITY-B}


Biquaternions analyticity only --- $\mathbb{B}$ alone.

See also ANALYTICITY-MAXWELL, Sec.~\ref{ANALYTICITY-MAXWELL}, and ANALYTICITY-DIRAC, Sec.~\ref{ANALYTICITY-DIRAC}.

\begin{enumerate}

\bibitem{SOUCE1983A} V. Soucek, \emph{Complex-quaternionic analysis applied to spin 1/2 massless fields}, Complex Variables {\bf 1} (1983) 327--346. %%QUATERNION, %%ANALYTICITY-B, %%DIRAC, %$D29052001.

\bibitem{SOUCE1983B} V. Soucek, \emph{(I) Holomorphicity in quaternionic analysis. (II) Complex quaternionic analysis, connections to mathematical Physics. (III) Cauchy integral formula}, in: Seminari di Geometria 1982-1983 (Universita di Bologna, Bologna, 1983) 147--171. %%QUATERNION, %%ANALYTICITY-B, %%TWISTOR, %$D29052001.

\bibitem{MESKA1984-} J. Meska, \emph{Regular functions of complex quaternionic variable}, Czech. Math. J. {\bf 34} (1984) 130--145. %%QUATERNION, %%ANALYTICITY-B, %$D16052001.

\bibitem{KRAVC1993-} V.G. Kravchenko and M.V. Shapiro, \emph{On the generalized system of Cauchy-Riemann equations with a quaternion parameter}, Russian Acad. Sci. Dokl. Math. {\bf 47} (1993) 315--319.  %%QUATERNION, %%ANALYTICITY-B, %$D29052001.

\bibitem{KRAVC1994A} V.V. Kravchenko and M.V. Shapiro, \emph{Helmholtz operator with a quaternionic wave number and associated function theory}, in: J. Lawrynowicz, ed., Deformations of mathematical structures II (Kluwer, Dordrecht, 1994) 101--128.  %%QUATERNION, %%ANALYTICITY-B, %$D01062001

\bibitem{KRAVC1995B} V.V. Kravchenko, \emph{Direct sum expansion of the Kernel of the Laplace operator with the aid of biquaternion zero divisors}, Diff. Equations. {\bf 31} (1995) 462--465. %%QUATERNION, %%ANALYTICITY-B, %%HARTOGS, %%LAPLACE-OPERATOR, %$D17052001.

\bibitem{SABAD2002A} I. Sabadini and D.C. Struppa, \emph{First order differential operators in real dimension eight}, Complex Variables  
{\bf 47} (2002) 953--968. %%QUATERNION, %%ANALYTICITY-B, %%OCTONION, %%CLIFFORD-R, %$D04022002.


\end{enumerate}

\subsection{ANALYTICITY-CLIFFORD-R}
%----------------------------------
\label{ANALYTICITY-CLIFFORD-R}


Real Clifford analysis --- but not $\mathbb{H}$ alone

\begin{enumerate}

\bibitem{NEF--1942-} W. Nef, \emph{Uber die singularen Gebilde der regul\"aren Funktionen einer Quaternionenvariabeln},  Comm. Math. Helv. {\bf 15} (1942/1943) 144--174. %%QUATERNION, %%ANALYTICITY-CLIFFORD-R, %%QUADRIPERIODIC-FUNCTION, %$D17052001.

\bibitem{NEF--1944-} W. Nef, \emph{Funktionentheorie einer Klasse von hyperbolischen und ultrahyperbolischen Differentialgleichungen zweiter Ordnung},  Comm. Math. Helv. {\bf 17} (1944/1945) 83--107. %%QUATERNION, %%ANALYTICITY-CLIFFORD-R, %$D17052001.

\bibitem{HAEFE1947-} H.G. Haefeli, \emph{Hypercomplexe Differentiale}, Comm. Math. Helv. {\bf 20} (1947/1948) 382--420. %%QUATERNION, %%ANALYTICITY-CLIFFORD-R, %$D13032002.

\bibitem{BRACK1982-} F. Brackx, R. Delanghe and F. Sommen, Clifford Analysis (Pitman Books, London, 1982) 308~pp.  %%BOOK, %%QUATERNION, %%ANALYTICITY-CLIFFORD-R, %$D01062001.

\bibitem{GOLDS1982A} B. Goldschmidt, \emph{Existence and representation of solutions of a class of elliptic systems of partial differential equations of first order in the space}, Math. Nachr. {\bf 108} (1982) 159--166. %%QUATERNION, %%ANALYTICITY-CLIFFORD-R, %$D01062001.

\bibitem{GOLDS1982B} B. Goldschmidt, \emph{A Cauchy integral formula for a class of elliptic systems of partial differential equations of first order in the space}, Math. Nachr. {\bf 108} (1982) 167--178. %%QUATERNION, %%ANALYTICITY-CLIFFORD-R, %$D01062001. 

\bibitem{RYAN-1982B} J. Ryan, \emph{Clifford analysis with generalized elliptic and quasi elliptic functions}, Applicable Anal. {\bf 13} (1982) 151--171.   %%QUATERNION, %%ANALYTICITY-CLIFFORD-R, %$D01062001. %$C Includes an historical overview.

\bibitem{LOUNE1983-} P. Lounesto and P. Bergh, \emph{Axially symmetric vector fields and their complex potentials}, Complex Variables {\bf 2} (1983) 139--150.  %%QUATERNION, %%ANALYTICITY-CLIFFORD-R, %$D01062001

\bibitem{RYAN-1984A} J. Ryan, \emph{Properties of isolated singularities of some functions taking values in real Clifford algebras}, Math. Proc. Cambridge Soc. {\bf A 84} (1984) 37--50. %%QUATERNION, %%ANALYTICITY-CLIFFORD-R, %$D01062001.

\bibitem{SOMME1984A} F. Sommen, \emph{Monogenic differential forms and homology theory}, Proc. Roy. Irish Acad. {\bf A 84} (1984) 87--109. %%QUATERNION, %%ANALYTICITY-CLIFFORD-R, %$D01062001.

\bibitem{LAWRY1986-} J. Lawrynowicz and J. Rembielinski, \emph{Pseudo-Euclidiean Hurwitz pairs and Generalized Fueter equations}, in: J.S.R. Chisholm and A.K. Common, eds.,  Clifford Algebras and Their Applications in Mathematical Physics (Reidel, Dordrecht, 1986)  39--48. %%QUATERNION, %%ANALYTICITY-CLIFFORD-R, %$D01062001.

\bibitem{RYAN-1986-} J. Ryan, \emph{Left regular polynomials in even dimensions, and tensor product of Clifford algebras},  in NATO Adv. Sci. Inst. ser. C. Math. Phys. Sci. {\bf 183} (Reidel, Dordrecht, 1986) 133--147. %%QUATERNION, %%ANALYTICITY-CLIFFORD-R, %$D01062001.

\bibitem{SOMME1986-} F. Sommen, \emph{Microfunctions with values in Clifford algebra II}, Sci. Papers of the Coll. of Art and Sci. {\bf 36} (University of Tokyo, 1986) 15--37. %%QUATERNION, %%ANALYTICITY-CLIFFORD-R, %%MAXWELL, %%FOURIER-TRANSFORM, %$D04022002.

\bibitem{SEMME1989-} S.W. Semmes, \emph{A criterion for the boundness of singular integrals on hypersurfaces}, Trans. Amer. Math. Soc. {\bf 311} (1989) 501--513. %%QUATERNION, %%ANALYTICITY-CLIFFORD-R, %$D01062001.

\bibitem{ZOLL-1989-} G. Z\"oll, \emph{Regular $n$-forms in Clifford analysis, their behavior under change of variables and their residues}, Complex Variables {\bf 11} (1989) 25--38. %%QUATERNION, %%ANALYTICITY-CLIFFORD-R, %$D01062001.

\bibitem{DELAN1990-}  R. Delanghe, F. Sommen, and X. Zhenyuan, \emph{Half Dirichlet problems for powers of the Dirac operator in the unit ball of R$^m (m \geq 3)$}, Bull. Soc. Math. Belg. {\bf 42} (1990) 409--429. %%QUATERNION, %%ANALYTICITY-CLIFFORD-R, %$D01062001.

\bibitem{XU---1990A} Z. Xu, \emph{On linear and nonlinear Riemann-Hilbert problems for regular function with values in a Clifford algebra}, Suppl. Chin. Ann. of Math. {\bf 11B} (1990) 349--358. %%QUATERNION, %%ANALYTICITY-CLIFFORD-R, %$D01062001.

\bibitem{XU---1990B} Z. Xu, \emph{On boundary value problem of Neumann type for hypercomplex function with values in a Clifford algebra}, Suppl. Rendiconti Circ. Math. Palermo {\bf 22} (1990) 213--226. %%QUATERNION, %%ANALYTICITY-CLIFFORD-R, %$D01062001.

\bibitem{GILBE1991-} J.E. Gilbert and M.A.M. Murray, Clifford Algebras and Dirac Analysis (Cambridge University Press, Cambridge, 1991) 334~pp. %%BOOK, %%QUATERNION, %%ANALYTICITY-CLIFFORD-R, %%WEISS, %$D20032002.

\bibitem{XU---1991-} Z. Xu, \emph{A function theory for the operator $D-\lambda$}, Complex Variables {\bf 16} (1991) 27--42. %%QUATERNION, %%ANALYTICITY-CLIFFORD-C, %$D01062001.

\bibitem{DELAN1992B} R. Delanghe, F. Sommen, and V. Soucek, \emph{Residues in Clifford analysis}, in: H. Begehr and A. Jeffrey, eds., \emph{Partial differential equations with complex analysis}, Pitnam Research Notes in Math. {\bf 262} (Longman, Burnt Hill, 1992) 61--92. %%QUATERNION, %%ANALYTICITY-CLIFFORD-R, %$D01062001.

\bibitem{RYAN-1992A} J. Ryan, \emph{Generalized Schwarzian derivatives for generalized fractional linear transformations}, Ann. Polonici Math. {\bf 57} (1992) 29--44. %%QUATERNION, %%ANALYTICITY-CLIFFORD-R, %$D01062001.

\bibitem{XU---1992-} Z. Xu, \emph{Helmoltz equations and boundary value problems}, in: H. Begehr and A. Jeffrey, eds., Partial Differential Equations With Complex Analysis, Pitnam Research Notes in Math. {\bf 262} (Longman, Burnt Hill, 1992) 204--214.  %%QUATERNION, %%ANALYTICITY-CLIFFORD-R, %$D01062001. %$C name:  Xu Zhenyuan.

\bibitem{BEGEH1993-}\emph{Chap IX:  Clifford analysis}, in: H. Begehr and R. Gilbert, Transformations, Transmutations, and Kernel Functions, Vol.~2 (Longman, New York, 1993) 215--240. %%QUATERNION, %%ANALYTICITY-CLIFFORD-R, %$D01062001.

\bibitem{SOMME1993-} F. Sommen and N. vanAcker, \emph{SO(m)-invariant differential operators on Clifford algebra-valued functions}, Found. Phys. {\bf 23} (1993) 1491--1519. %%QUATERNION, %%ANALYTICITY-CLIFFORD-R, %$D04022002.

\bibitem{GURLE1995-} K. G\"urlebeck and F. Kippig, \emph{Complex Clifford-analysis and elliptic boundary value problems}, Adv. Appl. Clifford Alg. {\bf 5} (1995) 51--62. %%QUATERNION, %%ANALYTICITY-CLIFFORD-C, $D24062002.

\bibitem{SPROS1996-} W. Spr\"ossig and K. G\"urlebeck, eds., Proc. of the Symp. ``Analytical and Numerical Methods in Quaternionic and Clifford Analysis,'' Seiffen, June 5--7, 1996 (TU Bergakademie Freiberg, 1996) 228~pp.  %%BOOK, %%QUATERNION, %%ANALYTICITY-CLIFFORD-R, %$D20032002.

\bibitem{SOMME1997B} F. Sommen and P. vanLancker, \emph{A product for special classes of monogenic functions and tensors}, Z. f\"ur Anal. und ihre Anwendungen {\bf 16} (1997) 1013--1026. %%QUATERNION, %%ANALYTICITY-CLIFFORD-R, %$D04022002.

\bibitem{VANLA1999A} P. vanLancker, \emph{Approximation theorems for spherical monogenics of complex degree}, Bull. Belg. Math. Soc. {\bf 6} (1999) 279--293. %%QUATERNION, %%ANALYTICITY-CLIFFORD-R, %$D04022002.

\bibitem{VANLA1999B} P. vanLancker, \emph{Taylor and Laurent series on the sphere}, Complex Var. {\bf 38} (1999) 321--365. %%QUATERNION, %%ANALYTICITY-CLIFFORD-R, %$D04022002.

\bibitem{SOMME1999A} F. Sommen, \emph{Clifford analysis in two and several variables}, Appl. Anal. {\bf 73} (1999) 225--253. %%QUATERNION, %%ANALYTICITY-CLIFFORD-R, %$D04022002.

\bibitem{SOMME2000-} F. Sommen, \emph{On a generalization of Fueter's theorem}, Z. f\"ur Anal. und ihre Anwendungen {\bf 19} (2000) 899--902. %%QUATERNION, %%ANALYTICITY-CLIFFORD-R, %$D04022002.

\bibitem{BERNS2001B} S. Bernstein, \emph{Integralgleichungen und Funktionenr\"aume f\"ur Randwerte monogener Funktionen}, Habilitation thesis (TU Bergakademie, Freiberg, 30 April 2001) 97~pp. %%QUATERNION, %%ANALYTICITY-CLIFFORD-R, %%CLIFFORD-C, %%THESIS, %$D28082001.

\bibitem{SOMME2001A} F. Sommen, \emph{Clifford analysis on the level of abstract vector variables}, in: F. Brackx et al., eds., Clifford Analysis and its Applications (Kluwer acad. publ., 2001) 303--322. %%QUATERNION, %%ANALYTICITY-CLIFFORD-R, %$D04022002.

\bibitem{SABAD2002B} I. Sabadini, F. Sommen, and D.C. Struppa, \emph{Series and integral representations for the biregular exponential function}, J. Natural Geom.  {\bf 21} (2002) 1--16. %%QUATERNION, %%ANALYTICITY-CLIFFORD-R, %%FOURIER-TRANSFORM, %$D04022002.

\bibitem{SABAD2002C} I. Sabadini, F. Sommen, D.C. Struppa, and P. vanLancker, \emph{Complexes of Dirac Operators in Clifford Algebras},
Math. ZS {\bf 239} (2002) 293--320 %%QUATERNION, %%ANALYTICITY-CLIFFORD-R, %$D04022002.

\bibitem{LAVIL2005A} G. Laville and E. Lehmann, \emph{Holomorphic Cliffordian product} (4 February 2005) 18~pp.;  e-print \underline{ arXiv:math.CV/0502088 }.  %%QUATERNION, %%ANALYTICITY-CLIFFORD-R, %$D30112005.

\bibitem{LAVIL2005B} G. Laville and E. Lehmann, \emph{Analytic Cliffordian functions} (4 February 2005) 19~pp.;  e-print \underline{ arXiv:math.CV/0502090 }.  %%QUATERNION, %%ANALYTICITY-CLIFFORD-R, %$D30112005.

\bibitem{COLOM2007D} F. Colombo, I. Sabadini and D.C. Struppa, \emph{Slice monogenic functions} (2007) 14~pp.; e-print \underline{ arXiv:0708.3595 }.  %%QUATERNION, %%ANALYTICITY-CLIFFORD-R, %$D06072008.

\bibitem{SEMME2007-} S. Semmes, \emph{Some remarks about Clifford analysis and fractal sets} (2007) 5~pp.; e-print \underline{ arXiv:0709.2356 }.   %%QUATERNION, %%ANALYTICITY-CLIFFORD, %$D06072008.


\end{enumerate}

\subsection{ANALYTICITY-CLIFFORD-C}
%----------------------------------
\label{ANALYTICITY-CLIFFORD-C}

Complex Clifford analysis --- but not $\mathbb{B}$ alone

\begin{enumerate}

\bibitem{FUETE1941-} R. Fueter, \emph{Uber ein Hartogs'schen Satz in der Theorie der analytischen Funktionen von $n$\, komplexen Variablen}, Comm. Math. Helv. {\bf 14} (1941/1942) 394--400. %%QUATERNION, %%ANALYTICITY-CLIFFORD-C, %%FUETER, %%HARTOGS, %$D17052001.

\bibitem{RYAN-1982C} J. Ryan, \emph{Complexified Clifford analysis}, Complex Variables {\bf 1} (1982) 119--149.  %%QUATERNION, %%ANALYTICITY-CLIFFORD-C, %$D01062001.

\bibitem{SOMME1982-} F. Sommen, \emph{Some connections between Clifford analysis and complex analysis}, Complex Variables {\bf 1} (1982) 97--118.  %%QUATERNION, %%ANALYTICITY-CLIFFORD-C, %$D01062001.

\bibitem{RYAN-1983-} J. Ryan, \emph{Singularities and Laurent expansions in complex Clifford analysis}, Applicable Anal. {\bf 16} (1983) 33--49.  %%QUATERNION, %%ANALYTICITY-CLIFFORD-C, %$D01062001.

\bibitem{BURES1984-} J. Bures, \emph{Some integral formulas in complex Clifford analysis}, Rendiconti Circ. Mat. Palermo--Suppl. {\bf 3} (1984) 81--87. %%QUATERNION, %%ANALYTICITY-CLIFFORD-C, %$D01062001.

\bibitem{RYAN-1984B} J. Ryan, \emph{Extensions of Clifford analysis to complex, finite dimensional, associative algebras with identity}, Proc. Roy. Irish Acad. {\bf 95} (1984) 277--298. %%QUATERNION, %%ANALYTICITY-CLIFFORD-C, %$D01062001.

\bibitem{RYAN-1984C} J. Ryan, \emph{Cauchy-Kowalewski extension theorems and representations of analytic functionals acting over special classes of real $n$-dimensional submanifolds of $C^{n+1}$}, Rendiconti Circ. Mat. Palermo Suppl. {\bf 3} (1984) 249--262.  %%QUATERNION, %%ANALYTICITY-CLIFFORD-C, %$D01062001.

\bibitem{BURES1985-} J. Bures and V. Soucek, \emph{Generalized hypercomplex analysis and its integral formulas}, Complex variables {\bf 5} (1985) 53--70. %%QUATERNION, %%ANALYTICITY-CLIFFORD-C, %$D01062001.

\bibitem{RYAN-1985-} J. Ryan, \emph{Conformal Clifford arising in Clifford analysis}, Proc. Roy. Irish Acad. {\bf A 85} (1985) 1--23. %%QUATERNION, %%ANALYTICITY-CLIFFORD-C, %$D01062001.

\bibitem{RYAN-1987-} J. Ryan, \emph{Applications of complex Clifford analysis to the Study of solutions to generalized Dirac and Klein-Gordon equations with holomorphic potentials}, J. Differential Eq. {\bf 67} (1987) 295--329. %%QUATERNION, %%ANALYTICITY-CLIFFORD-C, %$D01062001.

\bibitem{SOMME1987-} F. Sommen, \emph{Martinelli-Bochner type formulae in complex Clifford analysis}, Z. f\"ur Anal. und ihre Anwendungen {\bf 6} (1987) 75--82. %%QUATERNION, %%ANALYTICITY-CLIFFORD-C, %$D04022002.

\bibitem{RYAN-1990-} J. Ryan, \emph{Cells of harmonicity and generalized Cauchy integral formulae}, Proc. London Math. Soc. {\bf 60} (1990) 295--318. %%QUATERNION, %%ANALYTICITY-CLIFFORD-C, %$D01062001.

\bibitem{RYAN-1992B} J. Ryan, \emph{Plemelj formulae and transformations associated to plane wave decompositions in complex Clifford analysis}, Proc. London Math. Soc. {\bf 64} (1992) 70--94. %%QUATERNION, %%ANALYTICITY-CLIFFORD-C, %$D01062001.

\bibitem{RYAN-1996-} J. Ryan, \emph{Intrinsic Dirac operator in $C^n$}, Advances in Mathematics (1996) 99--133. %%QUATERNION, %%ANALYTICITY-CLIFFORD-C, %$D24062002.

\bibitem{SANO-1997-} K. Sano, \emph{Another type of Cauchy's integral formula in complex Clifford analysis}, Tokyo J. Math. {\bf 20} (1997) 187--204. %%QUATERNION, %%ANALYTICITY-CLIFFORD-C, %$D03082005.

\bibitem{BERNS2000-} S. Bernstein, \emph{A Borel-Pompeiu formula in $C^n$ and its application to inverse scattering theory}, in: J. Ryan and W. Spr\"ossig, eds., Clifford Algebra and their Applications in Mathematical Physics, Vol.~2: \emph{Clifford Analysis} (Birkh\"auser, Boston, 2000) 117--134. %%QUATERNION, %%ANALYTICITY-CLIFFORD-C, %%INVERSE-SCATTERING, %$D12052001.

\bibitem{BERNS2002-} S. Bernstein, \emph{Multidimensional inverse scattering and Clifford analysis}, Applied Mathematical Letters {\bf 15} (2002) 1035--1041. %%QUATERNION, %%ANALYTICITY-CLIFFORD-C, %%INVERSE-SCATTERING, %$D23072001.

\end{enumerate}


\subsection{MANIFOLD}
%--------------------
\label{MANIFOLD}

Quaternionic manifolds, K\"ahlerian manifolds, symmetric-spaces, etc.

\begin{enumerate}

\bibitem{STRIN1901-} I. Stringham, \emph{On the geometry of planes in a parabolic space of four dimensions}, Trans. Amer. Math. Soc. {\bf 2} (1901) 183--214. %%QUATERNION, %%MANIFOLDS, %%GEOMETRY, %$D24062002.

\bibitem{HATHA1902-} A.S. Hathaway, \emph{Quaternion space}, Trans. Amer. Math. Soc. {\bf 3} (1902) 46--59. %%QUATERNION, %%MANIFOLD, %%GEOMETRY, %%TURN, %%CLIFFORD, %$D14032002.

\bibitem{LEMAI1948-} G. Lemaitre, \emph{Quaternions et espace elliptique}, Acta Pontifica Acad. Scientiarium {\bf 12} (1948) 57--80.  %%QUATERNION, %%MANIFOLD, %$D06072008.

\bibitem{WOLF-1965} J.A. Wolf, \emph{Complex homogeneous contact manifolds and quaternion symmetric spaces}, J. Math. Mech., i.e., Indiana Univ. Math. J., {\bf 14} (1965) 1033-1048. %%QUATERNION, %%MANIFOLD, %%WOLF-SPACE, %$D04032006.

\bibitem{KRAIN1966-} V.Y. Kraines, \emph{Topology of quaternionic manifolds}, Trans. Amer. Math. Soc. {\bf 122} (1966) 357--367.  %%QUATERNION, %%MANIFOLD, GROUP-THEORY, KAHLER, %$D24062002.

\bibitem{ALEKS1968-} D.V. Alekseevsky, \emph{Compact quaternion spaces}, Functional Analysis {\bf 2} (1968) 106--114.  %%QUATERNION, %%MANIFOLD, %%WOLF-SPACE, %%KAHLER, %SYMMETRIC-SPACE, %$D04032006, %$C Theorem: A compact  quaternion K\"ahler homogeneous space P/K is symmetric.

\bibitem{ISHIH1974-} S. Ishihara, \emph{Quaternion K\"ahlerian manifolds}, J. Differential geometry {\bf 9} (1974) 483--500.  %%QUATERNION, %%MANIFOLD, %%KAHLER, %$D14032002.

\bibitem{HAHL-1975A} H. H\"ahl, \emph{Automorphismengruppen von lokalkompakten zusammen\-h\"angen\-den Quasik\"orpern und Translationebenen}, Geom. Dedicata {\bf 4} (1975) 305--321. %%QUATERNION, %%MANIFOLD, %$D04032006.

\bibitem{HAHL-1975B} H. H\"ahl, \emph{Vierdimensionale reelle Divisionalgebren mit dreidimensionaler Automorphismengruppen}, Geom. Dedicata {\bf 4} (1975) 323--331. %%QUATERNION, %%MANIFOLD, %$D04032006.

\bibitem{HAHL-1975C} H. H\"ahl, \emph{Geometrische homogen vierdimensionale reelle Divisionalgebren}, Geom. Dedicata {\bf 4} (1975) 333--361. %%QUATERNION, %%MANIFOLD, %$D04032006.

\bibitem{SALAM1986-} S.M. Salamon, \emph{Differential geometry of quaternionic manifolds}, Ann. Sc. Ec. Norm. Sup. {\bf 19} (1986) 31--55.   %%QUATERNION, %%MANIFOLD, %%KAHLER, %$D14032002.

\bibitem{HITCH1987-} N.J. Hitchin, A. Karlhede, U. Lindstrom, and M. Rocek, \emph{Hyperk\"ahler metrics and supersymmetry}, Commun. Math. Phys. {\bf 108} (1987) 535--589.  %%QUATERNION, %%MANIFOLD, %%SUPERSYMMETRY, %%KAHLER, %%TWISTOR, %$D06072008. 

\bibitem{SWANN1989-} A. Swann, \emph{Symplectic aspects of quaternionic geometry}, C. R. Acad. Sci. Math. {\bf 308} (1989) 225--228. %%QUATERNION, %%MANIFOLDS, %%KAHLER, %$D30112005.

\bibitem{JOYCE1991-} D. Joyce, \emph{The hypercomplex quotient and the quaternionic quotient}, Math. Annalen {\bf 290} (1991) 323--340. %%QUATERNION, %%MANIFOLD, %%KAHLER, %$D14032002.

\bibitem{ALEKS1994-} D.V. Alekseevsky and S. Marchiafava, \emph{Gradient quaternionic vector fields and a characterization of the quaternionic projective space}, Preprint ESI-138, to appear in C. R. Acad. Sci. Paris (Erwin Schr\"odinger Institute, Vienna, 1994) 8 pp.  Available at\\ \underline{ http://www.esi.ac.at/preprints/ESI-Preprints.html }.  %%QUATERNION, %%MANIFOLD, %%KAHLER, %$D04032006.

\bibitem{BONAN1994} E. Bonan, \emph{Isomorphismes sur une vari\'et\'e presque Hermitienne quaternionique}, in: G. Gentili et al., Proc. of the Meeting on Quaternionic Structures in Mathematics and Physics (SISSA, Trieste, 1994) 1--6. %%QUATERNION, %%MANIFOLD, %$D04032006.

\bibitem{BOYER1994-} C.P. Boyer, K. Galicki and B.M. Mann, \emph{Quaternionic geometry and 3-Sasakian manifolds}, in: G. Gentili et al., Proc. of the Meeting on Quaternionic Structures in Mathematics and Physics (SISSA, Trieste, 1994) 7--24. %%QUATERNION, %%MANIFOLD, %$D04032006.

\bibitem{CNOPS1994-} J. Cnops, \emph{Stokes' formula and the Dirac operator on imbedded manifolds}, in: G. Gentili et al., Proc. of the Meeting on Quaternionic Structures in Mathematics and Physics (SISSA, Trieste, 1994) 26--38. %%QUATERNION, %%MANIFOLD, %%ANALYTICITY, %$D04032006.

\bibitem{CORTE1994-} V. Cortes, \emph{Alekseevskian Spaces}, in: G. Gentili et al., Proc. of the Meeting on Quaternionic Structures in Mathematics and Physics (SISSA, Trieste, 1994) 39--91.  %%QUATERNION, %%MANIFOLD, %%KAHLER, %$D04032006.

\bibitem{DEWIT1994B} B. deWit and A. vanProeyen, \emph{Isometries of special manifolds}, in: G. Gentili et al., Proc. of the Meeting on Quaternionic Structures in Mathematics and Physics (SISSA, Trieste, 1994) 92--118. %%QUATERNION, %%MANIFOLD, %%KAHLER, %$D04032006.

\bibitem{FERNA1994A} M. Fernandez, R. Ibanez, and DeLeon, \emph{On a Brylinski conjecture for compact symplectic manifolds}, in: G. Gentili et al., Proc. of the Meeting on Quaternionic Structures in Mathematics and Physics (SISSA, Trieste, 1994) 119--126. %%QUATERNION, %%MANIFOLD, %%KAHLER, %$D04032006.

\bibitem{FERNA1994B} M. Fernandez and L. Ugarte., \emph{Canonical cohomology of compact $G_2$-nilmanifolds}, in: G. Gentili et al., Proc. of the Meeting on Quaternionic Structures in Mathematics and Physics (SISSA, Trieste, 1994) 127--138.  %%QUATERNION, %%MANIFOLD, %%KAHLER, %$D04032006.

\bibitem{GALPE1994-} A. Galperin, E. Ivanov, and V. Ogievetsky, \emph{Harmonic space and quaternionic manifolds}, Ann. of Phys. {\bf 230} (1994) 201--249. %%QUATERNION, %%MANIFOLDS, %$D24062002.

\bibitem{HANGA1994-} T. Hangan, \emph{On infinitesimal automorphisms of quaternionic manifolds}, in: G. Gentili et al., Proc. of the Meeting on Quaternionic Structures in Mathematics and Physics (SISSA, Trieste, 1994) 147--150.  %%QUATERNION, %%MANIFOLD, %%KAHLER, %$D04032006.

\bibitem{HIJAZ1994-} O. Hijazi, \emph{Twistor operators and eigenvalues of the Dirac operator}, in: G. Gentili et al., Proc. of the Meeting on Quaternionic Structures in Mathematics and Physics (SISSA, Trieste, 1994) 151--174.  %%QUATERNION, %%MANIFOLD, %%KAHLER, %$D04032006.

\bibitem{LICHN1994-} A. Lichnerowicz, \emph{Complex contact homogenous spaces and quaternion-K\"ahler symmetric spaces}, in: G. Gentili et al., Proc. of the Meeting on Quaternionic Structures in Mathematics and Physics (SISSA, Trieste, 1994) 175--188. %%QUATERNION, %%MANIFOLD, %%KAHLER, %$D04032006.

\bibitem{NAGAN1994-} T. Nagano, \emph{Symmetric spaces and quaternionic structures}, in: G. Gentili et al., Proc. of the Meeting on Quaternionic Structures in Mathematics and Physics (SISSA, Trieste, 1994) 203--218.  %%QUATERNION, %%MANIFOLD, %%KAHLER, %$D04032006.

\bibitem{NAGAT1994-} Y. Nagatomo, \emph{Instantons on quaternion-K\"ahler manifolds}, in: G. Gentili et al., Proc. of the Meeting on Quaternionic Structures in Mathematics and Physics (SISSA, Trieste, 1994) 219--230. %%QUATERNION, %%MANIFOLD, %%KAHLER, %%INSTANTONS, %$D04032006.

\bibitem{ORNEA1994-} L. Orena and P. Piccinni, \emph{Weyl structures on quaternionic manifolds}, in: G. Gentili et al., Proc. of the Meeting on Quaternionic Structures in Mathematics and Physics (SISSA, Trieste, 1994) 231--236. %%QUATERNION, %%MANIFOLD, %%KAHLER, %$D04032006.

\bibitem{SWANN1994-} A. Swann, \emph{Quaternionic K\"ahler metrics and nilpotent orbits}, in: G. Gentili et al., Proc. of the Meeting on Quaternionic Structures in Mathematics and Physics (SISSA, Trieste, 1994) 259--267.  %%QUATERNION, %%MANIFOLD, %%KAHLER, %$D04032006.

\bibitem{ALEKS1995-} D.V. Alekseevsky and S. Marchiafava, \emph{Quaternionic transformations and the first eigenvalues of laplacian on a quaternionic K\"ahler manifold}, C. R. Acad. Sci. Paris Ser. I Math. {\bf 320} (1995) 703--708.   %%QUATERNION, %%MANIFOLD, %%KAHLER, %$D04032006.

\bibitem{BATTA1996A} F. Battaglia, \emph{$S^1$-quotients of quaternion-K\"ahler manifolds}, Proc. Amer. Math. Soc. {\bf 124} (1996) 2185--2192.  %%QUATERNION, %%MANIFOLD, %%KAHLER, %$D04032006.

\bibitem{BATTA1996B} F. Battaglia, \emph{A hypercomplex Stiefel manifold}, Differ. Geom. Appl. {\bf 6} (1996) 121--128.  %%QUATERNION, %%MANIFOLD, %$D04032006.

\bibitem{ALEKS1997-} D.V. Alekseevsky and F. Podest\`a, \emph{Compact cohomogeneity one Riemannian manifolds of positive Euler characteristic and quaternionic K\"ahler manifolds}, in: de Gruyter, ed., Geometry, Topology and Physics (Campinas University, 1997) 1--33.   %%QUATERNION, %%MANIFOLD, %$D04032006.

\bibitem{ORNEA1997-} L. Ornea and P. Piccinni, \emph{Locally conformal K\"ahler structures in quaternionic geometry}, Trans. Am. Math. Soc. {\bf 349} (1997) 641--655.  %%QUATERNION, %%MANIFOLD, %%KAHLER, %$D04032006.

\bibitem{ORTEG1997-} M. Ortega and J. deDiosPeres, \emph{On the Ricci tensor of the real hypersurface of quaternionic hyperbolic space}, Manuscripta Math. {\bf 93} (1997) 49--57. %%QUATERNION, %%MANIFOLDS, %$D30112005, %$C EINSTEIN-MANIFOLD.

\bibitem{SWANN1997-} A. Swann, \emph{Some remarks on quaternion-Hermitian manifolds}, Archivum Mathematicum (Brno) {\bf 33} (1997) 349--354.  %%QUATERNION, %%MANIFOLD, %$D04032006.

\bibitem{ALEKS1998-} D.V. Alekseevsky, S. Marchiafava, and M. Pontecorvo, \emph{Compatible almost complex structures on quaternion-K\"ahler manifolds}, Ann. Global Anal. Geom. {\bf 16} (1998) 419--444.  %%QUATERNION, %%MANIFOLD, %%KAHLER, %$D04032006.

\bibitem{CADEK1998-} M. Cadek and J. Vanzura, \emph{Almost quaternionic structures on eight-manifolds}, Osaka J. Math. {\bf 35} (1998) 165--190.  %%QUATERNION, %%MANIFOLD, %%KAHLER, %$D04032006.

\bibitem{HONG-1998-} Y. Hong and C.S. Hou, \emph{Lagrangian submanifolds of quaternion Kaehlerian manifolds satisfying Chen's equality}, Contrib. to Algebra and Geometry {\bf 39} (1998) 413--421.  %%QUATERNION, %%MANIFOLD, %%KAHLER, %$D04032006.

\bibitem{PEDER1998-} H. Pedersen, Y.S. Poon, and A. Swann, \emph{Hypercomplex structures associated to quaternionic manifolds}, Differential Geom. Appl. {\bf 9} (1998) 273--292.  %%QUATERNION, %%MANIFOLD, %%KAHLER, %$D04032006.

\bibitem{TANIG1998-} T. Taniguchi, \emph{Isolation phenomena for quaternionic Yang-Mills connections}, Osaka J. Math. {\bf 35} (1998) 147--164.  %%QUATERNION, %%MANIFOLD, %%KAHLER, %%GAUGE-THEORY, %$D04032006.

\bibitem{ALEKS1999A} D.V. Alekseevsky and V. Cort\'es, \emph{Isometry groups of homogeneous quaternionic K\"ahler manifolds}, J. Geom. Anal. {\bf 9} (1999) 513--545.  %%QUATERNION, %%MANIFOLD, %%KAHLER, %$D04032006.

\bibitem{ALEKS1999B} D.V. Alekseevsky, S. Marchiafava, and M. Pontecorvo, \emph{Compatible complex structures on almost quaternionic manifolds}, Trans. Amer. Math. Soc. {\bf 351} (1999) 997--1014.  %%QUATERNION, %%MANIFOLD, %%KAHLER, %$D04032006.

\bibitem{POON-1999-} Y.S. Poon, \emph{Examples of hyper-K\"ahler connections with torsion},  in: S. Marchiafava et al., eds., Proceedings of the 2nd meeting on ``Quaternionic structures in mathematics and physics'' (World Scientific, Singapore, 2001) 321--327.  %%QUATERNION, %%MANIFOLD, %%KAHLER, %$D04032006.

\bibitem{ALEKS2001-} D.V. Alekseevsky and S. Marchiafava, \emph{Hermitian and K\"ahler submanifolds of a quaternionic K\"ahler manifold}, Osaka J. Math. {\bf 38} (2001) 869--904.  %%QUATERNION, %%MANIFOLD, %%KAHLER, %$D04032006.

\bibitem{PEDER2001-} H. Pedersen, \emph{Hypercomplex geometry}, in: S. Marchiafava et al., eds., Proceedings of the 2nd Meeting on Quaternionic Structures in Mathematics and Physics (World  
Scientific, Singapore, 2001) 313--319. %%QUATERNION, %%MANIFOLD, %$D04032006.

\bibitem{PICCI2001-} P. Piccinni and I. Vaisman, \emph{Foliations with transversal quaternionic structures}, Ann. Mat. Pura Appl. {\bf 180} (2001) 303--330.  %%QUATERNION, %%MANIFOLD %$D04032006.

\bibitem{SAWON2001-} J. Sawon, \emph{A new weight system on chord diagrams via hyperk�hler geometry},  in: S. Marchiafava et al., eds., Proceedings of the 2nd Meeting on Quaternionic Structures in Mathematics and Physics (World Scientific, Singapore, 2001) 349--363. %%QUATERNION, %%MANIFOLD, %%KAHLER, %$D04032006.

\bibitem{SWANN2001-} A. Swann, \emph{Weakening holonomy},  in: S. Marchiafava et al., eds., Proceedings of the 2nd Meeting on Quaternionic Structures in Mathematics and Physics (World Scientific, Singapore, 2001) 405--415.  %%QUATERNION, %%MANIFOLD, %$D04032006.

\bibitem{WOLF-2001-} J.A. Wolf, \emph{Complex forms of quaternionic symmetric spaces}, Progress in Mathematics {\bf 234} (2005) 265--277. Available at\\ \underline{ http://www.esi.ac.at/preprints/ESI-Preprints.html }.  %%QUATERNION, %%MANIFOLD, %%GROUP-THEORY, %%WOLF-SPACE, %$D04032006.

\bibitem{KIM--2003-} I. Kim and J.R. Parker, \emph{Geometry of quaternionic hyperbolic manifolds}, Math. Proc. Camb. Phil. Soc. {\bf 135}  (2003) 291--320. %%QUATERNION, %%MANIFOLD, %$D04032006.

\bibitem{LOWE-2003-} H. L\"owe, \emph{Sixteen-dimensional locally compact translation planes admitting $Sl_2(\mathbb{H})$ as a group of collineations}, Pacific J. Math. {\bf 209}  (2003) 325--337. %%QUATERNION, %%MANIFOLD, %$D04032006.

\bibitem{PEROT2005-} A. Perotti, \emph{Holomorphic functions and regular quaternionic functions on the hyperk\"ahler space $\mathbf{H}$}, to appear in: Proceedings V ISAAC Congress Catania (2005) 8~pp.; e-print \underline{ arXiv:0711.4440 }.  %%QUATERNION, %%MANIFOLD, %$D06072008.

\end{enumerate}



\section{\Huge RELATIVISTICS}
%============================
\label{RELATIVITICS}

Mathematical aspects of relativistics, i.e., relativity theory, and mathematical methods associated with that theory.  Papers dealing with specific applications to fields and physics are collected in Chaps.~\ref{FIELDS} and~\ref{PHYSICS}.

The exception to this rule is Sec.~\ref{GENERAL-RELATIVITY} on general relativity theory which does not allow for such a separation since according to it gravitation is a purely geometrical effect.


\subsection{SPECIAL-RELATIVITY}
%------------------------------
\label{SPECIAL-RELATIVITY}

Lorentz transformations and representations of the Lorentz group (rotations, boosts, and improper transformations), as well as more general rotations in Minkowskian or Euclidian four-space, including discrete transformations.

See also ELECTRODYNAMICS, Sec.~\ref{ELECTRODYNAMICS}, in particular for the seminal papers of A.W.\ Conway and L.\ Silberstein.

\begin{enumerate}

\bibitem{KLEIN1911-} F. Klein, \emph{Uber die geometrischen Grundlagen der Lorentzgruppe}, Phys. Zeitschr. {\bf 12} (1911) 17--27. %%QUATERNION, %%SPECIAL-RELATIVITY, %$D06022002.

\bibitem{SILBE1914-} L. Silberstein, The Theory of Relativity (MacMillan, 1914) 295~pp.  %%BOOK, %%QUATERNION, %%SPECIAL-RELATIVITY, %$D20032002.

%%\bibitem{CONWA1915-} A.W. Conway, Relativity (G. Bell and Son, London, 1915) %%MISS

\bibitem{CAILL1917-} C. Cailler, \emph{Sur quelques formules de la th\'eorie de la relativit\'e}, Arch. Sci. Phys. Nat. Gen\`eve {\bf 44} (1917) 237--255.  %%QUATERNION, %%SPECIAL-RELATIVITY, %%MAXWELL, %$D09022002.

\bibitem{HITCH1930-} F.L. Hitchcock, \emph{An analysis of rotations in Euclidean four-space by sedenions}, J. of Math. and Phys. {\bf 9} (1930) 188--193. %%QUATERNION, %%SPECIAL-RELATIVITY, %%SEDENION, %$D04032006.

\bibitem{SCHOU1930-} J.A. Schouten, \emph{Die Darstellung der Lorentzgruppe in der komplexen $E_2$ abgeleitet aus den Diracschen Zahlen}, Proc. Royal Acad. Amsterdam {\bf 38} (1930) 189--197. %%QUATERNION, %%SPECIAL-RELATIVITY, %%DIRAC, %%CLIFFORD, %%EDDINGTON, %$D30112005.

\bibitem{ROSEN1930-} N. Rosen, \emph{Note on the general Lorentz transformation}, J. of Math. and Phys. {\bf 9} (1930) 181--187. %%QUATERNION, %%SPECIAL-RELATIVITY, %$D04032006.

\bibitem{SOHON1930-} F.W. Sohon, \emph{Rotation and perversion groups in Euclidean space of four dimensions}, J. of Math. and Phys. {\bf 9} (1930) 194--260. %%QUATERNION, %%SPECIAL-RELATIVITY, %%GROUP-THEORY, %$D04032006.

\bibitem{GUTH-1933B} E. Guth, \emph{Einfache Ableitung der Darstellung der orthogonalen Transformationen in drei und vier reellen Ver\"anderlichen durch Quaternionen}, Anz. Akad. Wiss. Wien {\bf 70} (1933) 207--210. %%QUATERNION, %%SPECIAL-RELATIVITY,  %$D30112005.

\bibitem{MERCI1934-} A. Mercier, \emph{Application des nombres de Clifford \`a l'\'etablissement du th\'eor\`eme de relativit\'e de Lorentz}, Helv. Phys. Acta {\bf 7} (1934) 649--650.  %%QUATERNION, %%SPECIAL-RELATIVITY, %%MAXWELL, %%CLIFFORD, %$D09022002.

\bibitem{JUVET1935-} G. Juvet, \emph{Les rotations de l'espace Euclidien \`a quatre dimensions, leur expression au moyen des nombres de Clifford et leurs relations avec la th\'eorie des spineurs}, Comm. Math. Helv. {\bf 8} (1935/1936) 264--304.  %%QUATERNION, %%SPECIAL-RELATIVITY, %%SPINOR, %%SEMIVECTOR, %%CLIFFORD, %$D09022002.

\bibitem{SOMME1936-} A. Sommerfeld, \emph{Uber die Klein'schen Parameter $\alpha$, $\beta$, $\gamma$, $\delta$ und ihre Bedeutung f\"ur die Dirac-Theorie}, Sitz. Akad. Wissensch. Wien, IIa {\bf 145} (1936) 639--650. Reproduced in: A. Sommerfeld, Gesammelte Schrieften, Band IV (Friedr. Vieweg, Braunschweig, 1968). %%QUATERNION, %%SPECIAL-RELATIVITY, %%SPINOR, %%DIRAC, %$D05022002. %$C Null-sphere, two-valuedness, reflection, etc.

\bibitem{RAO--1938S} H.S.S. Rao, \emph{Eulerian parameters and Lorentz transformations}, Proc. Indian Acad. Sci.  {\bf 7} (1938) 339-342.  %%QUATERNION, %%SPECIAL-RELATIVITY, %$D06022002.

\bibitem{FISCH1940-} O.F. Fischer, \emph{Lorentz transformation and Hamilton's quaternions}, Phil. Mag. {\bf 30} (1940) 135--150. %%QUATERNION, %%SPECIAL-RELATIVITY, %$D06022002.

\bibitem{DIRAC1945-} P.A.M. Dirac, \emph{Application of quaternions to Lorentz transformations}, Proc. Roy. Irish Acad. {\bf A 50} (1945) 261--270. %%QUATERNION, %%SPECIAL-RELATIVITY, %%CONFORMALITY, %$D06022002.

\bibitem{COXET1946-} H.S.M. Coxeter, \emph{Quaternions and reflections}, Amer. Math. Monthly {\bf 53} (1946) 136--146. %%QUATERNION, %%SPECIAL-RELATIVITY, %%DISCRETE-SYMMETRY, %%REFLECTION, %$D06022002.

\bibitem{GURSE1956A} F. G\"ursey, \emph{Contribution to the quaternion formalism in special relativity}, Rev. Fac. Sci. Istanbul {\bf A 20} (1956) 149--171.  %%QUATERNION, %%SPECIAL-RELATIVITY, %$D06022002.

\bibitem{CONWA1947-} A.W. Conway, \emph{Applications of quaternions to rotations in hyperbolic space of four dimensions}, Proc. Roy. Soc. {\bf A 191} (1947) 137--145. %%QUATERNION, %%SPECIAL-RELATIVITY, %%SPIN-3/2, %$D13032002.

\bibitem{LANCZ1949-} C. Lanczos, The Variational Principles of Mechanics (Dover, New-York, 1949, 1986) 418~pp.  Quaternions pages 303--314. %%BOOK, %%QUATERNION, %%SPECIAL-RELATIVITY, %$D06022002.

\bibitem{MACDU1949-} C.C. MacDuffee, \emph{Orthogonal matrices in four-space}, Can. J. Math. {\bf 1} (1949) 69--72. %%QUATERNION, %%SPECIAL-RELATIVITY, %$D06022002.

\bibitem{LAMBE1950-} J. Lambek, \emph{Biquaternion vector fields over Minkowski space}, Thesis (McGill University, 1950). %%QUATERNION, %%SPECIAL-RELATIVITY,  %$D06022002. %$C Page 1: ``With the theory of relativity also came the fourth dimension. (However) Minkowski himself rejected the quaternion calculus as "too narrow and clumsy for the purpose in question".''  Page 1, again:  ``Professor Dirac believes (as stated by him in conversation) that, some day, Hamiltonian quaternions, as opposed to Hermitian quaternions, will re-assert themselves in relativity theory; but I do not see how this can be.''

\bibitem{SCHRE1950-} E.J. Schremp, \emph{On the interpretation of the parameters of the proper Lorentz group}, Proceedings of the 1950 International Congress of Mathematicians (Cambridge, Massachusetts, 1950) Vol.~I, p. 654--655. %%QUATERNION, %%SPECIAL-RELATIVITY, %$D24062002.

\bibitem{WAGNE1951-} H. Wagner, \emph{Zur mathematischen Behandlung von Spiegelungen}, Optik {\bf 8} (1951) 456--472.  %%QUATERNION, %%SPECIAL-RELATIVITY, %%DISCRETE-SYMMETRY, %%REFLECTION, %$D06022002.

\bibitem{SCHRE1952-} E.J. Schremp, \emph{On the geometry of the group-space of the proper Lorentz group}, Phys. Rev. {\bf 85} (1952) 721. %%QUATERNION, %%SPECIAL-RELATIVITY, %$D24062002.

\bibitem{WALKE1956-} M.J. Walker, \emph{Quaternions as 4-Vectors}, Am. J. Phys. {\bf 24} (1956) 515--522. %%QUATERNION, %%SPECIAL-RELATIVITY, %%POLARIZATION, %%STOKES-PARAMETERS, %$D23022002.

\bibitem{JOST-1957-} R. Jost, \emph{Eine Bemerkung zum CTP-Theorem}, Helv. Phys. Acta {\bf 30} (1957) 409--416.   %%QUATERNION, %%SPECIAL-RELATIVITY, %%QUANTUM-PHYSICS, %$D06072008.

\bibitem{JORDA1961-} P. Jordan, \emph{Uber die Darstellung der Lorentzgruppe mit Quaternionen}, in: Werner Heisenberg und die Physik unsere Zeit (Braunschweig, 1961) 84--89. %%QUATERNION, %%SPECIAL-RELATIVITY, %$D06022002, %$C-Hexavectors form a Lie algebra.

\bibitem{MACFA1962-} A.J. Macfarlane, \emph{On the restricted Lorentz group and groups homomorphically related to it},  J. Math. Phys. {\bf 3} (1962) 1116--1129. %%QUATERNION, %%SPECIAL-RELATIVITY, %%EINSTEIN-MAYER, %$D06022002, %$C p.116: "The Pauli matrices do not provide an adequate basis for the discussion of two-dimensional unimodular matrices".

\bibitem{BARUT1965-} A.O. Barut, \emph{Analyticity, complex and quaternionic Lorentz groups and internal quantum numbers}, in: W.E. Brittin and A.O. Barut, eds., Lect. in Th. Phys. {\bf 7A}, {Lorentz Group} (University of Colorado, Boulder, 1965) 121--131. %%QUATERNION, %%SPECIAL-RELATIVITY, %%ANALYTICITY, %%PARTICLE-PHYSICS, %$D24062002.

\bibitem{EHLER1966-} J. Ehlers, W. Rindler, and I. Robinson, \emph{Quaternions, bivectors, and the Lorentz group}, in: B. Hoffmann, ed., Perspectives in Geometry and Relativity (Indiana University Press, Bloomington, 1966) 134--149. %%QUATERNION, %%SPECIAL-RELATIVITY, %$D06022002.

\bibitem{KYRAL1967C} A. Kyrala, \emph{Four-dimensional vector analysis}, Chapter 8 of Theoretical Physics: Applications of Vectors, Matrices, Tensors and Quaternions (W.B. Saunders, Philadelphia, 1967) 374~pp.  %%QUATERNION, %%SPECIAL-RELATIVITY, %$D15032002.

\bibitem{SACHS1971-} M. Sachs, \emph{A resolution of the clock paradox}, Phys. Today (September 1971) 23--29.  %%QUATERNION, %%SPECIAL-RELATIVITY, %$09022002.

\bibitem{NEWCO1972-} W.A. Newcomb, \emph{Quaternionic clocks and odometers}, Lawrence Livermore Laboratory report UCRL-74016 (5 July 1972) 25~pp. %%QUATERNION, %%SPECIAL-RELATIVITY, %$D09022002.

\bibitem{SYNGE1972-} J.L. Synge, \emph{Quaternions, Lorentz transformations, and the Conway-Dirac-Eddington matrices}, Communications of the Dublin Institute for Advanced Studies {\bf A 21} (1972) 67~pp.  %%BOOK, %%QUATERNION, %%SPECIAL-RELATIVITY, %%REFLECTION, %%TURN, %%LANCZOS, %$D06022002.

\bibitem{MIGNA1975-} R. Mignani, \emph{Quaternionic form of superluminal Lorentz transformations}, Nuov. Cim. Lett. {\bf 13} (1975) 134--138. %%QUATERNION, %%SPECIAL-RELATIVITY, %$D06022002.

\bibitem{IMAED1976B} K. Imaeda, \emph{On ``quaternionic form of superluminal transformations,''} Nuov. Cim. Lett. {\bf 15} (1976) 91--92. %%QUATERNION, %%SPECIAL-RELATIVITY, %$D06022002.

\bibitem{HARTU1979-} R.W. Hartung, \emph{Pauli principle in Euclidean geometry}, Am. J. Phys. {\bf 47} (1979) 900--910. %%QUATERNION, %%SPECIAL-RELATIVITY, %%TURN, %%REFLECTION, %%DIRAC, %%PAULI, %$D06022002.

\bibitem{IMAED1979-} K. Imaeda, \emph{Quaternionic formulation of tachyons, superluminal transformations and a complex space-time}, Nuov. Cim. {\bf 50 B} (1979) 271--293. %%QUATERNION, %%SPECIAL-RELATIVITY, %$D06022002.

\bibitem{BAYLI1980-}  W.E. Baylis, \emph{Special relativity with $2 \times 2$ matrices}, Am. J. Phys. {\bf 48} (1980) 918--925. %%QUATERNION, %%SPECIAL-RELATIVITY, %%MAXWELL, %%DIRAC, %$D16032002.

\bibitem{TELI-1980-} M.T. Teli, \emph{Quaternionic form of unified Lorentz transformations}, Phys. Lett. {\bf A75} (1980) 460--462. %%QUATERNION, %%SPECIAL-RELATIVITY, %$D06022002.

\bibitem{BIEDE1981A} L.C. Biedenharn and J.D. Louck, \emph{The theory of turns adapted from Hamilton}, in: G.-C. Rota, ed., Encyclopedia of Mathematics and its Applications (Addison-Wesley, Reading, 1981) Vol.~8, Chap.~4, 180--204. %%QUATERNION, %%SPECIAL-RELATIVITY, %%TURN, %$D19032002.

\bibitem{RAO--1981-} K.N.S. Rao, D. Saroja, and A.V.G. Rao, \emph{On rotations in a pseudo-Euclidian space and proper Lorentz transformations}, J. Math. Phys.  {\bf 22} (1981) 2167--2179. %%QUATERNION, %%SPECIAL-RELATIVITY, %$D23022002.

\bibitem{RAO--1983-} K.N.S. Rao, A.V.G. Rao, and B.S. Narhari, \emph{On the quaternion representation of the proper Lorentz group SO(3,1)}, J. Math. Phys.  {\bf 24} (1983) 1945--1954. %%QUATERNION, %%SPECIAL-RELATIVITY, %$D06022002.

\bibitem{VANWY--1984-} C.B. vanWyk, \emph{Rotation associated with the product of two Lorentz transformations}, Am. J. Phys.  {\bf 52} (1984) 853--854. %%QUATERNION, %%SPECIAL-RELATIVITY, %$D06022002.

\bibitem{VANWY--1986A} C.B. vanWyk, \emph{General Lorentz transformations and applications}, J. Math. Phys.  {\bf 27} (1986) 1306--1310. %%QUATERNION, %%SPECIAL-RELATIVITY, %$D06022002.

\bibitem{VANWY--1986B} C.B. vanWyk, \emph{Lorentz transformations in terms of initial and final vectors}, J. Math. Phys.  {\bf 27} (1986) 1311--1314. %%QUATERNION, %%SPECIAL-RELATIVITY, %$D06022002.

\bibitem{BAYLI1988-}  W.E. Baylis and G. Jones, \emph{Special relativity with Clifford algebras and $2 \times 2$ matrices, and the exact product of two boosts}, J. Math. Phys. {\bf 29} (1988) 57--62. %%QUATERNION, %%SPECIAL-RELATIVITY, %$D16032002, %$C Critique of Salingaros.

\bibitem{BAYLI1989A}  W.E. Baylis, \emph{The Pauli-algebra approach to special relativity}, Nucl. Phys. B (Proc. Suppl.) {\bf 6} (1989) 129--131. %%QUATERNION, %%SPECIAL-RELATIVITY, %$D16032002.

\bibitem{BAYLI1989B}  W.E. Baylis, \emph{The Pauli-algebra approach to special relativity}, J. Phys. A: Math. Gen. {\bf 22} (1989) 1--15. %%QUATERNION, %%SPECIAL-RELATIVITY, %$D16032002.

\bibitem{SHARM1989A} C.S. Sharma, \emph{Representations of the general Lorentz group by $2\times 2$ matrices}, Nuov. Cim. {\bf B 103} (1989) 431--434. %%QUATERNION, %%SPECIAL-RELATIVITY, %$D06022002.

\bibitem{ABONY1991-} I. Abonyi, J.F. Bito, and J.K. Tar, \emph{A quaternion representation of the Lorentz group for classical physical applications}, J. Phys. {\bf A 24} (1991) 3245--3254. %%QUATERNION, %%SPECIAL-RELATIVITY, %$D06022002.  

\bibitem{ZENI-1992-} J.R. Zeni and W.A. Rodrigues, Jr., \emph{A thoughtful study of Lorentz transformations by Clifford algebras}, Int. J. Mod. Phys. {\bf 7} (1992) 1793--1817. %%QUATERNION, %%SPECIAL-RELATIVITY, %$D06022002.

\bibitem{LOUCK1993-} J.D. Louck, \emph{From the rotation group to the Poincar\'e group}, in: B. Gruber, ed., Symmetries in Science {\bf VI} (Plenum, New-York, 1993) 455--468. %%QUATERNION, %%SPECIAL-RELATIVITY, %$D05022002.

\bibitem{MANOG1993-} C.A. Manogue and J. Schray, \emph{Finite Lorentz transformations, automorphisms, and division algebras}, J. Math. Phys. {\bf 34} (1993) 3746--3767. %%QUATERNION, %%SPECIAL-RELATIVITY, %$D06022002.

\bibitem{REED-1993-} I.S. Reed, \emph{Generalized de Moivre's theorem, quaternions, and Lorentz transformations on a Minkowski space}, Linear Alg. and its Appl. {\bf 191} (1993) 15--40. %%QUATERNION, %%SPECIAL-RELATIVITY, %$D06022002.

\bibitem{REUSE1993-}  F. Reuse and J. Keller, \emph{Construction of a faithful vector representation of the Newtonian description of space-time and the Galilei group}, Adv. Appl. Clifford Alg. {\bf 3} (1993) 55--74. %%QUATERNION, %%SPECIAL-RELATIVITY, %%GALILEI, %%CLIFFORD, %%GRASSMANN, %$D24062002.

\bibitem{DELEO1996A} S. DeLeo, \emph{Quaternions and special relativity}, J. Math. Phys. {\bf 37} (1996) 2955--2968. %%QUATERNION, %%SPECIAL-RELATIVITY, %$D06022002.

\bibitem{YEFRE1996A} A.P. Yefremov, \emph{Quaternionic relativity. I. Inertial motion}, Gravit. \& Cosmology {\bf 2} (1996) 77--83. %%QUATERNION, %%SPECIAL-RELATIVITY, %%MECHANICS, %$D09022002.

\bibitem{YEFRE1996B} A.P. Yefremov, \emph{Quaternionic relativity. II. Non-inertial motion}, Gravit. \& Cosmology {\bf 2} (1996) 335--341. %%QUATERNION, %%SPECIAL-RELATIVITY, %%MECHANICS, %$D09022002.

\bibitem{CONTE1997-} E. Conte, \emph{On the generalization of the physical laws by biquaternions: an application to the generalization of Minkowski space-time}, Physics Essays {\bf 10} (1997) 437--441.  %%QUATERNION, %%SPECIAL-RELATIVITY, %$D24062002.

\bibitem{KRUGE1997-}  H. Kr\"uger, \emph{The electron as a self-interacting point charge. Classification of lightlike curves in spacetime under the group of SO(1,3) motions}, Adv. Appl. Clifford Alg. {\bf 7 (S)} (1997) 145--162. %%QUATERNION, %%SPECIAL-RELATIVITY, %$D17032002.

\bibitem{RINDL1999-} W. Rindler and I. Robinson, \emph{A plain man's guide to bivectors, biquaternions, and the algebra and geometry of Lorentz transformations}, in: A. Harvey, ed., On Einstein's Path --- Essays in Honor of Engelbert Schucking (Springer, New York, 1999) 407--433. %%QUATERNION, %%SPECIAL-RELATIVITY, %%BIVECTORS,  %$D13052001.

\bibitem{PEZZA2000-} W.M. Pezzaglia, Jr., \emph{Dimensionally democratic calculus of polydimensional physics}, in: R. Ablamowicz and B. Fauser, eds., Clifford Algebra and their Applications in Mathematical Physics, Vol.~1: \emph{Algebra and Physics} (Birkh\"auser, Boston, 2000)  101--123.  %%QUATERNION, %%SPECIAL-RELATIVITY, %%ANALYTICAL-MECHANICS, %%QUANTIZATION, %%CLIFFORD, %$D12052001.

\bibitem{PIAZZ2000-} F.M. Piazzese, \emph{A Pythagorean metric in relativity}, in: R. Ablamowicz and B. Fauser, eds., Clifford Algebra and their Applications in Mathematical Physics, Vol.~1: \emph{Algebra and Physics} (Birkh\"auser, Boston, 2000)  126--133. %%QUATERNION, %%SPECIAL-RELATIVITY, %%ANALYTICAL-MECHANICS, %$D12052001.

\bibitem{DELEO2001A} S. DeLeo, \emph{Quaternionic Lorentz group and Dirac equation}, Found. Phys. Lett. {\bf 14} (2001) 37--50. %%QUATERNION, %%SPECIAL-RELATIVITY, %%DIRAC, %%GSPONER, %$D06022002.

\bibitem{SILVA2002-} C.C. Silva and R. de Andrade Martins, \emph{Polar and axial vectors versus quaternions} Am. J. Phys. {\bf 70} (2002) 958--963.  %%QUATERNION, %%SPECIAL-RELATIVITY, %%DISCRETE-SYMMETRY, %%HISTORY, %$D30112005.

\bibitem{GIRAR2004-} P.R. Girard, Quaternions, Alg\`ebre de Clifford et Physique Relativiste (Presses Polytechniques et Universitaires Romandes, Lausanne, 2004) 165~~pp. %%QUATERNION, %%SPECIAL-RELATIVITY, %%GENERAL-RELATIVITY, %%CLIFFORD, %%BOOK, %$D30112005.

\bibitem{CASTR2006-} C. Castro and M. Pavsic, \emph{The extended relativity theory in Clifford spaces: reply to a review by W.A. Rodrigues, Jr.}, Prog. in Phys. {\bf 3} (2006) 27--29. %%QUATERNION, %%SPECIAL-RELATIVITY, %$D06072008.

\bibitem{MAJER2006-} V. Majernick, \emph{Quaternion formulation of the Galilean space-time transformation}, Acta. Phys. Slovaca {\bf 56} (2006) 9--14. %%QUATERNION, %%SPECIAL-RELATIVITY, %$D06072008.

\bibitem{CHRIS2007-} V. Christianto and F. Smarandache, \emph{Reply to ``Notes on Pioneer anomaly explanation by satellite-shift formula of quaternion relativity''}, Prog. in Phys. {\bf 3} (2007) 24--26. %%QUATERNION, %%SPECIAL-RELATIVITY, %$D06072008.

\bibitem{GIRAR2007-}  P.R. Girard, Quaternions, Clifford Algebras and Relativistic Physics (Birkhauser, Basel, 2007) 179~pp. %%QUATERNION, %%SPECIAL-RELATIVITY, %$D06072008. 

\bibitem{SAA--2007-} D. Saa, \emph{Fourvector algebra} (2007) 24~pp.; e-print \underline{ arXiv:0711.3220 }. %%QUATERNION, %%SPECIAL-RELATIVITY, %$D06072008.  

\end{enumerate}

\subsection{CONFORMALITY}
%------------------------
\label{CONFORMALITY}

Mostly mathematical papers dealing with conformal-transformations and conformal-invariance in hyperbolic four-space.  

See also TWISTOR, Sec.~\ref{TWISTOR}.

\begin{enumerate}


\bibitem{FUETE1931-} R. Fueter, \emph{Ueber automorph Funktionen der Picard'schen Gruppe I}, Comm. Math. Helv. {\bf 3} (1931) 42--68. %%QUATERNION, %%CONFORMALITY, %$D24062002.

\bibitem{FUETE1932B} R. Fueter, \emph{Formes d'Hermite, groupe de Picard et th\'eorie des ideaux de quaternions}, C. R. Acad. Sci. Paris. {\bf 194} (1932) 2009--2011.  %%QUATERNION, %%CONFORMALITY, %$D24062002.

\bibitem{WACHS1936-} S. Wachs, \emph{Essai sur la g\'eom\'etrie projective quaternionienne}, M\'e\-moires de l'Acad. Royale de Belgique --- classe des sciences {\bf 15} (1936) 134~pp. %%QUATERNION, %%CONFORMALITY, %$D13032002.

\bibitem{HAEFE1944-} H. Haefeli, \emph{Quaternionengeometrie und das Abbildungsproblem der regul\"aren Quaternionenfunktionen}, Comm. Math. Helv. {\bf 17} (1944) 135--164. %%QUATERNION, %%CONFORMALITY, %$D13032002.

\bibitem{GORML1947-} P.G. Gormley, \emph{Stereographic projection and the linear fractional group transformations of quaternions}, Proc. Roy. Irish Acad. {\bf A 51} (1947) 67--85. %%QUATERNION, %%CONFORMALITY, %%DIRAC, %$D13032002.

\bibitem{GURSE1956E} F. G\"ursey, \emph{On some conform invariant world-lines}, Rev. Fac. Sci. Univ. Istanbul  {\bf A 21} (1956) 129--142.  %%QUATERNION, %%CONFORMALITY, %%ELECTRODYNAMICS, %$D04032006, %$C Conformal invariant generalization of the Lorentz-Dirac equation. G\"ursey credits Synge (1937) and Eisenhart (1924) for the introduction of the relativistic Frenet formula.

\bibitem{HEPNE1962-} W.A. Hepner, \emph{The inhomogeneous Lorentz group and the conformal group}, Nuovo. Cimento {\bf 26} (1962) 351--367. %%QUATERNION, %%CONFORMALITY, %%LEPTODYNAMICS, %%DIRAC, %%TWISTOR, %$D30112005. %$C Justified the Feynman-Gell-Mann two-component interaction \cite{FEYNM1958-}.  Link with twistors according to \cite{PENRO1973-}.

\bibitem{BAKKE1983-} K.E. Bakkesa-Swamy and M. Nagaraj, \emph{Conformality, differentiability, and regularity of quaternion functions}, J. Indian Math. Soc. {\bf 47} (1983) 21--30. %%QUATERNION, %%CONFORMALITY, %%ANALYTICITY-H, %$D30112005.

\bibitem{ONDER1985-} T. Onder, \emph{Non-existence of almost-quaternion substructures on the complex projective space}, Can. Math. Bull. {\bf 28} (1985) 231--232. %%QUATERNION, %%CONFORMALITY, %%MANIFOLD, %%PROJECTIVE-GEOMETRY, %$D30112005.

\bibitem{HESTE1988-} D. Hestenes, \emph{Universal Geometric Algebra}, SIMON STEVIN, A Quarterly Journal of Pure and Applied Mathematics {\bf 62} (1988) 15~pp. %%QUATERNION, %%CONFORMALITY, %%CLIFFORD, %%HESTENES, %$D10022002.

\bibitem{ROBIN1991-} D.C. Robinson, \emph{Four-dimensional conformal and quaternionic structures}, J. Math. Phys. {\bf 32} (1991) 1259--1262. %%QUATERNION, %%CONFORMALITY, %$D30112005.

\bibitem{WILKE1993-} J.B. Wilker, \emph{The quaternion formalism for M\"obius groups in four or fewer dimensions}, Linear Algebra and its Appl. {\bf 190} (1993) 99--136. %%QUATERNION, %%CONFORMALITY, %%SPECIAL-RELATIVITY, %$D13032002.

\bibitem{RYAN-1994A} J. Ryan, \emph{Some applications of conformal covariance in Clifford analysis}, Chap.~4 in: J.Ryan, ed., Clifford Algebras in Analysis and Related Topics, Studies in Adv. Math (CRC Press Boca Raton, 1994) 129--156. %%QUATERNION, %%CONFORMALITY, %%CLIFFORD, %$D30112005.

\bibitem{RYAN-1994B} J. Ryan, \emph{The Fourier transform on the sphere}, in: G. Gentili et al., Proc. of the Meeting on Quaternionic Structures in Mathematics and Physics (SISSA, Trieste, 1994) 247--258.  %%QUATERNION, %%CONFORMALITY, %%ANALYTICITY-VARIA, %%MANIFOLD, %$D04032006.

\bibitem{WADA-2000-} M. Wada and O. Kobayashi, \emph{The Schwarzian and M\"obius transformations in higher dimensions}, in: J. Ryan and W. Spr\"ossig, eds., Clifford Algebra and their Applications in Mathematical Physics, Vol.~2: \emph{Clifford Analysis} (Birkh\"auser, Boston, 2000) 239--246. %%QUATERNION, %%CONFORMALITY, %%MOBIUS-TRANSFORMATION, %%SCHWARZIAN-DERIVATIVE, %%CLIFFORD, %$D12052001.

\bibitem{POZO-2001-} J.M. Pozo and G. Sobczyk, \emph{Realizations of the conformal group}, in: E.B. Corrochano and G. Sobczyk, eds., Geometric Algebra with Applications in Science and Engineering (Birkhauser, Boston, 2001) 43--59.  %%QUATERNION, %%CONFORMALITY, %$D13032002.

\bibitem{IVANO2007-} S. Ivanov and D. Vassilev, \emph{Conformal quaternionic contact curvature and the local sphere theorem} (2007) 30~pp.; e-print \underline{ arXiv:0707.1289 }.  %%QUATERNION, %%CONFORMALITY, %$D06072008.


\end{enumerate}


\subsection{TENSOR}
%------------------
\label{TENSOR}

Relations between quaternions and tensors.

\begin{enumerate}


\bibitem{JOHNS1926-} W.J. Johnston, \emph{A quaternion substitute for the theory of tensors}, Proc. Roy. Irish Acad. {\bf A 37} (1926) 13--27.  %%QUATERNION, %%TENSOR, %$D09022002.

\bibitem{RUSE-1936-} H.S. Ruse, \emph{On the geometry of Dirac's equations and their expression in tensor form}, Proc. Roy. Soc. Edinburgh {\bf 57} (1936/1937) 97--127.  %%QUATERNION, %%TENSOR, %%DIRAC, %$D09022002.

\bibitem{WHITT1937-} E.T. Whittaker, \emph{On the relations of the tensor-calculus to the spinor-calculus}, Proc. Roy. Soc. {\bf A 158} (1937) 38--46.  %%QUATERNION, %%TENSOR, %%DIRAC, %$D09022002.

\bibitem{KILMI1955-} C.W. Kilmister, \emph{The application of certain linear quaternion functions of quaternions to tensor analysis}, Proc. Roy. Irish Acad. {\bf A 57} (1955) 37--52.  %%QUATERNION, %%TENSOR, %%SPECIAL-RELATIVITY, %%LANCZOS, %$D06022002.

\end{enumerate}


\subsection{SPINOR}
%------------------
\label{SPINOR}

Relations between quaternions and spinors.

\begin{enumerate}


\bibitem{VEBLE1933A} O. Veblen, \emph{Geometry of two-component spinors}, Proc. Natl. Acad. Sci. {\bf 19} (1933) 462--474. %%QUATERNION, %%SPINOR, %%DIRAC, %$D04032002.

\bibitem{VEBLE1933B} O. Veblen, \emph{Geometry of four-component spinors}, Proc. Natl. Acad. Sci. {\bf 19} (1933) 503--517. %%QUATERNION, %%SPINOR, %%DIRAC, %$D04032002.

\bibitem{VEBLE1933C} O. Veblen, \emph{Spinors in projective relativity}, Proc. Natl. Acad. Sci. {\bf 19} (1933) 979--999. %%QUATERNION, %%SPINOR, %%DIRAC, %$D04032002.

\bibitem{MERCI1941-} A. Mercier, \emph{Beziehungen zwischen des Clifford'schen Zahlen und den Spinoren}, Helv. Phys. Acta {\bf 14} (1941) 565--573. %%QUATERNION, %%SPINOR, %%CLIFFORD, %$D09022002.

\bibitem{PAYNE1952-} W.T. Payne, \emph{Elementary spinor theory}, Am. J. Phys. {\bf 20} (1952) 253--262. %%QUATERNION, %%SPINOR, %$D17022002.

\bibitem{PAYNE1955-} W.T. Payne, \emph{Spinor theory and relativity I}, Am. J. Phys. {\bf 23} (1955) 526--536. %%QUATERNION, %%SPINOR, %$D23022002.

\bibitem{GURSE1956B} F. G\"ursey, \emph{Correspondence between quaternions and four-spinors}, Rev. Fac. Sci. Istanbul {\bf A 21} (1956) 33--54.  %%QUATERNION, %%SPINOR, %%DIRAC, %%LANCZOS, %$D09022002, %$C-G\"ursey claims priority over Lanczos.

\bibitem{RIESZ1958-} Marcel Riesz, \emph{Clifford numbers and spinors}, Lect. Series No 38 (Inst. for Fluid Dynamics and Appl. Math, Univ. Maryland, 1958). Reprinted in: M. Riesz, E.F. Bolinder and P. Lounesto, ed., Clifford Numbers and Spinors (Kluwer, Dordrecht, 1993) 245~pp.  %%BOOK, %%QUATERNION, %%SPINORS, %%CLIFFORD, %$D20032002. %$C Lectures given between 1957 et 1959. Reprint with commentaries and historical perspective by P. Lounesto in 1993.

\bibitem{PAYNE1959-} W.T. Payne, \emph{Spinor theory and relativity II}, Am. J. Phys. {\bf 27} (1959) 318--328. %%QUATERNION, %%SPINOR, %$D23022002.

\bibitem{PENRO1963-} R. Penrose, \emph{Null hypersurface initial data for classical electrodynamics} in: P.G. Bergmann's Aeronautical Res. Lab. Tech. Documentary Rept. 63--56; Reprinted in: Gen. Relat. \& Grav. {\bf 12} (1980) 225--264. %%QUATERNION, %%SPINOR, %%GENERAL-RELATIVITY, %%MAXWELL, %$D05022002.

\bibitem{ELLIS1966-} J.R. Ellis, \emph{A spinor approach to quaternion methods in relativity} Proc. Roy. Irish Acad. {\bf A 64} (1966) 127--142. %%QUATERNION, %%SPINOR, %%LINEAR-FUNCTION, %$D09022002.

\bibitem{TEITL1966-} S. Teitler, \emph{The structure of 4-spinors}, J. Math. Phys. {\bf 7} (1966) 1730--1738. %%QUATERNION, %%SPINOR, %%CLIFFORD, %$D09022002.

\bibitem{CRUME1969-} A. Crumeyrolle, \emph{Structure spinorielles}, Ann. Inst. H. Poincar\'e {\bf A11} (1969) 19--55. %%QUATERNION, %%SPINOR, %%CLIFFORD, %$D30112005.

\bibitem{CRUME1971-} A. Crumeyrolle, \emph{Groupes de spinoralit\'e}, Ann. Inst. H. Poincar\'e {\bf A14} (1971) 309--323. %%QUATERNION, %%SPINOR, %%CLIFFORD, %$D30112005.

\bibitem{HESTE1971A} D. Hestenes, \emph{Vectors, spinors, and complex numbers in classical and quantum physics}, Am. J. Phys. {\bf 39} (1971) 1013--1027. %%QUATERNION, %%SPINOR, %%QUANTICS, %%MAXWELL, %%MECHANICS, %%HESTENES, %$D10022002.

\bibitem{DELAH1972-} P. delaHarpe, \emph{The Clifford algebra and the spinor group of a Hilbert space}, Compositio Mathematica {\bf 25} (1972) 245--261. %%QUATERNION, %%SPINOR, %%CLIFFORD, %$D30112005.

\bibitem{BOLKE1973-} E.D. Bolker, \emph{The spinor spanner}, Amer. Math. Monthly {\bf 80} (1973) 977--984. %%QUATERNION, %%SPINOR, %$D04032006, %$C NO-QUATERNION.

\bibitem{BUGAJ1979-} K. Bugajska, \emph{Spinor structure of space-time}, Int. J. Theor. Phys. {\bf 18} (1979) 77--93. %%QUATERNION, %%SPINOR, %%CLIFFORD, %%SPECIAL-RELATIVITY, %%LORENTZ, %$D09022002.

\bibitem{LOUNE1980-} P. Lounesto, \emph{Sur les id\'eaux \`a gauche des alg\`ebres de Clifford et les produits scalaires des spineurs}, Ann. Inst. Henri Poincar\'e {\bf 33} (1980) 53--61. %%QUATERNION, %%SPINOR, %%CLIFFORD, %$D09022002.

\bibitem{LOUNE1981-} P. Lounesto, \emph{Scalar products of spinors and an extension of Brauer-Wall groups}, Found. Phys. {\bf 11} (1981) 72100740. %%QUATERNION, %%SPINOR, %%CLIFFORD, %$D09022002.

\bibitem{KUGO-1983-} T. Kugo and P. Townsend, \emph{Supersymmetry and the division algebras}, Nucl. Phys. {\bf B 221} (1983) 357--380. %%QUATERNION, %%SPINOR %%GROUP-THEORY, %%SUPERSYMMETRY, %$D13032002.  
\bibitem{FIGUE1990-} V.L. Figueiredo, E. Capelas de Oliviera, and W.A. Rodrigues Jr., \emph{Covariant, algebraic, and operator spinors}, Int. J. Theor. Phys. {\bf 29} (1990) 371--395. %%QUATERNION, %%SPINOR, %%CLIFFORD, %$D09022002.

\bibitem{DELAN1992A} R. Delanghe, F. Sommen ,and V. Soucek, \emph{Clifford algebra and spinor-valued functions}, Mathematics and its Appl. {\bf 53} (Kluwer, Dordrecht, 1992) 485~pp.  %%QUATERNION, %%SPINOR, %%CLIFFORD, %$D30112005.

\bibitem{BASRI1993-} S.A. Basri and A.O. Barut, \emph{Spinors, the Dirac formalism, and correct complex conjugation}, J. Modern. Phys. {\bf A 8} (1993) 3631--3648. %%QUATERNION, %%SPINOR, %%DIRAC, %%CLIFFORD, %%EINSTEIN-MAYER, %$D30112005.

\bibitem{LOUNE1993A} P. Lounesto, \emph{Clifford algebras and Hestenes spinors}, Found. of Phys. {\bf 23} (1993) 1203--1237. %%QUATERNION, %%SPINOR, %%CLIFFORD, %$D09022002.

\bibitem{LOUNE1993C} P. Lounesto, \emph{Marcel Riesz's work on Clifford algebras}, in: M. Riesz, E.F. Bolinder and P. Lounesto, ed., Clifford Numbers and Spinors  (Kluwer, Dordrecht, 1993) 215--241. %%BOOK, %%QUATERNION, %%SPINORS, %%CLIFFORD, %$D20032002.

\bibitem{NORTO1995-} A.H. Norton, \emph{Spinors and entanglement}, The Mathematica Journal {\bf 5}, Issue 2 (1995) 24--27.  %%QUATERNION, %%SPINOR, %$D04032006, %$C "Entanglement" is used in the ordinary, not quantical, sense.

\bibitem{RODRI1996A} W.A. Rodrigues, Jr., and Q.A.G. de Souza, \emph{Dirac-Hestenes spinor fields on Riemann-Cartan manifolds}, Int. J. Th. Phys. {\bf 35} (1996) 1849--1900. %%QUATERNION, %%SPINOR, %%CLIFFORD, %$D09022002.

\bibitem{BAYLI1997-}  W.E. Baylis, \emph{Eigenspinors and electron spin}, Adv. Appl. Clifford Alg. {\bf 7 (S)} (1997) 197--213. %%QUATERNION, %%SPINOR, %$D16032002.

\bibitem{HAMIL1997-} J.J. Hamilton, \emph{Hypercomplex numbers and the prescription of spin states}, J. Math. Phys. {\bf 38} (1997) 4914--4928. %%QUATERNION, %%SPINOR, %%SPIN-3/2, %$D17052001.

\bibitem{LOUNE1997-} P. Lounesto, Clifford Algebras and Spinors (Cambridge Univ. Press, Cambridge, 1997) 306~pp. %%BOOK, %%QUATERNION, %%SPINOR, %%CLIFFORD, %$D20032002.

\bibitem{KAMBE2003-} G. Kamberov, P. Norman, F. Pedit, and U. Pinkall, \emph{Surfaces, quaternions, and spinors} (American Mathematical Society, 2003) 150~pp.  %%QUATERNION, %%SPINORS, %$D06072008.

\bibitem{ANGLE2005-} P. Angles, \emph{Structure spinorielle associ\'ee \`a un espace vectoriel quaternionien \`a droite $E$ sur $\mathbb{H}$, muni d'une forme sesquilin\'eaire $b$ non d\'eg\'en\'er\'ee $\mathbb{H}$-antihermitienne (Spin-structures over $n$-dimensional skew-Hermitian $\mathbb{H}$-spaces)}, Adv. Appl. Clifford Alg. {\bf 15} (2005) 291--316. %%QUATERNION, %%SPINOR, %$D06072008.


\end{enumerate}

 
\subsection{TWISTOR}
%-------------------
\label{TWISTOR}

Twistors may be regarded as spinors of the $O(4,2)$ group, which is two-to-one isomorphic with the full 15-parameter conformal group in Minkowski space, including the full Poincar\'e group.  As 8-dimensional points in complexified space-time, twistors are able to encode angular-momentum/spin in addition to position/translation. They were introduced by Roger Penrose as possibly more fundamental physical objects than 4-dimensional points in Minkowski space-time.

\begin{enumerate}

\bibitem{PENRO1967-} R. Penrose, \emph{Twistor algebra}, J. Math. Phys. {\bf 8} (1967) 345--366. %%QUATERNION, %%TWISTOR, %$D05022002.

\bibitem{PENRO1968-} R. Penrose, \emph{Twistor quantization and curved space-time}, Int. J. Theor. Phys. {\bf 1} (1968) 61--99. %%QUATERNION, %%TWISTOR, %$D17022002.

\bibitem{PENRO1972-} R. Penrose, \emph{Twistor theory: an approach to the quantisation of fields and space-time}, Phys. Rep. {\bf 6} (1972) 241--316.  %%BOOK, %%QUATERNION, %%TWISTOR, %$D24062002.

\bibitem{LORD-1975-} E.A. Lord, \emph{Generalized quaternion methods in conformal geometry}, Int. J. Theor. Phys. {\bf 13} (1975) 89--102. %%QUATERNION, %%TWISTOR, %$D05022002.

\bibitem{PERJE1975-} Z. Perjes, \emph{Twistor variables of relativistic mechanics}, Phys. Rev. {\bf D 11} (1975) 2031--2041.  %%QUATERNION, %%TWISTOR, %$D30112005, %$C Structure of the angular momentum twistor.

\bibitem{WARD-1977-} R.S. Ward, \emph{On self-dual gauge fields}, Phys. Lett. {\bf 61A} (1977) 81--82. %%QUATERNION, %%TWISTOR, %%GAUGE-THEORY, %$D05022002.

\bibitem{ABLAM1982-} R. Ablamowicz, Z. Oziewicz, and J. Rzewuski, \emph{Clifford algebra approach to twistors}, J. Math. Phys. {\bf 23} (1982) 231--242. %%QUATERNION, %%TWISTOR, %%CLIFFORD, %$D05022002.

\bibitem{SOUCE1982-} V. Soucek, \emph{Complex-quaternions, their connection to twistor theory}, Czech. J. Phys. {\bf B 32} (1982) 688--691. %%QUATERNION, %%TWISTOR, %%ANALYTICITY-B, %$D29052001.

\bibitem{GLAZE1984-} J.F. Glazebrook, \emph{The construction of a class of harmonic maps to quaternionic projective-space}, J. London Math Soc. {\bf 30} (1984) 151--159.   %%QUATERNION, %%TWISTOR, %%PROJECTIVE-GEOMETRY, %$D30112005. 

\bibitem{PENRO1990-} R. Penrose, \emph{Twistors, particles, strings and links}, in: D.G. Quillen et al., eds., The Interface of Mathematics and Particle Physics (Clarendon Press, Oxford, 1990) 49--58. %%QUATERNION, %%TWISTOR, %$D05022002.

\bibitem{BASTO1992-} R.J. Baston, \emph{Quaternionic complexes}, J. Geom. Phys. {\bf 8} (1992) 29--52. %%QUATERNION, %%TWISTOR, %%ANALYTICITY-H, %$D05022002.

\bibitem{CARDO1992A} G.J. Cardoso, \emph{Twistors-diagram representation of mass-scattering integral expressions for Dirac fields}, Acta Phys. Pol. B {\bf 23} (1992) 887--906. %%QUATERNION, %%TWISTOR, %$D05022002.

\bibitem{FUJIK1994-} A. Fujiki, \emph{Nagata threefold and twistor space}, in: G. Gentili et al., Proc. of the Meeting on Quaternionic Structures in Mathematics and Physics (SISSA, Trieste, 1994) 139--146.  %%QUATERNION, %%TWISTOR, %%MANIFOLD, %$D04032006.

\bibitem{MOROI1994-} A. Moroianu and U. Semmelmann, \emph{K\"ahlerian Killing spinors, complex contact structures and twistor spaces}, in: G. Gentili et al., Proc. of the Meeting on Quaternionic Structures in Mathematics and Physics (SISSA, Trieste, 1994) 197--202; This note has appeared in C. R. Acad. Sci. Paris Ser. I Math. {\bf 323} (1996) 57--61. %%QUATERNION, %%TWISTOR, %%SPINOR, %%MANIFOLD, %%KAHLER, %$D04032006.

\bibitem{KELLE1997B}  J. Keller, \emph{Spinors, twistors, screws, mexors, and the massive spinning electron}, Adv. Appl. Clifford Alg. {\bf 7 (S)} (1997) 439--455. %%QUATERNION, %%TWISTOR, %$D17032002.

\bibitem{GOVER1999-} A.R. Gover and J. Slovak, \emph{Invariant local twistor calculus for quaternionic structures and related geometries}, J. Geom. Phys. {\bf 32} (1999) 14--56.  %%QUATERNION, %%TWISTOR, %%PROJECTIVE-GEOMETRY, %$D30112005.

\bibitem{PENRO1999-} R. Penrose, \emph{Some remarks on twistor theory}, in: A. Harvey, ed., On Einstein's Path --- Essays in Honor of Engelbert Schucking (Springer, New York, 1999) 353--366. %%QUATERNION, %%TWISTOR, %%HISTORY, %%GENERAL-RELATIVITY, %%SPIN-3/2,  %$D13052001.

\bibitem{ALEKS2000-} D.V. Alekseevsky, S. Marchiafava, and M. Pontecorvo, \emph{Spectral properties of the twistor fibration of a quaternion K\"ahler manifold}, J. Math. Pures Appl. {\bf 79} (2000) 95--110.  %%QUATERNION, %%TWISTOR, %%MANIFOLD, %%KAHLER, %$D04032006.

\bibitem{BETTE2000-} A. Bette, \emph{Twistor approach to relativistic dynamics and to Dirac equation --- A review}, in: R. Ablamowicz and B. Fauser, eds., Clifford Algebra and their Applications in Mathematical Physics, Vol.~1: \emph{Algebra and Physics} (Birkh\"auser, Boston, 2000)  75--92. %%QUATERNION, %%TWISTOR, %%SPECIAL-RELATIVITY, %%QUANTIZATION, %$D12052001.

\bibitem{BETTE2001-} A. Bette, \emph{Twistor dynamics of a massless spinning particle}, Int. J. Theor. Phys. {\bf 40} (2001) 377--385. %%QUATERNION, %%TWISTOR, %%QUANTIZATION, %$D12052001.

\bibitem{LAWRY2001-} J. Lawrynowicz and O. Suzuki, \emph{An introduction to pseudotwistors basic constructions}, in: S. Marchiafava et al., eds., Proceedings of the 2nd meeting on  
``Quaternionic structures in mathematics and physics'' (World Scientific, Singapore, 2001) 241--251. %%QUATERNION, %%TWISTOR, %$D04032006.

\bibitem{NAGAT2001-} Y. Nagatomo, \emph{Generalized ADHM-construction on Wolf spaces}, in: S. Marchiafava et al., eds., Proceedings of the 2nd Meeting on Quaternionic Structures in Mathematics and Physics (World Scientific, Singapore, 2001) 285--293. %%QUATERNION, %%TWISTOR, %%MANIFOLD, %%KAHLER, %$D04032006.

\bibitem{AGNEW2003-} A.F. Agnew, \emph{The twistor structure of the biquaternionic projective point}, Adv. Appl. Clifford Alg. {\bf 13} (2003) 231--240. %%QUATERNION, %%TWISTOR, %$D06072008.

\bibitem{ESPOS2005-} G. Esposito, \emph{From spinor geometry to complex general relativity}, Int. J. Geom. Meth. Mod. Phys. {\bf 2} (2005) 675--731; e-print \underline{ arXiv:hep-th/0504089 }.  %%QUATERNION, %%TWISTOR, %%GENERAL-RELATIVITY, %%SPIN-3/2, %$D06072008.

\bibitem{KASSA2005B} V.V. Kassandrov, \emph{Twistor algebraic dynamics in complex space-time and physical meaning of hidden dimensions}, in:  M.C.Duffy et al., eds., Proc. of the Int. Conf. on the Physical Interpretation of Relativity Theory,  PIRT-05  (Bauman Univ. Press, Moscow, 2005) 42--53; e-print arXiv:gr-qc/0602064. %%QUATERNION, %%TWISTOR, %%PARTICLE-PHYSICS, %%GENERAL-RELATIVITY, %$D04032006.

\bibitem{DAROC2007-} R. da Rocha and J. Vaz Jr., \emph{Conformal structures and twistors in the paravector model of spacetime}, Int. J. Geom. Meth. Mod. Phys. {\bf 4} (2007) 547--576; e-print \underline{ arXiv:math-ph/0412074 }.  %%QUATERNION, %%TWISTOR, %$D06072008.

\end{enumerate}


\subsection{GENERAL-RELATIVITY}
%------------------------------
\label{GENERAL-RELATIVITY}

Papers on general relativity, and papers in which fields (e.g., Dirac, Maxwell) are expressed in curved space-time.

% ==> Move papers ??

\begin{enumerate}


\bibitem{LARMO1919-} J. Larmor, \emph{On generalized relativity in connection with W.J. Johnston's symbolic calculus}, Proc. Roy. Soc. {\bf A 96} (1919) 334--363. %%QUATERNION, %%GENERAL-RELATIVITY, %%MAXWELL, %%CLIFFORD, %$D09022002. %$C Refers to \cite{JOHNS1919-}.

\bibitem{DEDON1932-} Th. DeDonder et Y. Dupont, \emph{G\'en\'eralisation relativiste des equations de Dirac}, Bull. de l'Acad. Roy. de Belg. Cl. Sc.  {\bf 18} (1932) 596--602.  %%QUATERNION, %%GENERAL-RELATIVITY, %%DIRAC, %%CLIFFORD, %$D01022002.

\bibitem{DEDON1933A} Th. DeDonder et Y. Dupont, \emph{G\'en\'eralisation relativiste des equations de Dirac (2)}, Bull. de l'Acad. Roy. de Belg. Cl. Sc.  {\bf 19} (1933) 472--478.  %%QUATERNION, %%GENERAL-RELATIVITY, %%DIRAC, %%CLIFFORD, %$D01022002.

\bibitem{DEDON1933B} Th. DeDonder et Y. Dupont, \emph{G\'en\'eralisation relativiste des equations de Dirac (3)}, Bull. de l'Acad. Roy. de Belg. Cl. Sc.  {\bf 19} (1933) 593--598.  %%QUATERNION, %%GENERAL-RELATIVITY, %%DIRAC, %%CLIFFORD, %$D01022002.

\bibitem{GODEL1949-} K. G\"odel, \emph{An example of a new type of cosmological solutions of Einstein's field equations of gravitation}, Rev. Mod. Phys. {\bf 21} (1949) 447--450. %%QUATERNION, %%GENERAL-RELATIVITY, %%PSEUDOQUATERNION, %$D24062002, %$C This is the often quoted paper in which pseudoquaternions are used under the name "hyperbolic quaternions" while there is nothing new with these "G\"odel quaternions" as they are now sometimes called. 

\bibitem{INGRA1953-} R.L. Ingraham, \emph{Spinor relativity}, Nuovo Cim. {\bf 10} (1953) 27-41. %%QUATERNION, %%GENERAL-RELATIVITY, %$D07022002.

\bibitem{GURSE1957A} F. G\"ursey, \emph{General relativistic interpretation of some spinor wave equations}, Nuov. Cim. {\bf 5} (1957) 154--171. %%QUATERNION, %%GENERAL-RELATIVITY, %%SEMIVECTOR, %C$ Essentially chapters VI to VIII of G\"ursey's Ph.D. thesis, %$D07022002.

\bibitem{PENRO1960-} R. Penrose, \emph{A spinor approach to general relativity}, Ann. Phys. {\bf 10} (1960) 171--201. %%QUATERNION, %%GENERAL-RELATIVITY, %%SPINOR, %$D17022002.

\bibitem{LANCZ1962A} C. Lanczos, \emph{The splitting of the Riemann tensor}, Rev. Mod. Phys. {\bf 34} (1962) 379--389. Reprinted in: W.R. Davis et al., eds., Cornelius Lanczos collected published papers with commentaries {\bf IV} (North Ca\-rolina State University, Raleigh NC, 1998) 2-1896 to 2-1906. %%QUATERNION, %%GENERAL-RELATIVITY, %%DIRAC, %$D05022002.

\bibitem{RASTA1964-} P. Rastall, \emph{Quaternions in relativity}, Rev. Mod. Phys. {\bf 36} (1964) 820--832. %%QUATERNION, %%GENERAL-RELATIVITY, %$D07022002.

\bibitem{DORIA1973-} F.A. Doria, \emph{Equations for a spin-two field from a Dirac-like equation}, Nuovo Cim. {\bf 7} (1973) 153--154. %%QUATERNION, %%GENERAL-RELATIVITY, %%CLIFFORD, %$D07022002.

\bibitem{EDMON1973B} J.D. Edmonds, Jr., \emph{Hypermass generalization of Einstein's gravitation theory}, Int. J. Th. Phys. {\bf 7} (1973) 475--482. %%QUATERNION, %%GENERAL-RELATIVITY, %$D07022002.

\bibitem{EDMON1974D} J.D. Edmonds, Jr., \emph{Quaternion wave equations in curved space-time}, Int. J. Th. Phys. {\bf 10} (1974) 115--122. %%QUATERNION, %%GENERAL-RELATIVITY, %$D07022002.

\bibitem{EDMON1974A} J.D. Edmonds, Jr., \emph{Five- and eight-vectors extensions of relativistic quantum theory: the preferred reference frame}, Int. J. Th. Phys. {\bf 10} (1974) 273--290. %%QUATERNION, %%GENERAL-RELATIVITY, %$D07022002.

\bibitem{CRUME1975-} A. Crumeyrolle, \emph{Une th�orie de Einstein-Dirac en spin maximum 1}, Ann. Inst. H. Poincar\'e {\bf A22} (1975) 43--41. %%QUATERNION, %%GENERAL-RELATIVITY, %%CLIFFORD, %%SPIN-1, %%SPIN-1/2, %$D30112005. 

\bibitem{DORIA1975-} F.A. Doria, \emph{A Weyl-like equations for the gravitational field}, Nuovo Cim. {\bf147} (1975) 480--482. %%QUATERNION, %%GENERAL-RELATIVITY, %%CLIFFORD, %$D07022002.

\bibitem{GALLO1975-} J.W. Gallop, \emph{Outline of a classical theory of quantum physics and gravitation}, Int. J. Theor. Phys. {\bf 14} (1975) 237--275. %%QUATERNION, %%GENERAL-RELATIVITY, %$D07022002.

\bibitem{EDMON1977-} J.D. Edmonds, Jr., \emph{Generalized quaternion formulation of relativistic quantum theory in curved space}, Found. Phys. {\bf 7} (1977) 835--879. %%QUATERNION, %%GENERAL-RELATIVITY, %$D07022002.

\bibitem{DOLAN1982-} B.P. Dolan, \emph{Quaternionic metrics and SU(2) Yang-Mills}, J. Phys. A: Math. Gen. {\bf 15} (1982) 2191--2200. %%QUATERNION, %%GENERAL-RELATIVITY, %%GAUGE-THEORY, %%COSMOLOGICAL-CONSTANT, %%LANCZOS, %%GURSEY, %$D30112005. 

\bibitem{SACHS1982-} M. Sachs, General Relativity and Matter (Reidel, Dordrecht, 1982) 208~pp. %%BOOK, %%QUATERNION, %%GENERAL-RELATIVITY, %%SPINOR, %%MAXWELL, %$09022002.

\bibitem{SINGH1982-} A. Singh, \emph{On the quaternionic form of linear equations for the gravitational field}, Nuov. Cim. Lett. {\bf 33} (1982) 457--459. %%QUATERNION, %%GENERAL-RELATIVITY, %$D07022002,

\bibitem{MORIT1983-} K. Morita, \emph{Quaternionic formulation of Dirac theory in special and general relativity}, Prog. Th. Phys. {\bf 70} (1983) 1648--1665. %%QUATERNION, %%GENERAL-RELATIVITY, %%DIRAC, %%LANCZOS,  %$D07022002, %$C Errata at page 670E.

\bibitem{GURSE1984-} F. G\"ursey and H.C. Tze, \emph{Quaternion analyticity and conformally K\"ahlerian structures in Euclidian gravity}, Lett. Math. Phys. {\bf 8} (1984) 387--395.  %%QUATERNION, %%GENERAL-RELATIVITY, %%ANALYTICITY, %%FIELD-THEORY, %%CONFORMALITY, %%KAHLER, %%H, %$D11102005.

\bibitem{MANN-1984-} R.B. Mann, \emph{Q-gravity}, Nucl. Phys. {\bf 39} (1984) 481--492. %%QUATERNION, %%GENERAL-RELATIVITY, %$D07022002,

\bibitem{HESTE1986C} D. Hestenes, \emph{Curvature calculations with spacetime algebra}, Int. J. Theor. Phys. {\bf 25} (1986) 581--588. %%QUATERNION, %%GENERAL-RELATIVITY, %%HESTENES, %$D07022002.

\bibitem{HESTE1986D} D. Hestenes, \emph{Spinor approach to gravitational motion and precession}, Int. J. Theor. Phys. {\bf 25} (1986) 589--598. %%QUATERNION, %%GENERAL-RELATIVITY, %%HESTENES, %$D07022002.

\bibitem{MARQU1991B} S. Marques, \emph{A new way to interpret the Dirac equation in a non-Riemannian manifold}, Preprint CBPF-NASA/Fermilab (1991) 4~pp. %%QUATERNION, %%GENERAL-RELATIVITY, %%DIRAC, %$D19032002.

\bibitem{PEDER1993-} H. Pedersen, Y.S. Poon, and A. Swann, \emph{The Einstein-Weyl equations in complex and quaternionic geometry}, Differ. Geom. Appl. {\bf 3} (1993) 309--321. %%QUATERNION, %%GENERAL-RELATIVITY, %%MANIFOLDS, %%CONFORMALITY, %%KAHLER, %$D30112005.

\bibitem{RODRI1993A} W.A. Rodrigues, Jr., and Q.A.G. deSouza, \emph{The Clifford bundle and the nature of the gravitational field}, Found. Phys. {\bf 23} (1993) 1465--1490. %%QUATERNION, %%GENERAL-RELATIVITY, %%CLIFFORD, %$D07022002.

\bibitem{KASTL1995-} D. Kastler, \emph{The Dirac operator and gravitation}, Comm. Math. Phys. {\bf 166} (1995) 633--643. %%QUATERNION, %%GENERAL-RELATIVITY, %%NONCOMMUTATIVE-GEOMETRY, %%CLIFFORD, %%ANALYTICITY-DIRAC, %$D06072008. 

%%\bibitem{FERNA2000A} V.V. Fernandez, A.M. Moya, and W.A. Rodrigues, Jr.,  \emph{A critical review of a gauge theory of gravitation}, Extended ``mathematical review'' record of Ref.~\cite{LASEN1998-} (2000) 7~pp. %%QUATERNION, %%GENERAL-RELATIVITY, %%CLIFFORD, %$D07022002. %%INCLUDE??

\bibitem{KASSA2005A} V.V. Kassandrov, \emph{Algebrodynamics in complex space-time and the complex-quaternionic origin of Minkowski geometry}, Gravit. \& Cosmol. {\bf 11} (2005) 354--358;  e-print arXiv:gr-qc/0405046. %%QUATERNION, %%GENERAL-RELATIVITY, %%SPECIAL-RELATIVITY, %%ANALYTICITY-MAXWELL, %%PARTICLE-PHYSICS, %$D04032006.

\bibitem{ULRYC2006-} S. Ulrych, \emph{Gravitoelectromagnetism in a complex Clifford algebra}, Phys. Lett. {\bf B 633} (2006) 631--635; e-print \underline{ arXiv:gr-qc/0602018 }.   %%QUATERNION, %%GENERAL-RELATIVITY, %$D06072008.



\end{enumerate}


\section{\Huge FIELDS}
%=====================
\label{FIELDS}

This chapter contains papers related to classical (i.e., ``non-quantized'') fields which have a stronger emphasis on mathematics and theory than physics and applications. 

Papers in which fields are expressed in curved space-time are collected in Sec.~\ref{GENERAL-RELATIVITY}.

Papers dealing with specific applications are collected in Chap.~\ref{PHYSICS} or~\ref{QUANTICS}.  For example, Sec.~\ref{SPIN-1/2} in the present chapter contains theoretical papers on Dirac's equation and field, while Sec.~\ref{QUANTUM-PHYSICS} contains papers in which Dirac's equation is applied to atomic physics.

The use of biquaternions naturally implies formulations in which Minkowski's metric and Einstein's relativity are automatically implemented.  However, papers in which Galilean relativity and non-relativistic limits are considered can also be written using quaternions or biquaternions in such a way that the relations to Minkowski's metric and Einstein's relativity are hidden or lost.  For example, Pauli's equation for a non-relativistic spin~$\Oh$ field can be written as a purely real quaternion equation.  Such papers are also collected in this chapter. 



\subsection{SPIN-1 (MAXWELL, PROCA)}
%-----------------------------------
\label{SPIN-1}

This section contains papers on Maxwell's and Proca's equations and fields, as well as papers on spin~0 fields when they are discussed in conjunction with spin~1 fields.

Applications are in sections ELECTRODYNAMICS and QUANTUM-ELEC\-TRODYNAMICS, i.e., Secs.~\ref{ELECTRODYNAMICS} and~\ref{QUANTUM-ELECTRODYNAMICS}.

See also ELECTRODYNAMICS, Sec.~\ref{ELECTRODYNAMICS}, for the seminal papers of A.W.\ Conway and L.\ Silberstein.

\begin{enumerate}


\bibitem{JOHNS1919-} W.J. Johnston, \emph{A linear associative algebra suitable for electromagnetic relations and the theory of relativity}, Proc. Roy. Soc. {\bf A 96} (1919) 331--333. %%QUATERNION, %%SPIN-1, %%MAXWELL, %%CLIFFORD, %$D09022002.

\bibitem{RUMER1930-} G. Rumer, \emph{Zur Wellentheorie des Lichtquants}, Zeits. f. Physik {\bf 65} (1930) 244--252. %%QUATERNION, %%SPIN-1, %%SPIN-1/2, %%DIRAC+MAXWELL, %%LANCZOS, %$D06072008.  %$C Cites and discusses \cite{LANCZ1929B, LANCZ1929C, LANCZ1929D}

\bibitem{PROCA1936-} A. Proca, \emph{Sur la th\'eorie ondulatoire des \'electrons positifs et n\'egatifs}, J. Phys. Radium {\bf 7} (1936) 347--353.  %%NO-QUATERNION, %%SPIN-1, %%PROCA, %$D05022002.  %$C Proca derives the Lanczos-Proca equation by means of a Lagrangian formalism.

\bibitem{PROCA1938-} A. Proca, \emph{Th\'eorie non relativiste des particules \`a spin entier}, J. Phys. Radium {\bf 9} (1938) 61-66.  %%QUATERNION, %%SPIN-1, %%PROCA, %$D05022002.  %$C In this paper, page 63, Proca qualifies the quadrivector variance ${\mathcal{L}}\A{\mathcal{L}}^+ $ as `quaternionic,' and the spinor variance ${\mathcal{L}}\A$ as `spinorial.'

\bibitem{CONWA1945A} A.W. Conway, \emph{Quaternions and matrices}, Proc. Roy. Irish Acad. {\bf A 50} (1945) 98--103.  %%QUATERNION, %%SPIN-1, %%PROCA, %$D19032002.

\bibitem{CONWA1945B} A.W. Conway, \emph{Cuaternios y matrices}, Revista Union Mat. Argentina {\bf 11} (1945) 11--17.  %%QUATERNION, %%SPIN-1, %%PROCA, %$D19032002. %C Spanish translation of \cite{CONWA1945A}.

\bibitem{ELLIS1964-} J.R. Ellis, \emph{Maxwell's equations and theories of Maxwell form},  (Ph.D. thesis, University of London, 1964) 417~pp.   %%QUATERNION, %%SPIN-1, %%MAXWELL, %%LANCZOS, %$D09022002.

\bibitem{MAJER1976-} V. Majernik and M. Nagy, \emph{Quaternionic form of Maxwell's equations with sources}, Nuov. Cim. Lett. {\bf 16} (1976) 265--268.  %%QUATERNION, %%SPIN-1, %%MAXWELL, %$D09022002.

\bibitem{HONIG1977-} W.M. Honig, \emph{Quaternionic electromagnetic wave equation and a dual charge-filled space}, Nuovo. Cim. Lett. {\bf 19} (1977) 137--140. %%QUATERNION, %%SPIN-1, %%MAXWELL, %$D09022002.

\bibitem{SINGH1981A} A. Singh, \emph{Quaternionic form of the electromagnetic-current equations with magnetic monopoles}, Nuov. Cim. Lett. {\bf 31} (1981) 145--148.  %%QUATERNION, %%SPIN-1, %%MAXWELL, %%MONOPOLE, %$D09022002.

\bibitem{SINGH1981B} A. Singh, \emph{On the quaternion form of the electromagnetic-current equations}, Nuov. Cim. Lett. {\bf 31} (1981) 97--98.  %%QUATERNION, %%SPIN-1, %%MAXWELL, %$D09022002.

\bibitem{JOLLY1984-} D.C. Jolly, \emph{Isomorphic $8 \times 8$ matrix representations of quaternion field theories}, Nuovo Cim. Lett. {\bf 39} (1984) 185--188.  %%QUATERNION, %%SPIN-1, %%MAXWELL, %%GROUP-THEORY, %$D13032002.

\bibitem{DASIL1985-} A. DaSilveira, \emph{Isomorphism between matrices and quaternions}, Lett. Nuovo Cim. {\bf 44} (1985) 80--82. %%QUATERNION, %%SPIN-1, %%MAXWELL, %$D09022002.

\bibitem{KRIST1992-} S. Kristyan and J. Szamosi, \emph{Quaternionic treatment of the electromagnetic wave equation}, Acta Phys. Hungarica {\bf 72} (1992) 243--248.  %%QUATERNION, %%SPIN-1, %%MAXWELL, %$D09022002.

\bibitem{NGUYE1992-} D.B. Nguyen, \emph{Pl\"ucker's relations and the electromagnetic field}, Am. J. Phys. {\bf 60} (1992) 1145--1147. %%QUATERNION, %%SPIN-1, %%MAXWELL, %$D09022002.

\bibitem{NEGI-1998-} O.P.S. Negi, S. Bisht, and P.S. Bisht, \emph{Revisiting quaternion formulation of electromagnetism}, Nuovo Cim. {\bf 113 B} (1998) 1449--1467. %%QUATERNION, %%SPIN-1, %%MAXWELL, %$D09022002.

\bibitem{GSPON2002-} A. Gsponer, \emph{On the ``equivalence'' of the Maxwell and Dirac equations}, Int. J. Theor. Phys. {\bf 41} (2002) 689--694; e-print \underline{ arXiv:math-ph/0201053 }. %%QUATERNION, %%SPIN-1, %%MAXWELL, %%DIRAC, %%LANCZOS, %$D13052001.

\bibitem{KRAVC2002B} V. Kravchenko \emph{On the relation between the Maxwell system and the Dirac equation}, WSEAS Transactions on systems {\bf 1} (2002) 115--118; e-print \underline{ arXiv:math-ph/0202009 }. %%QUATERNION, %%SPIN-1, %%MAXWELL, %%DIRAC, %$D30112005.

\bibitem{VARLA2005-} V.V. Varlamov, \emph{Maxwell field on the Poincar\'e group}, Int. J. of Modern Physics {\bf A 20} (2005) 4095--4112.  %%QUATERNION, %%SPIN-1, %%RUMER, %%GSPONER, %$D06072008.

\end{enumerate}


\subsection{SPIN-1/2 (DIRAC, LANCZOS, PAULI, WEYL)}
%--------------------------------------------------
\label{SPIN-1/2}

Dirac's, Lanczos's, and Weyl's relativistic equations and fields, as well as non-relativistic equations and fields such as Pauli's.  

Papers such as Lanczos's which relates to fields with spin~0 to~$\Th$ but whose main emphasis is spin $\Oh$ are included in this section.

Applications are in sections QUANTUM-PHYSICS and QUANTUM-EL\-EC\-TRODYNAMICS, i.e., Secs.~\ref{QUANTUM-PHYSICS} and~\ref{QUANTUM-ELECTRODYNAMICS}. \ref{QUANTUM-ELECTRODYNAMICS}.

See also EDDINGTON and SEMIVECTOR, Secs.~\ref{EDDINGTON} and~\ref{SEMIVECTOR}.

\begin{enumerate}

\bibitem{LANCZ1929B} C. Lanczos, \emph{Die tensoranalytischen Beziehungen der Diracschen Glei\-chung [The tensor analytical relationships of Dirac's equation]}, Zeits. f. Phys. {\bf 57} (1929) 447--473. Reprinted and translated in: W.R. Davis \emph{et al.}, eds., Cornelius Lanczos Collected Published Papers With Commentaries {\bf III} (North Ca\-rolina State University, Raleigh, 1998) pages 2-1132 to 2-1185; e-print \underline{ arXiv:physics/0508002 }. %%QUATERNION, %%SPIN-1/2, %%DIRAC, %%ISOSPIN, %%PROCA, %$D05022002.

\bibitem{LANCZ1929C} C. Lanczos, \emph{Zur kovarianten Formulierung der Diracschen Glei\-chung [On the covariant formulation of Dirac's equation]}, Zeits. f. Phys. {\bf 57} (1929) 474--483. Reprinted and translated in: W.R. Davis \emph{et al.}, eds., Cornelius Lanczos Collected Published Papers With Commentaries {\bf III} (North Ca\-rolina State University, Raleigh, 1998) pages 2-1186 to 2-1205; e-print \underline{ arXiv:physics/0508012 }. %%QUATERNION, %%SPIN-1/2, %%DIRAC, %%ISOSPIN, %$D05022002.

\bibitem{LANCZ1929D} C. Lanczos, \emph{Die Erhaltungss\"atze in der feldm\"assigen Darstellungen der Diracschen Theorie [The conservation law in the field theoretical representation of Dirac's theory]}, Zeits. f. Phys. {\bf 57} (1929) 484--493. Reprinted and translated  in: W.R. Davis \emph{et al.}, eds., Cornelius Lanczos Collected Published Papers With Commentaries {\bf III} (North Ca\-rolina State University, Raleigh, 1998) pages 2-1206 to 2-1225; e-print \underline{ arXiv:physics/0508013 }. %%QUATERNION, %%SPIN-1/2, %%DIRAC, %%ISOSPIN, %%PROCA, %$D05022002.

\bibitem{SCHOU1929-} J.A. Schouten, \emph{Ueber die in der Wellengleichung verwendeten hyperkomplexen Zahlen}, Proc. Royal Acad. Amsterdam {\bf 32} (1929) 105--108. %%QUATERNION, %%SPIN-1/2, %%DIRAC, %%CLIFFORD, %%EDDINGTON, %$D11032002.

\bibitem{LANCZ1930A} C. Lanczos, \emph{Dirac's wellenmechanische Theorie des Elektrons und ihre feldtheorische Ausgestaltung [Dirac's wave mechanical theory of the electron and its field-theoretical interpretation]}, Physikalische Zeits. {\bf 31} (1930) 120--130. Reprinted and translated in: W.R. Davis \emph{et al.}, eds., Cornelius Lanczos Collected Published Papers With Commentaries {\bf III} (North Ca\-rolina State University, Raleigh, 1998) 2-1226 to 2-1247; e-print \underline{ arXiv:physics/0508009 }. %%QUATERNION, %%SPIN-1/2, %%DIRAC, %%ISOSPIN, %%PROCA, %$D05022002.

\bibitem{DEDON1930-} Th. DeDonder, \emph{Th\'eorie invariante des fonctions hypercomplexes et des matrices de Dirac g\'en\'eralis\'ees}, Bull. de l'Acad. Roy. de Belg. Cl. Sc. {\bf 16} (1930) 1092--1097. %%QUATERNION, %%SPIN-1/2, %%DIRAC, %%CLIFFORD, %$D01022002.

\bibitem{JUVET1930-} G. Juvet, \emph{Op\'erateurs de Dirac et \'equations de Maxwell}, Comm. Math. Helv. {\bf 2} (1930) 225--235. %%QUATERNION, %%SPIN-1/2, %%DIRAC, %%MAXWELL, %%LANCZOS, %$D05022002. %$C Applies "Dirac's operator" to Maxwell's field.  Defines the "maxwellian" of a tensor.

\bibitem{PROCA1930A} A. Proca, \emph{Sur l'\'equation de Dirac}, C.R. Acad. Sci. Paris {\bf 190} (1930) 1377--1379. %%QUATERNION, %%SPIN-1/2, %%DIRAC, %%CLIFFORD, %$D05022002.

\bibitem{PROCA1930B} A. Proca, \emph{Sur l'\'equation de Dirac}, C.R. Acad. Sci. Paris {\bf 191} (1930) 26--27. %%QUATERNION, %%SPIN-1/2, %%DIRAC, %%CLIFFORD, %$D05022002.

\bibitem{PROCA1930C} A. Proca, \emph{Sur l'\'equation de Dirac}, J. Phys. Radium {\bf 1} (1930) 235--248. %%QUATERNION, %%SPIN-1/2, %%DIRAC, %%CLIFFORD, %$D05022002.

\bibitem{SAUTE1930A} F. Sauter, \emph{L\"osung der Diracschen Gleichungen ohne Spezialisierung der Diracschen Operatoren}, Zeitschr. f\"ur Phys. {\bf 63} (1930) 803--814. %%QUATERNION, %%SPIN-1/2, %%DIRAC, %%LANCZOS, %$D05022002.

\bibitem{SAUTE1930B} F. Sauter, \emph{Zur L\"osung der Diracschen Gleichungen ohne Spezialisierung der Diracschen Operatoren II}, Zeitschr. f\"ur Phys. {\bf 64} (1930) 295--303. %%QUATERNION, %%SPIN-1/2, %%DIRAC, %%LANCZOS, %$D23022002.

\bibitem{LANCZ1932-} C. Lanczos, \emph{The equation of Dirac for the electron}, in: C. Lanczos, Wave Mechanics, Part II, Lecture notes (Purdue University, 1931, 1932) 340--388. %%QUATERNION, %%SPIN-1/2, %%DIRAC, %$D05022002.

\bibitem{PROCA1932A} A. Proca, \emph{Sur une explication possible de la diff\'erence de masse entre  le proton et l'\'electron}, J. Phys. Radium  {\bf 3} (1932) 83--101.  %%QUATERNION, %%SPIN-1/2, %%DIRAC, %%LANCZOS, %%CLIFFORD, %$08082005.

\bibitem{PROCA1932B} A. Proca, \emph{Quelques observations concernant un article intitul\'e ``Sur l'\'equation de Dirac,''} J. de. Phys. Radium {\bf 3} (1932) 172--184.  %%QUATERNION, %%SPIN-1/2, %%DIRAC, %%CLIFFORD, %$D05022002.

\bibitem{FRANZ1935-} W. Franz, \emph{Zur Methodik der Dirac-Gleichung}, Sitzungsber. d. Bayrischen Akad. d. Wiss. {\bf III} (1935) 379--435. %%BOOK, %%QUATERNION, %%SPIN-1/2, %%DIRAC, %%SPINORS, %%CLIFFORD, %$D20032002.

\bibitem{CONWA1936-} A.W. Conway, \emph{A quaternion view of the electron wave equation}, in: Compte Rendus Congr. Intern. des Math\'ematiciens, Oslo 1936 (Broggers Boktrykkeri, Oslo, 1937) 233. %%QUATERNION, %%SPIN-1/2, %%DIRAC, %$D05022002.

\bibitem{WATSO1947-} A.G.D. Watson, \emph{On the geometry of the wave equation}, Proc. Cambridge Phil. Soc. {\bf 43} (1947) 491--505. %%QUATERNION, %%SPIN-1/2, %%DIRAC, %$D05022002.

\bibitem{JEHLE1949-} H. Jehle, \emph{Two-component wave equations}, Phys. Rev. {\bf 75} (1949) 1609. %%QUATERNION, %%SPIN-1/2, %%DIRAC, %$D05022002.

\bibitem{KILMI1949A} C.W. Kilmister, \emph{Two-component wave equations}, Phys. Rev. {\bf 76} (1949) 568. %%QUATERNION, %%SPIN-1/2, %%DIRAC, %$D05022002.

\bibitem{SERPE1949-} J. Serpe, \emph{Two-component wave equations}, Phys. Rev. {\bf 76} (1949) 1538. %%QUATERNION, %%SPIN-1/2, %%DIRAC, %$D05022002.

\bibitem{GURSE1950A} F. G\"ursey, \emph{Applications of quaternions to field equations}, Ph.D. thesis (University of London, 1950) 204~pp.  %%BOOK, %%QUATERNION, %%SPIN-1/2, %%DIRAC, %%GENERAL-RELATIVITY, %%LANCZOS, %%PROCA, %$D20032002. %$C G\"ursey generalizes Lanczos's equation to curved space-time. Two subsections of G\"ursey's thesis specifically deal with --- and clarify some aspects --- of Lanczos's papers of 1929, i.e., \cite{LANCZ1929B, LANCZ1929C, LANCZ1929D}. For a biography of Feza G\"ursey (1921-1992) and for a list of publications until 1984, see: \emph{Symmetries in Particle Physics}, edited by I. Bars, A. Chodos and C.-H. Tze (Plenum Press, New York, 1985).

\bibitem{GURSE1950B} F. G\"ursey, \emph{On two-component wave equation}, Phys. Rev. {\bf 77} (1950) 844--845. %%QUATERNION, %%SPIN-1/2, %%DIRAC, %$D05022002.

\bibitem{SOMME1951-} A. Sommerfeld, \emph{Die Diracsche Theorie des Elektrons},  in: Atombau und Spektrallinien (Friedr. Vieweg, Braunschweig, 1951) Vol.~II, 239--341. %%BOOK, %%QUATERNION, %%SPIN-1/2, %%DIRAC, %$D05022002.

\bibitem{GURSE1955-} F. G\"ursey, \emph{Connection between Dirac's electron and a classical spinning particle}, Phys. Rev. {\bf 97} (1955) 1712--1713. %%QUATERNION, %%SPIN-1/2, %%DIRAC, %%SPECIAL-RELATIVITY, %$D05022002. %$C G\"ursey credits de Broglie for two new conserved currents, and makes the Gordon decomposition of the current for the Dirac-Lanczos equation.  Moreover, on page 46, he wrongly claims "to have first written" the Dirac-Lanczos equation!

\bibitem{GURSE1956D} F. G\"ursey, \emph{New algebraic identities and divergence equations for the Dirac electron}, Rev. Fac. Sci. Univ. Istanbul  {\bf A 21} (1956) 85--95.  %%QUATERNION, %%SPIN-1/2, %%DIRAC-LANCZOS, %$C Discusses the properties of the "de Broglie" currents, %$D04032006.

\bibitem{EBERL1962-} W.F. Eberlein, \emph{The spin model of Euclidian 3-space}, Amer. Math. Monthly {\bf 69} (1962) 587--598; errata p.960. %%QUATERNION, %%SPIN-1/2, %$D04032006.

\bibitem{TEITL1965-} S. Teitler, \emph{``Vector'' Clifford algebras and the classical theory of fields}, Nuovo Cim. Suppl. {\bf 3} (1965) 1--14. %%QUATERNION, %%SPIN-1/2, %%DIRAC, %%PROCA, %%SPINOR, %%CLIFFORD, %$D09022002.

\bibitem{HESTE1967A} D. Hestenes, \emph{Real spinor fields}, J. Math. Phys. {\bf 8} (1967) 798--808.  %%QUATERNION, %%SPIN-1/2, %%DIRAC, %%SPINOR, %%HESTENES, %$D10022002.

\bibitem{CASAN1968A} G. Casanova, \emph{Sur les th\'eories de D. Hestenes et de Dirac}, C.R. Acad. Sci. A {\bf 266} (1968) 1551--1554. %%QUATERNION, %%SPIN-1/2, %%DIRAC, %%GURSEY, %%HESTENES, %$D04032006, %$D Casanova qualifies the "Dirac-Hestenes" equation as "G�rsey's equation interpreted by Hestenes." 

\bibitem{BOUDE1971-}  R. Boudet, \emph{Sur une forme intrins\`eque de l'\'equation de Dirac et son interpr\'etation g\'eom\'etrique}, C.R. Acad. Sci. Paris. {\bf 272} (1971) 767--768.  %%QUATERNION, %%SPIN-1/2, %%DIRAC, %$D16032002.

\bibitem{EDMON1972-} J.D. Edmonds, Jr., \emph{Nature's natural numbers: relativistic quantum theory over the ring of complex quaternions}, Int. J. Th. Phys. {\bf 6} (1972) 205--224. %%QUATERNION, %%SPIN-1/2, %%DIRAC, %%LANCZOS, %$D06022002.

\bibitem{EDMON1973A} J.D. Edmonds, Jr., \emph{Generalized charge in the eight-component spin 1/2 wave equation}, Found. Phys. {\bf 3} (1973) 313--319. %%QUATERNION, %%SPIN-1/2, %%DIRAC, %$D06022002.

\bibitem{EDMON1974C} J.D. Edmonds, Jr., \emph{Quaternion quantum theory: new physics or number mysticism\,?}, Am. J. Phys. {\bf 42} (1974) 220--223. %%QUATERNION, %%SPIN-1/2, %%DIRAC, %%LANCZOS, %$D06022002.

\bibitem{BOUDE1974-}  R. Boudet, \emph{Sur le tenseur de Tetrode et l'angle de Takabayasi. Cas du potentiel central}, C.R. Acad. Sci. Paris. {\bf 278} (1974) 1063--1065.  %%QUATERNION, %%SPIN-1/2, %%DIRAC, %%TETRODE, %%HESTENES, %$D16032002.

\bibitem{EDMON1975A} J.D. Edmonds, Jr., \emph{Mass term variation in the Dirac hydrogen atom}, Int. J. Th. Phys. {\bf 13} (1975) 431--435. %%QUATERNION, %%SPIN-1/2, %%DIRAC, %$D0602200, %$ AG.

\bibitem{EDMON1975B} J.D. Edmonds, Jr., \emph{Comment on the Dirac-like equation for the photon}, Nuov. Cim. Lett. {\bf 13} (1975) 185--186.  %%QUATERNION, %%SPIN-1/2, %%DIRAC, %$D06022002.

\bibitem{EDMON1978-} J.D. Edmonds, Jr., \emph{Yet another formulation of the Dirac equation}, Found. Phys. {\bf 8} (1978) 439--444.  %%QUATERNION, %%SPIN-1/2, %%DIRAC, %$D06022002.

\bibitem{GREIDE1980-} K. Greider, \emph{Relativistic quantum theory with correct conservation laws}, Phys. Rev. Lett. {\bf 44} (1980) 1718--1721. %%QUATERNION, %%SPIN-1/2, %%DIRAC, %%CLIFFORD, %%EDDINGTON, %%HESTENES, %$D24062002.  %$C Inconsistencies of standard Dirac formalism are avoided.

\bibitem{HESTE1981-} D. Hestenes, \emph{Geometry of the Dirac Theory}, in: A Symposium on the Mathematics of Physical Space-Time (Facultad de Quimica, Universidad Nacional Autonoma de Mexico, Mexico City, Mexico, 1981) 67--96. %%QUATERNION, %%SPIN-1/2, %%DIRAC, %%QED, %%HESTENES, %$D10022002.

\bibitem{GREIDE1984-} K. Greider, \emph{A unifying Clifford formalism for relativistic fields}, Found. Phys. {\bf 14} (1984) 467--506. %%QUATERNION, %%SPIN-1/2, %%DIRAC, %%MAXWELL, %%CLIFFORD, %%EDDINGTON, %%HESTENES, %$D24062002.  %$C Inconsistencies of standard Dirac formalism are avoided.

\bibitem{BOUDE1985-}  R. Boudet, \emph{Conservation laws in the Dirac theory}, J. Math. Phys. {\bf 26} (1985) 718--724.  %%QUATERNION, %%SPIN-1/2, %%DIRAC, %%TETRODE, %%HESTENES, %$D16032002.

\bibitem{MORIT1986-} K. Morita, \emph{A role of quaternions in the Dirac theory}, Prog. Th. Phys. {\bf 75} (1986) 220--223. %%QUATERNION, %%SPIN-1/2, %%DIRAC, %%SPECIAL-RELATIVITY, %%CONFORMALITY, %$D16032002, %$C-Real quaternion Lorentz transformation like \cite{DIRAC1945-}.

\bibitem{ROTEL1989A} P. Rotelli, \emph{The Dirac equation on the quaternion field}, Mod. Phys. Lett. {\bf A 4} (1989) 933--940. %%QUATERNION, %%SPIN-1/2, %%DIRAC, %$D15032002.

\bibitem{RODRI1990-} W.A. Rodrigues, Jr., and E. Capelas de Oliviera, \emph{Dirac and Maxwell equations in the Clifford and spin-Clifford bundles}, Int. J. Theor. Phys. {\bf 29} (1990) 397--412.  %%QUATERNION, %%SPIN-1/2, %%DIRAC, %%MAXWELL, %%WRONG, %$D17032002.

\bibitem{HESTE1991A} D. Hestenes, \emph{Zitterbewegung in Radiative Processes}, in: D. Hestenes and A. Weingartshofer, eds., The Electron (Kluwer Academic Publishers, Dordrecht, 1991) 21--36. %%QUATERNION, %%SPIN-1/2, %%DIRAC, %%ZITTERBEWEGUNG, %%HESTENES, %$D10022002.

\bibitem{BAYLI1992A}  W.E. Baylis, \emph{Classical eigenspinors and the Dirac equation}, Phys. Rev. {\bf A 45} (1992) 4293--4302. %%QUATERNION, %%SPIN-1/2, %%DIRAC, %%SPINOR, %$D16032002.

\bibitem{KELLE1992-} J. Keller and A. Rodrigues, \emph{Geometric superalgebra and the Dirac equation}, J. Math. Phys. {\bf 33} (1992) 161--170.  %%QUATERNION, %%SPIN-1/2, %%DIRAC, %$D05022002.

\bibitem{PARRA1992-} J.M. Parra, \emph{On Dirac and Dirac-Darwin-Hestenes equation}, in: A. Micalli et al., eds., Clifford Algebras and their Applications in Mathematical Physics (Kluwer, Dordrecht, 1992) 463--477.  %%QUATERNION, %%SPIN-1/2, %%DIRAC, %$D05022002.

\bibitem{DAVIA1993-}  C. Daviau, \emph{Linear and nonlinear Dirac equation},  Found. Phys.  {\bf 23 } (1993) 1431--1443. %%QUATERNION, %%SPIN-1/2, %%DIRAC, %%NONLINEAR, %%HESTENES, %$D16032002.

\bibitem{HESTE1993C} D. Hestenes, \emph{The kinematic origin of complex wave functions  physics and probability}, in: W.T. Grandy and P.W. Miloni, ed., Essays in Honor of Edwin T. Jaynes (Cambridge U. Press, Cambridge, 1993) 153-160. %%QUATERNION, %%SPIN-1/2, %%DIRAC, %%HESTENES, %$D10022002.

\bibitem{HESTE1993D} D. Hestenes, \emph{Zitterbewegung modeling}, Foundations of Physics {\bf 23} (1993) 365--368. %%QUATERNION, %%SPIN-1/2, %%DIRAC, %%HESTENES, %$D10022002.

\bibitem{PAVSIC1993-} M. Pavsic, E. Recami, W.A. Rodrigues, Jr., and G. Salesi, \emph{Spin and electron structure}, Phys. Lett. B {\bf 318} (1993) 481--488. %%QUATERNION, %%SPIN-1/2, %%DIRAC, %%ZITTERBEWEGUNG, %%CLIFFORD, %$D07022002.

\bibitem{RODRI1993B} W.A. Rodrigues, Jr., J. Vaz, Jr., E. Recami, and G. Salesi, \emph{About zitterbewegung and electron structure}, Phys. Lett. B {\bf 318} (1993) 623--628. %%QUATERNION, %%SPIN-1/2, %%DIRAC, %%ZITTERBEWEGUNG, %%CLIFFORD, %$D07022002.

\bibitem{VAZ--1993A}  J. Vaz, Jr., and W.A. Rodrigues, Jr., \emph{Equivalence of Dirac and Maxwell equations and quantum mechanics} Int. J. Theor. Phys. {\bf 32} (1993) 945--959.  %%QUATERNION, %%SPIN-1/2, %%DIRAC, %%MAXWELL, %%WRONG, %$D17032002.

\bibitem{VAZ--1993B}  J. Vaz, Jr., and W.A. Rodrigues, Jr., \emph{Zitterbewegung and the electromagnetic field of the electron}, Phys. Lett. B {\bf 319} (1993) 203--208. %%QUATERNION, %%SPIN-1/2, %%DIRAC, %%ZITTERBEWEGUNG, %%CLIFFORD, %$D07022002.

\bibitem{GSPON1994A}  A. Gsponer and J.-P. Hurni,  \emph{Lanczos' equation to replace Dirac's equation?}, in: J.D. Brown et al., eds., Proceedings of the Cornelius Lanczos International Centenary Conference (SIAM Publishers, Philadelphia, 1994) 509--512; e-print \underline{ arXiv:hep-ph/0112317 }. %%QUATERNION, %%SPIN-1/2, %%DIRAC, %%LANCZOS, %$D20032002.

\bibitem{VRBIK1994B} J. Vrbik, \emph{Dirac equation and Clifford algebra}, J. Math. Phys. {\bf 35} (1994) 2309--2314.  %%QUATERNION, %%SPIN-1/2, %%DIRAC, %%CLIFFORD, %%HESTENES, %%GAUGE-THEORY, %$D30112005. %$C Critique of Hestenes. 

\bibitem{DELEO1996B} S. DeLeo, \emph{A one-component Dirac equation}, Mod. Phys. Lett. {\bf A11} (1996) 3973--3985. %%QUATERNION, %%SPIN-1/2, %%DIRAC, %$D15032002.

\bibitem{DELEO1996C} S. DeLeo and P. Rotelli, \emph{The quaternionic Dirac Lagrangian}, Mod. Phys. Lett. {\bf A11} (1996) 357--366. %%QUATERNION, %%SPIN-1/2, %%DIRAC, %$D15032002. 

\bibitem{SRINI1996-} S.K. Srinivasan and E.C.G. Sudarshan, \emph{A direct derivation of the Dirac equation via quaternion measures}, J. Phys. {\bf A 29} (1996) 5181--5186. %%QUATERNION, %%SPIN-1/2, %%DIRAC, %$D06022002.

\bibitem{CAMPO1997-}  A. Campolataro, \emph{Classical electrodynamics and relativistic quantum mechanics}, Adv. Appl. Clifford Alg. {\bf 7 (S)} (1997) 167--173. %%QUATERNION, %%SPIN-1/2, %%DIRAC, %%MAXWELL, %$D17032002.

\bibitem{PENRO1997-} R. Penrose, \emph{The mathematics of the electron's spin}, Eur. J. Phys. {\bf 18} (1997) 164--168.  %%QUATERNION, %%SPIN-1/2, %%DIRAC, %$D06022002.

\bibitem{VAZ--1997A}  J. Vaz, Jr., and W.A. Rodrigues, Jr., \emph{Maxwell and Dirac's theories as an already unified theory}, Adv. Appl. Clifford Alg. {\bf 7 (S)} (1997) 369--386. %%QUATERNION, %%SPIN-1/2, %%DIRAC, %%MAXWELL, %%WRONG, %$D17032002.

\bibitem{DELEO1998A} S. DeLeo and W.A. Rodrigues, Jr., \emph{Quaternionic electron theory: Dirac's equation}, Int. J. Th. Phys. {\bf 37} (1998) 1511--1529. %%QUATERNION, %%SPIN-1/2, %%DIRAC, %$D15032002.

\bibitem{DELEO1998B} S. DeLeo and W.A. Rodrigues, Jr., \emph{Quaternionic electron theory: geometry, algebra and Dirac's spinors}, Int. J. Th. Phys. {\bf 37} (1998) 1707--1720. %%QUATERNION, %%SPIN-1/2, %%DIRAC, %$D15032002.

\bibitem{DELEO1998C} S. DeLeo, W.A. Rodrigues, Jr., and J. Vaz, Jr., \emph{Complex geometry and Dirac equation}, Int. J. Th. Phys. {\bf 37} (1998) 12415--2431. %%QUATERNION, %%SPIN-1/2, %%DIRAC, %$D15032002.

\bibitem{GSPON1998B} A. Gsponer and J.-P. Hurni, \emph{Lanczos-Einstein-Petiau: From Dirac's equation to nonlinear wave mechanics,} in: W.R. Davis et al., eds., Cornelius Lanczos Collected Published Papers With Commentaries {\bf III} (North Carolina State University, Raleigh, 1998) 2-1248 to 2-1277; e-print \underline{ arXiv:physics/0508036 }. %%QUATERNION, %%SPIN-1/2, %%DIRAC, %%LANCZOS, %%MAXWELL, %%PROCA, %%EINSTEIN-MAYER, %%PETIAU, %$D20032002.

\bibitem{RODRI1998A} W.A. Rodrigues, Jr., and J. Vaz, Jr., \emph{From electromagnetism to relativistic quantum mechanics}, Found. Phys. {\bf 28} (1998) 789--814.  %%QUATERNION, %%SPIN-1/2, %%DIRAC, %%MAXWELL, %%MUON, %$D17032002.

\bibitem{DELEO1999B} S. DeLeo, Z. Oziewicz, W.A. Rodrigues, Jr., and J. Vaz, Jr., \emph{Dirac-Hestenes Lagrangian}, Int. J. Th. Phys. {\bf 38} (1999) 2349--2369. %%QUATERNION, %%SPIN-1/2, %%DIRAC, %$D15032002. 

\bibitem{EDMON1999-} J.D. Edmonds, Jr., \emph{Dirac's equation in half of his algebra}, Eur. J. Phys. {\bf 20} (1999) 461--467.  %%QUATERNION, %%SPIN-1/2, %%DIRAC, %%LANCZOS, %$D06022002.

\bibitem{HEHL-1999-} F.W. Hehl et al., \emph{On the structure of the energy-momentum and spin currents in Dirac's electron theory}, in: A. Harvey, ed., On Einstein's Path --- Essays in Honor of Engelbert Schucking (Springer, New York, 1999) 257--273. %%QUATERNION, %%SPIN-1/2, %%DIRAC, %%ENERGY-MOMENTUM-TENSOR, %%GRAVITATIONAL-MOMENTS, %%DIFFERENTIAL-FORM,  %$D13052001.

\bibitem{BELL-2000-} S.B.M. Bell, J.P. Cullerne, and B.M. Diaz, \emph{Classical behavior of the Dirac bispinor}, Found. Phys. {\bf 30} (2000) 35--57.  %%QUATERNION, %%SPIN-1/2, %%DIRAC, %$D04032006.

\bibitem{DRAY-2000-} T. Dray and C. Manogue, \emph{Quaternionic spin}, in: R. Ablamowicz and B. Fauser, eds., Clifford Algebra and their Applications in Mathematical Physics, Vol.~1: \emph{Algebra and Physics} (Birkh\"auser, Boston, 2000)  21--37. %%QUATERNION, %%SPIN-1/2, %%DIRAC, %%SPINOR, %%OCTONION, %$D12052001.

\bibitem{CELER2001-} M.-N. C\'el\'erier and L. Nottale, \emph{Dirac equation in scale relativity} (21 December 2001) 33~pp.;  e-print \underline{ arXiv:hep-th/0112213 }.  %%QUATERNION, %%SPIN-1/2, %%DIRAC, %$D30112005.

\bibitem{FAUSE2001-} B. Fauser, \emph{Equivalence of Daviau's, Hestenes', and Parra's formulations of Dirac theory}, Int. J. Theor. Phys. {\bf 40} (2001) 399--411.  %%QUATERNION, %%SPIN-1/2, %%DIRAC, %%HESTENES, %$D30112005. 

\bibitem{GSPON2001-} A. Gsponer and J.-P. Hurni, \emph{Comment on formulating and generalizing Dirac's, Proca's, and Maxwell's equations with biquaternions or Clifford numbers}, Found. Phys. Lett. {\bf 14} (2001) 77--85; e-print \underline{ arXiv:math-ph/0201049 }. %%QUATERNION, %%SPIN-1/2, %%DIRAC, %%LANCZOS, %%PROCA, %%MAXWELL, %%EINSTEIN-MAYER, %$D13052001.

\bibitem{JOYCE2001-}  W.P. Joyce, \emph{Dirac theory in spacetime algebra: I. The generalized bivector Dirac equation}, J. Phys. A: Math. Gen. {\bf 34} (2001) 1991--2005. %%QUATERNION, %%SPIN-1/2, %%DIRAC, %$D16032002, %$C See comment by W.E. Baylis, Ref.~\cite{BAYLI2002-}.

\bibitem{BAYLI2002-}  W.E. Baylis, \emph{Comment on ``Dirac theory in spacetime algebra''}, J. Phys. A: Math. Gen. {\bf 35} (2002) 4791--4796. %%QUATERNION, %%SPIN-1/2, %%DIRAC, %%GSPONER, %$D16032002.

\bibitem{JOYCE2002A}  W.P. Joyce and J.G. Martin, \emph{Equivalence of Dirac formulations}, J. Phys. A: Math. Gen. {\bf 35} (2002) 4729--4736. %%QUATERNION, %%SPIN-1/2, %%DIRAC, %$D16032002.

\bibitem{JOYCE2002B}  W.P. Joyce, \emph{Gauge freedom of Dirac theory in complexified spacetime algebra}, J. Phys. A: Math. Gen. {\bf 35} (2002) 4737--4747. %%QUATERNION, %%SPIN-1/2, %%DIRAC, %$D16032002.

\bibitem{JOYCE2002C}  W.P. Joyce, \emph{Reply to comments on `Dirac theory in spacetime algebra'}, J. Phys. A: Math. Gen. {\bf 35} (2002) 4797--4798. %%QUATERNION, %%SPIN-1/2, %%DIRAC, %$D30112005.

\bibitem{CELER2002-} M.-N. C\'el\'erier and L. Nottale, \emph{A scale-relativistic derivation of the Dirac equation}, Electromagn. Phenom. {\bf 3} (2003) 70--80;  e-print \underline{ arXiv:hep-th/0210027 }.  %%QUATERNION, %%SPIN-1/2, %%DIRAC, %$D30112005.

\bibitem{YUFEN2002-} L. Yu-Fen, \emph{Triality, biquaternion and vector representation of the Dirac equation}, Adv. Appl. Clifford Alg. {\bf 12} (2002) 109--124.  %%QUATERNION, %%SPIN-1/2, %%DIRAC, %%C The "Dirac-Lanczos" equation of 1929 is rediscovered: Eq. (33) p.9, %$D30112005.

\bibitem{RODRI2003-} Rodrigues, W. A. Jr., \emph{Maxwell-Dirac equivalences of the first and second kinds and the Seiberg-Witten equations}, Int. J. Math. and Math. Sci. (2003) 2707--2734.  %%QUATERNION, %%SPIN-1/2, %%DIRAC, %%MAXWELL, %%WRONG, %$D06072008. 

\bibitem{TANIS2003A} M. Tanisli and G. \"Ozg\"ur, \emph{Biquaternionic representations of angular momentum and Dirac equation}, Acta Physica Slovaca {\bf 53} (2003) 243--252. %%QUATERNION, %%SPIN-1/2, %%DIRAC, %$D04032006.

\bibitem{RODRI2004-} Rodrigues, W. A. Jr., \emph{Algebraic and Dirac-Hestenes spinors and spinor fields}, J. Math. Phys. {\bf 45} (2004) 2908--2944.  %%QUATERNION, %%SPIN-1/2, %%DIRAC, %%MAXWELL, %%WRONG, %$D06072008.

\bibitem{CHRIS2006-}  V. Christianto, \emph{A New Wave Quantum Relativistic Equation from Quaternionic Representation of Maxwell-Dirac Isomorphism as an Alternative to Barut-Dirac Equation}, Electronic Journal of Theoretical Physics (EJTP) {\bf 3} (2006) 117--144. %%QUATERNION, %%SPIN-1/2, %%BARUT, %$D06072008. 

\bibitem{MORIT2007-} K. Morita, \emph{Quaternions, Lorentz group and the Dirac equation}, Progress Theor. Phys. {\bf 117} (2007) 501--532; e-print \underline{ arXiv:hep-th/0701074 }. %%QUATERNION, %%SPIN-1/2,  %%GSPONER, %%HURNI, %$D06072008.

\bibitem{REDIN2007-} N. Redington and M.A.K. Lodhi,  \emph{A simple five-dimensional wave equation for a Dirac particle}, J. Math. Phys. {\bf 48} (2007) 1--18; e-print \underline{ arXiv:quant-ph/0512140 }. %%QUATERNION, %%SPIN-1/2, %$D06072008. 

\end{enumerate}


\subsection{SPIN-3/2}
%--------------------
\label{SPIN-3/2}


\begin{enumerate}

\bibitem{EDMON1976B} J.D. Edmonds, Jr., \emph{A relativistic ``higher spin'' quaternion wave equation giving a variation on the Pauli equation}, Found. Phys. {\bf 6} (1976) 33--36.  %%QUATERNION, %%SPIN-3/2, %$D24062002.

\bibitem{MORIT1984-} K. Morita, \emph{Quaternions and simple $D=4$ supergravity}, Prog. Th. Phys. {\bf 72} (1984) 1056--1059.  %%QUATERNION, %%SPIN-3/2, %%SUPERSYMMETRY, %$D24062002.

\bibitem{MORIT1985-} K. Morita, \emph{Quaternionic variational formalism for Poincar\'e gauge theory and supergravity}, Prog. Th. Phys. {\bf 73} (1985) 999--1015.  %%QUATERNION, %%SPIN-3/2, %%SUPERSYMMETRY, %$D24062002.

\bibitem{PENRO1991-} R. Penrose, \emph{Twistors as spin 3/2 charges}, in: A. Zichichi et al., eds., Gravitation and Modern Cosmology (Plenum Press, New York, 1991) 129--137. %%QUATERNION, %%SPIN-3/2, %%TWISTOR, %$D30112005.

\bibitem{BURES2001-} J. Bures, F. Sommen, V. Soucek and P. vanLancker, \emph{Rarita-Schwinger type operators in Clifford analysis} J. Funct. Anal. {\bf 185} (2001) 425--455.  %%QUATERNION, %%SPIN-3/2, %$D24062002.
 
\bibitem{GSPON2002G} A. Gsponer and J.-P. Hurni, \emph{Lanczos's equation as a way out of the spin 3/2 crisis?}, Hadronic Journal {\bf 26} (2003) 327--350; e-print \underline{ arXiv:math-ph/00210055 }.  %%QUATERNION, %%SPIN-3/2, %%PARTICLE-PHYSICS, %%DIRAC, %%LANCZOS, %$D13052001.

\end{enumerate}
   

\subsection{ANALYTICITY-MAXWELL}
%-------------------------------
\label{ANALYTICITY-MAXWELL}


Hypercomplex analysis applied to Maxwell's field, with emphasis on physics.

\begin{enumerate}

\bibitem{LANCZ1919A} Kornel Lanczos, \emph{Die Funktionentheoretischen Beziehungen der Max\-well\-schen  Aethergleichungen --- Ein Beitrag zur  Relativit\"ats- und Elektronentheorie [The relations of the homogeneous Maxwell's equations to the theory of functions --- A contribution to the theory of relativity and electrons]} (Verlagsbuchhandlung Josef N\'emeth, Budapest, 1919) 80~pp. Handwritten and lithographed in 50 copies. Reprinted in: W.R.\ Davis \emph{et al.}, eds., Cornelius Lanczos Collected Published Papers With Commentaries (North Carolina State University, Raleigh, 1998) Volume {\bf VI}, pages A-1 to A-82. %%BOOK, %%QUATERNION, %%ANALYTICITY-MAXWELL, %$D01062001. 

For a typeseted version, see \cite{LANCZ1919B}.

\bibitem{LANCZ1919B} Cornelius Lanczos, \emph{The relations of the homogeneous Maxwell's equations to the theory of functions --- A contribution to the theory of relativity and electrons} (1919, Typeseted by Jean-Pierre Hurni with a preface by Andre Gsponer, 2004) 58~pp.; e-print \underline{ arXiv:physics/0408079 }. %%QUATERNION, %%ANALYTICITY-MAXWELL, %$D03082005.

\bibitem{RAINI1925-} G.Y. Rainich, \emph{Electrodynamics in the general relativity theory}, Trans. Am. Math. Soc. {\bf 27} (1925) 106--136. See Part III, ``Integral properties and singularities.'' %%QUATERNION, %%ANALYTICITY-MAXWELL, %%SILBERSTEIN, %$D04032006.

\bibitem{EICHL1939-} M. Eichler, \emph{Allgemeine Integration einiger partieller Differentialgleichungen der mathematischen Physik durch Quaternionenfunktionen}, Comment. Math. Helv. {\bf 12} (1939) 212--224. %%QUATERNION, %%ANALYTICITY-MAXWELL, %$D01062001.

\bibitem{IMAED1950-} K. Imaeda, \emph{Linearization of Minkowski space and five-dimensional space}, Prog. Th. Phys. {\bf 5 } (1950) 133--134. %%QUATERNION, %%ANALYTICITY-MAXWELL, %$D01062001.

\bibitem{IMAED1951-} K. Imaeda, \emph{A study of field equations and spaces by means of hypercomplex numbers}, Memoirs of the Faculty of Liberal Arts and Education {\bf 2 } (Yamanashi University, Kofu, Japan, 1951) 111--118. %%QUATERNION, %%ANALYTICITY-MAXWELL, %%PROCA, %$D02562001.

\bibitem{SNEER1968-} M.S. Sneerson, \emph{Maxwell's equation, and functionally invariant solutions of the wave equation}, Differential Equations {\bf 4} (1968) 386--394. %%QUATERNION, %%ANALYTICITY-MAXWELL, %$D30112005.

\bibitem{SNEER1971-} M.S. Sneerson, \emph{Linear nonseparable transformations and the Hall effect}, Differential Equations {\bf 7} (1971) 294--295. %%QUATERNION, %%ANALYTICITY-MAXWELL, %$D30112005.

\bibitem{SNEER1972-} M.S. Sneerson, \emph{A problem involving the directional derivative for a harmonic function of three independent variables}, Differential Equations {\bf 8} (1972) 1479--1481. %%QUATERNION, %%ANALYTICITY-MAXWELL, %$D30112005.

\bibitem{EVANS1976-} D.D. Evans, \emph{Complex variable theory generalized to electromagnetics: The theory of functions of a quaternion variable}, Ph.D. Thesis (Univ. of California, 1976). %%QUATERNION, %%ANALYTICITY-MAXWELL, %$D04032006.

\bibitem{IMAED1976A} K. Imaeda, \emph{A new formulation of classical electrodynamics}, Nuov. Cim. {\bf 32 B} (1976) 138--162. %%QUATERNION, %%ANALYTICITY-MAXWELL, %$D01062001.

\bibitem{IMAED1986B} K. Imaeda, \emph{Quaternionic formulation of classical electromagnetic fields and theory of functions of a biquaternionic variable},   in: J.S.R. Chisholm and A.K. Common, eds.,  Clifford Algebras and Their Applications in Mathematical Physics (Reidel, Dordrecht, 1986) 495--500. %%QUATERNION, %%ANALYTICITY-MAXWELL, %$D01062001.

\bibitem{WEING1973-} D. Weingarten, \emph{Complex symmetries of electrodynamics}, Ann. Phys. {\bf 76} (1973) 510--548. %%QUATERNION, %%ANALYTICITY-MAXWELL, %%CONFORMALITY, %$D16052001.

\bibitem{KASSA1995-} V.V. Kassandrov, \emph{Biquaternion electrodynamics and Weyl-Cartan geometry of space-time}, Gravitation \& Cosmology {\bf 3} (1995) 216--222. %%QUATERNION, %%ANALYTICITY-MAXWELL, %$D02962001.

\bibitem{KRAVC1995C} V.V. Kravchenko and M.V. Shapiro, \emph{Quaternionic time-harmonic Maxwell operator}, J. Phys. {\bf A 28} (1995) 5017--5031.  %%QUATERNION, %%ANALYTICITY-MAXWELL, %$D01062001.

\bibitem{COLOM1998-} F. Colombo, P. Loutaunau, I. Sabadini, and D.C. Struppa, \emph{Regular functions of biquaternionic variables and Maxwell's equations}, J. Geom. Phys. {\bf 26} (1998) 183--201. %%QUATERNION, %%ANALYTICITY-MAXWELL, %$D01062001.

\bibitem{KASSA1998-} V.V. Kassandrov, \emph{Conformal mappings, hyperanalyticity and field dynamics}, Acta Applicandae Math. {\bf 50} (1998) 197--206.  %%QUATERNION, %%ANALYTICITY-MAXWELL, %%CONFORMALITY, %$D01062001.

\bibitem{GSPON1998A} A. Gsponer and J.-P. Hurni, \emph{Lanczos's functional theory of electrodynamics --- A commentary on Lanczos's Ph.D. dissertation,} in: W.R. Davis et al., eds., Cornelius Lanczos Collected Published Papers With Commentaries,~{\bf I} (North Carolina State University, Raleigh, 1998) 2-15 to 2-23; e-print \underline{ arXiv:math-ph/0402012 }. %%QUATERNION, %%ANALYTICITY-MAXWELL, %$D01062001.

\bibitem{KASSA1999-} V.V. Kassandrov and V.N. Trishin, \emph{``Particle-like'' singular solutions in Einstein-Maxwell theory and in algebraic dynamics}, Gravit \&  Cosmol.  {\bf 5} (1999) 272--276.  %%QUATERNION, %%ANALYTICITY-MAXWELL, %%GAUGE-THEORY, %%TWISTOR, %$D01062001.

\bibitem{KASSA2000-} V.V. Kassandrov and J.A. Rizcallah, \emph{Twistor and ``weak'' gauge structures in the framework of quaternionic analysis} (29 Dec 2000) 21~pp.; e-print \underline{ arXiv:gr-qc/0012109 }.  %%QUATERNION, %%ANALYTICITY-MAXWELL, %%PARTICLE-PHYSICS, %%GAUGE-THEORY, %%TWISTOR,  %$D30112005.

\bibitem{KRAVC2001A} V.V. Kravchenko, \emph{Applied quaternionic analysis: Maxwell's system and Dirac's equation}, in: W. Tutschke, ed., Functional-Analytic and Complex Methods, their Interactions, and Applications to Partial Differential Equations (World Scientific, 2001) 143--160. %%QUATERNION, %%ANALYTICITY-MAXWELL, %%CONSTITUTIVE-TENSOR, %%DIRAC, %%GSPONER, %$D06022002.

\bibitem{KASSA2002-} V.V. Kassandrov, \emph{General solution of the complex 4-eikonal equation and the ``algebrodynamical'' field theory}, Grav. \& Cosmol. {\bf 8} (2002) 57--62.  %%QUATERNION, %%ANALYTICITY-MAXWELL, %%PARTICLE-PHYSICS, %%GENERAL-RELATIVITY, %$D04032006.

\bibitem{KRAVC2004-} V. Kravchenko \emph{On the reduction of the multidimensional Schr\"odinger equation to a first order equation and its relation to the pseudoanalytic function theory}, J. Physics A: Mathematical and General {\bf 38} (2005) 851--868; e-print \underline{ arXiv:math.AP/0408172 }. %%QUATERNION, %%ANALYTICITY-MAXWELL, %$D30112005.

\bibitem{KASSA2004A} V.V. Kassandrov, \emph{Singular sources of maxwell fields with self-quantized electric charge}, in: A. Chubykalo, V. Onoochin, A. Espinoza, and V. Smirnov-Rueda, eds., Has the Last Word been Said on Classical Electrodynamics? (Rinton Press, 2004) 42--67.  %%QUATERNION, %%ANALYTICITY-MAXWELL, %%PARTICLE-PHYSICS, %%GENERAL-RELATIVITY, %$D04032006.

\bibitem{KASSA2004B} V.V. Kassandrov, \emph{Nature of time and particles-caustics: physical world in algebrodynamics and in twistor theory}, Hypercomplex Num. Geom. Phys. {\bf 1} (2004) 89--105.  %%QUATERNION, %%ANALYTICITY-MAXWELL, %%PARTICLE-PHYSICS, %%GENERAL-RELATIVITY, %%TWISTOR, %$D04032006.

\bibitem{GSPON2004C} A. Gsponer, \emph{On the physical nature of the Lanczos-Newman ``circle electron'' singularity}, Report ISRI-04-04 (19 May 2004) 36~pp.; e-print \underline{ arXiv:gr-qc/0405046 }.  %%QUATERNION, %%ANALYTICITY-MAXWELL, %%PARTICLE-PHYSICS, %%RENORMALIZATION, %%LANCZOS, %%HADRODYNAMICS, %$D03082005.

\bibitem{KASSA2006-} V.V. Kassandrov, \emph{On the structure of general solution to the equations of shear-free null congruences}, in: Proceedings of the Int. School-Seminar on geometry and analysis in memory of N.V.Efimov. (Rostov-na-Donu Univ. Press, 2004) 65--68; e-print \underline{ arXiv:gr-qc/0602046 }.  %%QUATERNION, %%ANALYTICITY-MAXWELL, %%PARTICLE-PHYSICS, %%GENERAL-RELATIVITY, %$D04032006.


\end{enumerate}

\subsection{ANALYTICITY-DIRAC}
%-----------------------------
\label{ANALYTICITY-DIRAC}

Hypercomplex analysis applied to Dirac's field.

\begin{enumerate}


\bibitem{IWANE1930-} D. Iwanenko and K. Nikolsky, \emph{�ber den Zusammenhang zwischen den Cauchy-Riemannschen und Diracschen Differentialgleichungen},   Zeits. f. Phys. {\bf 63} (1930) 129--137. %%QUATERNION, %%ANALYTICITY-DIRAC, %%MAXWELL, %$D08082001.

\bibitem{BOSSH1940-} P. Bosshard, \emph{Die Cliffordschen Zahlen, ihre Algebra, und ihre Funktionentheorie} (Ph.D. thesis, Universit\"at Zurich, 1940) 48~pp. %%QUATERNION, %%ANALYTICITY-DIRAC, %%CLIFFORD-C, %$D01022002.

\bibitem{FUETE1943-} R. Fueter, \emph{Die Funktionentheorie der Diracschen Differentialgleichungen}, Comm. Math. Helv. {\bf 16} (1943) 19--28. %%QUATERNION, %%ANALYTICITY-DIRAC, %$D17052001.

\bibitem{KRISZ1947-} A. Kriszten, \emph{Funktionentheorie und Randwertproblem der Diracschen Differentialgleichungen}, Comm. Math. Helv. {\bf 20} (1947) 333--365. %%QUATERNION, %%ANALYTICITY-DIRAC, %%CLIFFORD-C, %$D17052001.

\bibitem{SNEER1964-} M.S. Sneerson, \emph{A class of solutions of a system of differential equations of Moisil and Dirac}, Amer. Math. Soc. Transl. {\bf 42} (1964) 195--198. %%QUATERNION, %%ANALYTICITY-DIRAC, %Cites the paper of D. Iwanenko and K. Nikolsky, %$D11052004.

\bibitem{KRAVC1996-} V.V. Kravchenko, H.R. Malonek, and G. Santana, \emph{Biquaternionic integral representations for massive Dirac spinors in a magnetic field and generalized biquaternionic differentiability}, Math. Meth. in the Appl. Sci. {\bf 19} (1996) 1415--1431. %%QUATERNION, %%ANALYTICITY-DIRAC, %$D29052001.

\bibitem{PALAM1999-} V.P. Palamodov, \emph{Holomorphic synthesis of monogenic functions of several quaternionic variables}, J. Anal. Math. {\bf 78} (1999) 177--204. %%QUATERNION, %%ANALYTICITY-DIRAC, %$D29052001.

\bibitem{KRAVC2000-} V.V. Kravchenko, \emph{On a new approach for solving the Dirac equations with some potentials and Maxwell's system in inhomogeneous media}, in: J. Elschner, I. Gohberg and B. Silbermann, eds., Operator Theory: Advances and Applications {\bf 121} (Birkh\"auser Verlag, 2001) 278--306. %%QUATERNION, %%ANALYTICITY-DIRAC, %%MAXWELL, %$D29052001.

\bibitem{KRAVC2002A} V.V. Kravchenko, R. Castillo, \emph{An analogue of the Sommerfeld radiation condition for the Dirac operator},  Mathematical  
Methods in the Applied Sciences  {\bf 25} (2002) 1383--1394. %%QUATERNION, %%ANALYTICITY-DIRAC, %$D29052001.

\bibitem{SABAD2003-} I. Sabadini, F. Sommen and D.C. Struppa, \emph{The Dirac complex on abstract vector variables: megaforms}, Experimental Math. {\bf 12} (2003) 351--364.  %%QUATERNION, %%ANALYTICITY-DIRAC, %$D06072008.

\bibitem{CASTA2005-} A. Castaneda and V.V. Kravchenko, \emph{New applications of pseudoanalytic function theory to the Dirac equation}, J. Phys. A: Math. Gen. {\bf 38} (2005) 9207--9219. %%QUATERNION, %%ANALYTICITY-DIRAC, %%SPIN-1/2, %%LANCZOS, %%GSPONER, %$D30112005.

\end{enumerate}




\section{\Huge PHYSICS}
%======================
\label{PHYSICS}

This chapter contains papers in which the quaternion formalism and quaternion methods are applied to physics, with an emphasis on applications of a fundamental rather than practical character.

However, papers in which quaternions are applied to GENERAL-RELATIVITY are in Chap.~\ref{RELATIVITICS} on ``relativistics,'' and papers in which quaternions are applied to QUANTUM-PHYSICS are in Chap.~\ref{QUANTICS} on ``quantics.''

Moreover, as was explained in Chap.~\ref{Introduction}, the bibliography also includes selected papers in which a formalism allied to quaternions is used (e.g., semivectors or Clifford numbers), especially if these papers could have been written using quaternions rather than the closely related formalism. 

On the other hand, the numerous papers which use a standard tensor or matrix formalism (e.g., the Pauli- or Dirac-matrices, and the corresponding two- or four-component formalisms) are excluded from the bibliography.  The reason for this exclusion is conceptual rather than conventional (because Pauli's matrices are actually very closely related to Hamilton's quaternion units).  The reason is that the emphasis of standard formalisms is on the vector-space structure supported by them, while the emphasis of the biquaternion formalism is on the full algebraic structure (vector-space and multiplication ring) provided by the biquaternion algebra.

The instances in which selected papers using Pauli-matrices and two-component spinors are included in the bibliography (e.g., by F.\ G\"ursey, L.M.\ Brown, and even by Pauli himself) correspond to cases in which their use is essentially equivalent to that of biquaternions due to the isomorphism $\mathbb{B} \equiv  
M_2(\mathbb{C})$. 


\subsection{PHYSICS-VARIA}
%------------------------
\label{PHYSICS-VARIA}

Papers and books dealing with several applications to physics, and few selected references dealing with applications in computer graphics, modeling, etc. 


See also the books by W.R.\ Hamilton, P.G.\ Tait, and C.J.\ Joly listed in the MATH-VARIA section, Sec.~\ref{MATH-VARIA}, which contain chapters or sections on the applications of quaternions to classical physics topics such as mechanics, hydrodynamics, astronomy, etc..


\begin{enumerate}

\bibitem{HAMIL1862-} W. R. Hamilton, \emph{On some quaternion equations connected with Fresnel's wave surface for biaxial crystals}, Proc. Roy. Irish Acad. {\bf 7} (1862) 122--124, 163.  %%QUATERNION, %%PHYSICS-VARIA, %$D06072008.

\bibitem{TAIT-1890-} P.G. Tait, \emph{On the importance of quaternions in physics}, Phil. Mag. {\bf } (January 1890) SP-2:297--308.  %%QUATERNION, %%PHYSICS-VARIA, %$D24062002.

\bibitem{MERCI1935B} A. Mercier, \emph{Expression du second principe de la thermodynamique au moyen des nombres de Clifford}, Suppl. Arch. Sci. Phys. Nat. Gen\`eve (1935) 112--113.  %%QUATERNION, %%PHYSICS-VARIA, %%THERMODYNAMICS, %%CLIFFORD, %$D20032002.

\bibitem{FISCH1951-} O.F. Fischer, \emph{Universal mechanics and Hamiltons quaternions} (Axion Institute, Stockholm, 1951) 356~pp. %%BOOK, %%QUATERNION, %%PHYSICS-VARIA, %%MECHANICS, %%ELECTRODYNAMICS, %%MAXWELL, %%DIRAC, %%GENERAL-RELATIVITY, %$D20032002.

\bibitem{HAYES1964-} M.V. Hayes, A Unified Field Theory (The Stinghour Press, Lunenburg, Vermont, 1964) 70~pp. %%BOOK, %%QUATERNION, %%PHYSICS-VARIA, %$D04032006.

\bibitem{WINAN1977-} J.G. Winans, \emph{Quaternion physical quantities}, Found. of Phys. {\bf 7} (1977) 341--349. Errata, ibid {\bf 11} (1981) 651. %%QUATERNION, %%PHYSICS-VARIA, %%UNITS, %$D14032002.

\bibitem{KOMKO1979-} V. Komkov, \emph{Quaternions, Frechet differentiation, and some equations of mathematical physics. 1: Critical point theory}, J. Math. Anal. Appl. {\bf 71} (1979) 187--209.  %%QUATERNION, %%PHYSICS-VARIA, %%ANALYSIS, %%MECHANICS, %%ELECTRODYNAMICS, %$D30112005.

\bibitem{VESEL1982-} F.J. Vesely, \emph{Angular Monte-Carlo integration using quaternion parameters: A spherical reference potential for $CCl_4$}, J. Computational Phys. {\bf 47} (1982) 291-296.  %%QUATERNION, %%PHYSICS-VARIA, %%MONTE-CARLO, %$D20032002.

\bibitem{FIRNE1986-} M. Firneis and F. Firneis, \emph{On the use of quaternions in spherical and positional astronomy}, Astronomical Journal {\bf 91} (1986) 177--178. %%QUATERNION, %%PHYSICS-VARIA, %$D30112005. %$C Cite E. Study and W. Blaschke.

\bibitem{HORN-1987-} B.K.P. Horn, \emph{Closed-form solution of absolute orientation using unit quaternions}, J. Opt. Soc. America {\bf A 4} (1987) 629--642. %%QUATERNION, %%PHYSICS-VARIA, %$D20032002.
 
\bibitem{DAVIE1991-} A.J. Davies, R. Foot, G.C. Joshi, and B.H.J. McKellar, \emph{Quaternionic methods in integral transforms of geophysical interest}, Geophys. J. Int. {\bf 99} (1991) 579--582. %%QUATERNION, %%PHYSICS-VARIA, %%ANALYTICITY-H, %%GEOPHYSICS, %$D29052001.

\bibitem{LOSCO1991-} L. Losco, F. Pelletier, and J.P. Taillard, \emph{Modeling a chain of rigid bodies by biquaternions}, Eur. J. Mech. {\bf A 10} (1991) 433--451.  %%QUATERNION, %%PHYSICS-VARIA, %$D20032002. %$C Biquaternions of Clifford: $i^2=0$.

\bibitem{GSPON1993B} A. Gsponer and J.-P. Hurni,  \emph{The physical heritage of Sir W.R. Hamilton}. Presented at the Conference The Mathematical Heritage of Sir William Rowan Hamilton (Trinity College, Dublin, 17-20 August, 1993) 35~pp.; e-print \underline{ arXiv:math-ph/0201058 }. %%QUATERNION, %%PHYSICS-VARIA, %%HISTORY, %$D19032002. 

\bibitem{HASLW1995-} T. Haslwanter, \emph{Mathematics of the three dimensional eye rotations}, Vision Res. {\bf 35} (1995) 1727--1739.  %%QUATERNION, %%PHYSICS-VARIA, %$D06072008.  

\bibitem{ALEAR1996-} J.E. Al\'e Araneda, \emph{Dimensional-directional analysis by a quaternionic representation of physical quantities}, J. Franklin Inst. {\bf 333} (1996) 113--126. %%QUATERNION, %%PHYSICS-VARIA, %$D06072008.

\bibitem{MEIST1997-} L. Meister, \emph{Quaternions and their applications in photogrammetry and navigation}, Doctor rerum naturalium habilitatus of the Fakult\"at f\"ur Mathematik und Informatik der TU Bergakademie Freiberg (1997) 64~pp.  %%QUATERNION, %%PHYSICS-VARIA, %$D06072008.

\bibitem{SIMIN1997-} D.J. Siminovitch, \emph{An NMR rotation operator disentanglement strategy for establishing properties of the Euler-Rodrigues parameters}, J. of Physics A: Math. Gen. {\bf 30} (1997) 2577--2584.  %%QUATERNION, %%PHYSICS-VARIA, %$D04032006.  

\bibitem{KUIPE1999-} J.B. Kuipers, Quaternions and rotation sequences --- A Primer with Applications to Orbits, Aerospace, and Virtual Reality (Princeton University Press, 1999) 371~pp.  %%BOOK, %%QUATERNION, %%PHYSICS-VARIA, %$D20032002.

\bibitem{BAR-I2000-} I.-Y. Bar-Itzhack, \emph{New method for extracting the quaternion from a rotation matrix}, J. Guidance, Control, and Dynamics {\bf 23} (2000) 1085-1087.%%QUATERNION, %%PHYSICS-VARIA, %$D06072008.

\bibitem{GEBRE2000-} D. Gebre-Egziabher, G.H. Elkaim, J.D. Powell and B.W. Parkinson, \emph{A gyro-free quaternion-based attitude determination system suitable for implementation using low cost sensors}, in: Position Location and Navigation Symposium (IEEE, 2000) 185--192.   %%QUATERNION, %%PHYSICS-VARIA, %$D06072008.

\bibitem{NADLE2001-} A. Nadler, I.-Y. Bar-Itzhack and H. Weiss, \emph{Iterative algorithms for attitude estimation using global positioning system phase measurements}, J. Guidance, Control, and Dynamics {\bf 24} (2001) 983--990. %%QUATERNION, %%PHYSICS-VARIA, %$D06072008.

\bibitem{BAR-I2002A} I.-Y. Bar-Itzhack and R.R. Harman, \emph{Optimal fusion of a given quaternion with vector measurements}, J. Guidance, Control, and Dynamics {\bf 25} (2002) 188-190. %%QUATERNION, %%PHYSICS-VARIA, %$D06072008.

\bibitem{BAR-I2002B} I.-Y. Bar-Itzhack and R.R. Harman, \emph{In-space calibration of a skewed gyro quadruplet}, J. Guidance, Control, and Dynamics {\bf 25} (2002) 852--859. %%QUATERNION, %%PHYSICS-VARIA, %$D06072008.

\bibitem{HORN2002-} M.E. Horn, \emph{Quaternions in university-level physics: Considering special relativity}, German Physical Society Spring Conference  (2002) 6~pp.;  e-print \underline{ arXiv:physics/0308017 }.  %%QUATERNION, %%PHYSICS-VARIA, %%SPECIAL-RELATIVITY, %$D03082005.

\bibitem{MUKUN2002-} R. Mukundan, \emph{Quaternions:  from classical mechanics to computer graphics, and beyond}, in: Proceedings of the 7th Asian Technology Conference in Mathematics (2002) 97--106.  %%QUATERNION, %%PHYSICS-VARIA, %$D06072008.

\bibitem{JACK-2003-} P.M. Jack, \emph{Physical space as a quaternion structure, I: Maxwell equations. A brief Note}, Report hypcx-20001015e (July 18, 2003) 6~pp; e-print \underline{ arXiv:math-ph/0307038 }.  %%QUATERNION, %%PHYSICS-VARIA, %%ELECTRODYNAMICS, %%THERMODYNAMICS, %$D30112005.

\bibitem{CAPOZ2005-} S. Capozziello and A. Lattanzi, \emph{Chiral tetrahedrons as unitary quaternions: molecules and particles under the same standard?}, Int. J. Quantum Chem. {\bf 104} (2005) 885--893; e-print \underline{ arXiv:physics/0502092 }. %%QUATERNION, %%PHYSICS-VARIA, %%CHEMISTRY, %%LANCZOS, %%CONWAY, %$D04032006, %$C Reference to Lanczos and Conway identical as in {CELER2004-}.

\bibitem{KARNE2007-} C.F.F. Karney, \emph{Quaternions in molecular modeling}, J. Molecular Graphics and Modeling {\bf 25} (2007) 595--604; e-print \underline{ arXiv:physics/0506177 }.  %%QUATERNION, %%PHYSICS-VARIA, %$D06072008.

\end{enumerate}


\subsection{MECHANICS}
%---------------------
\label{MECHANICS}

This section also includes papers on ``analytical mechanics,'' and ``generalized dynamics,'' etc.

\begin{enumerate}

\bibitem{HAMIL1850-} W. R. Hamilton, \emph{On quaternions and the rotation of a solid body}, Proc. Roy. Irish Acad. {\bf 4} (1850) 38--56.   %%QUATERNION, %%MECHANICS, %$D06072008.

\bibitem{HAMIL1845-} W.R. Hamilton, \emph{On the application of the method of quaternions to some dynamical systems}, Proc. Roy. Irish Acad. {\bf 3} (1847) Appendix, xxxvi--l. %%QUATERNION, %%MECHANICS, %%KEPLER, %%LAPLACE-RUNGE-LENZ, %$D30112005, %$C Communicated on July 14 and July 21, 1845.

\bibitem{JOLY-1899A} C.J. Joly, \emph{Astatics and quaternion functions}, Proc. Roy. Irish Acad. {\bf 5} (1899) 366--369. %%QUATERNION, %%MECHANICS, %$D15032002.

\bibitem{JOLY-1902A} C.J. Joly, \emph{The interpretation of quaternion as a point symbol}, Trans. Roy. Irish Acad. {\bf A 32} (1902-1904) 1--16. %%QUATERNION, %%MECHANICS, %%PROJECTIVE-GEOMETRY, %$D15032002.

\bibitem{JOLY-1904-} C.J. Joly, \emph{Some new relations in the theory of screws}, Proc. Roy. Irish Acad. {\bf A 8} (1904) 69--70. %%QUATERNION, %%MECHANICS, %%MATRIX, %%LINEAR-FUNCTION, %$D14032002.

\bibitem{FLINT1920-} H.T. Flint, \emph{Applications of quaternions to the theory of relativity}, Phil. Mag. {\bf 39} (1920) 439--449. %%QUATERNION, %%MECHANICS, %%ELECTRDYNAMICS, %$D09022002.

\bibitem{HILL-1945-} E.L. Hill, \emph{Rotations of a rigid body about a fixed point}, Amer. J. Phys. {\bf 13} (1945) 137--140.  %%QUATERNION, %%MECHANICS, %$D06072008. %$C Juvet+Schidlof

\bibitem{PROCA1946-} A. Proca, \emph{Sur les \'equations relativistes des particules \'el\'ementaires}, Cr. Acad. Sci. Paris {\bf 223} (1946) 270--272.  %%QUATERNION, %%MECHANICS, %%SPINOR, %$D09022002.

\bibitem{PROCA1947-} A. Proca, \emph{New possible equations for fundamental particles}, Phys. Soc. Cambridge Conf. Rep. (1947) 180--181.  %%QUATERNION, %%MECHANICS, %%SPINOR, %$D09022002.

\bibitem{PROCA1952-} A. Proca, \emph{Sur l'espace-temps des particules fondamentales et les espaces spinoriels sous-jacents}, Bull. Scientifique Roumain {\bf 1} (1952) 18--24.  %%QUATERNION, %%MECHANICS, %%SPINOR, %$D09022002.

\bibitem{PROCA1954-} A. Proca, \emph{M\'ecanique du point}, J. Phys. Radium {\bf 15} (1954) 65--72.  %%QUATERNION, %%MECHANICS, %%SPINOR, %$D09022002.

\bibitem{PROCA1955A} A. Proca, \emph{Particules de tr\`es grandes vitesses en m\'ecanique spinorielle}, Nuovo Cim. {\bf 2} (1955) 962--971. %%QUATERNION, %%MECHANICS, %%SPINOR, %$D09022002.

\bibitem{PROCA1955B} A. Proca, \emph{Interf\'erences en m\'ecanique spinorielle}, Nuovo Cim. {\bf 2} (1955) 972--979. %%QUATERNION, %%MECHANICS, %%SPINOR, %$D09022002.

\bibitem{PROCA1956A} A. Proca, \emph{Sur la m\'ecanique spinorielle du point charg\'e}, J. Phys. Radium {\bf 17} (1956) 81--82.  %%QUATERNION, %%MECHANICS, %%SPINOR, %$D09022002.

\bibitem{PROCA1956B} A. Proca, \emph{Sur un nouveau principe d'\'equivalence sugg\'er\'e par les m\'ecaniques spinorielles}, J. Phys. Radium {\bf 17} (1956) 81--82.  %%QUATERNION, %%MECHANICS, %%SPINOR, %$D09022002.

\bibitem{GURSE1957B} F. G\"ursey, \emph{Relativistic kinematics of a classical point particle in spinor form}, Nuov. Cim. {\bf 5} (1957) 784--809.  %%QUATERNION, %%MECHANICS, %$D09022002.

\bibitem{BLASC1960-} W. Blaschke, Kinematik und Quaternionen (Deutscher Verlag der Wissenschaften, Berlin, 1960) 84~pp. %%BOOK, %%QUATERNION, %%MECHANICS, %$D20032002.

\bibitem{KUSTA1965-} P. Kustaanheimo and E. Stiefel, \emph{Perturbation theory of Kepler motion based on spinor regularization}, J. f\"ur die reine und  angew. Math. {\bf 218} (1965) 204--219.  %%QUATERNION, %%MECHANICS, %%SPINOR, %$D15032002. 

\bibitem{KYRAL1967B} A. Kyrala, \emph{An alternative derivation of relativistic mechanics}, Section 8.9 of Theoretical Physics: Applications of Vectors, Matrices, Tensors and Quaternions (W.B. Saunders, Philadelphia, 1967) 270--276.  %%QUATERNION, %%MECHANICS, %$D09022002.

\bibitem{ROCHE1972-} E.Y. Rocher, \emph{Noumeon: elementary entity of a new mechanics}, J. Math. Phys. {\bf 13} (1972) 1919--1925.  %%QUATERNION, %%MECHANICS, %$D06072008. %$C Exemple d'une aberration autoris�e (provoqu�e?) par l'arithm�tique mais sans signification physique. D'un autre c�t�, c'est cet article  qui m'a mis sur la voie du Rumer-quaternionien en le citant...

\bibitem{HESTE1974A} D. Hestenes, \emph{Proper particle mechanics}, J. Math. Phys. {\bf 15} (1974) 1768--1777. %%QUATERNION, %%MECHANICS, %%HESTENES, %$D09022002.

\bibitem{HESTE1974B} D. Hestenes, \emph{Proper dynamics of a rigid point particle}, J. Math. Phys. {\bf 15} (1974) 1778--1786. %%QUATERNION, %%MECHANICS, %%HESTENES, %$D09022002.

\bibitem{GOLDS1975-} H. Goldstein, \emph{Prehistory of the ``Runge-Lenz'' vector}, Am. J. Phys. {\bf 43} (1975) 737--738. %%QUATERNION, %%MECHANICS, %%KEPLER, %%LAPLACE-RUNGE-LENZ, %$D30112005.

\bibitem{GOLDS1976-} H. Goldstein, \emph{More on the prehistory of the Laplace or Runge-Lenz vector}, Am. J. Phys. {\bf 44} (1976) 1123--1124. %%QUATERNION, %%MECHANICS, %%KEPLER, %%LAPLACE-RUNGE-LENZ, %$D30112005.

\bibitem{CHELN1980-} Iu. N. Chelnokov, \emph{On integration of kinematic equations of a rigid body's screw-motion}, PMM: J. of Appl. Math. \& Mechanics {\bf 44} (1980) 19--23. %%QUATERNION, %%MECHANICS, %$D09022002, %$C Biquaternions of Clifford, not Hamilton's complexified quaternions, see \cite{CLIFF1873-}.

\bibitem{CORNI1983-} F.H.J. Cornish, \emph{Kepler orbits and the harmonic oscillator}, J. Phys. A: Math. Gen. {\bf 17} (1984) 2191--2197. %%QUATERNION, %%MECHANICS, %$D09022002.

\bibitem{VIVAR1983-} M.D. Vivarelli, \emph{The KS (Kustaanheimo-Stiefel) transformation in hypercomplex form}, Celestial Mechanics {\bf 29} (1983) 45--50.  %%QUATERNION, %%MECHANICS, %$D09022002.

\bibitem{CICOG1985A-} G. Cicogna, \emph{On the quaternionic bifurcation}, J. Phys. A: Math. Gen. {\bf 18} (1985) L829--L832. %%QUATERNION, %%MECHANICS, %%GROUP-THEORY, %$D30112005.

\bibitem{VIVAR1985-} M.D. Vivarelli, \emph{The KS (Kustaanheimo-Stiefel) transformation in hypercomplex form and the quantization of the negative-energy orbit manifold in the Kepler problem}, Celestial Mechanics {\bf 36} (1985) 349--364.  %%QUATERNION, %%MECHANICS, %%SPINORS, QUANTUM-PHYSICS, %$D24062002

\bibitem{CICOG1985B-} G. Cicogna and G. Gaeta, \emph{Periodic-solutions from quaternionic bifurcation}, Lett. Nuovo Cim. {\bf 44} (1985) 65--68. %%QUATERNION, %%MECHANICS, %%GROUP-THEORY, %$D30112005.

\bibitem{HESTE1986E} D. Hestenes, \emph{Spinor mechanics and perturbation theory}, 1--23, in: D Hestenes.  New Foundation for Classical Mechanics (Reidel, Dordrecht, 1986) 564--573. %%QUATERNION, %%MECHANICS, %%SPINOR, %%HESTENES, %$D10022002.

\bibitem{HESTE1986F} D. Hestenes,  New Foundation for Classical Mechanics (Reidel, Dordrecht, 1986) 644~pp. %%BOOK, %%QUATERNION, %%MECHANICS, %%HESTENES, %%TABLE-OF-CONTENT, %$D10022002.

\bibitem{AGRAW1987A} O.P. Agrawal, \emph{Hamilton operators and dual-number quaternions in spatial kinematics}, Mechanisms and Machine Theory {\bf 22} (1987) 569--575. %%QUATERNION, %%MECHANICS, %$D06072008.   

\bibitem{AGRAW1987B} O.P. Agrawal, \emph{Quaternions, Hamilton operators, and kinematics of mechanical systems},  Journal of Mechanisms, Transmission, and Automation: Advances in Design Automation (Robotics, Mechanisms, and Machine Systems), edited by S.S. Ras, {\bf 2}  (ASME, New York,1987) 317--322. %%QUATERNION, %%MECHANICS, %$D06072008.   

\bibitem{CICOG1987-} G. Cicogna and G. Gaeta, \emph{Quaternionic-like bifurcation in the absence of symmetry}, J. Phys. A: Math. gen. {\bf 20} (1987) 79--89. %%QUATERNION, %%MECHANICS, %%GROUP-THEORY, %$D30112005.

\bibitem{FIRNE1988-} F. Firneis, M. Firneis, L. Dimitrov, G. Frank, and R. Thaller, \emph{On some applications of quaternions in geometry}, Proceedings ICEGDG {\bf 1} (Techn. University, Vienna, 1988) 158--164.  %%QUATERNION, %%MECHANICS, %%KINEMATICS, %$D30112005.

\bibitem{CAMPI1991-} C. Campigotto, \emph{The Kustaanheimo-Stiefel transformation, the hydrogen-oscillator connection and orthogonal polynomial}, in: A. Ronvaux and D. Lambert, Le Probl\`eme de factorisation de Hurwitz (Universit\'e de Namur, 1991) 29~pp. %%QUATERNION, %%MECHANICS, %$D09022002.

\bibitem{SMOLI1991-} A.L. Smolin, \emph{Hypercomplex equations of dynamics}, Sov. Phys. J. {\bf 34} (1991) 79--81. %%QUATERNION, %%MECHANICS, %$D09022002.

\bibitem{KUTRU1992-} V.N. Kutrunov, \emph{The quaternion method of regularizing integral equations of the theory of elasticity}, J. Appl. Maths. Mechs. {\bf 56} (1992) 765--770.  %%QUATERNION, MECHANICS, %%ANALYTICITY-H, %%ELASTICITY, %$D29052001.

\bibitem{HESTE1993B} D. Hestenes, \emph{Hamiltonian mechanics with geometric calculus}, in: Z. Oziewicz et al., eds., Spinors, Twistors, Clifford Algebras and Quantum Deformations (Kluwer Academic Publishers, Dordrecht, 1993) 203-214. %%QUATERNION, %%MECHANICS, %%HESTENES, %$D10022002.

\bibitem{RUSSO1993-} F. Russo Spena, \emph{A note on quaternion algebra and finite rotations}, Nuovo Cim. {\bf 108 B} (1993) 689--698.  %%QUATERNION, %%MECHANICS, %%KINEMATICS, %$D30112005. %$C Kinematics of deformable bodies.

\bibitem{VOLD-1993A} T.G. Vold, \emph{An introduction to geometric algebra with an application in rigid body mechanics}, Am. J. Phys. {\bf 61} (1992) 491--504. %%QUATERNION, %%MECHANICS, %%HESTENES, %$D17022002.

\bibitem{VRBIK1994A} J. Vrbik, \emph{Celestial mechanics via quaternions}, Can. J. Phys. {\bf 72} (1994) 141--146.  %%QUATERNION, %%MECHANICS, %%KEPLER, %$D30112005.

\bibitem{VRBIK1995-} J. Vrbik, \emph{Perturbed Kepler-problem in quaternionic form}, J. Phys. A: Math. Gen. {\bf 28} (1995) 6245--6252. %%QUATERNION, %%MECHANICS, %%KEPLER, %$D30112005.

\bibitem{HESTE1998A} D. Hestenes, \emph{Spinor particle mechanics}, Fundamental Theories of Physics {\bf 94} (1998) 129--143. %%QUATERNION, %%MECHANICS, %%HESTENES, %$D09022002.

\bibitem{KOSEN1998-} I.I. Kosenko, \emph{Integration of the equations of a rotational motion of a rigid body in quaternion algebra. The Euler case}, J. of Applied Mathematics and Mechanics {\bf 62} (1998) 193--200.  %%QUATERNION, %%MECHANICS, %$D04032006. 

\bibitem{HESTE2002B} D. Hestenes and E.D. Fasse, \emph{Homogeneous rigid body mechanics with elastic coupling}, in: L. Dorst et al., eds, Applications of Geometric Algebra in Computer Science and Engineering (Birkh\"auser, Boston, 2002) 197--212. %%QUATERNION, %%MECHANICS, %%HESTENES, %$D10022002. %$C Preprint title \emph{Modeling elastically coupled rigid bodies with geometric algebra}.

\bibitem{GAETA2002A} G. Gaeta and P. Morando, \emph{Hyper-Hamiltonian dynamics},  J. Phys. A {\bf 35} (2002) 3925--3943. %%QUATERNION, %%MECHANICS, %%MANIFOLD, %%KAHLER, %$D04032006.

\bibitem{GAETA2002B} G. Gaeta and P. Morando, \emph{Quaternionic integrable systems}, in:  S. Abenda, G. Gaeta, and S. Walcher, eds.,  Symmetry and perturbation theory, SPT 2002, Cala Gonone (World Sci. Publishing, River Edge, NJ, 2002) 72--81. %%QUATERNION, %%MECHANICS, %%MANIFOLD, %%KAHLER, %$D04032006.

\bibitem{MORAN2003-} P. Morando and M. Tarallo, \emph{Hyper-Hamiltonian dynamics and quaternionic regularity}, Mod. Phys. Lett. A {\bf 18} (2003) 1841--1847. %%QUATERNION, %%MECHANICS, %%MANIFOLD, %%KAHLER, %$D04032006.

\bibitem{VRBIK2003-} J. Vrbik, \emph{A novel solution to Kepler's problem}, Eur. J. Phys. {\bf 24}  (2003) 575--583. %%QUATERNION, %%MECHANICS, %%KEPLER, %$D04032006.

\bibitem{MEIST2005-} L. Meister and H. Schaeben, \emph{A concise quaternion geometry of rotations}, Math. Meth. Appl. Sci. {\bf 28} (2005) 101--126.  %%QUATERNION, %%MECHANICS, %$D06072008.


\end{enumerate}


\subsection{HYDRODYNAMICS}
%-------------------------
\label{HYDRODYNAMICS}

\begin{enumerate}

\bibitem{ROSE-1950-} A. Rose, \emph{On the use of a complex (quaternion) velocity potential in the three dimensions}, Comm. Math. Helv. {\bf 24} (1950) 135--148. %%QUATERNION, %%HYDRODYNAMICS, %%ANALYTICITY-H, %$D29052001.

\bibitem{GURLE1990A} K. G\"urlebeck and W. Spr\"ossig, \emph{A quaternionic treatment of Navier-Stokes equations}, Rendiconti Circ. Mat Palermo--Suppl. {\bf 22} (1990) 77--95. %%QUATERNION, %%HYDRODYNAMICS, %%ANALYTICITY-H, %$D29052001.

\bibitem{SPROS1992-} W. Spr\"ossig and K. G\"urlebeck, \emph{Application of quaternionic analysis on generalized non-linear Stokes eigenvalue problems}, in: H. Begehr and A. Jeffrey, eds., Partial Differential Equations With Complex Analysis, Pitnam Research Notes in Math. {\bf 262} (Longman, Burnt Hill, 1992) 52--60.  %%QUATERNION, %%HYDRODYNAMICS, %%ANALYTICITY-H, %$D29052001.

\bibitem{NAGEM1998-} R.J. Nagem, C. Rebbi, G. Sandri, and S. Shei, \emph{Gauge transformations and local conservation equations for linear acoustics and for Maxwell's equations}, Nuovo Cim.  {\bf 113B} (1998) 1509--1517. %%QUATERNION, %%HYDRODYNAMICS, %%ACOUSTICS, %%PHYSICS-VARIA, %%ELECTRODYNAMICS, %$D04032006, %$C Rediscovery of Lanczos's Lagrangian of electrodynamics, i.e., \cite{LANCZ1949-} page 384.

\bibitem{SPROS2000-} W. Spr\"ossig, \emph{Quaternionic analysis in fluid mechanics}, in: J. Ryan and W. Spr\"ossig, eds., Clifford Algebra and their Applications in Mathematical Physics, Vol.~2: \emph{Clifford Analysis} (Birkh\"auser, Boston, 2000) 37--53. %%QUATERNION, %%HYDRODYNAMICS, %%ANALYTICITY-H, %$D29052001.

\bibitem{GIBBO2002-} J.D. Gibbon, \emph{A quaternionic structure in the three-dimensional Euler and equations for ideal MHD}, Physica D {\bf 166}, (2002) 17--28.  %%QUATERNION, %%HYDRODYNAMICS, %$D06072008. 

\bibitem{TANIS2003B} M. Tanisli, \emph{The quaternionic energy conservation equation for acoustics}, Acta Physica Slovaca {\bf 53} (2003) 253--258. %%QUATERNION, %%HYDRODYNAMICS, %%ACOUSTICS, %%PHYSICS-VARIA, %$D04032006.

\bibitem{GIBBO2006-} J.D. Gibbon, D.D. Holm, R.M. Kerr and I. Roulstone, \emph{Quaternions and particle dynamics in the Euler fluid equations}, Nonlinearity {\bf 19} (2006) 1969--1983. %%QUATERNION, %%HYDRODYNAMICS, %$D06072008.  


 
\end{enumerate}


\subsection{ELECTRODYNAMICS}
%---------------------------
\label{ELECTRODYNAMICS}

Classical electrodynamics.

\begin{enumerate}

\bibitem{MAXWE1873-} J.C. Maxwell, {Treatise on Electricity and Magnetism} (1873).  %%QUATERNION, %%ELECTRODYNAMICS, %$D09022002.

\bibitem{MAXWE1885-} J.C. Maxwell, \emph{Trait\'e d'\'electricit\'e et de magn\'etisme}, 2 volumes (Gauthier-Villars, Paris, 1885, 1887).  %%QUATERNION, %%ELECTRODYNAMICS, %$D09022002.

\bibitem{SARRA1887-} M. Sarrau, \emph{Note sur la th\'eorie des quaternions}, in: J.C. Maxwell, Trait\'e d'\'Electricit\'e et de Magn\'etisme, Vol.~II (1887) 591--632.  %%QUATERNION, %%ELECTRODYNAMICS, %$D09022002.

\bibitem{SILBE1907A} L. Silberstein, \emph{Elektromagnetische Grundgleichungen in bivectorieller Behandlung}, Ann. der Phys {\bf 22} (1907) 579--586. %%QUATERNION, %%ELECTRODYNAMICS, %$D04032006.

\bibitem{SILBE1907B} L. Silberstein, \emph{Nachtrag sur Abhandlung \"uber ``Elektromagnetische Grundgleichungen in bivectorieller Behandlung,''} Ann. der Phys {\bf 22} (1907) 783--784. %%QUATERNION, ELECTRODYNAMICS, %$D04032006.

\bibitem{CONWA1907-} A.W. Conway, \emph{A theorem on moving distributions of electricity}, Proc. Roy. Irish Acad. {\bf 27} (1907) 1--8.  %%QUATERNION, %%ELECTRODYNAMICS, %$D09022002.

\bibitem{CONWA1908-} A.W. Conway, \emph{The dynamics of a rigid electron}, Proc. Roy. Irish Acad. {\bf 27} (1908) 169--181.  %%QUATERNION, %%ELECTRODYNAMICS, %$D09022002.

\bibitem{CONWA1910-} A.W. Conway, \emph{On the motion of an electrified sphere}, Proc. Roy. Irish Acad. {\bf 28} (1910) 1--15.  %%QUATERNION, %%ELECTRODYNAMICS, %$D09022002.

\bibitem{CONWA1911-} A.W. Conway, \emph{On the application of quaternions to some recent developments of electrical theory}, Proc. Roy. Irish Acad. {\bf 29} (1911) 1--9. %%QUATERNION, %%ELECTRODYNAMICS, %%MAXWELL, %%SPECIAL-RELATIVITY, %$D09022002.

\bibitem{SILBE1912-} L. Silberstein, \emph{Quaternionic form of relativity}, Phil. Mag. {\bf 23} (1912) 790--809.  %%QUATERNION, %%ELECTRODYNAMICS, %%SPECIAL-RELATIVITY, %%MAXWELL, %$D09022002.

\bibitem{CONWA1912-} A.W. Conway, \emph{The quaternion form of relativity}, Phil. Mag. {\bf 24} (1912) 208.  %%QUATERNION, %%ELECTRODYNAMICS, %%SPECIAL-RELATIVITY, %%MAXWELL, %$D09022002.

\bibitem{SILBE1913-} L. Silberstein, \emph{Second memoir on quaternionic relativity}, Phil. Mag. {\bf 25} (1913) 135--144.  %%QUATERNION, %%ELECTRODYNAMICS, %%SPECIAL-RELATIVITY, %%MAXWELL, %$D09022002. %$C Definition of the correct energy-momentum tensor of the enelctromagnetic field.

\bibitem{JUVET1932-} G. Juvet and A. Schidlof, \emph{Sur les nombres hypercomplexes de Clifford et leurs applications \`a l'analyse vectorielle ordinaire, \`a l'\'electromagn\'etisme de Minkowski et \`a la th\'eorie de Dirac}, Bull. Soc. Sci. Nat. Neuch\^atel {\bf 57} (1932) 127--141.  %%QUATERNION, %%ELECTRODYNAMICS, %%MAXWELL, %%MONOPOLE, %%DIRAC, %%CLIFFORD, %%PROCA, %$D05022002.

\bibitem{WATSO1936-} W.H. Watson, \emph{Note on the representation of electromagnetic fields by biquaternions}, Trans. Roy. Soc. Canada {\bf 30} (1936) 105--113. %%QUATERNION, %%ELECTRODYNAMICS, %%BORN-ELECTRODYNAMICS, %$D01022002.

\bibitem{WATSO1937-} W.H. Watson, \emph{On a system of functional dynamics and optics}, Phil. Trans. Roy. Soc. {\bf A 236} (1937) 155--190. %%QUATERNION, %%ELECTRODYNAMICS, %%MECHANICS, %%BORN-ELECTRODYNAMICS, %%DIFFERENTIAL-FORM, %$D30112005.

\bibitem{RAO--1938M} B.S.M. Rao, \emph{Biquaternions in Born's electrodynamics}, Proc. Indian Acad. Sci.  {\bf 7} (1938) 333--338. %%QUATERNION, %%ELECTRODYNAMICS, %%BORN-ELECTRODYNAMICS, %%CONSTITUTIVE-TENSOR, %$D01022002.

\bibitem{WEISS1941-} P. Weiss, \emph{On some applications of quaternions to restricted relativity and classical radiation theory}, Proc. Roy. Irish Acad. {\bf A 46} (1941) 129--168. %%QUATERNION, %%ELECTRODYNAMICS, %%SPINOR, %%LORENTZ-DIRAC, %$D09022002.

\bibitem{MERCI1949-} A. Mercier, \emph{Sur les fondements de l'\'electrodynamique classique (m\'ethode axiomatique)}, Arch. Sci. Phys. Nat. Gen\`eve {\bf 2} (1949) 584--588.  %%QUATERNION, %%ELECTRODYNAMICS, %%CLIFFORD, %%MONOPOLE, %$D09022002. %$C No magnetic monopoles!

\bibitem{GURSE1954-} F. G\"ursey, \emph{Dual invariance of Maxwell's tensor}, Rev. Fac. Sci. Istanbul {\bf A 19} (1954) 154--160.   %%QUATERNION, %%ELECTRODYNAMICS, %$D09022002. %$C Energy-momentum tensor.

\bibitem{KYRAL1961-} A. Kyrala, \emph{An approach to the unification of classical, quantum and relativistic formulations of electromagnetics and dynamics}, Acta Phys. Austriaca {\bf 14} (1961) 448--459.  %%QUATERNION, %%ELECTRODYNAMICS, %%MECHANICS, %$D09022002.

\bibitem{BAYLI1989C}  W.E. Baylis, \emph{Relativistic dynamics of charges in external fields: The Pauli-algebra approach}, J. Phys. A: Math. Gen. {\bf 22} (1989) 17--29. %%QUATERNION, %%ELECTRODYNAMICS, %%SPECIAL-RELATIVITY, %$D16032002.

\bibitem{CORNB1992-} S. Cornbleet, \emph{An electromagnetic theory of rays in a nonuniform medium},  in: H. Blok et al., eds., Stud. in Math. Phys. {\bf 3} \emph{Huygens's principle 1690--1990. Theory and applications} (North Holland, Amsterdam, 1992)  451--459.  %%QUATERNION, %%ELECTRODYNAMICS, %%CONFORMALITY, %%MAXWELL, %$D13032002.

\bibitem{BAYLI1993-}  W.E. Baylis, \emph{Light polarization: A geometric-algebra approach}, Am. J. Phys. {\bf 61} (1993) 534--545. %%QUATERNION, %%ELECTRODYNAMICS, %%POLARIZATION, %$D16032002.

\bibitem{HILLI1993B}  P. Hillion, \emph{Spinor electromagnetism in isotropic chiral media}, Adv. Appl. Clifford Alg. {\bf 3} (1993) 107--120. %%QUATERNION, %%ELECTRODYNAMICS, %%CONSTITUTIVE-TENSOR, %%EINSTEIN-MAYER, %$D24062002.

\bibitem{KAUFF1993-} T. Kauffmann and W.Y. Sun, \emph{Quaternion mechanics and electromagnetism}, Ann. Fond. L. de Broglie {\bf 18} (1993) 213--291. %%QUATERNION, %%ELECTRODYNAMICS, %%MECHANICS, %$D09022002.  %$C Inspired by \cite{KYRAL1961-, KYRAL1967C}. 

\bibitem{VOLD-1993B} T.G. Vold, \emph{An introduction to geometric algebra and its application to electrodynamics}, Am. J. Phys. {\bf 61} (1992) 505--513. %%QUATERNION, %%ELECTRODYNAMICS, %%HESTENES, %$D17022002.

\bibitem{HILLI1995-}  P. Hillion, \emph{Constitutive relations and Clifford algebra in electromagnetism}, Adv. Appl. Clifford Alg. {\bf 5} (1995) 141--158. %%QUATERNION, %%ELECTRODYNAMICS, %%CONSTITUTIVE-TENSOR, %$D24062002, %$C Formulas (16a,16b) are the order reversal covariant form of Maxwell's equations.

\bibitem{SWEET2001-} D. Sweetser and G. Sandri, \emph{Maxwell's vision: electromagnetism with Hamilton's quaternions}, in: S. Marchiafava et al., eds., Proceedings of the 2nd meeting ``Quaternionic structures in mathematics and physics'' (World Scientific, Singapore, 2001) 417--420. %%QUATERNION, %%ELECTRODYNAMICS, %$D04032006.

\bibitem{CHERN2002-} A.A. Chernitskii, \emph{Born-Infeld electrodynamics: Clifford number and spinor representations}, Int. J. of Math. and Math. Sci. {\bf 31} (2002) 77--84.  %%QUATERNION, %%ELECTRODYNAMICS, %%BORN-ELECTRODYNAMICS, %%CLIFFORD, %$D30112005.

\bibitem{KRAVC2002C} V.V. Kravchenko, \emph{On a quaternionic reformulation of Maxwell's equations for inhomogeneous media and new solutions}, Z. f. Anal. u. ihre Anwend. {\bf 21} (2002) 21--26. %%QUATERNION, %%ELECTRODYNAMICS, %%MAXWELL, %%CONSTITUTIVE-TENSOR, %$D06022002.

\bibitem{CHERN2003-} A.A. Chernitskii, \emph{Source function and dyon's field in Clifford number representation for electrodynamics}, Adv. Appl. Clifford Alg. {\bf 13} (2003) 219--230.  %%QUATERNION, %%ELECTRODYNAMICS, %%MONOPOLE, %%CLIFFORD, %$D30112005.

\bibitem{GSPON2003J} A. Gsponer, \emph{What is spin?}, Report ISRI-03-10 (10 September 2003) 6~pp.; e-print \underline{ arXiv:physics/0308027 }.  %%QUATERNION, %%ELECTRODYNAMICS, %%QUANTUM-PHYSICS, %%DIRAC, %%PROCA, %%LANCZOS, %$D03082005.

\bibitem{KHMEL2003-} K.V. Khmelnytskaya, V. Kravchenko, and V.S. Rabinovich, \emph{Quaternionic fundamental solutions for the numerical analysis of electromagnetic scattering problems}, Z. f\"ur Anal. und ihhre Anwendungen {\bf 22} (2003) 659--589. %%QUATERNION, %%ELECTRODYNAMICS, %$D30112005.

\bibitem{KRAVC2003A} V. Kravchenko, \emph{Quaternionic equation for electromagnetic fields in inhomogeneous media}, Progress in Analysis {\bf I,II} (World Scientific, River Edge, NJ, 2003) 361--366. %%QUATERNION, %%ELECTRODYNAMICS, %$D30112005.

\bibitem{GRUDS2004-} S.M. Grudsky, K.V. Khmelnytskaya, and V.V. Kravchenko, \emph{On a quaternionic Maxwell equation for the time-dependent electromagnetic field in a chiral medium}, J. Phys. A: Math. Gen. {\bf 37} (2004) 4641--4647. %%QUATERNION, %%ELECTRODYNAMICS, %%MAXWELL, %%CONSTITUTIVE-TENSOR, %$D30112005.

\bibitem{GSPON2004E} A. Gsponer and J.-P. Hurni, \emph{Cornelius Lanczos's derivation of the usual action integral of classical electrodynamics}, Foundations of Physics {\bf  35} (2005) 865--880; e-print \underline{ arXiv:math-ph/0408100 }.  %%QUATERNION, %%ELECTRODYNAMICS, %%RENORMALIZATION, %%LANCZOS, %$D03082005.

\bibitem{ACEVE2005-} M. Acevedo, J. Lopez-Bonilla, and M. Sanchez-Meraz, \emph{Quaternion, Maxwell equations and Lorentz transformations}, Apeiron {\bf 12} (2005) 371--384. %%QUATERNION, %%ELECTRODYNAMICS, %$D30112005. %$C Debye potential.

\bibitem{GSPON2006C} A. Gsponer, \emph{The locally-conserved current density of the Lienard-Wiechert field}, Report ISRI-06-03 (29 January 2007) 7~pp.; e-print \underline{ arXiv:physics/0612090 }. %%QUATERNION, %%ELECTRODYNAMICS, %$D06072008.
   
\bibitem{GSPON2006D} A. Gsponer, \emph{Derivation of the potential, field, and locally-conserved charge-current density of an arbitrarily moving point-charge}, Report ISRI-06-04 (29 January 2007) 19~pp.; e-print \underline{ arXiv:physics/0612232 }. %%QUATERNION, %%ELECTRODYNAMICS, %$D06072008.

\bibitem{TANIS2006-} M. Tanisli, \emph{Gauge transformation and electromagnetism with biquaternions}, Europhys. Lett. {\bf 74} (2006) 569--573.   %%QUATERNION, %%ELECTRODYNAMICS, %%MAXWELL, %%LANCZOS, %$D06072008.

\bibitem{PRIVA2007-} A.N. Privalchuk and E.A. Tolkachev, \emph{Reduction of nonlinear equations of the non-commutative electrodynamics in  quaternion formulation}, submitted to SIGMA: Symmetry, Integrability and Geometry: Methods and Applications (2007) 4~pp. %%QUATERNION, %%ELECTRODYNAMICS, %$D06072008.

\bibitem{BISHT2007-} P.S. Bisht and O.P.S. Negi, \emph{Revisiting quaternionic dual electrodynamics} (2007) 15~pp.; e-print \underline{ arXiv:0709.0088 }. %%QUATERNION, %%ELECTRODYNAMICS, %$D06072008.

\bibitem{CATON2007-} F. Catoni, \emph{Commutative (Segre's) Quaternion Fields and Relation with Maxwell Equation},  Advances in Applied Clifford Algebras, published on-line (SpringerLink, August 28, 2007). %%QUATERNION, %%ELECTRODYNAMICS, %$D06072008. 

 
\end{enumerate}


\subsection{LEPTODYNAMICS}
%-------------------------
\label{LEPTODYNAMICS}

Weak interactions of leptons and quarks, electroweak models and interactions.


\begin{enumerate}

\bibitem{FEYNM1958-} R.P. Feynman and M. Gell-Mann, \emph{Theory of the Fermi interaction}, Phys. Rev. {\bf 109} (1958) 193--198. %%NO-QUATERNION, %%LEPTODYNAMICS, %%DIRAC, %%TWO-COMPONENT-SPINOR, %%$D30112005. %$C Reference 2 is citing \cite{FEYNM1951-}. 

\bibitem{GURSE1961-} F. G\"ursey, \emph{On the structure and parity of weak interaction currents}, Ann. Phys. {\bf 12} (1961) 91--117. %%QUATERNION, %%LEPTODYNAMICS, %%PARTICLE-PHYSICS, %%ISOSPIN, %%SPIN-3/2, %$D16032002.

\bibitem{NDILI1976-} F.N. Ndili, \emph{Spontaneous symmetry breaking with quaternionic scalar fields and electron-muon mass ratio}, Int. J. Theor. Phys. {\bf 15} (1976) 265--268.   %%QUATERNION, %%LEPTODYNAMICS, %$D30112005.

\bibitem{CHKAR1978-}  J.L. Chkareuli, \emph{Leptons and quarks in the quaternion model}, JETP Lett. {\bf 27} (1978) 557--561. %%QUATERNION, %%LEPTODYNAMICS, %$D27062001.

\bibitem{WARD-1978-}  B.F.L. Ward, \emph{Weakly coupled fields: a more active view},  Nuovo Cim.  {\bf 38 A } (1978) 299--313. %%QUATERNION, %%LEPTODYNAMICS, %$D16032002.

\bibitem{CHKAR1979-}  J.L. Chkareuli, \emph{CP violation and the Cabibbo angle in the quaternion model}, JETP Lett. {\bf 29} (1979) 148--151.  %%QUATERNION, %%LEPTODYNAMICS, %$D27062001.

\bibitem{WARD-1979-}  B.F.L. Ward, \emph{Quarks, quaternions and weakly coupled fields},  Nuovo Cim.  {\bf 51 A } (1979) 208--218. %%QUATERNION, %%LEPTODYNAMICS, %$D16032002. %$C Both "real Minkowski space" and "Euclidian 4-space" may be regarded as "real slices" of complexified space-time.

\bibitem{CHKAR1981-}  J.L. Chkareuli, \emph{The weak interaction of leptons and quarks in the quaternionic model}, Sov. J. Nucl. Phys. {\bf 34} (1981) 258--265.  %%QUATERNION, %%LEPTODYNAMICS, %$D27062001. 

\bibitem{MORIT1981-} K. Morita, \emph{Gauge theories over quaternions and Weinberg-Salam}, Prog. Th. Phys. {\bf 65} (1981) 2071--2074. %%QUATERNION, %%LEPTODYNAMICS, %$D16032002.

\bibitem{HESTE1982-} D. Hestenes, \emph{Space-time structure of weak and electromagnetic interactions}, Found. Phys. {\bf 12} (1982) 153--168. %%QUATERNION, %%LEPTODYNAMICS, %%HESTENES, %$D10022002.

\bibitem{MORIT1982A} K. Morita, \emph{Quaternionic Weinberg-Salam theory}, Prog. Th. Phys. {\bf 67} (1982) 1860--1876. %%QUATERNION, %%LEPTODYNAMICS, %$D15032002.

\bibitem{MORIT1993-} K. Morita, \emph{Quaternion and non-commutative geometry}, Prog. Th. Phys. {\bf 90} (1993) 219--236. %%QUATERNION, %%LEPTODYNAMICS, %%NONCOMMUTATIVE-GEOMETRY, %$D16032002.

\bibitem{MORIT1994-} K. Morita, \emph{Weinberg-Salam theory in non-commutative geometry}, Prog. Th. Phys. {\bf 91} (1994) 959--974. %%QUATERNION, %%LEPTODYNAMICS, %%NONCOMMUTATIVE-GEOMETRY, %$D16032002.

\bibitem{DELEO1995B} S. DeLeo and P. Rotelli, \emph{Quaternion Higgs and the electroweak gauge group}, Int. J. Mod. Phys. {\bf A 10} (1995) 4359--4370. %%QUATERNION, %%LEPTODYNAMICS, %$D16032002.

\bibitem{DELEO1996F} S. DeLeo and P. Rotelli, \emph{Quaternion electroweak theory}, J. Phys. G: Nucl. part. Phys. {\bf 22} (1966) 1137--1150. %%QUATERNION, %%LEPTODYNAMICS, %$D16032002.

\bibitem{BARUT1997-}  A.O. Barut, \emph{Neutrinos and electromagnetic fields}, Adv. Appl. Clifford Alg. {\bf 7 (S)} (1997) 357--367. %%QUATERNION, %%LEPTODYNAMICS, %%PARTICLE-PHYSICS, %$D17032002.

\bibitem{CASAN1991-} G. Casanova, \emph{Masses des neutrinos}, Unpublished report (Paris, March 1991) 15~pp. %%QUATERNION, %%LEPTODYNAMICS, %%PARTICLE-PHYSICS,  %$D16032002.

\bibitem{BOUDE1997A}  R. Boudet, \emph{The Glashow-Salam-Weinberg electroweak theory in the real algebra of spacetime}, Adv. Appl. Clifford Alg. {\bf 7 (S)} (1997) 321--336. %%QUATERNION, %%LEPTODYNAMICS, %$D16032002.

\bibitem{BOUDE1997B}  R. Boudet, \emph{The Takabayasi moving frame, from the A potential to the Z boson}, in: S. Jeffers et al., eds., The Present Status of the Theory of Light (Kluwer Acad. Pub., 1997) 471--481. %%QUATERNION, %%LEPTODYNAMICS, %$D16032002.

\bibitem{DELEO1997A} S. DeLeo, \emph{Quaternionic electroweak theory and Cabibbo-Kobayashi-Maskawa matrix}, Int. J. Th. Phys. {\bf 36} (1997) 1165--1177. %%QUATERNION, %%LEPTODYNAMICS, %$D16032002.

\bibitem{VELTM1997-} M. Veltman, \emph{Reflexions on the Higgs system}, Report 97-05 (CERN, 1997) 63~pp. %%BOOK, %%QUATERNION, %%LEPTODYNAMICS, %%GAUGE-THEORY,  %$D14052001, %$C Veltman uses quaternions in the form of 2x2 matrices.

\bibitem{BOUDE2003-} R. Boudet, \emph{Identification de la jauge $SU(2) \otimes U(1)$ de l'�lectrofaible � un produit de sous-groupes orthogonaux de l'espace-temps}, Ann. Fond. L. de Broglie {\bf 28} (2003) 315--330. %%QUATERNION, %%LEPTODYNAMICS, %$D06072008.

\end{enumerate}


\subsection{HADRODYNAMICS}
%-------------------------
\label{HADRODYNAMICS}

Strong interactions, piondynamics, quantum-chromodynamics (QCD).


\begin{enumerate}

\bibitem{YANG1954-} C.N. Yang and R.L. Mills, \emph{Conservation of isotopic spin and isotopic spin gauge invariance}, Phys. Rev. {\bf 96} (1954) 191--195. %%NO-QUATERNION, %%HADRODYNAMICS, %%PARTICLE-PHYSICS, %%GAUGE-THEORY, %$D30112005, %$C Perfect example of a paper that could have been written using quaternions instead of $2 \times 2$ unitary matrices. %%PARTICLE-PHYSICS??

\bibitem{SCHRE1955-} E.J. Schremp, \emph{Isotopic spin and the group space of the Lorentz group}, Phys. Rev. {\bf 99} (1955) 1603. %%QUATERNION, %%HADRODYNAMICS, %%DIRAC, %%PARTICLE-PHYSICS, %%ISOSPIN, %%LANCZOS, %$D05022002.

\bibitem{SCHRE1957-} E.J. Schremp, \emph{Parity nonconservation and the group-space of the proper Lorentz group}, Phys. Rev. {\bf 108} (1957) 1076-1077.  %%QUATERNION, %%HADRODYNAMICS, %%PARTICLE-PHYSICS, %%ISOSPIN, %%GURSEY, %%C Summary published in Bull. Am. Phys. Soc. Ser. II {\bf 2} (1957) 191, %$D24062002.

\bibitem{SCHWI1957-} J. Schwinger, \emph{A theory of fundamental interactions}, Ann. of Phys. {\bf 2} (1957) 407--434  %%NO-QUATERNION, %%HADRODYNAMICS, %%LEPTODYNAMICS, %%SIGMA-MODEL, %%PARTICLE-PHYSICS, %$D30112005. %$C This paper was received on July 31, 1957, while G\"ursey's paper \cite{GURSEY1958A} was received on November 5, 1957.  Both Schwinger an G\"ursey refer to B. Touscheck and W. Pauli's papers of 1957 in Nuovo cimento. %%PARTICLE-PHYSICS??

\bibitem{YANG1957-} C.N. Yang, \emph{Comments about symmetry laws}, Proc. 7th Rochester Conf. (April 15--19, 1957) IX-25 -- IX-26. %%QUATERNION, %%HADRODYNAMICS, %%PARTICLE-PHYSICS, %%COMPLEX-NUMBERS, %%CPT-THEOREM, %$D24062002.

\bibitem{RUELL1958-} D. Ruelle, \emph{Repr\'esentation du spin isobarique des particules \`a interactions fortes}, Nucl. Phys. {\bf 7} (1958) 443--450.  %%QUATERNION, %%HADRODYNAMICS, %%PARTICLE-PHYSICS, %$D15032002.

\bibitem{FEINB1959-} G. Feinberg and F. G\"ursey, \emph{Space-time properties and internal symmetries of strong interactions}, Phys. Rev. {\bf 114} (1959) 1153--1170. %%QUATERNION, %%HADRODYNAMICS, %$D03082005.

\bibitem{SCHRE1959-} E.J. Schremp, \emph{$G$-Conjugation and the group-space of the proper Lorentz group}, Phys. Rev. {\bf 113} (1959) 936--943. %%QUATERNION, %%HADRODYNAMICS, %%PARTICLE-PHYSICS, %%ISOSPIN, %%GURSEY, %$D05022002.

\bibitem{GELLM1960-} M. Gell-Mann and M. L\'evy, \emph{The axial vector current in beta decay}, Nuovo Cim. {\bf 16} (1960) 705--725. %%NO-QUATERNION, %%HADRODYNAMICS, %%LEPTODYNAMICS, %%SIGMA-MODEL, %%GURSEY, %$D30112005, %$C First renormalizable chiral-invariant version of the sigma-model.  %%LEPTODYNAMICS??

\bibitem{GURSE1960-} F. G\"ursey, \emph{On the symmetries of strong and weak interactions}, Nuovo Cim. {\bf 16} (1960) 230--240.  %%QUATERNION, %%HADRODYNAMICS, %%PARTICLE-PHYSICS, %%ISOSPIN, %%SIGMA-MODEL, %$D16032002, %$C First chiral-invariant exponential model of multiple-pion production/interaction. 

\bibitem{NAMBU1961A} Y. Nambu and G. Jona-Lasinio, \emph{Dynamical model of elementary particles based on an analogy with superconductivity. I}, Phys. Rev. {\bf 122} (1961) 345--358. %%NO-QUATERNION, %%HADRODYNAMICS, %%$D30112005.  %%PARTICLE-PHYSICS??

\bibitem{NAMBU1961B} Y. Nambu and G. Jona-Lasinio, \emph{Dynamical model of elementary particles based on an analogy with superconductivity. II}, Phys. Rev. {\bf 124} (1961) 246-254. %%NO-QUATERNION, %%HADRODYNAMICS, %%GURSEY %%$D30112005.  %%PARTICLE-PHYSICS??

\bibitem{SCHRE1962-} E.J. Schremp, \emph{Quaternion approach to elementary particle theory I}, NRL Quarterly on Nuclear Science and Technology (October 1962) 7--21.  %%QUATERNION, %%HADRODYNAMICS, %%PARTICLE-PHYSICS, %%ISOSPIN, %%SPECIAL-RELATIVITY, %%LANCZOS, %%CONWAY, $D24062002.

\bibitem{SCHRE1963-} E.J. Schremp, \emph{Quaternion approach to elementary particle theory II}, NRL Quarterly on Nuclear Science and Technology (January 1963) 1--21.  %%QUATERNION, %%HADRODYNAMICS, %%PARTICLE-PHYSICS, %%ISOSPIN, %%SPECIAL-RELATIVITY, %%LANCZOS, $D24062002.

\bibitem{HESTE1967B} D. Hestenes, \emph{Spin and isospin}, J. Math. Phys. {\bf 8} (1967) 809--812.  %%QUATERNION, %%HADRODYNAMICS, %%PARTICLE-PHYSICS, %%HESTENES, %$D10022002.

\bibitem{GURSEY1968-} F. G\"ursey, \emph{Effective Lagrangians in particle physics}, Acta. Phys. Austr. Suppl. {\bf 5} (1968) 185--225.  %%QUATERNION, %%HADRODYNAMICS, %%SIGMA-MODEL, %$D03082005.

\bibitem{SKYRM1971-} T.H.R. Skyrme, \emph{Kinks and the Dirac equation}, J. Math. Phys. {\bf 12} (1971) 1735--1743. Reprinted in \cite{BROWN1994-}. %%QUATERNION, %%HADRODYNAMICS, %%SKYRME, %%DIRAC, %$D30112005.

\bibitem{EDMON1975C} J.D. Edmonds, Jr., \emph{Six bits for nine colored quarks}, Int. J. Th. Phys. {\bf 13} (1975) 431--435. %%QUATERNION, %%HADRODYNAMICS, %%PARTICLE-PHYSICS, %$D06022002.

\bibitem{ADLER1978-} S.L. Adler, \emph{Classical algebraic chromodynamics}, Phys. Rev. {\bf D 17} (1978) 3212--3224.  %%QUATERNION, %%HADRODYNAMICS, %%PARTICLE-PHYSICS, %$D15032002.

\bibitem{ADLER1979-} S.L. Adler, \emph{Algebraic chromodynamics}, Phys. Lett. {\bf 86 B} (1979) 203--205.  %%QUATERNION, %%HADRODYNAMICS, %%PARTICLE-PHYSICS, %$D15032002.

\bibitem{ADLER1980-} S.L. Adler, \emph{Quaternion chromodynamics as a theory of composite quarks and leptons}, Phys. Rev. {\bf D 21} (1980) 2903--2915. %%QUATERNION, %%HADRODYNAMICS, %%PARTICLE-PHYSICS, %$D15032002.

\bibitem{BIEDE1980-}  L.C. Biedenharn, D. Sepaneru, and L.P. Horwitz, \emph{Quaternionic quantum mechanics and Adler's chromostatics}, in: K.B. Wolf, ed., Lect. Notes in Phys. {\bf 135} (Springer, Berlin, 1980) 51--66. %%QUATERNION, %%HADRODYNAMICS, $$QQPH, %%PARTICLE-PHYSICS, %$D30112005.

\bibitem{REMBI1980A} J. Rembielinski, \emph{Quaternionic Hilbert space and colour confinement: I.}, J. Phys. {\bf A 13} (1980) 15--22. %%QUATERNION, %%HADRODYNAMICS, %%PARTICLE-PHYSICS, %$D15032002.

\bibitem{REMBI1980B} J. Rembielinski, \emph{Quaternionic Hilbert space and colour confinement: II. The admissible symmetry groups}, J. Phys. {\bf A 13} (1980) 23--30. %%QUATERNION, %%HADRODYNAMICS, %%PARTICLE-PHYSICS, %$D15032002.

\bibitem{REMBI1981-} J. Rembielinski, \emph{Algebraic confinement of coloured states}, J. Phys. {\bf A 14} (1981) 2609--2624. %%QUATERNION, %%HADRODYNAMICS, %%PARTICLE-PHYSICS, %$D15032002.

\bibitem{ZAHED1986-} I. Zahed and G.E. Brown, \emph{The Skyrme model}, Phys. Rep. {\bf 142} (1986) 1--102. Reprinted in \cite{BROWN1994-}. %%QUATERNION, %%HADRODYNAMICS, %%SKYRME, %$D04032006.

\bibitem{SANYU1990-} V.I. Sanyuk, \emph{Genesis and evolution of the Skyrme model from 1954 to the present}, Int. J. Mod. Phys. {\bf A7} (1992) 1--40.  Reprinted in \cite{BROWN1994-}. %%QUATERNION, %%HADRODYNAMICS, %%SKYRME, %$D30112005. 

\bibitem{CASAN1992-} G. Casanova, \emph{Th\'eorie relativiste du nucl\'eon et du doublet $\Xi$}, in: A. Micali et al., eds., Clifford Algebras and their Applications in Mathematical Physics (Kluwer Academic Publishers, Dordrecht, 1992) 353--361. %%QUATERNION, %%HADRODYNAMICS, %%PARTICLE-PHYSICS, %%ISOSPIN, %$D16032002.

\bibitem{BROWN1994-} G.E. Brown, ed., Selected Papers, with Commentary, of Tony Hilton Royle Skyrme (World Scientific, Singapore, 1994) 438~pp. %%BOOK, %%QUATERNION, %%HADRODYNAMICS, %%SKYRME, %$D04032006.

\bibitem{KRAVC1995A} V.V. Kravchenko, \emph{On a biquaternionic bag model}, Zeitschr. f\"ur Anal. und ihre Anwend. {\bf 14} (1995) 3--14.  %%QUATERNION, %%HADRODYNAMICS, %%PARTICLE-PHYSICS, %%ANALYTICITY-B, %$D15032002.

\bibitem{DAHM1997-} R. Dahm, \emph{Relativistic SU(4) and quaternions}, in: J. Keller and Z. Oziewicz, eds., \emph{The Theory of the Electron}, Adv. Appl. Clifford Alg.  {\bf 7 (S)}, 337--356. %%QUATERNION, %%HADRODYNAMICS, %%PARTICLE-PHYSICS, %$D15032002. 

\bibitem{DAHM1998-} R. Dahm, \emph{Complex quaternions in spacetime symmetry and relativistic spin-flavour supermultiplets}, Phys. of Atomic Nuclei {\bf 61} (1998) 1885--1891.  %%QUATERNION, %%HADRODYNAMICS, %%PARTICLE-PHYSICS, %$D15032002. 

\bibitem{KENNE2002-}  S. Kennedy, \emph{Geometric-algebra approach to the Weyl-Lanczos equation under the influence of the Li\'enard-Wiechert potential at nucleon distances}, Submitted to Electromagnetic Phenomena, Ukraine (9 September 2002) 13~pp.  %%QUATERNION, %%HADRODYNAMICS, %%LANCZOS, %%GSPONER, %$D30112005.


\end{enumerate}


\subsection{PARTICLE-PHYSICS}
%----------------------------
\label{PARTICLE-PHYSICS}
 
Papers dealing with elementary-particles physics and field-theory, as well as mathematical-physics papers primarily motivated by their relevance to elementary particle physics. 

Papers using OCTONION structures and fields are collected in Sec.~\ref{OCTONION}.

\begin{enumerate}

\bibitem{GURSE1956C} F. G\"ursey, \emph{On a conform-invariant spinor wave equation}, Nuovo Cim. {\bf 3} (1956) 988--1006. %%QUATERNION, %%PARTICLE-PHYSICS, %%CONFORMALITY, %%DIRAC, %$D13032002.

\bibitem{PAULI1957-} W. Pauli, \emph{On the conservation of the lepton charge}, Nuovo. Cim. {\bf 1} (1957) 204--215.  %%QUATERNION, %%PARTICLE-PHYSICS, %%LEPTODYNAMICS, %$D30112005. %$C This paper makes use of implicit properties of the Pauli algebra of complex 2x2 matrices which are explicit in the biquaternion algebera, as was shown by G\"ursey in \cite{GURSEY1958A, GURSEY1958B}.  %%LEPTODYNAMICS??

\bibitem{GURSE1958A} F. G\"ursey, \emph{Relation of charge independence and baryon conservation to Pauli's transformation}, Nuovo Cim. {\bf 7} (1958) 411--415.  %%QUATERNION, %%PARTICLE-PHYSICS, %%HADRODYNAMICS, %%LEPTODYNAMICS, %%ISOSPIN, %%DIRAC-LANCZOS, %%DYSON, %$D16032002, %$C G\"ursey makes the isospin decomposition of Lanczos's fundamental equation, Ref. \cite{LANCZ1929A}, without citing Lanczos (as he did in his Ph.D.\ thesis \cite{GURSE1950-}).  In fact G\"ursey shows how simple it is to rewrite and generalize Pauli's paper \cite{PAULI1957-} using biquaternions (in the form of ``Pauli matrices'').

\bibitem{GURSE1958B} F. G\"ursey, \emph{On the group structure of elementary particles}, Nucl. Phys. {\bf 8} (1958) 675--691.  %%QUATERNION, %%PARTICLE-PHYSICS, %%HADRODYNAMICS, %$D16032002.

\bibitem{HEISE1958-} W. Heisenberg and W. Pauli, \emph{On the isospin group in the theory of elementary particles}, Unpublished preprint (March 1958) Reprinted in: W. Blum, H.-P. D\"urr, and H. Rechenberg, eds., Werner Heisenberg Collected Works, Serie A / Part III (Springer Verlag, 1993) 337--351, with a postscript by W. Pauli on page 351.  %%QUATERNION, %%PARTICLE-PHYSICS, %%ISOSPIN, %%GURSEY, %$D30112005.  %$C Heisenberg and Pauli refer to G\"ursey's use of biquaternions (in the form of 2x2 matrices, in \cite{GURSE1958A}, to show how isospin arises from Lanczos's fundamental equation) with the sentence: ``G\"ursey has shown how the Dirac equation for a system of free neutrons and protons and neutrons can be written in a nonconventional way...'' In Ref. \cite{MARX1958-} G. Marx explains (page 346) that Heisenberg and Pauli are actually working with biquaternions.

\bibitem{MARX-1958-} G. Marx, \emph{On the second order wave equation of the fermions}, Nucl. Phys. {\bf 9} (1958/1959) 337--346.  %%QUATERNION, %%PARTICLE-PHYSICS, %%ELD, %%LED, %%HAD, %$D20032002. %$C On page 346 Marx refers to Ref. \cite{HEISE1958-} with the comment: ``Heisenberg considers $\xi$ not as complex but as quaternion operator.''

\bibitem{DEBRO1963A} L. deBroglie, D. Bohm, P. Hillion, F. Halbwachs, T. Takabayasi, and J.-P. Vigier, \emph{Rotator model of elementary particles considered as relativistic extended structures in Minkowski space}, Phys. Rev. {\bf 129} (1963) 438--450.  %%QUATERNION, %%PARTICLE-PHYSICS, %$D30112005.

\bibitem{DEBRO1963B} L. deBroglie, F. Halbwachs, P. Hillion, T. Takabayasi, and J.-P. Vigier, \emph{Space-time model of extended elementary particles in Minkowski space.  II. Free particles and interaction theory}, Phys. Rev. {\bf 129} (1963) 451--466.  %%QUATERNION, %%PARTICLE-PHYSICS, %$D30112005.

\bibitem{PAULI1959-} W. Pauli$^\dagger$ and B. Touschek, \emph{Report and comment on F.\ G\"ursey's ``Group structure of elementary particles''}, Supp. Nuovo. Cim. {\bf 14} (1959) 207--211.  %%QUATERNION, %%PARTICLE-PHYSICS, %%GURSEY, %$D30112005. %$C This is the last paper drafted by Pauli, published posthumously. Its aim, with reference to \cite{GUSEY1958B}, is to answer the questions: ``How do spin and isospin interplay, and is it reasonable to extrapolate isospin to the leptons?''

\bibitem{EBRAH1974A} A. Ebrahim and F. G\"ursey, \emph{General chiral $SU_2 \times SU_2$ Lagrangian and representation mixing}, Lett. Nuovo Cim. {\bf 9} (1974) 9--14, Errata 716.  %%QUATERNION, %%PARTICLE-PHYSICS, %$D03082005.  %%HADRODYNAMICS??

\bibitem{EBRAH1974B} A. Ebrahim and F. G\"ursey, \emph{Current-current Sugawara form and pion-kaon scattering}, Nuovo Cim. {\bf 21 A} (1974) 249--263.  %%QUATERNION, %%PARTICLE-PHYSICS, %$D03082005. %%HADRODYNAMICS??

\bibitem{CASAN1975B} G. Casanova, \emph{Equation relativiste du nucl\'eon et du doublet $\Xi$}, C.R. Acad. Sci. A {\bf 280} (1975) 1321--1324. %%QUATERNION, %%PARTICLE-PHYSICS, %%GURSEY, %%HESTENES, %$D04032006. 

\bibitem{CASAN1975C} G. Casanova, \emph{Existence et classification des particules}, C.R. Acad. Sci. A {\bf 281} (1975) 257--260. %%QUATERNION, %%PARTICLE-PHYSICS, %%HESTENES, %$D04032006. 

\bibitem{MINAM1975-} M. Minami, \emph{Quaternionic gauge-fields on $S_7$ and Yang's $SU(2)$ monopole}, Prog. Theor. Phys. {\bf 63} (1980) 303--321. %%QUATERNION, %%PARTICLE-PHYSICS, %%GAUGE-THEORY, %%MONOPOLE, %$D30112005.

\bibitem{WU---1975-} T.T. Wu and C.N. Yang, \emph{Concept of nonintegrable phase factors and global formulation of gauge fields}, Phys. Rev. {\bf D 12} (1975) 3845--3857.  %%QUATERNION, %%PARTICLE-PHYSICS, %%MONOPOLE, %%GAUGE-THEORY, %$D30112005.

\bibitem{CASAN1976A} G. Casanova, \emph{Diff\'erences de masse des multiplets baryoniques fondamentaux}, C.R. Acad. Sci. A {\bf 282} (1976) 349--351. %%QUATERNION, %%PARTICLE-PHYSICS, %%HESTENES, %$D04032006. 

\bibitem{CASAN1976B} G. Casanova, \emph{Diff\'erences de masse des multiplets m\'esoniques fondamentaux}, C.R. Acad. Sci. A {\bf 282} (1976) 665--667. %%QUATERNION, %%PARTICLE-PHYSICS, %%HESTENES, %$D04032006. 

\bibitem{CASAN1976-} G. Casanova, L'Alg\`ebre Vectorielle (Presses universitaires de France, Paris, 1976) 128~pp. %%BOOK, %%QUATERNION, %%PARTICLE-PHYSICS, %%SPECIAL-RELATIVITY, %%DIRAC, %%HADRODYNAMICS, %%GURSEY, %$D16032002.

\bibitem{EDMON1976A} J.D. Edmonds, Jr., \emph{Hypercomplex number approach to Schwinger's quantum source theory}, Int. J. Th. Phys. {\bf 15} (1976) 911--925. %%QUATERNION, %%PARTICLE-PHYSICS, %%GENERAL-RELATIVITY,  %$D24062002. 

\bibitem{ATIYA1978A} M.F. Atiyah and  N.J. Hitchin \emph{Construction of instantons}, Phys. Lett. {\bf 65 A} (1978) 185--187. %%QUATERNION, %%PARTICLE-PHYSICS, %%GAUGE-THEORY, %%INSTANTON, %$D05022002.

\bibitem{DEALF1978-} V. deAlfaro, S. Fubini, and G. Furlan, \emph{Classical solutions of generally invariant gauge theories}, Phys. Lett. {\bf 73 B} (1978) 463--467. %%QUATERNION, %%PARTICLE-PHYSICS, %%GAUGE-THEORY, %$D30112005.

\bibitem{GURSE1978A} F. G\"ursey, \emph{Some algebraic structures in particle theory}, in: Proc. 2nd John Hopkins Workshop on Current Problems in High Energy Particle Physics (John Hopkins Univ, Baltimore, 1978) 3--25.  %%QUATERNION, %%PARTICLE-PHYSICS, %$D15032002.

\bibitem{GURSE1978B} F. G\"ursey, \emph{Quaternion analyticity in field theory}, in: Proc. 2nd John Hopkins Workshop on Current Problems in High Energy Particle Physics (John Hopkins Univ, Baltimore, 1978) 179--221.   %%QUATERNION, %%PARTICLE-PHYSICS, %%ANALYTICITY-H, %%GAUGE-THEORY, %%INSTANTON, %$D02062001.

\bibitem{TUCKE1978-} R.W. Tucker and W.J. Zakrzweski, \emph{M\"obius invariance and classical solutions of $SU(2)$ gauge theory}, Nucl. Phys. {\bf B 143} (1978) 428--430.  %%QUATERNION, %%PARTICLE-PHYSICS, %%GAUGE-THEORY, %%CONFORMALITY,  %$D15032002.

\bibitem{CASAL1979-} R. Casalbuoni, G. Demokos, and S. Kovesi-Domokos, \emph{A new class of solutions to classical Yang-Mills equations}, Phys. Lett. {\bf 81 B} (1979) 34--36. %%QUATERNION, %%PARTICLE-PHYSICS, %%GAUGE-THEORY, %$D30112005.

%%\bibitem{GURSE1979-} F. G\"ursey, M.A. Jafarizadeh, and H.C. Tze, \emph{Quaternionic multi $S^4 \approx HP}(1)$ gravitational and chiral instantons}, Phys. Lett. {\bf B 88} (1979) 282--286. %%QUATERNION, %%PARTICLE-PHYSICS, %%FIELD-THEORY, %%INSTANTONS, %%ANALYTICITY, %$D30112005.

\bibitem{KAFIE1979-} Y.N. Kafiev, \emph{4-dimensional sigma model on quaternionic projective space}, Phys. Lett. {\bf B 87} (1979) 219--221. %%QUATERNION, %%PARTICLE-PHYSICS, %%SIGMA-MODEL, %%SPROJECTIVE-GEOMETRY, %%INSTANTON, %%KAHLER, %$D30112005.

\bibitem{LUKIE1979-} J. Lukierski, \emph{Complex and quaternionic supergeometry}, in: P. vanNieuwenhuizen and D.Z. Freeman, eds., \emph{Supergravity} (North-Holland, Amsterdam, 1979) 301.  %%QUATERNION, %%PARTICLE-PHYSICS, %$D05022002.

\bibitem{LUKIE1980-} J. Lukierski, \emph{Four dimensional quaternionic $\sigma$ models}, in: W. Ruehl, ed., Field Theoretical Methods in Particle Physics, Proceedings NATO, Kaiserslautern (Plenum, New York, 1980) 349--359. %%QUATERNION, %%PARTICLE-PHYSICS, %%MANIFOLD, %%SKYRME, %$D24062002.

\bibitem{GURSE1980A} F. G\"ursey and H.C. Tze, \emph{Complex and quaternionic analyticity in chiral and gauge theories, I}, Annals of Phys. {\bf 128} (1980) 29--130. %%C Part II has apparently not been published.  %%QUATERNION, %%PARTICLE-PHYSICS, %%ANALYTICITY-H, %%SIGMA-MODEL, %%GAUGE-THEORY, %%KAHLER, %%DIFFERENTIAL-FORMS, %$D02062001.

\bibitem{GURSE1980B} F. G\"ursey, \emph{Quaternion methods in field theory}, in: Proc. 4th John Hopkins workshop on current problems in particle theory (John Hopkins Univ, Baltimore, 1980) 255--288.  %%QUATERNION, %%PARTICLE-PHYSICS, %%ANALYTICITY-H, %%GAUGE-THEORY, %%INSTANTON, %$D02062001.

\bibitem{GLEBO1981-} A.L. Glebov, \emph{Classical particle with spin and Clifford algebra}, Theor. Math. Phys. {\bf 48} (1981) 786--790.   %%QUATERNION, %%PARTICLE-PHYSICS, %$D06072008.

\bibitem{WITTE1982-} E. Witten, \emph{An $SU(2)$ anomaly}, Phys. Lett. {\bf 117B} (1982) 324--328. %%QUATERNION, %%PARTICLE-PHYSICS, %$D04032006.

\bibitem{KELLE1984-} J. Keller, \emph{Space-time dual geometry theory of elementary particles and their interaction fields}, Int. J. Th. Phys. {\bf 23} (1984) 817--837. %%QUATERNION, %%PARTICLE-PHYSICS, %$D24062002, %$C Cite Juvet and Mercier.

\bibitem{BUGAJ1985-} K. Bugajska, \emph{Internal structure of fermions}, J. Math. Phys. {\bf 26} (1985) 77--93. %%QUATERNION, %%PARTICLE-PHYSICS, %%DIRAC, %%MAXWELL, %%CLIFFORD, %%KAHLER, %%DIFFERENTIAL-FORM, %$D09022002.

\bibitem{AOYAM1986-} S. Aoyama and T.W. vanHolten, \emph{Sigma-models on quaternionic manifolds and anomalies}, Z. Phys. C. Part. Fields. {\bf 31} (1986) 487--489. %%QUATERNION, %%PARTICLE-PHYSICS, %%SIGMA-MODEL, %$D30112005.

\bibitem{PERNI1986-} M. Pernici and P. vanNeuwenhuizen, \emph{A covariant action for the SU(2) spinning string as a hyper-K\"ahler or quaternionic nonlinear sigma model}, Phys. Lett. {\bf B 169} (1986) 381--385. %%QUATERNION, %%PARTICLE-PHYSICS, %%SIGMA-MODEL, %%SUPERSYMMETRY, %%KAHLER, %$D30112005.

\bibitem{ALBEV1987-} S. Albeveiro and R. Hoeghrohn, \emph{Quaternionic non-Abelian relativistic quantum fields in 4 space-time dimensions}, Phys. Lett. {\bf B 189} (1987) 329--336.  %%QUATERNION, %%PARTICLE-PHYSICS, %%ANALYTICITY-H, %%FIELD-THEORY, %$D30112005.

\bibitem{GOVOR1987A} A.B. Govorkov, \emph{Quaternion gauge fields. Pseudocolor}, Theor. Math. Phys. {\bf 68} (1987) 893--900. %%QUATERNION, %%PARTICLE-PHYSICS, %$D15032002.

\bibitem{GOVOR1987B} A.B. Govorkov, \emph{Fock representation for quaternion fields}, Theor. Math. Phys. {\bf 69} (1987) 1007--1013. %%QUATERNION, %%PARTICLE-PHYSICS, %$D15032002.

\bibitem{GURSE1987A} F. G\"ursey, \emph{Quaternionic and octonionic structures in physics}, in: M. G. Dancel et al., eds, Symmetries in Physics (1600-1980),  Proceedings of the 1st international Meeting on the History of Scientific Ideas, Barcelona, 1983, (Univ. Autonoma Barcelona, Barcelona, 1987) 557--592. %%QUATERNION, %%OCTONION, %%HISTORY, %$D19032002, %$D19032002.

\bibitem{GURSE1987B} F. G\"ursey, \emph{Super Poincar� groups and division algebras}, Modern Phys. Lett. {\bf A 2} (1987) 967--976. %%QUATERNION, %%PARTICLE-PHYSICS, %%GROUP-THEORY, %%SUPERSYMMETRY, %%OCTONION, %$D30112005.

\bibitem{ITOH-1987-} M. Itoh, \emph{Quaternion structure on the moduli space of Yang-Mills connections}, Math. Ann. {\bf 276} (1987) 581--593. %%QUATERNION, %%PARTICLE-PHYSICS, %%GAUGE-THEORY, %%KAHLER, %$D30112005.

\bibitem{SCHEU1987-} H. Scheurich, \emph{Principles of quaternionic vacuum thermodynamics and a unified gravistrong interaction model}, Ann. der Phys. {\bf 44} (1987) 473--487.  %%QUATERNION, %%PARTICLE-PHYSICS, %%EDDINGTON, %$D15032002.

\bibitem{KIMUR1988-} T. Kimura and I. Oda, \emph{Superparticles and division algebras}, Prog. Theor. Phys.  {\bf 80} (1988) 1--6.  %%QUATERNION, %%PARTICLE-PHYSICS, %%OCTONION, %$D15032002. 

\bibitem{LAMBE1988-} D. Lambert and J. Rembielinski, \emph{From G\"odel quaternions to non-linear sigma models}, J. Phys. {\bf A 21} (1988) 2677--2691.  %%QUATERNION, %%PARTICLE-PHYSICS, %%SIGMA-MODEL, %%GROUP-THEORY, %%PSEUDOQUATERNION, %$D14032002.

\bibitem{TZE--1988-} C.H. Tze and S. Nam, \emph{Global dynamics of electric and magnetic membranes on the complex, quaternionic and octonionic Hopf bundles}, Phys. Lett. {\bf B 210} (1988) 76--84.  %%QUATERNION, %%PARTICLE-PHYSICS, %%MONOPOLE, %%HOPF-FIBRATION, %$D21032002.

\bibitem{CHEUN1989-} H.Y. Cheung and F. G\"ursey, \emph{Hadronic superalgebra from skyrmion operators}, Phys. Lett. {\bf B 219} (1989) 127--129. %%QUATERNION, %%PARTICLE-PHYSICS, %%SKYRME, %$D03082005.

\bibitem{TZE--1989-} C.H. Tze and S. Nam, \emph{Topological phase entanglements of membrane solitons in division algebras sigma models with a Hopf term}, Ann. Phys. {\bf 193} (1989) 419--471.  %%QUATERNION, %%PARTICLE-PHYSICS, %%HOPF-FIBRATION, %$D04032006, %$C Fractional spin and statistics.

\bibitem{ALBEV1990-} S. Albeverio, K. Iwata, and T. Kolsrud, \emph{Random fields as solutions of the inhomogeneous quaternionic Cauchy-Riemann equation.  I.  Invariance and analytic continuation}, Commun. Math. Phys. {\bf 132} (1990) 555--580. %%QUATERNION, %%PARTICLE-PHYSICS, %%ANALYTICITY-H, %%FIELD-THEORY, %$D01062001.

\bibitem{CHEUN1990-} H.Y. Cheung and F. G\"ursey, \emph{Composite Skyrme model}, Mod. Phys. Lett. {\bf A 5} (1990) 1685--1691. %%QUATERNION, %%PARTICLE-PHYSICS, %%SKYRME, %$D03082005.

\bibitem{CHISH1990-} J.S.R. Chisholm and R.S. Farwell, \emph{Unified spin gauge theories of the four fundamental forces}, in: D.G. Quillen et al., eds., The Interface of Mathematics and Particle Physics (Clarendon Press, Oxford, 1990) 193--202. %%QUATERNION, %%PARTICLE-PHYSICS, %%SU(3), %%HESTENES, %$D16032002, 

\bibitem{BISHT1991-} P.S. Bisht, O.P.S. Negi, and B.S. Rajput, \emph{Quaternion gauge theory of dyon fields}, Prog. Theor. Phys. {\bf 85} (1991) 157--168.  %%QUATERNION, %%PARTICLE-PHYSICS, %%MONOPOLE, %$D30112005.

\bibitem{GALIC1991-} K. Galicki and Y.S. Poon, \emph{Duality and Yang-Mills fields of quaternionic K\"ahler manifolds}, J. Math. Phys. {\bf 32} (1991) 1263--1268.  %%QUATERNION, %%PARTICLE-PHYSICS, %%GAUGE-THEORY, %%KAHLER, %$D30112005.

\bibitem{CHISH1992A} J.S.R. Chisholm and R.S. Farwell, \emph{Tetrahedral structure of idempotents of the Clifford algebra $\cl_{1,3}$}, in: A. Micali et al., eds., Clifford Algebras and their Applications in Mathematical Physics (Kluwer Academic Publishers, Dordrecht, 1992) 27--32. %%QUATERNION, %%PARTICLE-PHYSICS, %%CLIFFORD, %%SU(3), %%HESTENES, %$D16032002. 

\bibitem{CHISH1992B} J.S.R. Chisholm and R.S. Farwell, \emph{Unified spin gauge theories of the four fundamental forces}, in: A. Micali et al., eds., Clifford Algebras and their Applications in Mathematical Physics (Kluwer Academic Publishers, Dordrecht, 1992) 363--370. %%QUATERNION, %%PARTICLE-PHYSICS, %%SU(3), %%HESTENES, %$D16032002. 

\bibitem{GURSE1992-} F. G\"ursey, W. Jiang, \emph{Euclidean space-time diffeomorphism and their Fueter subgroups}, J. Math. Phys. {\bf 33} (1992) 682--700.  %%QUATERNION, %%PARTICLE-PHYSICS, %%ANALYTICITY, %%FIELD-THEORY, %%H, CONFORMALITY, %$D02062001.

\bibitem{DIXON1993-}  G. Dixon, \emph{Particle families and the division algebras},  J. Phys. G: Nucl. Phys. {\bf 12 } (1986) 561--570. %%QUATERNION, %%PARTICLE-PHYSICS, %%OCTONION, %$D19032002.

\bibitem{EVANS1993-} M. Evans, F. G\"ursey, and V. Ogievetsky, \emph{From two-dimensional conformal to four-dimensional self-dual theories: Quaternionic analyticity}, Phys. Rev. D {\bf 47} (1993) 3496-3508.  %%QUATERNION, %%PARTICLE-PHYSICS, %%ANALYTICITY-H, %%FIELD-THEORY, CONFORMALITY, %$D02062001.

\bibitem{GALPE1993-} A. Galperin and V. Ogievetsky, \emph{Harmonic potentials for quaternionic symmetrical sigma-models}, Phys. lett. {\bf B 301} (1993) 61--71. %%QUATERNION, %%PARTICLE-PHYSICS, %%SIGMA-MODEL, %%SUPERSYMMETRY, %%KAHLER, %$D30112005.

\bibitem{VROEG1993-}  P.G. Vroegindewij, \emph{The equations of Dirac and the $M_2(H)$--representation of $\cl_{1,3}$ },  Found. Phys.  {\bf 23 } (1993) 1445--1463. %%QUATERNION, %%PARTICLE-PHYSICS, %%DIRAC, %%SU(3), %%HESTENES, %$D16032002.

\bibitem{DEWIT1994A} B. DeWit and A. vanProeyen, \emph{Hidden symmetries, spectral geometry and quaternionic manifolds}, Int. J. Mod. Phys. {\bf D3} (1994) 31--47. %%QUATERNION, %%PARTICLE-PHYSICS, %%SIGMA-MODEL, %%SUPERSYMMETRY, %%KAHLER, %$D30112005.

\bibitem{MERKU1994-} S. Merkulov, H. Pedersen, and A. Swann, \emph{Topological quantum-field theory in quaternionic geometry}, J. Geom. Phys. {\bf 14} (1994) 121--145. %%QUATERNION, %%PARTICLE-PHYSICS, %%SUPERSYMMETRY, %%KAHLER, %%INSTANTON, %$D30112005.

\bibitem{DELEO1995C} S. DeLeo, \emph{Duffin-Kemmer-Petiau equation on the quaternion field}, Prog. Theor. Phys. {\bf 94} (1995) 1109--1120.  %%QUATERNION, %%PARTICLE-PHYSICS, %$D15032002.

\bibitem{BRUMB1996C}  S.P. Brumby, R. Foot and R.R. Volkas, \emph{Quaternionic formulation of the exact parity model} (1996) 30~pp.; e-print \underline{ arXiv:hep-th/9602139 }. %%QUATERNION, %%PARTICLE-PHYSICS, %%PARITY, %$D06072008.

\bibitem{DELEO1996D} S. DeLeo, \emph{Quaternions for GUTs}, Int. J. Th. Phys. {\bf 35} (1996) 1821--1837. %%QUATERNION, %%PARTICLE-PHYSICS, %$D15032002.

\bibitem{GURSEY1996-}  F. G\"ursey, H.C. Tze, On the Role of Division, Jordan and Related Algebras in Particle Physics (World Scientific, 1996) 462~pp.  %%BOOK, %%QUATERNION, %%PARTICLE-PHYSICS, %%ISOSPIN-page-36, %$D20032002. 

\bibitem{OKONE1996-} C. Okonek and A. Teleman, \emph{Quaternionic monopoles}, Commun. Math. Phys. {\bf 180} (1996) 363--388. %%QUATERNION, %%PARTICLE-PHYSICS, %%MONOPOLE, %%KAHLER, %$D30112005.

\bibitem{LAMBE1996-} J. Lambek, \emph{Quaternions and the particles of nature}, Report from the department of mathematics and statistics 96-03 (McGill University, October 1, 1996) 14~pp.. %%QUATERNION, %%PARTICLE-PHYSICS, %$D06022002.

\bibitem{LIM--1997-} C.S. Lim, \emph{Quaternionic mass matrices and CP-symmetry}, Mod. Phys. Lett. {\bf A12} (1997) 2829--2835. %%QUATERNION, %%PARTICLE-PHYSICS, %$D30112005.

\bibitem{ZUCCH1998-} R. Zucchini, \emph{The quaternionic geometry of four-dimensional conformal field theory}, J. Geom. Phys. {\bf 27} (1998) 113--153. %%QUATERNION, %%PARTICLE-PHYSICS, %%ANALYTICITY, %%CONFORMALITY, %$D13032002.

\bibitem{MAIA-1999-} M.D. Maia, \emph{Spin and isospin in quaternions quantum mechanics} (8 Apr 1999) 7~pp; e-print \underline{ arXiv:hep-th/9904067 }. %%QUATERNION, %%PARTICLE-PHYSICS, %%ISOSPIN, %%QQPH, %$D24062002.

\bibitem{LAMBE2000-} J. Lambek, \emph{Four-vector representation of fundamental particles}, Int. J. Theor. Phys. {\bf 39} (2000) 2253--2258.  %%QUATERNION, %%PARTICLE-PHYSICS, %$D06022002. 

\bibitem{BUDIN2002A} P. Budinich, \emph{From the geometry of pure spinors with their division algebras to fermion physics}, Found. Phys. {\bf 32} (2002) 1347--1398; e-print \underline{ arXiv:hep-th/0107158 }. %%QUATERNION, %%PARTICLE-PHYSICS, %%CLIFFORD, %%HESTENES, %$D06072008.    

\bibitem{BUDIN2002B} P. Budinich, \emph{The possible role of pure spinors in some sectors of particle physics} (2002) 20~pp.; e-print \underline{ arXiv:hep-th/0207216 }.  %%QUATERNION, %%PARTICLE-PHYSICS, %%CLIFFORD, %%HESTENES, %$D06072008.   

\bibitem{BUDIN2003-} P. Budinich, \emph{Internal symmetry from division algebras in pure spinor geometry} (2003) 12~pp.; e-print \underline{ arXiv:hep-th/0311045 }. %%QUATERNION, %PARTICLE-PHYSICS, %%CLIFFORD, %%HESTENES, %$D06072008.  

\bibitem{BEIL-2003-} R.G. Beil and K.L. Ketner, \emph{Peirce, Clifford, and quantum theory}, Int. J. Theor. Phys. {\bf 42} (2003) 1957--1972.  %%QUATERNION, %%PARTICLE-PHYSICS, %%CLIFFORD, %$D30112005.

\bibitem{NOTTA2003-} L. Nottale, M.-N. C\'el\'erier, and T. Lehner, \emph{Gauge field theory in scale relativity} (10 July 2003) 17~pp.;  e-print \underline{ arXiv:hep-th/0307093 }.  %%QUATERNION, %%PARTICLE-PHYSICS, %%GAUGE-THEORY, %$D30112005.

\bibitem{KASSA2004C} V.V. Kassandrov, \emph{The algebrodynamics: primordial light, particles-caustics and the flow of time}, Hypercomplex Numbers in Geometry and Physics {\bf 1} (2004) 84--99.  %%QUATERNION, %%PARTICLE-PHYSICS, %$D06072008.

\bibitem{TOPPA2004-} F. Toppan, \emph{Hermitian versus holomorphic complex and quaternionic generalized supersymmetries of the M-theory. A classification}, J. of High Energy Physics {\bf 0409} (2004) 016, 25~pp.; e-print \underline{ arXiv:hep-th/0406022 }. %%QUATERNION, %%PARTICLE-PHYSICS, %%SUPERSYMMETRY, %$D04032006.

\bibitem{GSPON2005-} A. Gsponer, \emph{Integral-quaternion formulation of Lambek's representation of fundamental particles and their interactions}, Report ISRI-02-03 (13 November 2005) 10~pp; e-print \underline{ arXiv:math-ph/0511047 }.  %%QUATERNION, %%PARTICLE-PHYSICS, %%INTEGRAL-QUATERNION, %$D30112005.

\bibitem{FRIGE2005-} M. Frigerio, S. Kaneko, E. Ma, and M. Tanimoto, \emph{Quaternion family symmetry of quarks and leptons}, Phys. Rev. {\bf D71} (2005) 011901(R), 5~pp.; e-print \underline{ arXiv:hep-ph/0409187 }.  %%QUATERNION, %%PARTICLE-PHYSICS, %$D06072008.

\bibitem{NOTTA2006-} L. Nottale, M.-N. C�l�rier and T. Lehner, \emph{Non-Abelian gauge field theory in scale relativity}, J. Math. Phys. {\bf 47} (2006) 032303, 19~pp.; e-print \underline{ arXiv:hep-th/0605280 }.  %%QUATERNION, %%PARTICLE-PHYSICS, %$D06072008.

\bibitem{GU---2007-} Y.-Q. Gu and T.-T. Li, \emph{Eigen Equation of the Nonlinear Spinor} (3 March 2007) 8~pp.; e-print \underline{ arXiv:0704.0436 }. %%QUATERNION, %%PARTICLES, %%GSPONER, %%HURNI, %$D06072008. 

\bibitem{YEFRE2007-} A. Yefremov, F. Smarandache, and V. Christianto, \emph{Yang-Mills field from quaternion space geometry, and its Klein-Gordon representation}, Prog. in Phys. {\bf 3} (2007) 42--50.  %%QUATERNION, %%PARTICLE-PHYSICS, %$D06072008.

\bibitem{RADCH2008-} V. Radchenko, \emph{Nonlinear classical fields} (21 January 2007) 45~pp.; e-print \underline{ arXiv:math-ph/0701054 }.  %%QUATERNION,   %%PARTICLE-PHYSICS,  %%GSPONER, %%HURNI, %$D06072008.

\end{enumerate}


\section{\Huge QUANTICS}
%=======================
\label{QUANTICS}

Quantum theory and its generalizations.

\subsection{QUANTUM-PHYSICS}
%---------------------------
\label{QUANTUM-PHYSICS}

Standard relativistic and non-relativistic quantum physics in which the scalars belong to the field of complex numbers. 


\begin{enumerate}

\bibitem{JORDA1927-} P. Jordan, \emph{Zur Quantenmechanik der Gasentartung}, Zeits. f. Phys. {\bf 44} (1927) 473--480. %%QUATERNION, %%QUANTUM-PHYSICS, %%SPIN-1/2, %%PAULI, %$D06022002. %$C Jordan uses quaternions in 2x2 matrix form (without referring to Pauli's paper \cite{PAULI1927-}) to formulate the anticommutations rules of the second-quantized electron field.

\bibitem{PAULI1927-} W. Pauli, \emph{Zur Quantenmechanik des magnetischen Electrons}, Z. Phys. {\bf 43} (1927) 601--623. %%QUATERNION, %%QUANTUM-PHYSICS, %%SPIN-1/2, %%PAULI, %$D05022002. %$C This is the paper in which Pauli defines his ``Pauli matrices.''  In footnote 2 Pauli acknowledges that it was Jordan who showed him the commutation relations of the spin 1/2 operators, and their relations to quaternions (see \cite{JORDA1927-}).  It is apparently the only published paper of Pauli in which the word ``quaternion'' is printed.

\bibitem{WEYL-1927-} H. Weyl, \emph{Quantenmechanik und Gruppentheorie}, Zeits. f. Phys. {\bf 46} (1927) 1--46. %%QUATERNION, %%QUANTUM-PHYSICS, %%GROUP-THEORY, %$D30112005, %$C Quaternions p.20.

\bibitem{JORDA1928-} P. Jordan and E. Wigner, \emph{\"Uber das Paulische \"Aquivalentzverbot}, Zeits. f. Phys. {\bf 47} (1927) 631--651. %%QUATERNION, %%QUANTUM-PHYSICS, %%SPIN-1/2, %%PAULI, %$D14052007. %$C Jordan and Wigner  refer to \cite{JORDA1927-} and use quaternions to formulate the anticommutations rules of the second-quantized electron field (Pauli-Fermi gas).

\bibitem{CONWA1932A} A.W. Conway, \emph{The radiation of angular momentum}, Proc. Roy. Irish Acad. {\bf A 41} (1932) 8--17. %%QUATERNION, %%QUANTUM-PHYSICS, %%SPHERICAL-HARMONICS, %%ANGULAR-MOMENTUM, $C Cited by W. Heitler, The Quantum Theory of radiation (Oxford University Press, 1954) p.402, %$D13032002.

\bibitem{CONWA1932B} A.W. Conway, \emph{The radiation of angular momentum (Abstract)}, Nature {\bf 129} (25 June 1932) 950. %%QUATERNION, %%QUANTUM-PHYSICS, %%SPHERICAL-HARMONICS, %%ANGULAR-MOMENTUM, %C Abstract of \cite{CONWA1932A}, %$D30112005.

\bibitem{CONWA1937-} A.W. Conway, \emph{Quaternion treatment of the electron wave equation}, Proc. Roy. Soc. {\bf A 162} (1937) 145--154. %%QUATERNION, %%QUANTUM-PHYSICS, %%DIRAC, %%LANCZOS, %$D05022002.

\bibitem{CONWA1948-} A.W. Conway, \emph{Quaternions and quantum mechanics}, Acta Pontifica Acad. Scientiarium {\bf 12} (1948) 259--277. %%QUATERNION, %%QUANTUM-PHYSICS, %%DIRAC, %$D05022002.

\bibitem{DYSON1962A} F. Dyson, \emph{Statistical theory of the energy levels of complex systems}, J. Math. Phys. {\bf 3} (1962) 140--175. %%QUATERNION, %%QUANTUM-PHYSICS, %$D14032002.

\bibitem{CASAN1968B} G. Casanova, \emph{Sur certaines solutions de l'\'equation de Dirac-Hestenes}, C.R. Acad. Sci. A {\bf 267} (1968) 661--663. %%QUATERNION, %%QUANTUM-PHYSICS, %%HESTENES, %$D04032006. 

\bibitem{CASAN1968C} G. Casanova, \emph{Sur certains champs magn\'etiques en th\'eorie de Dirac-Hestenes}, C.R. Acad. Sci. A {\bf 267} (1968) 674--676. %%QUATERNION, %%QUANTUM-PHYSICS, %%HESTENES, %$D04032006. 

\bibitem{CASAN1969A} G. Casanova, \emph{Principe de superposition en th\'eorie de Dirac-Hestenes}, C.R. Acad. Sci. A {\bf 268} (1969) 437--440. %%QUATERNION, %%QUANTUM-PHYSICS, %%HESTENES, %$D04032006. 

\bibitem{CASAN1969B} G. Casanova, \emph{Particules neutre de spin $1$}, C.R. Acad. Sci. A {\bf 268} (1969) 673--676. %%QUATERNION, %%PARTICLE-PHYSICS, %%HESTENES, %$D04032006, %$C De Broglie's method of "fusion." 

\bibitem{CASAN1970A} G. Casanova, \emph{L'atome d'hydrog\`ene en th\'eorie de Dirac-Hestenes}, C.R. Acad. Sci. A {\bf 270} (1970) 1202--1204. %%QUATERNION, %%QUANTUM-PHYSICS, %%HESTENES, %$D04032006. 

\bibitem{CASAN1970B} G. Casanova, \emph{Solutions planes de l'\'equation de Dirac-Hestenes dans un champ central}, C.R. Acad. Sci. A {\bf 270} (1970) 1470--1472. %%QUATERNION, %%QUANTUM-PHYSICS,  %%HESTENES, %$D04032006, %$C The solutions differ from Dirac's equation... 

\bibitem{CASAN1970C} G. Casanova, \emph{Moments cin\'etiques et magn\'etiques en th\'eorie de Dirac-Hestenes}, C.R. Acad. Sci. A {\bf 271} (1970) 817--820. %%QUATERNION, %%QUANTUM-PHYSICS, %%HESTENES, %$D04032006. 

\bibitem{HAUTO1970-} A.P. Hautot, \emph{Sur une m\'ethode quaternionique de s\'eparation des variables}, Physica {\bf 48} (1970) 609--619. %%QUATERNION, %%QUANTUM-PHYSICS, %%LANCZOS, %%CONWAY, %$D03082005, %$C Hydrogen atom.

\bibitem{HAUTO1971-} A.P. Hautot, \emph{Sur la compl\'etude de l'ensemble des fonctions propres de l'atome d'hydrog\`ene relativiste}, Physica {\bf 53} (1971) 154--156. %%QUATERNION, %%QUANTUM-PHYSICS, %$D04032006, %$C Hydrogen atom.

\bibitem{HESTE1971B} D. Hestenes, R. Gurtler, \emph{Local observables in quantum theory}, Am. J. Phys. {\bf 39} (1971) 1028--1038. %%QUATERNION, %%QUANTUM-PHYSICS, %%DIRAC, %%PAULI, %%HESTENES, %$D10022002.

\bibitem{HAUTO1972-} A.P. Hautot, \emph{The exact motion of a charged particle in the magnetic field $B= (x^2+y^2)^{-\frac{1}{2}} ({-\gamma y}/{x^2+y^2},{\gamma x}/{x^2+y^2},\alpha)$}, Physica {\bf 58} (1972) 37--44. %%QUATERNION, %$D04032006.

\bibitem{HESTE1973-} D. Hestenes, \emph{Local observables in the Dirac theory}, J. Math. Phys. {\bf 14} (1973) 893--905. %%QUATERNION, %%QUANTUM-PHYSICS, %%DIRAC, %%HESTENES, %$D10022002.

\bibitem{CRUBE1974-} A. Crubellier and S. Feneuille, \emph{Application de la m\'ethode de factorisation et de la th\'eorie des groupes au probl\`eme de l'\'electron de Landau relativiste}, J. Phys. A: Math., Nucl. Gen {\bf 7} (1974) 1051--1060. %%QUATERNION, %%QUANTUM-PHYSICS, %$D04032006.

\bibitem{CASAN1975A} G. Casanova, \emph{Sur la m\'ethode du rep\`ere mobile instantan\'e en th\'eorie de Dirac}, C.R. Acad. Sci. A {\bf 280} (1975) 299--302. %%QUATERNION, %%QUANTUM-PHYSICS, %%HESTENES, %$D04032006. 

\bibitem{GURTL1975-}  R. Gurtler and D. Hestenes, \emph{Consistency in the Formulation of the Dirac, Pauli and Schroedinger Theories},  J. Math. Phys. {\bf 16} (1975) 573-583. %%QUATERNION, %%QUANTUM-PHYSICS, %%DIRAC, %%PAULI, %%SCHRODINGER, %%HESTENES, %$D10022002.

\bibitem{GREUB1975-} W. Greub and H.-R. Petry, \emph{Minimal coupling and complex line bundles}, J. Math. Phys. {\bf 16} (1975) 1347--1351. %%QUATERNION, %%QUANTUM-PHYSICS, %%SCHRODINGER, %%MAGNETIC-MONOPOLE %$D30112005.

\bibitem{HESTE1975-} D. Hestenes, \emph{Observables, operators, and complex numbers in the Dirac theory}, J. Math. Phys. {\bf 16} (1975) 556--572. %%QUATERNION, %%QUANTUM-PHYSICS, %%DIRAC, %%HESTENES, %$D10022002.

\bibitem{KENNE1976-}  S. Kennedy and R. Gamache, \emph{A geometric-algebra treatment of the Feynman-Vernon-Hellewarth space of the two-state problem}, Am. J. Phys. {\bf 64} (1976) 1475-1482. %%QUATERNION, %%QUANTUM-PHYSICS, %%SPIN-1/2, %$D30112005.

\bibitem{BARRE1978-} T.W. Barrett, \emph{A deterministic interpretation of the commutation and uncertainty relations of quantum theory and a redefinition of Planck's constant as a coupling condition}, Nuovo Cim. {\bf 45 B} (1978) 297--309.  %%QUATERNION, %%QUANTUM-PHYSICS, %$D30112005. 

\bibitem{HESTE1979-} D. Hestenes, \emph{Spin and uncertainty in the interpretation of quantum mechanics}, Am. J. Phys. {\bf 47} (1979) 399--415. %%QUATERNION, %%QUANTUM-PHYSICS, %%DIRAC, %%PAULI, %%SCHRODINGER, %%HESTENES, %$D10022002.

\bibitem{ARAKI1980-} H. Araki, \emph{On a characterization of the state space of quantum mechanics}, Comm. Math. Phys. {\bf 75} (1980) 1--24. %%QUATERNION, %%QUANTUM-PHYSICS, %%QQM, %$D04032006.

\bibitem{BEREZ1981-} A.V. Berezin, E.A. Tolkachev, and I. Fedorov, \emph{Solution of the Dirac equation in quaternions}, Sov. Phys. Journal. {\bf 24} (1981) 935--937. %%QUATERNION, %%QUANTUM-PHYSICS, %%DIRAC, %$D06022002. %$C Solution with an anomalous magnetic moment in the field of a plane wave.

\bibitem{GOUGH1984-} W. Gough, \emph{Quaternions and spherical harmonics}, Eur. J. Phys. {\bf 5} (1984) 163--171. %%QUATERNION, %%QUANTUM-PHYSICS, %%SPHERICAL-HARMONICS, %$D13032002.

\bibitem{HESTE1985-} D. Hestenes, \emph{Quantum mechanics from self-interaction}, Found. Phys. {\bf 15} (1985) 63--87. %%QUATERNION, %%QUANTUM-PHYSICS, %%DIRAC, %%ZITTERBEWEGUNG, %%HESTENES, %$D10022002.

\bibitem{GOUGH1986-} W. Gough, \emph{The analysis of spin and spin-orbit coupling in quantum and classical physics by quaternions}, Eur. J. Phys. {\bf 7} (1986) 35--42. %%QUATERNION, %%QUANTUM-PHYSICS, %%SPHERICAL-HARMONICS, %$D13032002.

\bibitem{HESTE1986B} D. Hestenes, \emph{Clifford algebras and the interpretation of quantum mechanics}, in: J.S.R. Chisholm and A.K. Common, eds.,  Clifford Algebras and Their Applications in Mathematical Physics (Reidel, Dordrecht, 1986) 321--346. %%QUATERNION, %%QUANTUM-PHYSICS, %%DIRAC, %%LEPTODYNAMICS, %%GAUGE-THEORY, %%HESTENES, %%KAHLER, %$D10022002.

\bibitem{GOUGH1987-} W. Gough, \emph{Quaternion quantum mechanical treatment of an electron in a magnetic field}, Eur. J. Phys. {\bf 8} (1987) 164--170. %%QUATERNION, %%QUANTUM-PHYSICS, %$D17042002.

\bibitem{YANG-1987-} C.N. Yang, \emph{Square root of minus one, complex phases and Erwin Schr\"odinger}, in: C.W. Kilmister, ed., Schr\"odinger:  Centenary Celebration of a Polymath (Cambridge University Press, 1987) 53--64.  %%QUATERNION, %%QUANTUM-PHYSICS, %%COMPLEX-NUMBERS, %$D04032006.

\bibitem{SIMON1988A} R. Simon and N. Kumar, \emph{A note on the Berry phase for systems having one degree of freedom}, J. Phys. A: Math. gen. {\bf 21} (1988) 1725--1727. %%QUATERNION, %%QUANTUM-PHYSICS, %%BERRY-PHASE, %$D11032002.

\bibitem{GOUGH1989A} W. Gough, \emph{A quaternion expression for the quantum mechanical probability and current densities}, Eur. J. Phys. {\bf 10} (1989) 188--193. %%QUATERNION, %%QUANTUM-PHYSICS, %$D17042002.

\bibitem{AVRON1989-} J.E. Avron, L. Sadun, J. Segert, and B. Simon, \emph{Chern numbers, quaternions, and Berry's phases in Fermi systems}, Commun. Math. Phys. {\bf 129} (1989) 595--627.  %%QUATERNION, %%QUANTUM-PHYSICS, %%BERRY-PHASE, %$D20032002.

\bibitem{GOUGH1989B} W. Gough, \emph{On the probability of a relativistic free electron}, Eur. J. Phys. {\bf 10} (1989) 318--319. %%QUATERNION, %%QUANTUM-PHYSICS, %$D17042002.

\bibitem{SIMON1989A} R. Simon, N. Mukunda, and E.C.G. Sudarshan, \emph{Hamilton's theory of turns generalized to SP(2,R)}, Phys. Rev. Lett. {\bf 62} (1989) 1331--1334. %%QUATERNION, %%QUANTUM-PHYSICS, %%BERRY-PHASE, %$D11032002.

\bibitem{SIMON1989B} R. Simon, N. Mukunda, and E.C.G. Sudarshan, \emph{The theory of screws: A new geometric representation for the group SU(1,1)}, J. Math. Phys. {\bf 30} (1989) 1000--1006. %%QUATERNION, %%QUANTUM-PHYSICS, %%BERRY-PHASE, %$D11032002.

\bibitem{SIMON1989C} R. Simon, N. Mukunda, and E.C.G. Sudarshan, \emph{Hamilton's theory of turns and a new geometrical representation for polarization optics}, Pramana -- J. Phys. {\bf 32} (1989) 769--792.  %%QUATERNION, %%QUANTUM-PHYSICS, %%BERRY-PHASE, %$D11032002.

\bibitem{GOUGH1990-} W. Gough, \emph{Mixing scalars and vectors---an elegant view of physics}, Eur. J. Phys. {\bf 11} (1990) 326--333. %%QUATERNION, %%QUANTUM-PHYSICS, %$D17042002.

\bibitem{HESTE1990A} D. Hestenes, \emph{The Zitterbewegung interpretation of quantum mechanics}, Found. Phys. {\bf 20} (1990) 1213--1232.  %%QUATERNION, %%QUANTUM-PHYSICS, %%DIRAC, %%ZITTERBEWEGUNG, %%HESTENES, %$D10022002.

\bibitem{HESTE1990B} D. Hestenes, \emph{On Decoupling Probability from Kinematics in Quantum Mechanics}, in: P.F. Fougere, ed., Maximum Entropy and Bayseian Methods (Kluwer Academic Publishers, Dordrecht, 1990) 161--183. %%QUATERNION, %%QUANTUM-PHYSICS, %%DIRAC, %%HESTENES, %$D10022002.

\bibitem{BOUDE1991-} R. Boudet, \emph{The role of the duality rotation in the Dirac theory.  Comparison between the Darwin and the Kr\"uger solutions for the central potential problem}, in: D. Hestenes and A. Weingartshofer, eds., The Electron (Kluwer Academic Publishers, Dordrecht, 1991) 83--104.  %%QUATERNION, %%QUANTUM-PHYSICS, %%DIRAC, %%HESTENES, %$D16032002. 

\bibitem{KRUGE1991-} H. Kr\"uger, \emph{New solutions of the Dirac equation for central fields}, in: D. Hestenes and A. Weingartshofer, eds., The Electron (Kluwer Academic Publishers, Dordrecht, 1991) 49--81.  %%QUATERNION, %%QUANTUM-PHYSICS, %%DIRAC, %$D15032002. 

\bibitem{SIMON1990-} R. Simon, N. Mukunda, \emph{Minimal three-component SU(2) gadget for polarization optics}, Phys. Lett. A. {\bf 143} (1990) 165--169. %%QUATERNION, %%QUANTUM-PHYSICS, %%BERRY-PHASE, %$D11032002.

\bibitem{ADLER1992-} R.J. Adler and R.A. Martin, \emph{The electron $g$ factor and factorization of the Pauli equation}, Am. J. Phys. {\bf 60} (1992) 837--839. %%QUATERNION, %%QUANTUM-PHYSICS, %%DIRAC, %%PAULI, %$D06022002. 

\bibitem{BOUDE1992-} R. Boudet, \emph{Les alg\`ebres de Clifford et les transformations des multivecteurs.  L'alg\`ebre de Clifford $R(1,3)$ et la constante de Planck}, in: A. Micali et al., eds., Clifford Algebras and their Applications in Mathematical Physics (Kluwer Academic Publishers, Dordrecht, 1992) 343--352. %%QUATERNION, %%QUANTUM-PHYSICS, %%DIRAC, %%TETRODE, %%HESTENES, %$D16032002. 

\bibitem{LOUPI1992-}  G. Loupias, \emph{Alg\`ebres de Clifford et relations d'anticommutation canoniques}, Adv. Appl. Clifford Alg. {\bf 2} (1992) 9--52. %%QUATERNION, %%QUANTUM-PHYSICS, %$D24062002.

\bibitem{SIMON1992-} R. Simon and N. Mukunda, \emph{Hamilton's turns and geometric phase for two-level systems}, J. Phys. {\bf A25} (1992) 6135--6144. %%QUATERNION, %%QUANTUM-PHYSICS, %%BERRY-PHASE, %$D11032002.

\bibitem{CASAN1993-}  G. Casanova, \emph{Non-localisation des \'electrons dans leur onde}, Adv. Appl. Clifford Alg. {\bf 3} (1993) 127--132. %%QUATERNION, %%QUANTUM-PHYSICS, %%DIRAC, %$D24062002.

\bibitem{DORAN1993B} C. Doran, A. Lasenby, and S. Gull, \emph{States and operators in the spacetime algebra}, Found. Phys. {\bf 23}  (1993) 1239--1264.  %%QUATERNION, %%QUANTUM-PHYSICS, %%HESTENES, %$D17032002.

\bibitem{GULL-1993A} S. Gull, C. Doran, and A. Lasenby, \emph{Electron paths, tunneling and diffraction in the spacetime algebra}, Found. Phys. {\bf 23}  (1993) 1329--1356.  %%QUATERNION, %%QUANTUM-PHYSICS, %%PAULI %%HESTENES, %$D17032002.

\bibitem{JANCE1993-} B. Jancewicz, \emph{A Hilbert space for the classical electromagnetic field}, Found. Phys. {\bf 23} (1993) 1405--1421.  %%QUATERNION, %%QUANTUM-PHYSICS, %%MAXWELL, %%HESTENES, %$D24062002.

\bibitem{KRUGE1993-} H. Kr\"uger, \emph{Classical limit of real Dirac theory: quantization of relativistic central field orbits}, Found. of Phys. {\bf 23} (1993) 1265--1288.  %%QUATERNION, %%QUANTUM-PHYSICS, %%DIRAC, %%SPINOR, %$D15032002. 

\bibitem{LASEN1993A} A. Lasenby, C. Doran, and S. Gull, \emph{A multivector derivative approach to Lagrangian field theory}, Found. Phys. {\bf 23}  (1993) 1295--1327.  %%QUATERNION, %%QUANTUM-PHYSICS, %%HESTENES, %$D17032002.

\bibitem{CASAN1996-}  G. Casanova, \emph{Sur la longueur d'onde de l'\'electron et le potentiel de Yukawa}, Adv. Appl. Clifford Alg. {\bf 6} (1996) 143--150. %%QUATERNION, %%QUANTUM-PHYSICS, %%DIRAC, %%HADRODYNAMICS, %$D24062002.

\bibitem{CHALL1996-}  A. Challinor, A. Lasenby, S. Gull, and C. Doran, \emph{A relativistic, causal account of spin measurement}, Phys. Lett. A {\bf 218}  (1996) 128--138.  %%QUATERNION, %%QUANTUM-PHYSICS, %%DIRAC, %%PAULI, %%HESTENES, %$D17032002.

\bibitem{DORAN1996-} C. Doran, A. Lasenby, S. Gull, S. Somaroo, and A. Challinor, \emph{Spacetime algebra and electrophysics}, in: P.W. Hawbes, ed., Advances in Imaging and Electron Physics {\bf 95} (1996) 271--386.  %%QUATERNION, %%QUANTUM-PHYSICS, %%CLIFFORD, %%HESTENES, %%DIRAC, %%SPHERICAL-HARMONICS, %$D24062002.

\bibitem{SLATE1996-} P.B. Slater, \emph{Bayesian inference for complex and quaternionic two-level quantum systems}, Physica A {\bf 223} (1966) 167--174. %%QUATERNION, %%QUANTUM-PHYSICS, %%DENSITY-OPERATOR, %$D24062002.

\bibitem{CASAN1997-}  G. Casanova, \emph{The electron's double nature}, Adv. Appl. Clifford Alg. {\bf 7 (S)} (1997) 163--166. %%QUATERNION, %%QUANTUM-PHYSICS, %%DIRAC, %$D24062002.

\bibitem{CHALL1997-}  A. Challinor, A. Lasenby, S. Somaroo, C. Doran, and S. Gull, \emph{Tunneling times of electrons}, Phys. Lett. A {\bf 227}  (1997) 143--152.  %%QUATERNION, %%QUANTUM-PHYSICS, %$D17032002, %$C NO-HESTENES.

\bibitem{DAVIA1997-}  C. Daviau, \emph{Solutions of the Dirac equation and a nonlinear Dirac equation for the hydrogen atom}, Adv. Appl. Clifford Alg. {\bf 7 (S)} (1997) 175--194. %%QUATERNION, %%QUANTUM-PHYSICS, %%DIRAC, %%HESTENES, %$D17032002.

\bibitem{HESTE1997-}  D. Hestenes, \emph{Real Dirac theory}, Adv. Appl. Clifford Alg. {\bf 7 (S)} (1997) 97--144. %%QUATERNION, %%QUANTUM-PHYSICS, %%DIRAC, %%HESTENES, %$D17032002.

\bibitem{RACIT1997-} F. Raciti and E. Venturino, \emph{Quaternion methods for random matrices in quantum physics}, in: Applications of Clifford algebras and Clifford analysis in physics and engineering, Electronic Proceedings of IKM97, Weimar, Germany, Feb. 26 - March 1, 1997 (26 January 1997) 5~pp.  %%QUATERNION, %%QUANTUM-PHYSICS, %$D04032006.

\bibitem{KOCIK1999-} J. Kocik, \emph{Duplex numbers, diffusion systems, and generalized quantum mechanics}, Int. J. Phys. {\bf 38} (1999) 2221--2230. %%QUATERNION, %%QUANTUM-PHYSICS, %$D24062002, %$C Quantum mechanics cannot be reduced to diffusion theory. "i" is necessary.

\bibitem{RABEI1999-} E.M. Rabei, Arvind, N. Mukunda, and R. Simon, \emph{Bargmann invariants and geometric phase --- A generalized connection}, Phys. Rev. {\bf A60} (1999) 3397--3409. %%QUATERNION, %%QUANTUM-PHYSICS, %%BERRY-PHASE, %$D11032002.

\bibitem{SAUE-1999-} T. Saue and H.J. Jensen, \emph{Quaternion symmetry in relativistic molecular calculations}, J. Chem. Phys. {\bf 111} (1999) 6211--6222. %%QUATERNION, %%QUANTUM-PHYSICS, %%CHEMISTRY, %%DIRAC, %$D30052001.

\bibitem{SOMAR1999-}  S. Somaroo, A. Lasenby, and C. Doran, \emph{Geometric algebra and the causal approach to multiparticle quantum mechanics}, J. Math. Phys. {\bf 40}  (1999) 3327--3340.  %%QUATERNION, %%QUANTUM-PHYSICS, %%HESTENES, %$D17032002.

\bibitem{LEWIS2000-} A. Lewis, A. Lasenby, and C. Doran, \emph{Electron scattering in the spacetime algebra}, in: R. Ablamowicz and B. Fauser, eds., Clifford Algebra and their Applications in Mathematical Physics, Vol.~1: \emph{Algebra and Physics} (Birkhauser, Boston, 2000)  47--71. %%QUATERNION, %%QUANTUM-PHYSICS, %%QED, %%CLIFFORD, %$D12052001.

\bibitem{VISSC2000-} L. Visscher and T. Saue, \emph{Approximate relativistic electronic structure methods based on the quaternion modified Dirac equation}, J. Chem. Phys. {\bf 113} (2000) 3996--4002. %%QUATERNION, %%QUANTUM-PHYSICS, %%DIRAC, %%CHEMISTRY, %$D30052001.

\bibitem{HAVEL2002-}  T.F. Havel and C. Doran, \emph{Interaction and entanglement in the multiparticle spacetime algebra}, in: L. Dorst  
et al., eds., Applications of Geometric Algebra in Computer Science and Engineering (Birkh\"auser, Boston, 2002) 227--247  %%QUATERNION, %%QUANTUM-PHYSICS, %%HESTENES, %$D17032002.

\bibitem{PARKE2002-}  R. Parker and C. Doran, \emph{Analysis of 1 and 2 particle quantum systems using geometric algebra}, in: L. Dorst et al., eds, Applications of Geometric Algebra in Computer Science and Engineering (Birkh\"auser, Boston, 2002) 213--226.  %%QUATERNION, %%QUANTUM-PHYSICS, %%HESTENES, %$D17032002.

\bibitem{DREIS2003-} D.W. Dreisigmeyer, R. Clawson, R. Eykholt, and P.M. Young, \emph{Dynamic and geometric phase formulas in Hestenes-Dirac theory},  Found. of Phys. Letters {\bf 16} (2003) 429--445.  %%QUATERNION, %%QUANTUM-PHYSICS, %%BERRY-PHASE, %%HESTENES, %%GSPONER, %$D30112005.

\bibitem{MUKUN2003-} N. Mukunda, Arvind, S. Chaturvedi, and R. Simon, \emph{Bargmann invariants and off-diagonal geometric phases for multi-level quantum systems --- a unitary group approach}, Phys. Rev. {A 65} (2003) 012102--11. %%QUATERNION, %%QUANTUM-PHYSICS, %%BERRY-PHASE, %$D11032002.

\bibitem{GOTTL2003-} D.H. Gottlieb, \emph{Eigenbundles, quaternions, and Berry's phase} (20 April 2003) 22~pp.;  e-print \underline{ arXiv:math.AT/0304281 }.  %%QUATERNION, %%QUANTUM-PHYSICS, %%BERRY-PHASE, %%GSPONER, %$D03082005.

\bibitem{GUPTA2003-} R.C. Gupta, \emph{Concept of quaternion-mass for wave-particle duality: A novel approach}, Preprint (2003) 13~pp.;  e-print \underline{ arXiv:physics/0305024 }.  %%QUATERNION, %%QUANTUM-PHYSICS, %$D30112005.

\bibitem{SAUE-2003-} T. Saue and H.J. Jensen, \emph{Linear response at the 4-component relativistic-level: Application to the frequency-dependent dipole polarizabilities of the coinage metal dimers}, J. Chem. Phys. {\bf 118} (2003) 522--536. %%QUATERNION, %%QUANTUM-PHYSICS, %%CHEMISTRY, %%DIRAC, %$D30052001.

\bibitem{TCHOU2003-} A.L. Tchougr\'eeff, \emph{$SO(4)$ group and deductive molecular dynamics}, Molecular Structure (Theochem) {\bf 630} (2003) 243--263. %%QUATERNION, %%QUANTUM-PHYSICS, %%CHEMISTRY, %$D04032006.

\bibitem{CELER2004-} M.-N. C\'el\'erier and L. Nottale, \emph{Quantum-classical transition in scale relativity}, J. of Physics A: Math. Gen. {\bf 37} (2004) 931--955. %%QUATERNION, %%QUANTUM-PHYSICS, %%LANCZOS, %%CONWAY, %$D04032006, %$C Reference to Lanczos and Conway identical as in {CAPOZ2005-}.

\bibitem{TUCCI2005-} R.R. Tucci, \emph{An Introduction to Cartan's KAK Decomposition for QC Programmers} (18 July 2005) 12~pp.; e-print \underline{  arXiv:quant-ph/0507171 }.   %%QUATERNION, %%QUANTUM-PHYSICS, %%MATRIX, %%GROUP-THEORY, %%GSPONER, %$D06072008.

\bibitem{NOBIL2006-} R. Nobili, \emph{Fourteen steps into quantum mechanics}, HTML document (Posted in 2006) about 13 pp.;  available at\\ \underline{ http://www.pd.infn.it/~rnobili/qm/14steps/14steps.htm }. %%QUATERNION, %%QUANTUM-PHYSICS, %%QQM, %$D04032006.

\bibitem{NOTTA2007-} L. Nottale and M.-N. C�l�rier, \emph{Derivation of the postulates of quantum mechanics from the first principles of scale relativity}, J. Phys. A: Math. Theor. {\bf 40} (2007)  14471--14498; e-print \underline{ arXiv:0711.2418 }.  %%QUATERNION, %%QUANTUM-PHYSICS, %$D06072008.

\bibitem{FRENK2008-} I. Frenkel and M. Libine, \emph{Quaternionic analysis, representation theory, and physics} (2008) 63~pp.; e-print \underline{ arXiv:0711.2699 }.  %%QUATERNION, %%QUANTUM-PHYSICS, %%ANALYTICITY, %$D06072008.

  
\end{enumerate}

\subsection{QUANTUM-ELECTRODYNAMICS}
%-----------------------------------
\label{QUANTUM-ELECTRODYNAMICS}

Quantum relativistic perturbation-theory of electromagnetic interactions.

Most papers relate to the attempt of R.P. Feynman's student L.M. Brown to reformulate quantum electrodynamics using, instead of the usual four-component-spinor Dirac matrix formalism, the two-component-spinor Pauli-matrix formalism which is essentially equivalent to that of biquaternions due to the isomorphism $\mathbb{B} \sim M_2(\mathbb{C})$. 

\begin{enumerate}


\bibitem{FEYNM1951-} R.P. Feynman, \emph{An operator calculus having applications in quantum electrodynamics}, Phys. Rev. {\bf 84} (1951) 108--128. %%QUATERNION, %%QUANTUM-ELECTRODYNAMICS, %%$D30112005, %%$C Feynman mentions quaternions and hypercomplex numbers on page 126 and 128.

\bibitem{BROWN1958-} L.M. Brown, \emph{Two-component fermion theory}, Phys. Rev. {\bf 109} (1958) 193--198. %%NO-QUATERNION, %%QUANTUM-ELECTRODYNAMICS, %%TWO-COMPONENT-SPINOR, %%$D30112005.

\bibitem{KIBBL1958-} T.W.B. Kibble and J.C. Polkinghorne, \emph{Higher order spinor lagrangians}, Nuovo Cim. {\bf 8} (1958) 74--83. %%NO-QUATERNION, %%QUANTUM-ELECTRODYNAMICS, %%TWO-COMPONENT-SPINOR, %%$D30112005.

\bibitem{TONIN1959-} M. Tonin, \emph{Quantization of the two-component fermion theory}, Nuovo Cim. {\bf 14} (1959) 1108--1119. %%NO-QUATERNION, %%QUANTUM-ELECTRODYNAMICS, %%TWO-COMPONENT-SPINOR, %%$D30112005.

\bibitem{PIETS1961-} H. Pietschmann, \emph{Zur Renormierung der zweikomponentigen Quantenelektrodynamik}, Acta Physica Austriaca {\bf 14} (1961) 63--74. %%NO-QUATERNION, %%QUANTUM-ELECTRODYNAMICS, %%TWO-COMPONENT-SPINOR, %%$D30112005.

\bibitem{BARUT1962-} A.O. Barut and G.H. Mullen, \emph{Quantization of two-component higher order spinor equations}, Ann. Phys. {\bf 20} (New York, 1962) 184-202. %%NO-QUATERNION, %%QUANTUM-ELECTRODYNAMICS, %%TWO-COMPONENT-SPINOR, %%$D30112005.

\bibitem{BROWN1962-} L.M. Brown, \emph{Two-component fermion theory}, Lectures in Theoretical Physics {\bf 4} (Interscience, New York, 1962) 324--357. %%NO-QUATERNION, %%QUANTUM-ELECTRODYNAMICS, %%TWO-COMPONENT-SPINOR, %%$D30112005.

\bibitem{ROTEL1989B} P. Rotelli, \emph{Quaternion trace theorems and first order electron-muon scattering}, Mod. Phys. Lett. {\bf A 4} (1989) 1763--1771.  %%QUATERNION, %%QUANTUM-ELECTRODYNAMICS, %$D16032002.
 
\bibitem{CHO---1990-} H.T. Cho, A. Diek, and R. Kantowski, \emph{A Clifford algebra quantization of Dirac's electron-positron field}, J. Math. Phys. {\bf 31} (1990) 2192--2200. %%QUATERNION, %%QUANTUM-ELECTRODYNAMICS, %%CLIFFORD, %$D30112005.

\bibitem{FAUSE1997-} B. Fauser and H. Stumpf, \emph{Positronium as an example of algebraic composite calculations}, in: J. Keller and Z. Oziewicz, eds., The Theory of the Electron, Adv. Appl. Clifford Alg. {\bf 7 (S)} (1997) 399--418.  %%QUATERNION, %%QUANTUM-ELECTRODYNAMICS, %%CLIFFORD, %$D24062002.

\bibitem{FAUSE2000-} B. Fauser and R. Ablamowicz, \emph{On the decomposition of Clifford algebras of arbitrary bilinear form}, in: R. Ablamowicz and B. Fauser, eds., Clifford Algebra and their Applications in Mathematical Physics, Vol.~1: \emph{Algebra and Physics} (Birkh\"auser, Boston, 2000)  341--366.  %%QUATERNION, %%QUANTUM-ELECTRODYNAMICS, %%CLIFFORD, %%MULTI-PARTICLE-STATES, %%QUANTUM-FIELD-THEORY, %%NORMAL-ORDERING, %$D12052001.

\bibitem{VARLA2002-} V.V. Varlamov, \emph{About algebraic foundations of Majorana-Oppenheimer quantum electrodynamics and de Broglie-Jordan neutrino theory of light}, Ann. Fondation L. de Broglie {\bf 27} (2002) 273--286.   %%QUATERNION, %%QUANTUM-ELECTRODYNAMICS, %$D06072008.

\bibitem{VARLA2006-} V.V. Varlamov, \emph{A note on the Majorana-Oppenheimer quantum electrodynamics} (2006) 13~pp.; e-print \underline{ arXiv:math-ph/0206008 }.   %%QUATERNION, %%QUANTUM-ELECTRODYNAMICS, %$D06072008.


\end{enumerate}

\subsection{QUATERNIONIC-QUANTUM-PHYSICS}
%----------------------------------------
\label{QUATERNIONIC-QUANTUM-PHYSICS}


Generalized quantum physics in which real quaternions are used instead of complex numbers as the scalar field. 

In the \TeX source of the bibliography the keyword QUATERNIONIC-QUANTUM-PHYSICS is abreviated QQPH, and there two special keywords: QQPH/CERN for the publications of the CERN-University-of-Geneva school (led by J.M.\ Jauch) and QQPH/ALDER for the papers of S.L.\ Adler and his followers. 


\begin{enumerate}

\bibitem{JORDA1932-} P. Jordan, \emph{Uber eine Klasse nichassoziativer hyperkomplexer Algebren}, Nachr. Ges. Wiss. G\"ottingen. {\bf 33} (1932) 569--575. %%QUATERNION, %%QQPH, %$D1503200, %$C-First "Jordan algebra".

\bibitem{JORDA1933A} P. Jordan, \emph{Uber die Multiplikation  quantenmechanischer Grossen. I.}, Zeits. f. Phys. {\bf 80} (1933) 285--291. %%QUATERNION, %%QQPH, %$D06022002.

\bibitem{JORDA1933B} P. Jordan, \emph{Uber Verallgemeinerungsm\"oglichkeiten des Formalismus Quantenmechanik}, Nachr. Ges. Wiss. G\"ottingen. {\bf 39} (1933) 209--217. %%QUATERNION, %%QQPH, %$D15032002, %$C-First suggestion of an algebraic generalization of quantum theory.

\bibitem{JORDA1934A} P. Jordan, J. vonNeumann, and E. Wigner, \emph{On a generalization of the quantum mechanical formalism}, Ann. of Math. {\bf 35} (1934) 29--64.%%QUATERNION, %%QQPH, %$D06022002.

\bibitem{JORDA1934B} P. Jordan, \emph{Uber die Multiplikation  quantenmechanischer Grossen. II.}, Zeits. f. Phys. {\bf 81} (1934) 505--512. %%QUATERNION, %%QQPH, %$D06022002.

\bibitem{TEICH1935-} O. Teichm\"uller, \emph{Operatoren im Wachsschen Raum}, J. f\"ur Math. {\bf 174} (1935) 73--124. %%BOOK, %%QUATERNION, %%QQPH, %$D20032002, %$C The "Wachs-space" is the quaternionic generalization of the "Hilbert-space,"  see \cite{WACHS1936-}.

\bibitem{BIRKH1936-} G. Birkhoff and J. vonNeumann, \emph{The logic of quantum mechanics}, Ann. of Math. {\bf 37} (1936) 823--843. %%QUATERNION, %%QQPH, %$D15032002.

\bibitem{FINKE1959A} D. Finkelstein, J.M. Jauch, and D. Speiser, \emph{Zur Frage der Ladungsquantisierung}, Helv. Phys. Acta {\bf 32} (1959) 258--250. %%QUATERNION, %%QQPH/CERN, %$D15032002.

\bibitem{FINKE1959B} D. Finkelstein, J.M. Jauch, and D. Speiser, \emph{Notes on quaternion quantum mechanics. I, II, III.},  Reports 59-7, 59-9, 59-17 (CERN, 1959).  Published in: C.A. Hooker, ed., Logico-Algebraic Approach to Quantum Mechanics.~II. (Reidel, Dordrecht, 1979) 367--421. %%QUATERNION, %%QQPH/CERN, %$D15032002.

\bibitem{STUCK1959-} E.C.G. Stuckelberg, \emph{Field quantization and time reversal in real Hilbert space}, Helv. Phys. Acta {\bf 32} (1959) 254--256. %%QUATERNION, %%QQPH/CERN, %%PSEUDOQUATERNION, %%DYSON?, %$D15032002.

\bibitem{KANEN1960-} T. Kaneno, \emph{On a possible generalization of quantum mechanics}, Prog. Th. Phys. {\bf 23} (1960) 17--31. %%QUATERNION, %%QQPH, %%PARTICLE-PHYSICS, %$D15032002.

\bibitem{FINKE1962A} D. Finkelstein, J.M. Jauch, S. Schiminovich, and D. Speiser, \emph{Foundations of quaternion quantum mechanics}, J. Math. Phys. {\bf 3} (1962) 207--220. %%QUATERNION, %%QQPH/CERN, %$D15032002.

\bibitem{FINKE1962B} D. Finkelstein, J.M. Jauch, S. Schiminovich, and D. Speiser, \emph{Appendix: Quaternionic Hilbert space}, J. Math. Phys. {\bf 3} (1962) 218--220. %%QUATERNION, %%QQPH/CERN, %$D15032002, %$C Appendix to \cite{FINKE1962A}.

\bibitem{EMCH-1963A} G. Emch, \emph{M\'ecanique quantique quaternionienne et relativit\'e restreinte I}, Helv. Phys. Acta {\bf 36} (1963) 739--769. %%QUATERNION, %%QQPH/CERN,  %$D15032002.

\bibitem{EMCH-1963B} G. Emch, \emph{M\'ecanique quantique quaternionienne et relativit\'e restreinte II}, Helv. Phys. Acta {\bf 36} (1963) 770--788. %%QUATERNION, %%QQPH/CERN,  %$D15032002.

\bibitem{FINKE1963B} D. Finkelstein, J.M. Jauch, S. Schiminovich, and D. Speiser, \emph{Principle of general Q covariance}, J. Math. Phys. {\bf 4} (1963) 788--796. %%QUATERNION, %%QQPH/CERN, %$D14032002.

\bibitem{BARGM1964-} V. Bargmann, \emph{Appendix: Wigner's theorem in quaternion quantum theory}, J. Math. Phys. {\bf 5} (1964) 866--868. %%QUATERNION, %%QQPH, %$D30112005, %$C Appendix to: \emph{Note on Wigner's theorem on symmetry operations}.

\bibitem{EMCH-1965-} G. Emch, \emph{Representations of the Lorentz group in quaternionic quantum mechanics},  in: W.E. Brittin and A.O. Barut, eds., Lect. in Th. Phys. {\bf 7A}, {Lorentz Group} (University of Colorado, Boulder, 1965) 1--36.  %%QUATERNION, %%QQPH/CERN, %%GROUP-THEORY, %%QUANTUM-PHYSICS,  %$D15032002.

\bibitem{HORWI1965-} L.P. Horwitz and L.C. Biedenharn, \emph{Instrinsic superselection rules of algebraic Hilbert space}, Helv. Phys. Acta {\bf 38} (1965) 385--408.  %%QUATERNION, %%QQPH, %$D30112005.

\bibitem{TAVEL1965-} M. Tavel, D. Finkelstein, and S. Schiminovich, \emph{Weak and electromagnetic interactions in quaternion quantum mechanics}, Bull. Am. Phys. Soc. {\bf 9} (1965) 436. %%QUATERNION, %%QQPH/CERN,  %$D15032002.

\bibitem{KYRAL1967A} A. Kyrala, \emph{Quaternion form of wave-particle dualism}, Section 9.3 of Theoretical Physics: Applications of Vectors, Matrices, Tensors and Quaternions (W.B. Saunders, Philadelphia, 1967) 374~pp.  %%QUATERNION, %%QQPH, %$D15032002.

\bibitem{FINKE1979-} D. Finkelstein, J.M. Jauch, and D. Speiser, \emph{Notes on quaternion quantum mechanics. I, II, III.}, in: C.A. Hooker, ed., Logico-Algebraic Approach to Quantum Mechanics.~II. (Reidel, Dordrecht, 1979) 367--421. %%QUATERNION, %%QQPH/CERN, %$D15032002.

\bibitem{HORWI1979B}  L.P. Horwitz, D. Sepaneru, and L.C. Biedenharn, \emph{Quaternion quantum mechanics}, Ann. Israel. Phys. Soc. {\bf 3} (1980) 300--306.  %%QUATERNION, %%QQPH,  %$D15032002.

\bibitem{PERES1979-} A. Peres, \emph{Proposed test for complex versus quaternion quantum theory}, Phys. Rev. Lett. {\bf 42} (1979) 683--686.  %%QUATERNION, %%QQPH, %%GSPONER, %$D15032002.

\bibitem{REMBI1979-} J. Rembielinski, \emph{Notes on the proposed test for complex versus quaternionic quantum theory}, Phys. Lett. {\bf B 88} (1979) 279--281.  %%QUATERNION, %%QQPH,  %$D15032002.

\bibitem{SOUCE1979-} J. Soucek, \emph{Quaternion quantum mechanics as a description of tachyons and quarks}, Czech. J. Phys. {\bf B 29} (1979) 315--318.  %%QUATERNION, %%QQPH,  %$D15032002.

\bibitem{WOLFF1979-} U. Wolff, \emph{Some remarks on quantum theory with hypercomplex numbers and gauge theory}, Preprint MPI-PAE/PTh 26/79 (Max-Planck Institute, July 1979).  %%QUATERNION, %%QQPH, %$D15032002.

\bibitem{BIEDE1981B} L. C. Biedenharn and L. P. Horwitz, \emph{Nonassociative algebras and exceptional gauge groups}, in: J. Ehlers et al., Differential geometric methods in mathematical physics (Proc. Internat. Conf., Tech. Univ. Clausthal, Clausthal-Zellerfeld, 1978), Lecture Notes in Phys. {\bf 139} (Springer, Berlin-New York, 1981) 152--166. %%QUATERNION, %%QQPH, %%PARTICLE-PHYSICS, %%GAUGE-THEORY, %$D30112005.

\bibitem{SOUCE1981-} J. Soucek, \emph{Quaternion quantum mechanics as a true $3$+$1$-dimensional theory of tachyons}, J. Phys. {\bf A 14} (1981) 1629--1640.  %%QUATERNION, %%QQPH,  %$D15032002.

\bibitem{TRUIN1981-} P. Truini, L.C. Biedenharn, and G. Cassinelli, \emph{Imprimitivity theorem and quaternionic quantum mechanics}, Hadronic Journal {\bf 4} (1981) 981--994.  %%QUATERNION, %%QQPH,  %$D15032002.

\bibitem{WOLFF1981-} U. Wolff, \emph{A quaternion quantum system}, Phys. Lett. {\bf A 84} (1981) 89--92.   %%QUATERNION, %%QQPH,  %$D15032002.

\bibitem{SOFFE1983-} A. Soffer and L.P. Horwitz, \emph{$B^*$-algebra representations in a quaternionic Hilbert module}, J. Math. Phys. {\bf 24} (1983) 2780--2782. %%QUATERNION, %%QQPH, %$D30112005.

\bibitem{HORWI1984-} L.P. Horwitz and L.C. Biedenharn, \emph{Quaternion quantum mechanics: second quantization and gauge fields}, Ann. of Phys. {\bf 157} (1984) 432--488. Errata, Ann. of Phys. {\bf 159} (1985) 481.  %%QUATERNION, %%QQPH,  %$D15032002.

\bibitem{KAISE1984-} H. Kaiser, E.A. George, and S.A. Werner, \emph{Neutron interferomagnetic search for quaternions in quantum mechanics}, Phys. Rev. {\bf A 29} (1984) 2276--2279.  %%QUATERNION, %%QQPH,  %$D15032002.

\bibitem{ADLER1985-} S.L. Adler, \emph{Quaternionic quantum field theory}, Phys. Rev. Lett. {\bf 55} (1985) 783--786. %%QUATERNION, %%QQPH/ADLER, %$D15032002.

\bibitem{ADLER1986A} S.L. Adler, \emph{Quaternionic quantum field theory}, Comm. Math. Phys. {\bf 104} (1986) 611--656. %%QUATERNION, %%QQPH/ADLER, %$D15032002.

\bibitem{ADLER1986B} S.L. Adler, \emph{Time-dependent perturbation theory for quaternionic quantum mechanics, with application to CP non-conservation in K-meson decays}, Phys. Rev. {\bf D 34} (1986) 1871--1877. %%QUATERNION, %%QQPH/ADLER, %$D15032002.

\bibitem{ADLER1986C} S.L. Adler, \emph{Superweak $CP$\, nonconservation arising from an underlying quaternionic quantum dynamics}, Phys. Rev. Lett. {\bf 57} (1986) 167--169. %%QUATERNION, %%QQPH/ADLER, %$D15032002.

\bibitem{ADLER1986D} S.L. Adler, \emph{Quaternionic field theory and a possible dynamics for composite quarks and leptons}, Proceedings of the rencontres de Moriond (1986) 11~pp. %%QUATERNION, %%QQPH/ADLER, %$D15032002.

\bibitem{ADLER1987-} S.L. Adler, \emph{Quaternionic Gaussian multiple integrals}, 601--629, in: I. Batalin et al., eds., Quantum Field Theory and Quantum Statistics: Essays in Honor of the 60th Birthday of E.S. Fradkin, Vol.~1 (Adam Hilger, Bristol, 1987) 601--623. %%QUATERNION, %%QQPH/ADLER, %$D20032002.

\bibitem{NASH-1987A} C.G. Nash and G.C. Joshi, \emph{Spontaneous symmetry breaking and the Higgs  mechanism for quaternion fields},  J. Math. Phys. {\bf 28} (1987) 463--467.  %%QUATERNION, %%QQPH, %%GAUGE-THEORY, %$D30112005.

\bibitem{NASH-1987B} C.G. Nash and G.C. Joshi, \emph{Composite systems in quaternionic quantum mechanics},  J. Math. Phys. {\bf 28} (1987) 2883--2885.  %%QUATERNION, %%QQPH,  %$D15032002.

\bibitem{NASH-1987C} C.G. Nash and G.C. Joshi, \emph{Component states of a composite quaternionic system},  J. Math. Phys. {\bf 28} (1987) 2886--2890.  %%QUATERNION, %%QQPH,  %$D15032002.

\bibitem{SHARM1987-} C.S. Sharma and T.J. Coulson, \emph{Spectral theory for unitary operators on a quaternionic Hilbert space}, J. Math. Phys. {\bf 28} (1987) 1941--1946.   %%QUATERNION, %%QQPH,  %$D15032002.

\bibitem{ADLER1988-} S.L. Adler, \emph{Scattering and decay theory for quaternionic quantum mechanics, and the structure of induced T nonconservation}, Phys. Rev. {\bf D 37} (1988) 3654--3662. %%QUATERNION, %%QQPH/ADLER, %$D15032002.

\bibitem{KLEIN1988-} A.G. Klein, \emph{Schr\"odinger inviolate: Neutron optical search for violations of quantum mechanics}, Physics B {\bf 151} (1988) 44--49. %%QUATERNION, %%QQPH,  %$D15032002.

\bibitem{SHARM1988-} C.S. Sharma, \emph{Complex structure on a real Hilbert space and symplectic structure on a complex Hilbert space}, J. Math. Phys. {\bf 29} (1988) 1067--1078.  %%QUATERNION, %%QQPH,  %$D15032002.

\bibitem{ADLER1989-} S.L. Adler, \emph{A new embedding of quantum electrodynamics in a non-Abelian gauge structure}, Phys. Lett. {\bf B 221} (1989) 39--43. %%QUATERNION, %%QQPH/ADLER, %%QUANTUM-ELECTRODYNAMICS, %$D15032002.

\bibitem{DAVIE1989-} A.J. Davies and B.H.J. McKellar, \emph{Nonrelativistic quaternionic quantum mechanics in one dimension}, Phys. Rev. {\bf A 40} (1989) 4209--4214. %%QUATERNION, %%QQPH,  %$D15032002.

\bibitem{RAZON1989-} A. Razon, L.P. Horwitz, and L.C. Biedenharn, \emph{On a basic theorem of quaternion modules}, J. Math. Phys. {\bf 30} (1989) 59.  %%QUATERNION, %%QQPH, %$D30112005.

\bibitem{SHARM1989B} C.S. Sharma and D.F. Almedia, \emph{Additive functionals and operators on a quaternionic Hilbert space}, J. Math. Phys. {\bf 30} (1989) 369--375. %%QUATERNION, %%QQPH,  %$D15032002.

\bibitem{ADLER1990-} S.L. Adler, \emph{Scattering theory in quaternionic quantum mechanics}, in: A. Das., ed., From Symmetries to Strings: Forty Years of Rochester Conferences (World Scientific, River Edje NY, 1990) 37--56.  %%QUATERNION, %%QQPH/ADLER, %$D15032002.

\bibitem{DAVIE1990-} A.J. Davies, \emph{Quaternionic Dirac equation}, Phys. Rev. {\bf D 41} (1990) 2628--2630. %%QUATERNION, %%QQPH, %%DIRAC, %$D15032002.

\bibitem{SHARM1990-} C.S. Sharma and D.F. Almedia, \emph{Additive isometries on a quaternionic Hilbert space}, J. Math. Phys. {\bf 31} (1990) 1035--1041. %%QUATERNION, %%QQPH,  %$D15032002.

\bibitem{ADLER1991-} S.L. Adler, \emph{Linear momentum and angular momentum in quaternionic quantum mechanics}, in: M. Kaku et al., eds., Quarks, Symmetries and Strings (World Scientific, Singapore, 1991) 253--255. %%QUATERNION, %%QQPH/ADLER, %$D15032002.

\bibitem{HORWI1991A} L.P. Horwitz and A. Razon, \emph{Tensor product of quaternion Hilbert modules}, Acta Applicandae Mathematicae {\bf 24} (1991) 141--178.  %%QUATERNION, %%QQPH, %$D15032002.

%%\bibitem{HORWI1991B} L.P. Horwitz, A. Razon, \emph{Tensor product of quaternion Hilbert modules II}, Acta Appl. Math. {\bf 24} (1991) 179--.  %%MISS

\bibitem{LEVAY1991-} P. L\'evay, \emph{Quaternionic gauge-fields and the geometric phase}, J. Math. Phys. {\bf 32} (1991) 2347--2357. %%QUATERNION, %%QQPH, %%BERRY-PHASE, %$D30112005.

\bibitem{DAVIE1992-} A.J. Davies and B.H.J. McKellar, \emph{Observability of quaternionic quantum mechanics}, Phys. Rev. {\bf A 46} (1992) 3671--3675. %%QUATERNION, %%QQPH, %$D15032002.

\bibitem{DELEO1992-} S. DeLeo and P. Rotelli, \emph{The quaternion scalar field}, Phys. Rev. {\bf D 45} (1992) 580--585. %%QUATERNION, %%QQPH, %$D15032002.

\bibitem{MARCH1992-} S. Marchiafava and J. Rembielinski, \emph{Quantum quaternions}, J. Math. Phys. {\bf B 33} (1992) 171--173. %%QUATERNION, %%QQPH, %%NONCOMMUTATIVE-GEOMETRY, %$D30112005.

\bibitem{NASH-1992-} C.G. Nash and G.C. Joshi, \emph{Quaternionic quantum mechanics is consistent with complex quantum mechanics},  Int. J. Th. Phys. {\bf 31} (1992) 965--981. %%QUATERNION, %%QQPH, %$D15032002.

\bibitem{RAZON1992-} A. Razon and L.P. Horwitz, \emph{Uniqueness of the scalar product in the tensor product of quaternion Hilbert modules}, J. Math. Phys. {\bf 33} (1992) 3098--3104. %%QUATERNION, %%QQPH, %$D15032002.

\bibitem{CONTE1993A} E. Conte, \emph{On a generalization of quantum mechanics by biquaternions}, Hadronic Journal {\bf 16} (1993) 261--275.  %%QUATERNION, %%QQPH, %$D15032002.

\bibitem{CONTE1993B} E. Conte, \emph{An example of wave packet reduction using biquaternions}, Physics Essays {\bf 6} (1993) 532--535.  %%QUATERNION, %%QQPH, %$D24062002.

\bibitem{HANLO1993-} B.E. Hanlon and G.C. Joshi, \emph{Spontaneous CP violation from a quaternionic Kaluza-Klein theory}, Int. J. Mod. Phys. {\bf A 8} (1993) 3263--3283. %%QUATERNION, %%QQPH, %$D15032002.

\bibitem{HORWI1993-} L.P. Horwitz, \emph{Some spectral properties of anti-self-adjoint operators on a quaternionic Hilbert space}, J. Math. Phys. {\bf 34} (1993) 3405--3419. %%QUATERNION, %%QQPH, %$D15032002.

\bibitem{ADLER1994-} S.L. Adler, \emph{Generalized quantum dynamics}, Nucl. Phys. {\bf B 415} (1994) 195--242.  %%QUATERNION, %%QQPH/ADLER, %$D15032002.

\bibitem{CONTE1994-} E. Conte, \emph{Wave function collapse in biquaternion quantum mechanics}, Physics Essays {\bf 7} (1994) 429--435.   %%QUATERNION, %%QQPH, %$D24062002.

\bibitem{DELEO1994-} S. DeLeo and P. Rotelli, \emph{Translations between quaternion and complex quantum mechanics}, Prog. Th. Phys. {\bf 92} (1994) 917--926. %%QUATERNION, %%QQPH, %$D15032002.

\bibitem{HORWI1994A} L.P. Horwitz, \emph{A soluble model for scattering and decay in quaternionic quantum mechanics. I: Decay}, J. Math. Phys. {\bf 35} (1994) 2743--2760. %%QUATERTNION, %%QQPH, %$D30112005.

\bibitem{HORWI1994B} L.P. Horwitz, \emph{A soluble model for scattering and decay in quaternionic quantum mechanics. II: Scattering}, J. Math. Phys. {\bf 35} (1994) 2761--2771.  %%QUATERTNION, %%QQPH, %$D30112005.

\bibitem{ADLER1995-} S.L. Adler, Quaternionic Quantum Mechanics and Quantum Fields (Oxford University Press, Oxford, 1995) 586~pp. %%BOOK, %%QUATERNION, %%QQPH/ADLER, %$D15032002.

\bibitem{HOLLA1995-} S.S. Holland, Jr., \emph{Projections algebraically generate the bounded operators on real or quaternionic Hilbert-space}, Proc. Am. Math. Soc. {\bf 123} (1995) 3361--3362. %%QUATERNION, %%QQPH, $C Hilbert spaces with scalars in R or H, %$D30112005.

\bibitem{ADLER1996A} S.L. Adler, \emph{Projective group representations in quaternionic Hilbert space}, J. Math. Phys. {\bf 37} (1996) 2352--2360. %%QUATERNION, %%QQPH/ADLER, %$D15032002.

\bibitem{ADLER1996B} S.L. Adler, \emph{Response to the Comment by G. Emch on projective group representations in quaternionic Hilbert space}, J. Math. Phys. {\bf 37} (1996) 6586--6589. %%QUATERNION, %%QQPH/ADLER, %$D15032002.

\bibitem{ADLER1996C} S.L. Adler and J. Anandan, \emph{Nonadiabatic geometric phase in quaternionic Hilbert space}, Found. Phys. {\bf 26} (1996) 1579--1589. %%QQPH/ADLER, %%BERRY-PHASE, %$D11032002. 

\bibitem{BRUMB1996A} S.P. Brumby and G.C. Joshi, \emph{Experimental status of quaternionic quantum mechanics}, Chaos, Solitons \& Fractals {\bf 7} (1996) 747--752. %%QUATERNION, %%QQPH, %$D15032002.

\bibitem{BRUMB1996B} S.P. Brumby and G.C. Joshi, \emph{Global effects in quaternionic quantum field theory}, Found. of Phys. {\bf 26} (1996) 1591--1599. %%QUATERNION, %%QQPH, %$D15032002.

\bibitem{DELEO1996E} S. DeLeo and P. Rotelli, \emph{Odd dimensional translation between complex and quaternionic quantum mechanics}, Prog. Theor. Phys. {\bf 96} (1996) 247--255. %%QUATERNION, %%QQPH, %$D15032002.

\bibitem{EMCH-1996-} G.G. Emch, \emph{Comments on a recent paper by S. Adler on projective group representations in quaternionic Hilbert space}, J. Math. Phys. {\bf 37} (1996) 6582--6585. %%QUATERNION, %%QQPH/ADLER, %$D15032002.

\bibitem{HORWI1996A} L.P. Horwitz, \emph{Hypercomplex quantum mechanics}, Found. of Phys. {\bf 26} (1996) 851--862. %%QUATERNION, %%QQPH,  %$D15032002.

%%\bibitem{HORWI1996B} \emph{Bibliography of Lawrence P. Horwitz}, Found. of Phys. {\bf 26} (1996) 1711--1723. %%OK-11102005

\bibitem{PERES1996-} A. Peres, \emph{Quaternionic quantum interferometry}, in: F. DeMartini, G. Denardo and Y. Shih, eds., Quantum Interferometry, Proc. of an Adriatico Workshop (ICTP, Trieste, 1996) 431--437.  %%QUATERNION, %%QQPH,  %$D15032002.

\bibitem{TAO--1996-} T. Tao and A.C. Millard, \emph{On the structure of projective group representations in quaternionic Hilbert space}, J. Math. Phys. {\bf 37} (1996) 5848--5857.  %%QUATERNION, %%QQPH, %%GROUP-THEORY, %$D30112005.

\bibitem{ADLER1997-} S.L. Adler and G.G. Emch, \emph{A rejoinder on quaternionic projective representations}, J. Math. Phys. {\bf 38} (1997) 4758--4762. %%QUATERNION, %%QQPH/ADLER, %$D15032002.

\bibitem{BRUMB1997-} S.P. Brumby, B.E. Hanlon, and G.C. Joshi, \emph{Implications of quaternionic dark matter}, Phys. Lett. {\bf B 401} (1997) 247--253. %%QUATERNION, %%QQPH,  %$D15032002.

\bibitem{DELEO1997B} S. DeLeo and W.A. Rodrigues, Jr., \emph{Quantum mechanics: from complex to complexified quaternions}, Int. J. Th. Phys. {\bf 36} (1997) 2725--2757.  %%QUATERNION, %%QQPH, %$D15032002.

\bibitem{HORWI1997-} L.P. Horwitz, \emph{Schwinger algebra for quaternionic quantum mechanics}, Found. of Phys. {\bf 27} (1997) 1011--1034. %%QUATERNION, %%QQPH,  %$D15032002.

\bibitem{EMCH-1998-} G.G. Emch and A.Z. Jadczyk, \emph{On quaternions and monopoles}, in: F. Gesztesy et al. (eds) Stochastic Processes, Physics and  
Geometry; Can. Math. Soc. Conference Proceedings Series, ``Stochastic  
Processes, Physics and Geometry: New Interplays. A Volume in Honor of  
Sergio Albeverio.'' (Amer. Math. Soc., 2000) 333~pp; e-print \underline{ arXiv:quant-ph/9803002 }. %%QUATERNION, %%QQPH/ADLER, %$D24062002.

\bibitem{KAROW1999-} M. Karow, \emph{Self-adjoint operators and pairs of Hermitian forms over the quaternions}, Linear Alg. Appl. {\bf 299} (1999) 101--117.  %%QUATERNION, %%QQPH, %$D30112005.

\bibitem{DELEO2000-} S. DeLeo and G. Scolarici, \emph{Right eigenvalue in quaternionic quantum mechanics}, J. Phys. A: Math. Gen. {\bf 33} (2000) 2971--2995. %%QUATERNION, %%QQPH, %$D15032002.

\bibitem{SCOLA2000A} G. Scolarici and L. Solombrino, \emph{Quaternionic symmetry groups and particle multiplets}, J. Math. Phys. {\bf 41} (2000) 4950--4603. %%QUATERNION, %%QQPH, %%GROUP-THEORY, %$D04022002.

\bibitem{SCOLA2000B} G. Scolarici and L. Solombrino, \emph{$t$-violation and quaternionic state oscillations}, J. Phys. A: Math. Gen. {\bf 33} (2000) 7827--7838. %%QUATERNION, %%QQPH, %%CP-VIOLATION, %$D04022002.

\bibitem{MAIA-2001-} M.D. Maia and V.B. Bezerra, \emph{Geometric phase in quaternionic quantum mechanics}, Int. J. Theor. Phys. {\bf 40} (2001) 1283--1295. %%QUATERNION, %%QQPH, %%BERRY-PHASE, %$D11032002.

\bibitem{DELEO2002-} S. DeLeo, G. Ducati, and C.C. Nishi, \emph{Quaternionic potentials in non-relativistic quantum mechanics}, J. Phys. A: Math. Gen.  {\bf 35} (2002) 5411--5426. %%QUATERNION, %%QQPH, %%GSPONER, %$D30112005.

\bibitem{DAJKA2003-} J. Dajka and M. Szopa, \emph{Holonomy in quaternionic quantum mechanics}, Int. J. Theor. Phys. {\bf 42} (2003) 1053--1057. %%QUATERNION, %%QQPH, %%GEOMETRIC-PHASE, %$D04032006.

\bibitem{MULAS2003-} M. Mulase and A. Waldron, \emph{Duality of orthogonal and symplectic matrix integrals and quaternionic Feynman graphs}, Commun. Math. Phys. {\bf 240} (2003) 553--586.  %%QUATERNION, %%QQPH, %%MATRIX, %$D30112005.

\bibitem{BLASI2004-} A. Blasi, G. Scolarici and L. Solombrino, \emph{Alternative descriptions in quaternionic quantum mechanics}, J. Math. Phys. {\bf 46} (2005) 042104; e-print \underline{ arXiv:quant-ph/0407158 }.  %%QUATERNION, %%QUATERNIONIC-QUANTUM-PHYSICS, %$D06072008.

\bibitem{JIANG2004B} T. Jiang, \emph{An algorithm for quaternionic linear equations in quaternionic quantum theory}, J. Math. Phys. {\bf 45} (2004) 4218--4222.  %%QUATERNION, %%QUATERNIONIC-QUANTUM-PHYSICS, %$D06072008.

\bibitem{DELEO2005-} S. De Leo abd G. Ducati, \emph{Quaternionic bound states}, J. of Physics A {\bf 38} (2005) 3443--3454. %%QUATERNION, %%QUATERNIONIC-QUANTUM-PHYSICS, %%GSPONER, %%HURNI, %$D06072008.

\bibitem{JIANG2005-} T. Jiang, \emph{Algebraic methods for diagonalization of a quaternion matrix in quaternionic quantum theory}, J. Math. Phys. {\bf 46} (2005) 052106, 8~pp.  %%QUATERNION, %%QUATERNIONIC-QUANTUM-PHYSICS, %$D06072008.

\bibitem{JIANG2006-} T. Jiang, \emph{Cramer rule for quaternionic linear equations in quaternionic quantum theory}, Rep. Math. Phys. {\bf 57} (2006) 463--468.  %%QUATERNION, %%QUATERNIONIC-QUANTUM-PHYSICS, %$D06072008.

\bibitem{ADLER2006-} S.L. Adler, \emph{Quaternionic quantum mechanics, trace dynamics, and emergent quantum theory}, in: S.L. Adler, Adventures in Theoretical Physics --- Selected Papers with Commentaries,  World Scientific Series in 20th Century Physics {\bf 37} (World Scientific, Singapore, 2006) 107--114; e-print \underline{ arXiv:hep-ph/0505177 }.�%%QUATERNION, %%QQPH/ADLER, %$D04032006.  

\bibitem{JIANG2007-} T. Jiang and L. Chen, \emph{Algebraic algorithms for least squares problem in quaternionic quantum theory}, Computer Physics Communications {\bf 176} (2007) 481--485.   %%QUATERNION, %%QUATERNIONIC-QUANTUM-PHYSICS, %$D06072008.

\bibitem{SCOLA2007A} G. Scolarici, \emph{Complex projection of quasianti-Hermitian quaternionic Hamiltonian dynamics}, SIGMA (Symmetry, Integrability and Geometry: Methods and Applications) {\bf  
3} (2007) 088, 10~pp.; e-print \underline{ arXiv:0709.1198 }. %%QUATERNION, %%QUATERNIONIC-QUANTUM-PHYSICS, %$D06072008.


\bibitem{SCOLA2007B} G. Scolarici and L. Solombrino, \emph{Quasistationary quaternionic Hamiltonians and complex stochastic maps} (2007) 9~pp.; e-print \underline{ arXiv:0711.1244 }.   %%QUATERNION, %%QUATERNIONIC-QUANTUM-PHYSICS, %$D06072008.

\end{enumerate}


\section{\Huge ALLIED FORMALISMS}
%================================
\label{ALLIED FORMALISMS}

This chapter contains papers on algebraic formalisms that are closely related to quaternions and biquaternions, either because they are essentially equivalent to them (e.g., Einstein-Mayer's semivectors) or a generalization of them (e.g., Clifford numbers).

However, this chapter covers only the major formalisms allied to quaternions, i.e., only those which have gained some level of popularity, and not the many formalisms that have been, and continue to be, introduced.


Since these formalisms have been invented after the discovery of quaternions, the related sections are listed in the historical order in which they were introduced.\footnote{For an historical account, see Ref.~\cite{VANDE1985-} in Sec.~\ref{HISTORY}.}


\subsection{OCTONION (``Cayley's numbers'')}
%-------------------------------------------
\label{OCTONION}

Octonions were independently discovered, first in 1843 a few months after the discovery of quaternions by Charles Grave (who called them ``octaves''), and then in 1845 by Arthur Cayley (who called them ``biquaternions'').


\begin{enumerate}

\bibitem{MOUFA1933} R. Moufang, \emph{Alternativk\"orper und der Satz vom vollst\"andigen Vierseit $(D_9)$}, Abhandlung. aus dem Math. Seminar Univ. Hamburg {\bf 9} (1933) 207--222. %%QUATERNION, %%OCTONION, %%HILBERT-DESARGUE-THEOREM, %$D19032002.

\bibitem{JORDA1949-} P. Jordan, \emph{Uber die nicht-Desarguessche ebene projektive Geometrie}, Hamb. Abh. {\bf 16} (1949) 74--76. %%QUATERNION, %%OCTONIONS, %$D30112005.

\bibitem{PAIS-1961-} A. Pais, \emph{Remark on the algebra of interactions}, Phys. Rev. Lett. {\bf 7} (1961) 291--293. %%QUATERNION, %%OCTONION, %%PARTICLE-PHYSICS, %$D19032002.

\bibitem{GAMBA1967-} A. Gamba, \emph{Peculiarities of the eight-dimensional space}, J. Math. Phys. {\bf 8} (1967) 775--781. %%QUATERNION, %%OCTONION, %%PARTICLE-PHYSICS, %$D19032002.

\bibitem{PENNE1968-} R. Penney, \emph{Octonions and the Dirac equation}, Am. J. Phys. {\bf 36} (1968) 871--873. %%QUATERNION, %%OCTONION, %%PARTICLE-PHYSICS, %$D19032002.

\bibitem{PENNE1971-} R. Penney, \emph{Octonions and isospin}, Nuovo Cim. {\bf 3 B} (1971) 95--113. %%QUATERNION, %%OCTONION, %%PARTICLE-PHYSICS, %$D19032002.

\bibitem{GUNAY1973A} M. G\"unaydin and F. G\"ursey, \emph{An octonionic representation of the Poincar\'e group}, Lett. Nuovo Cim. {\bf 6} (1973) 401--406. %%QUATERNION, %%OCTONION, %%PARTICLE-PHYSICS, %$D19032002.

\bibitem{GUNAY1973B} M. G\"unaydin and F. G\"ursey, \emph{Quark structure and octonions}, J. Math. Phys. {\bf 14} (1973) 1651--1667. %%QUATERNION, %%OCTONION, %%PARTICLE-PHYSICS, %$D19032002.

\bibitem{GUNAY1974-} M. G\"unaydin and F. G\"ursey, \emph{Quark statistics and octonions}, Phys. Rev. D {\bf 9} (1974) 3387--3391. %%QUATERNION, %%OCTONION, %%PARTICLE-PHYSICS, %$D19032002.

\bibitem{GURSE1975A} F. G\"ursey, P. Ramond, and P. Sikivie, \emph{Six-quark model for the suppression of $\Delta S = 1$ neutral currents}, Phys. Rev. D {\bf 12} (1975) 2166--2168. %%QUATERNION, %%OCTONION, %%PARTICLE-PHYSICS, %%HADRODYNAMICS, %$D19032002.

\bibitem{GURSE1975B} F. G\"ursey, \emph{Algebraic methods and quark structure}, in: H. Araki, ed., Kyoto International Symposium on Mathematical Physics, Lect. Notes in Phys. {\bf 39} (Springer, New York, 1975) 189--195. %%QUATERNION, %%OCTONION, %%PARTICLE-PHYSICS, %%HADRODYNAMICS, %$D30112005.

\bibitem{GURSE1976A} F. G\"ursey, \emph{Charge space, exceptional observables and groups}, in: New Pathways in High-Energy Physics (Plenum Press, 1976) Vol.~1, 231--248. %%QUATERNION, %%OCTONION, %%PARTICLE-PHYSICS, %%OQPH, %$D19032002.

\bibitem{GURSE1976B} F. G\"ursey, P. Ramond, and P. Sikivie, \emph{A universal gauge theory based on $E_6$}, Phys. Lett. {\bf 60B} (1976) 177--180. %%QUATERNION, %%OCTONION, %%PARTICLE-PHYSICS, %$D19032002.

\bibitem{HAYAS1977-} M.J. Hayashi, \emph{A new approach to chromodynamics}, Preprint SLAC-PUB-1936 (May 1977) 16~pp. %%QUATERNION, %%OCTONION, %%PARTICLE-PHYSICS, %$D19032002.

\bibitem{GUNAY1978A} M. G\"unaydin, C. Piron, and H. Ruegg, \emph{Moufang plane and octonionic quantum mechanics}, Comm. Math. Phys. {\bf 61} (1978) 69--85. %%QUATERNION, %%OCTONION, %%OQPH, %$D19032002.

\bibitem{GUNAY1978B} M. G\"unaydin, \emph{Moufang plane and octonionic quantum mechanics}, in: G. Domokos, ed., Second John Hopkins Workshop on Current Problems in Particle Physics (John Hopkins University, Baltimore, 1978) 56--85. %%QUATERNION, %%OCTONION, %%OQPH, %$D19032002.

\bibitem{KOSIN1978-} P. Kosinski and J. Rembielinski, \emph{Difficulties with an octonionic Hilbert space description of the elementary particles}, Phys. Lett. {\bf 79 B} (1978) 309--310. %%QUATERNION, %%OCTONION, %%PARTICLE-PHYSICS, %$D19032002.

\bibitem{NAHM-1978-} W. Nahm, \emph{An octonionic generalization of Yang-Mills}, Preprint TH-2489 (CERN, April 1978) 6~pp. %%QUATERNION, %%OCTONION, %%PARTICLE-PHYSICS, %$D19032002.

\bibitem{REMBI1978-} J. Rembielinski, \emph{Tensor product of the octonionic Hilbert spaces and colour confinement}, J. Phys. {\bf A 11} (1978) 2323--2331. %%QUATERNION, %%OCTONION, %%PARTICLE-PHYSICS, %$D19032002.

\bibitem{RUEGG1978-} H. Ruegg, \emph{Octonionic quark confinement}, Acta Phys. Polonica {\bf B 9} (1978) 1037--1050.  %%QUATERNION, %%OCTONION, %%HADRODYNAMICS, %$D19032002.

\bibitem{HORWI1979A} L.P. Horwitz, D. Sepunaru, and L.C. Biedenharn, \emph{Some quantum aspects of theories with hypercomplex and non-associative structures}, Proc. of the Third Int. Workshop on Current Problems in High Energy Particle Theory (Physics department, John Hopkins University, 1965) 121--153.  %%QUATERNION, %%OCTONIONS, %%QQPH, %$D30112005.

\bibitem{MORIT1979-} K. Morita, \emph{Hypercomplex quark fields and quantum chromodynamics}, Lett. Nuovo Cim. {\bf 26} (1979) 50--54.  %%QUATERNION, %%OCTONION, %%HADRODYNAMICS, %$D30112005.

\bibitem{MORIT1982B} K. Morita, \emph{Algebraic gauge theory of quarks and leptons}, Prog. Th. Phys. {\bf 68} (1982) 2159--2175. %%QUATERNION, %%OCTONION, %%PARTICLE-PHYSICS, %$D15032002.

\bibitem{GUNAY1983-} M. G\"unaydin, G. Sierra, and P.K. Townsend, \emph{Exceptional supergravity theories and the magic square }, Phys. Lett. {\bf 133B} (1983) 72-76. %%QUATERNION, %%OCTONION, %%PARTICLE-PHYSICS, %$D19032002.

\bibitem{GURSE1983B} F. G\"ursey, and C.-H. Tze, \emph{Octonionic torsion on $S^7$ and Englert's compactification of $d=11$ supergravity}, Phys. Lett. {\bf 127B} (1983) 191--196. %%QUATERNION, %%OCTONION, %%PARTICLE-PHYSICS, %$D19032002.

\bibitem{DUNDA1984-} R. Dundarer, F. G\"ursey and C.H. Tze, \emph{Generalized vector products, duality, and octonionic identities in $D=8$ geometry}, J. Math. Phys. {\bf 25} (1984) 1496--1506. %%QUATERNION, %%OCTONION, %$D30112005.

\bibitem{MOFFA1984-} J.W. Moffat, \emph{Higher-dimensional Riemannian geometry and quaternion and octonion spaces}, J. Math. Phys. {\bf 25} (1984) 347--350. %%QUATERNION, %%OCTONION, %%GENERAL-RELATIVITY, %%PSEUDOQUATERNION, %$D17032002.

\bibitem{WENE-1984-} G.P. Wene, \emph{A construction relating Clifford algebras and Cayley-Dickson algebras}, J. Math. Phys. {\bf 25} (1984) 2323--2331. %%QUATERNION, %%OCTONION, %%PARTICLE-PHYSICS, %%REVIEW, %$D19032002.

\bibitem{FUBIN1985-} S. Fubini and H. Nicolai, \emph{The octonionic instanton}, Phys. Lett. {\bf 155B} (1985) 369--372. %%QUATERNION, %%OCTONION, %%PARTICLE-PHYSICS, %$D19032002.

\bibitem{CATTO1985-} S. Catto and F. G\"ursey, \emph{Algebraic treatment of effective supersymmetry}, Nuovo Cimento {\bf A 86} (1985) 201--218. %%QUATERNION, %%OCTONION, %%PARTICLE-PHYSICS, %%SUPERSYMMETRY, %$D04032006.

\bibitem{MARQU1985-} S. Marques, \emph{An extension of quaternionic metrics to octonions}, J. Math. Phys. {\bf 26} (1985) 3131--3139. %%QUATERNION, %%OCTONION, %%GENERAL-RELATIVITY, %$D19032002.

\bibitem{DEALF1986-} V. deAlfaro, S. Fubini, and G. Furlan, \emph{Why we like octonions}, Progr. Theor. Phys. Suppl. {\bf 86} (1986) 274--286. %%QUATERNION, %%OCTONION, %%PARTICLE-PHYSICS, %$D19032002.

\bibitem{DUNDA1986-} R. Dundarer, F. G\"ursey, and C.H. Tze, \emph{Self-duality and octonionic analyticity of $S^7$-valued antisymmetric fields in eight dimensions}, Nuclear Phys. {\bf B 266} (1986) 440--450. %%QUATERNION, %%OCTONION, %%ANALYTICITY, %%FIELD-THEORY, %$D30112005.

\bibitem{FAIRI1986-} D.B. Fairlie and C.A. Manogue, \emph{Lorentz invariance and the composite string}, Phys. Rev. {\bf D34} (1986) 1832--1834. %%QUATERNION, %%OCTONION, %%PARTICLE-PHYSICS, %%SUPERSYMMETRY, %$D30112005.

\bibitem{FAIRI1987-} D.B. Fairlie and C.A. Manogue, \emph{A parametrization of the covariant superstring}, Phys. Rev. {\bf D36} (1987) 475-489. %%QUATERNION, %%OCTONION, %%PARTICLE-PHYSICS, %%SUPERSYMMETRY, %$D30112005.

\bibitem{GODDA1987-} P. Goddard, W. Nahm, D.I. Olive, H. Ruegg, and A. Schwimmer, \emph{Fermions and octonions}, Comm. Math. Phys. {\bf 112} (1987) 385--408.  %%QUATERNION, %%OCTONION, %%PARTICLE-PHYSICS, %%REVIEW, %$D19032002.

\bibitem{IMAED1987-} K. Imaeda, H. Tachibana, M. Imaeda, and S. Ohta, \emph{Solutions of the octonion wave equation and the theory of functions of an octonion variable}, Nuovo Cim. {\bf 100 B} (1987) 53--71. %%QUATERNION, %%OCTONION, %%PARTICLE-PHYSICS, %%ANALYTICITY, %$D19032002.

\bibitem{MARQU1987-} S. Marques, \emph{Geometrical properties of an internal local octonionic space in curved spacetime}, Phys. Rev. D {\bf 36} (1987) 1716--1723. %%QUATERNION, %%OCTONION, %%GENERAL-RELATIVITY, %$D19032002.

\bibitem{MARQU1988-} S. Marques, \emph{The Dirac equation in a non-Riemannian manifold: I. An analysis using the complex algebra}, J. Math. Phys. {\bf 29} (1988) 2127--2131. %%QUATERNION, %%OCTONION, %%GENERAL-RELATIVITY, %%DIRAC, %$D19032002.

\bibitem{MARQU1989-} S. Marques, \emph{Geometrical properties of an internal local octonionic space in a non-Riemannian manifold}, Preprint CBPF-NF-030/1989 (1989) 18~pp. %%QUATERNION, %%OCTONION, %%GENERAL-RELATIVITY, %$D19032002.

\bibitem{MANOG1989-} C.A. Manogue and A. Sudbery, \emph{General solutions of covariant superstring equations of motion}, Phys. Rev. {\bf D40} (1989) 4073--4077. %%QUATERNION, %%OCTONION, %%STRING, %$D19032002.

\bibitem{TACHI1989-} H. Tachibana and K. Imaeda, \emph{Octonions, superstrings and ten-dimensional spinors}, Nuovo Cim. {\bf 104 B} (1989) 91--106. %%QUATERNION, %%OCTONION, %%PARTICLE-PHYSICS, %$D19032002.

\bibitem{MARQU1990-} S. Marques, \emph{The Dirac equation in a non-Riemannian manifold: II. An analysis using an internal local n-dimensional space of the Yang-Mills type}, J. Math. Phys. {\bf 31} (1990) 2127--2131. %%QUATERNION, %%OCTONION, %%GENERAL-RELATIVITY, %%DIRAC %$D19032002.

\bibitem{DUNDA1991-} A.R. Dundarer and F. G\"ursey, \emph{Octonionic representations of ${\rm SO}(8)$ and its subgroups and cosets}, J. Math. Phys. {\bf 32} (1991) 1176--1181. %%QUATERNION, %%OCTONION, %%GROUP-THEORY, %$D30112005.

\bibitem{MARQU1991A} S. Marques, \emph{The Dirac equation in a non-Riemannian manifold: III. An analysis using the algebra of quaternions and octonions}, J. Math. Phys. {\bf 32} (1991) 1383--1394. %%QUATERNION, %%OCTONION, %%GENERAL-RELATIVITY, %%DIRAC, %$D19032002.

\bibitem{WALDR1992-} A.K. Waldron and G.C. Joshi, \emph{Gauging the octonionic algebra}, Preprint UM-P-92/60 (University of Melbourne, 1992) 20~pp. %%QUATERNION, %%OCTONION, %%PARTICLE-PHYSICS, %$D19032002.

\bibitem{DRAY-1999-} T. Dray and C.A. Manogue, \emph{The exceptional Jordan algebra eigenvalue problem}, Int. J. Theor. Phys. {\bf 38} (1999) 2901--2916. %%QUATERNION, %%OCTONION, %%OQPH, %%EQUATION, %$D19032002.

\bibitem{COLOM2000-} F. Colombo, I. Sabadini, and D.C. Struppa, \emph{Dirac equation in the octonionic algebra}, Contemporary. Math. {\bf 251} (2000) 117--134. %%QUATERNION, %%OCTONION, %%DIRAC, %$D19032002.

\bibitem{BAEZ-2002-} J. Baez, \emph{The octonions}, Bull. Amer. Math. Soc. {\bf 39} (2002) 145--205; Errata Bull. Amer. Math. Soc. {\bf 42} (2005) 213; e-print \underline{ arXiv:math/0105155 }.  %%QUATERNION, %%OCTONION, %$D06072008.

\bibitem{LUKIE2002-} J. Lukierski and F. Toppan, \emph{Generalized space-time supersymmetries, division algebras and octonionic M-theory}, Phys. Lett. {\bf B 539} (2002) 266--276; e-print \underline{ arXiv:hep-th/0203149 }.  %%QUATERNION, %%OCTONION, %%PARTICLE-PHYSICS, %%SUPERSYMMETRY, %$D04032006.

\bibitem{LUKIE2003-} J. Lukierski and F. Toppan, \emph{Octonionic M-theory and $D=11$ generalized conformal and superconformal algebras}, Phys. Lett. {\bf B 567} (2003) 125--132; e-print\underline{ arXiv:hep-th/0212201 }.  %%QUATERNION, %%OCTONION, %%PARTICLE-PHYSICS, %%SUPERSYMMETRY, %$D04032006.

\bibitem{CARRI2003A} H.L. Carrion, M. Rojas and F. Toppan, \emph{Octonionic realizations of 1-dimensional extended supersymmetries. A classification}, Mod. Phys. Lett. {\bf A 18} (2003) 787--798; e-print \underline{ arXiv:hep-th/0212030 }.  %%QUATERNION, %%OCTONION, %%PARTICLE-PHYSICS, %%SUPERSYMMETRY, %$D04032006.

\bibitem{CARRI2003B} H.L. Carrion, M. Rojas and F. Toppan,  \emph{Quaternionic and octonionic spinors. A classification}, J. of High Energy Physics {\bf 0304} (2003) 040, 24~pp; e-print \underline{ arXiv:hep-th/0302113 }.  %%QUATERNION, %%OCTONION, %%SPINOR, %$D04032006.

\end{enumerate}


\subsection{GRASSMANN (``Grassmann's calculus of extensions'')}
%-----------------------------------------------
\label{GRASSMANN}

Hermann G. Grassmann's ``calculus of extensions'' of 1844/1862 is the first attempt to define explicitly the notions of ``$n$-dimensional vector space,'' and of ``internal'' and ``external'' products of vectors.

\begin{enumerate}

\bibitem{TAIT-1891-} P.G. Tait, \emph{Quaternions and the Ausdehnungslehre}, Nature {\bf } (4 June 1891) SP-2:456. %%QUATERNION, %%HISTORY, %%GRASSMANN, %$D24062002.

\bibitem{JOLY-1900-} C.J. Joly, \emph{On the place of the Ausdehnungslehre in the general associative algebra of the quaternion type}, Proc. Roy. Irish Acad. {\bf 6} (1900) 13--18. %%QUATERNION, %%GRASSMANN, %$D17032002.

\bibitem{H.C.P.1901-} Book review: \emph{Elements of Quaternions, by Sir W. Hamilton, 2nd edition, edited by C.J. Joly, Vol. II}, Nature {\bf 64}  (1901) 206.  %%QUATERNION, %%APPRECIATION, %$D04032006. 

\bibitem{CHRYS1901-} G. Chrystal, \emph{Obituary notice of Professor Tait}, Nature {\bf 64}  (1901) 305--307.  %%QUATERNION, %%HISTORY, %$D04032006. 

\bibitem{MARKI1936-} M. Markic, \emph{Transformantes nouveau v\'ehicule math\'ematique -- Synth\`ese des triquaternions de Combebiac et du syst\`eme g\'eom\'etrique de Grassmann -- Calcul des quadriquaternions}, Ann. Fac. Sci. Toulouse {\bf 28} (1936) 103--148.  %%QUATERNION, %%GRASSMANN, %%ALGEBRA, %$D06072008.

\bibitem{MARKI1937-} M. Markic, \emph{Transformantes nouveau v\'ehicule math\'ematique -- Synth\`ese des triquaternions de Combebiac et du syst\`eme g\'eom\'etrique de Grassmann -- Calcul des quadriquaternions. (suite)}, Ann. Fac. Sci. Toulouse {\bf 1} (1937) 201--248.  %%QUATERNION, %%GRASSMANN, %%ALGEBRA, %$D06072008.

\bibitem{FEARN1979-} D. Fearnley-Sander, \emph{Hermann Grassmann and the creation of linear algebra}, Am. Math. Monthly {\bf 86} (1979) 809--817. %%QUATERNION, %%GRASSMANN, %$D17032002.

\bibitem{FEARN1982-} D. Fearnley-Sander, \emph{Hermann Grassmann and the prehistory of universal algebra}, Am. Math. Monthly {\bf 89} (1982) 161--166. %%QUATERNION, %%GRASSMANN, %$D17032002, %$C-Page 163: "Universal algebra = commutativity and associativity of addition, and distributivity (left and right) and associativity of multiplication over addition".

\bibitem{BARNA1985-} M. Barnabei, A. Brini, and G.-C. Rota, \emph{On the exterior calculus of invariant theory}, J. of Algebra {\bf 96} (1985) 120--160. %%QUATERNION, %%GRASSMANN, %%PEANO, %$D17032002.

\bibitem{STEWA1986-} I. Stewart, \emph{Hermann Grassmann was right}, Nature {\bf 321} (1986) 17. %%QUATERNION, %%GRASSMANN, %$D17032002.

\bibitem{LASEN1993B} A. Lasenby, C. Doran, and S. Gull, \emph{Grassmann calculus, pseudoclassical mechanics, and geometric algebra}, J. Math. Phys. {\bf 34}  (1993) 3683--37127.   %%QUATERNION, %%GRASSMANN, %%MECHANICS, %$D17032002.

\bibitem{HESTE1996B} D. Hestenes, \emph{Grassmann's vision}, in: G. Schubring, ed., Hermann Gunther Grassmann (1809-1877): Visionary Mathematician, Scientist and Neohumanist Scholar (Kluwer Academic Publishers, Dordrecht, 1996). %%QUATERNION, %%GRASSMANN, %%HESTENES, %$D10022002.


\end{enumerate}


\subsection{CLIFFORD (``Clifford numbers'')}
%-------------------------------------------
\label{CLIFFORD}

William K. Clifford introduced in 1873 two different kind of 8-dimensional algebras that he called ``biquaternions,'' and then in 1878 a general theory of linear associative algebras of dimension $2^n$, denoted $\cl_{p,q}$ where $n=p+q$.

\begin{enumerate}

\bibitem{CLIFF1873-} W.K. Clifford, \emph{Preliminary sketch of biquaternions}, Proc. London Math. Soc. {\bf 4} (1873) 381--395. Reprinted in: R. Tucker, ed., Mathematical Papers by William Kingdon Clifford (MacMillan, London, 1882) 181--200. %%QUATERNION, %%CLIFFORD, %%MECHANICS, %$D24062002. %$C Clifford defines his own variety of ``biquaternions,'' which are neither Hamilton's complex quaternions, nor Cayley's octonions.

\bibitem{CLIFF1876-} W.K. Clifford, \emph{Further note on biquaternions}, in: R. Tucker, ed., Mathematical Papers by William Kingdon Clifford (MacMillan, London, 1882) 385--394. %%QUATERNION, %%CLIFFORD, %$D06072008. 
           
\bibitem{CLIFF1878-} W.K. Clifford, \emph{Applications of Grassmann's extensive algebras}, Am. J. Math. {\bf 1} (1878) 350--358. Reprinted in: R. Tucker, ed., Mathematical Papers by William Kingdon Clifford (MacMillan, London, 1882) 266--276. %%QUATERNION, %%CLIFFORD, %%APPRECIATION, %$D04032006. 

\bibitem{TUCKE1882-} R. Tucker, ed., Mathematical Papers by William Kingdon Clifford (Macmillan, London, 1882) 658~pp. %%QUATERNION, %%CLIFFORD, %$D04032006. 

\bibitem{STUDY1891-} E. Study, \emph{Von der bewegungen und Umlegungen}, Math. Annalen {\bf 39} (1891) 441--556. %%QUATERNION, %%CLIFFORD, %%MECHANICS, %$D30112005. %$C Clifford ``biquaternions'' are used to obtain a parametric representation of the Euclidian group of rigid motions, see \cite{VANDE1985-}.

\bibitem{JOLY-1897-} C.J. Joly, \emph{The associative algebra applicable to hyperspace}, Proc. Roy. Irish Acad. {\bf 5} (1897) 73--123. %%QUATERNION, %%CLIFFORD, %$D14032002.

\bibitem{COQUE1982-} R. Coquereaux, \emph{Modulo 8 periodicity of real Clifford algebras and particle physics}, Phys. Lett. {\bf 115B} (1982) 389--395.  %%QUATERNION, %%CLIFFORD, %%PARTICLE-PHYSICS, %$D14032002. %$C Mostly mathematical: a classification of Clifford algebras.

\bibitem{BLAU-1986-} M. Blau, \emph{Clifford algebras and K\"ahler-Dirac spinors}, Ph.D. dissertation, Report UWTHPh-1986-16 (Universitat Wien, 1986) 200~pp.   %%QUATERNION, %%CLIFFORD, %%KAHLER, %%DIRAC, %%SPINORS, %$D14032002. 

\bibitem{STEIN1987-} P.A.J. Steiner, \emph{Real Clifford algebras and their representations over the reals}, J. Phys. {\bf A 20} (1987) 3095--3098. %%QUATERNION, %%CLIFFORD, %$D14032002.

\bibitem{OKUBO1991A} S. Okubo, \emph{Real representations of finite Clifford algebras. I. Classification}, J. Math. Phys. {\bf 32} (1991) 1657--1668.  %%QUATERNION, %%CLIFFORD, %$D30112005.

\bibitem{OKUBO1991B} S. Okubo, \emph{Real representations of finite Clifford algebras. II. Explicit construction and pseudo-octonion}, J. Math. Phys. {\bf 32} (1991) 1669--1673. %%QUATERNION, %%CLIFFORD, %%OCTONION, %$D30112005.

\bibitem{DIEK-1995-} A. Diek and R. Kantowski, \emph{Some Clifford algebra history}, in: Clifford Algebras and Spinor Structures (Kluwer, Dordrecht, 1995) 3--12.  %%QUATERNION, %%CLIFFORD, %%HISTORY, %%BIBLIOGRAPHY, %$D30112005. %$C The representation of complex numbers as ordered pairs of real numbers due to Hamilton is wrongly attributed to Gauss: See \cite[p.~178]{VANDE1985-}, \cite{MACDU1944-}, or \cite{DICKS1918-}.

\bibitem{SOMME1997A} F. Sommen, \emph{The problem of defining abstract bivectors}, Result. Math. {\bf 31} (1997) 148--160. %%QUATERNION, %%CLIFFORD, %%HESTENES, %$D01062001. %$Critique of Hestenes. 

\bibitem{ABLAM2000B-} R. Ablamowicz and B. Fauser, \emph{Heck algebra representations in ideals generated by q-Young Clifford idempotents}, in: R. Ablamowicz and B. Fauser, eds., Clifford Algebra and their Applications in Mathematical Physics, Vol.~1: \emph{Algebra and Physics} (Birkh\"auser, Boston, 2000)  245--268. %%QUATERNION, %%CLIFFORD, %%MULTI-PARTICLE-STATES, %%REVERSION, %$D12052001.

\bibitem{VARLA2001-} V.V. Varlamov, \emph{Discrete symmetries and Clifford algebras}, Int. J. Theor. Phys. {\bf 40} (2001) 769--806.  %%QUATERNION, %%CLIFFORD, %%DISCRETE-SYMMETRY, %$D14032002. 

\bibitem{ULRYC2008-} S. Ulrych, \emph{Representations of Clifford algebras with hyperbolic numbers}, Adv. Appl. Cliff. Alg. {\bf 18} (2008) 93--114.   %%QUATERNION, %%CLIFFORD, %$D06072008.

\end{enumerate}



\subsection{EDDINGTON (``Eddington numbers'')}
%----------------------
\label{EDDINGTON}

 The ``E-numbers'' introduced by Arthur S. Eddington in 1928 are equivalent to elements of the algebra of the Dirac matrices, i.e., to elements of the Clifford algebra $\cl_{4,1}$.  Eddington's formalism corresponds to the first use of Clifford numbers in relation to Dirac's theory.


\begin{enumerate}


\bibitem{EDDIN1928-} A.S. Eddington, \emph{A symmetrical treatment of the wave equation}, Proc. Roy. Soc. A {\bf 121} (1928) 524--542. %%QUATERNION, %%EDDINGTON, %%DIRAC, %%CLIFFORD, %$D24062002.

\bibitem{EDDIN1929-} A.S. Eddington, \emph{The charge of an electron}, Proc. Roy. Soc. A {\bf 122} (1929) 359--369. %%QUATERNION, %%EDDINGTON, %%DIRAC, %%CLIFFORD, %$D24062002.

\bibitem{TEMPL1930-} G. Temple, \emph{The group properties of Dirac's operators}, Proc. Roy. Soc. {\bf A127} (1930) 339--348. %%QUATERNION, %%EDDINGTON, %%DIRAC, %$D17032002.

\bibitem{MCCRE1939-} W.H. McCrea, \emph{On matrices of quaternions and the representation of Eddington's E-numbers}, Proc. Roy. Irish Acad. {\bf A 45} (1939) 65--67. %%QUATERNION, %%EDDINGTON, %$17012002.

\bibitem{MCCRE1940-} W.H. McCrea, \emph{Quaternion analogy of wave-tensor calculus}, Phil. Mag. {\bf 30} (1940) 261--281. %%QUATERNION, %%EDDINGTON, %$17012002.

\bibitem{KILMI1949B} C.W. Kilmister, \emph{The use of quaternions in wave-tensor calculus}, Proc. Roy. Soc. {\bf A 199} (1949) 517--532. %%QUATERNION, %%EDDINGTON, %$17012002.

\bibitem{KILMI1951-} C.W. Kilmister, \emph{Tensor identities in wave-tensor calculus}, Proc. Roy. Soc. {\bf A 207} (1951) 402--415. %%QUATERNION, %%EDDINGTON, %$17012002.

\bibitem{KILMI1953A} C.W. Kilmister, \emph{A new quaternion approach to meson theory}, Proc. Roy. Irish Acad. {\bf A 55} (1953) 73--99.  %%QUATERNION, %%EDDINGTON, %$17012002.

\bibitem{KILMI1953B} C.W. Kilmister, \emph{A note on Milner's E-numbers}, Proc. Roy. Soc. {\bf A 218} (1953) 144--148. %%QUATERNION, %%EDDINGTON, %$17012002.

\end{enumerate}



\subsection{SEMIVECTOR (``Einstein-Mayer's spinors'')}
%----------------------------------------------------
\label{SEMIVECTOR}

Semivectors were introduced in 1932 by Einstein and his assistant Walter Mayer, possibly in reaction to the work on Dirac's equation by his previous assistant, Cornelius Lanczos.  Einstein-Mayer's formulation of Dirac's equation is based on $4 \times 4$ matrices which in flat spacetime reduces to Lanczos's biquaternionic formulation.

\begin{enumerate}


\bibitem{EINST1932-} A. Einstein and W. Mayer, \emph{Semi-Vektoren und Spinoren}, Sitzber. Preuss. Akad. Wiss. Physik.-Math. Kl. (1932) 522--550. %%QUATERNION, %%SEMIVECTOR, %$D05022002.

\bibitem{EINST1933A} A. Einstein and W. Mayer, \emph{Die Diracgleichungen f\"ur Semivektoren}, Proc. Roy. Acad. Amsterdam {\bf 36} (1933) 497--516. %%QUATERNION, %%SEMIVECTOR, %$D05022002.

\bibitem{EINST1933B} A. Einstein and W. Mayer, \emph{Spaltung der nat\"urlichsten Feldgleichungen f\"ur Semi-Vektoren in Spinor-Gleichungen vom Dirac'schen Typus}, Proc. Roy. Acad. Amsterdam {\bf 36} (1933) 615--619.  %%QUATERNION, %%SEMIVECTOR, %$D05022002.

\bibitem{GUTH-1933A} E. Guth, \emph{Semivektoren, Spinoren und Quaternionen}, Anz. Akad. Wiss. Wien {\bf 70} (1933) 200--207. %%QUATERNION, %%SEMIVECTOR, %%SPINOR, %$D30112005, %$C Guth makes the Pierce decomposition of a quaternion into two spinors that he calls ``semiquaternion'' or ``spinquaternion.'' 

\bibitem{SCHOU1933-} J.A. Schouten, \emph{Zur generellen Feldtheorie. Semivektoren und Spinraum}, Zeits. f\"ur Phys. {\bf 84} (1933) 92--111. %%QUATERNION, %%SEMIVECTOR, %$D05022002.

\bibitem{BARGM1934-} V. Bargmann, \emph{Uber den Zusammenhang zwischen Semivektoren und Spinoren und die Reduktion der Diracgleichungen f\"ur Semivektoren}, Helv. Phys. Acta {\bf 7} (1934) 57--82. %%QUATERNION, %%SEMIVECTOR, %$D05022002.

\bibitem{EINST1934-} A. Einstein and W. Mayer, \emph{Darstellung der Semi-Vektoren als gew\"ohnliche Vektoren von besonderem differentitions Charakter}, Ann. of Math. {\bf 35} (1934) 104--110. %%QUATERNION, %%SEMIVECTOR, %$D05022002.

\bibitem{ULLMO1934-} J. Ullmo, \emph{Quelques propri\'et\'es du groupe de Lorentz, semi-vecteurs et spineurs}, J. de Phys. {\bf 5} (1934) 230--240. %%QUATERNION, %%SEMIVECTOR, %$D05022002.

\bibitem{BLATO1935-} J. Blaton, \emph{Quaternionen, Semivektoren und Spinoren}, Zeitschr. f\"ur Phys. {\bf 95} (1935) 337--354. %%QUATERNION, %%SEMIVECTOR, %%SPINOR, %$D05022002.

\bibitem{SCHER1935-} W. Scherrer, \emph{Quaternionen und Semivektoren}, Comm. Math. Helv. {\bf 7} (1935) 141--149. %%QUATERNION, %%SEMIVECTOR, %$D05022002.

\bibitem{RAO--1936-} B.S.M. Rao, \emph{Semivectors in Born's field theory}, Proc. Indian Acad. Sci.  {\bf 4} (1936) 436--451. %%QUATERNION, %%SEMIVECTOR, %%BORN-ELECTRODYNAMICS, %$D05022002.

\bibitem{VANDO2004-} J. vanDongen, \emph{Einstein's methodology, semivectors, and the unification of electrons and protons}, Arch. Hist. Exact. Sci. {\bf 58} (2004) 219--254.  %%QUATERNION, %%SEMIVECTOR, %%HISTORY, %$D03082005, %$C Only proton-neutron, no isospin doublets in general.

\bibitem{GOENN2004-} H.F.M. Goenner, \emph{Einstein, spinors, and semivectors}, Sec.~7.3 of ``On the history of unified field theories,'' Living Reviews of Relativity {\bf 7} 2 (2004) 152~pp.  %%QUATERNION, %%SEMIVECTOR, %%HISTORY, %%LANCZOS, %$D03082005, %$C Four pages on semivectors.


\end{enumerate}

\subsection{HESTENES (``Hestenes's space-time algebra'')}
%--------------------
\label{HESTENES}

Hestenes's 16-dimensional ``space-time algebra'' formalism (based on the Clifford algebra $\cl_{1,3}$) was designed in 1966 to provide a substitute for the standard 32-dimensional Dirac formalism (equivalent to the Clifford algebra $\cl_{4,1}$) such that all references to complex numbers are avoided.


\begin{enumerate}


\bibitem{HESTE1966-} D. Hestenes, {Space-Time Algebra} (Gordon and Breach, New York, 1966, 1987, 1992) 93~pp. %%BOOK, %%QUATERNION, %%HESTENES, %$D10022002.

\bibitem{HESTE1984-} D. Hestenes and G. Sobczyk, Clifford Algebra to Geometric Calculus: A Unified Language for Mathematics and Physics (Reidel, Dordrecht, 1984) 314~pp. %%BOOK, %%QUATERNION, %%HESTENES, %$D10022002.

\bibitem{HESTE1986A} D. Hestenes, \emph{A unified language for mathematics and physics}, in: J.S.R. Chisholm and A.K. Common, eds., Clifford Algebras and Their Applications in Mathematical Physics (Reidel, Dordrecht, 1986) 1--23. %%QUATERNION, %%HESTENES, %$D10022002.

\bibitem{BOUDE1988-}  R. Boudet, \emph{La g\'eom\'etrie des particules du groupe $SU(2)$}, Ann. Fond. Louis de Broglie {\bf 13} (1988) 105--137. %%QUATERNION, %%HESTENES, %%PARTICLE-PHYSICS, %%ISOSPIN, %%LED, %%HAD, %%GURSEY, %$D16032002. %$C Plea for Hestenes's formalism.

\bibitem{HESTE1991B} D. Hestenes and R. Ziegler, \emph{Projective geometry with Clifford algebra}, Acta Applicandae Mathematicae {\bf 23} (1991) 25--63.  %%QUATERNION, %%HESTENES, %$D10022002.

\bibitem{HESTE1991-} D. Hestenes, \emph{The design of linear algebra and geometry}, Acta Applicandae Mathematicae {\bf 23} (1991) 65--93.  %%QUATERNION, %%HESTENES, %$D06072008.

\bibitem{BAYLI1992B}  W.E. Baylis, \emph{Why i?}, Am. J. Phys. {\bf 60} (1992) 788-797. %%QUATERNION, %%HESTENES, %%SPECIAL-RELATIVITY, %%MAXWELL, %%QPH, %$D16032002.

\bibitem{HESTE1992-} D. Hestenes, \emph{Mathematical viruses}, in: A. Micali et al., eds., Clifford Algebras and their Applications in Mathematical Physics (Kluwer Academic Publishers, Dordrecht, 1992) 3--16. %%QUATERNION, %%HESTENES, %$D10022002.

\bibitem{GULL-1993B} S. Gull, A. Lasenby, and C. Doran, \emph{Imaginary numbers are not real --- The geometric algebra of spacetime}, Found. Phys. {\bf 23}  (1993) 1175--1201.  %%QUATERNION, %%HESTENES, %$D17032002.

\bibitem{SOBCZ1993-} G. Sobczyk, \emph{David Hestenes: the early years}, Found. of Phys. {\bf 23} (1993) 1291--1293.  %%QUATERNION, %%HESTENES, %$D24062002.

\bibitem{HESTE1996} D. Hestenes, Spacetime Calculus for Gravitation Theory (Monograph, 1996) 74~pp. %%BOOK, %%QUATERNION, %%HESTENES, %$D10022002.

\bibitem{KELLE1997A}  J. Keller, \emph{On the electron theory}, Adv. Appl. Clifford Alg. {\bf 7 (S)} (1997) 3--26. %%QUATERNION, %%HESTENES, %$D17032002.

\bibitem{HESTE1998B}  D. Hestenes, Space Time Calculus, Draft of an overview of ``space time algebra,'' (1998) 73~pp.  %%BOOK, %%QUATERNION, %%HESTENES, %$D20032002.

\bibitem{HESTE2001C} D. Hestenes, \emph{Old wine in new bottles: A new algebraic framework for computational geometry}, in: E.B. Corrochano and G. Sobczyk, eds., Geometric Algebra with Applications in Science and Engineering (Birkhauser, Boston, 2001) 3--17.  %%QUATERNION, %%HESTENES, %$D17032002.

\bibitem{LASEN2001-} A. Lasenby and J. Lasenby, \emph{Applications of geometric algebra in physics and links with engineering}, in: E.B. Corrochano and G. Sobczyk, eds., Geometric Algebra with Applications in Science and Engineering (Birkhauser, Boston, 2001) 430--457.  %%QUATERNION, %%HESTENES, %$D17032002.

\bibitem{LOUNE2001-} P. Lounesto, \emph{Counterexamples for validation and discovering new theorems}, in: E.B. Corrochano and G. Sobczyk, eds., Geometric Algebra with Applications in Science and Engineering (Birkhauser, Boston, 2001) 477--490.  %%QUATERNION, %%HESTENES, %%PATHOLOGICAL-SCIENCE, %$D17032002.

\bibitem{SOBCZ2001-} G. Sobczyk, \emph{Universal geometric algebra}, in: E.B. Corrochano and G. Sobczyk, eds., Geometric Algebra with Applications in Science and Engineering (Birkhauser, Boston, 2001) 3--17.  %%QUATERNION, %%HESTENES, %$D17032002.

\bibitem{HORN-2007-} M.E. Horn, \emph{Quaternions and geometric algebra --- Quaternionen und Geometrische Algebra}, in: Volkhard Nordmeier, Arne Oberlaender (Eds.): Tagungs-CD des Fachverbandes Didaktik der Physik der DPG in Kassel, Beitrag 28.2, ISBN 978-3-86541-190-7, LOB - Lehmanns Media, Berlin 2006 (in German), 22~pp.; e-print \underline{ arXiv:0709.2238 }.   %%QUATERNION, %%HESTENES, %$D06072008.


\end{enumerate}




\section{\Huge MISCELLANEA}
%==========================
\label{MISCELLANEA}


\subsection{HISTORY and APPRECIATION}
%------------------------------------
\label{HISTORY}

Historical papers on quaternions and on Hamilton (1805-1865), as well as some appreciations of quaternions and Hamilton's devotion to them.

Historical papers on algebras in general. 

The historical papers on specific formalisms allied to quaternions are in the respective sections of Chap.~\ref{ALLIED FORMALISMS}.

\begin{enumerate}

\bibitem{TAIT-1866-} P.G. Tait, \emph{Sir William Rowan Hamilton}, North British Review {\bf 45}  (1866) 37--74. %%QUATERNION, %%HISTORY, %%APPRECIATION, %$D04032006, %$C Written the year following Hamilton(1805-1865)'s death.

\bibitem{MAXWE1870-} J.C. Maxwell, \emph{Address to the mathematical and physical sections of the British Association}, Brit. Ass. Rep. {\bf XL} (1870) 215--229. %%QUATERNION, %%HISTORY, %$D07112001.

\bibitem{MAXWE1872-} J.C. Maxwell, \emph{Letter to Lewis Campbell} (18 October, 1872). %%QUATERNION, %%HISTORY, %$D07112001.

\bibitem{TAIT-1880-} P.G. Tait, \emph{Hamilton}, Encyclopedia Britannica {\bf } (1880) SP-2:440-444.  %%QUATERNION, %%HISTORY, %%APPRECIATION, %$D24062002.

\bibitem{TAIT-1894-} P.G. Tait, \emph{On the intrinsic nature of the quaternion method}, Proc. Roy. Soc. Edinburgh {\bf 20} (1894) 276--284. %%QUATERNION, %%HISTORY, %$D19032002.

\bibitem{GIBBS1891A} J.W. Gibbs, \emph{On the role of quaternions in the algebra of vectors}, Nature {\bf 43} (2 April 1891) 511--513. %%QUATERNION, %%HISTORY, %%APPRECIATION, %$D24062002, %$C Reproduced in J.W. Gibbs, \emph{The scientific papers} (Dover, New York, 1961) 155--160.

\bibitem{GIBBS1891B} J.W. Gibbs, \emph{Quaternions and the ``Ausdehnungslehre,''} Nature {\bf 44} (28 May 1891) 79--82. %%QUATERNION, %%HISTORY, %%APPRECIATION, %$D24062002, %$C Reproduced in J.W. Gibbs, \emph{The scientific papers} (Dover, New York, 1961)  161--168.


\bibitem{GIBBS1893A} J.W. Gibbs, \emph{Quaternions and the algebra of vectors}, Nature {\bf 47} (16 March 1893) 463--464. %%QUATERNION, %%HISTORY, %%APPRECIATION, %$D24062002, %$C Reproduced in J.W. Gibbs, \emph{The scientific papers} (Dover, New York, 1961) 169--172.

\bibitem{MACFA1893A} A. MacFarlane, \emph{Vectors versus quaternions}, Nature {\bf 48} (25 May 1893) 75--76.  %%QUATERNION, %%APPRECIATION, %$D04032006. 

\bibitem{KNOTT1893A} C.G. Knott, \emph{Vectors and quaternions}, Nature {\bf 48} (15 June 1893) 148--149.  %%QUATERNION, %%APPRECIATION, %$D04032006. 

\bibitem{LODGE1893-} A. Lodge, \emph{Vectors and quaternions}, Nature {\bf 48} (29 June 1893) 198--199. %%QUATERNION, %%APPRECIATION, %$D04032006. 

\bibitem{GIBBS1893B} J.W. Gibbs, \emph{Quaternions and vector analysis}, Nature {\bf 48} (17 August 1893) 364--367.  %%QUATERNION, %%HISTORY, %%APPRECIATION, %$D24062002, %$C Reproduced in J.W. Gibbs \emph{The scientific papers} (Dover, New York, 1961) 173--181.

\bibitem{BALL-1893-} R.S. Ball, \emph{The discussion on quaternions}, Nature {\bf 48} (24 August 1893) 391.  %%QUATERNION, %%APPRECIATION, %$D04032006. 

\bibitem{KNOTT1893B} C.G. Knott, \emph{Quaternions and vectors}, Nature {\bf 48} (28 September 1893) 516--517. %%QUATERNION, %%APPRECIATION, %$D04032006. 

\bibitem{MACFA1893B} A. MacFarlane, \emph{Vectors and quaternions}, Nature {\bf 48} (5 October 1893) 540--541.  %%QUATERNION, %%APPRECIATION, %$D04032006. 

\bibitem{TAIT-1898-} P.G. Tait, {Scientific Papers}, 2 volumes (Cambridge University Press, 1898, 1900) 943~pp. The numbers ``SP-1:'' and ``SP-2:'' in the above papers correspond to pages in volumes 1 and 2. %%QUATERNION, %%HISTORY, %%VARIA, %$D19032002.

\bibitem{TAIT-1899B} P.G. Tait, \emph{On the claim recently made for Gauss to the invention (not the discovery) of quaternions}, Proc. Roy. Soc. Edinburgh {\bf 23} (1899/1900) 17--23. %%QUATERNION, %%HISTORY, %$D19032002.

\bibitem{BELL-1939-} E.T. Bell, \emph{Hamilton --- Une trag\'edie irlandaise}, Chap.~19 in:  Les grands math\'ematiciens (Payot, Paris, 1939) 368--390. Translation of  E.T. Bell, \emph{Men of Mathematics} (Simon \& Schuster, New York, 1937).  %%QUATERNION, %%HISTORY, %$D19032002.

\bibitem{WHITT1940-} E.T. Whittaker, \emph{The Hamiltonian revival}, The Math. Gazette {\bf 24} (1940) 153--158.  %%QUATERNION, %%HISTORY, %$D19032002.

\bibitem{PIAGG1943-} H.T.H. Piaggio, \emph{The significance and development of Hamilton's quaternions}, Nature {\bf 152} (1943) 553--555.  %%QUATERNION, %%APPRECIATION, %$D06072008. %$C Piaggio fut l'un des auteurs cit�s par Lanczos dans les Physikalische Berichte.

\bibitem{BATEM1944-} H. Bateman, \emph{Hamilton's work in dynamics and its influence on modern thought}, Scripta Math. {\bf 10} (1944) 51--63. %%QUATERNION, %%HISTORY, %%MECHANICS, %%LANCZOS, %$D19032002.

\bibitem{GINSB1944-} J. Ginsburg, \emph{Editorial to special issue dedicated to the memory of Hamilton's discovery of quaternions}, Scripta Math. {\bf 10} (1944) 7. %%QUATERNION, %%HISTORY, %$D19032002.

\bibitem{MACDU1944-} C.C. MacDuffee, \emph{Algebra's debt to Hamilton}, Scripta Math. {\bf 10} (1944) 25--36. %%QUATERNION, %%HISTORY, %$D19032002.

\bibitem{MURNA1944-} F.D. Murnaghan, \emph{An elementary presentation of quaternions}, Scripta Math. {\bf 10} (1944) 37--49.  %%QUATERNION, %%HISTORY, %$D19032002.

\bibitem{SMITH1944-} D.E. Smith, \emph{Sir William Rowan Hamilton}, Scripta Math. {\bf 10} (1944) 9--11. %%QUATERNION, %%HISTORY, %$D19032002.

\bibitem{SYNGE1944-} J.L. Synge, \emph{The life and early work of sir William Rowan Hamilton}, Scripta Math. {\bf 10} (1944) 13--25. %%QUATERNION, %%HISTORY, %$D19032002.

\bibitem{BEST-1945-} R.I. Best, J.L. Synge, D. Birkoff, A.J. McConnell, E.T. Whittaker, A.W. Conway, F.D. Murnaghan, and J. Riverdale Colthurst, \emph{Quaternion centenary celebration}, Proc. Roy. Irish Acad. {\bf A 50} (1945) 69--121. %%QUATERNION, %%HISTORY, %$D19032002.

\bibitem{BIRKH1945-} G. Birkhoff, \emph{Letter from George D. Birkoff}, Proc. Roy. Irish Acad. {\bf A 50} (1945) 72--75. %%QUATERNION, %%HISTORY, %$D19032002.

\bibitem{MURNA1945-} F.D. Murnaghan, \emph{A modern presentation of quaternions}, Proc. Roy. Irish Acad. {\bf A 50} (1945) 104--112. %%QUATERNION, %%HISTORY, %%SEMIVECTOR, %$D19032002.

\bibitem{SYNGE1945-} J.L. Synge, \emph{Message from J.L. Synge}, Proc. Roy. Irish Acad. {\bf A 50} (1945) 71--72. %%QUATERNION, %%HISTORY, %$D19032002.

\bibitem{WHITT1945-} E.T. Whittaker, \emph{The sequence of ideas in the discovery of quaternions}, Proc. Roy. Irish Acad. {\bf A 50} (1945) 93--98. %%QUATERNION, %%HISTORY, %$D19032002.

\bibitem{CONWA1951-} A.W. Conway, \emph{Hamilton, his life work and influence}, in: Proc. Second Canadian Math. Congress, Toronto (Univ. of Toronto Press, Toronto, 1951) 32--41.  %%QUATERNION, %%HISTORY, %%HAMILTON, %$D24062002.

\bibitem{BORK-1966-} A.M. Bork, \emph{``Vectors versus quaternions''---The letters in Nature}, Am. J. Phys. {\bf 34} (1966) 202--211. %%QUATERNION, %%HISTORY, %$D19032002.

\bibitem{STEPH1966-} R.J. Stephenson, \emph{Development of vector analysis from quaternions}, Am. J. Phys. {\bf 34} (1966) 194--201. %%QUATERNION, %%HISTORY, %$D19032002.

\bibitem{LANCZ1967B} C. Lanczos, \emph{William Rowan Hamilton---an appreciation}, Amer. Sci. {\bf 2} (1967) 129--143. Reprinted in: W.R. Davis et al., eds., Cornelius Lanczos collected published papers with commentaries {\bf IV} (North Ca\-rolina State University, Raleigh NC, 1998) 2-1859 to 2-1873.  %%QUATERNION, %%HISTORY, %$D19032002.

\bibitem{DYSON1972B} F. Dyson, \emph{Missed opportunities}, Bull. Am. Math. Soc. {\bf 78} (1972) 635--652. %%QUATERNION, %%HISTORY, %$D19032002. %$C See pages 643--644 dedicated to the difference between vecteurs et quaternions.

\bibitem{VANDE1976-} B.L. vanderWaerden, \emph{Hamilton's discovery of quaternions}, Mathematics Magazine {\bf 49} (1976) 227--234. %%QUATERNION, %%HISTORY, %$D04032006.

\bibitem{ODONN1983-} S. O'Donnel, {William Rowan Hamilton, Portrait of a Prodigy} (Boole Press Dublin, Dublin, 1983) 224~pp. %%BOOK, %%QUATERNION, %%HISTORY, %$D20032002.

\bibitem{GIRAR1984-} P.R. Girard, \emph{The quaternion group and modern physics}, Eur. J. Phys. {\bf 5} (1984) 25--32. %%QUATERNION, %%HISTORY, %%SPECIAL-RELATIVITY, %%GENERAL-RELATIVITY, %$D19032002.

\bibitem{VANDE1985-} B.L. vanderWaerden, \emph{The discovery of algebras}, Chap. 10 of A History of Algebra (Springer, Berlin, 1985) 177--201. %%QUATERNION, %%HISTORY, %%MATH-VARIA, %$D14032002.

\bibitem{SCHWE1986-} S.S. Schweber, \emph{Feynman and the visualization of space-time processes}, Rev. Mod. Phys. {\bf 58} (1986) 449--508. %%QUATERNION, %%HISTORY, %%$D30112005, %$C Page 472: ``Feynman tried to use quaternions and (biquaternions) to represent wave functions, but was not able to obtain a "natural" representation of the Dirac equation as an integral over paths.''

\bibitem{ALTMA1989-} S.L. Altmann, \emph{Hamilton, Rodrigues, and the quaternion scandal}, Math. Mag. {\bf 62} (1989) 291--308. %%QUATERNION, %%HISTORY, %$D19032002.

\bibitem{ANDER1992-} R. Anderson and G.C. Joshi, \emph{Quaternions and the heuristic role of mathematical structures in physics}, Physics Essays {\bf 6} (1993) 308--319. %%QUATERNION, %%HISTORY, %$D19032002.

\bibitem{SPEAR1993-} T.D. Spearman, \emph{William Rowan Hamilton 1805--1865}, Proc. Roy. Irish Acad. {\bf 95A} Suppl. (1993) 1--12. %%QUATERNION, %%HISTORY, %$D30112005.

\bibitem{MEHRA1994-} J. Mehra, The Beat of a Different Drum --- The Life and Science of Richard Feynman (Clarendon Press, Oxford, 1994) 630~pp. %%BOOK, %%QUATERNION, %%HISTORY, %$D30112005, %$C Page 174: ``With the help of quaternions and path integrals, Feynman tried to formulate the relativistic theory of an electron in one dimension.''

\bibitem{KOETS1995-} T. Koetsier, \emph{Explanation in the historiography of mathematics: The case of Hamilton's quaternions}, Studies in History and Philosophy of Science {\bf 26} (1995) 593--616.  %%QUATERNION, %%HISTORY, %$D04032006. 

\bibitem{LAMBE1995-} J. Lambek, \emph{If Hamilton had prevailed: quaternions in physics}, The Mathematical Intellig. {\bf 17} (1995) 7--15. Errata, ibid {\bf 18} (1996) 3. %%QUATERNION, %%HISTORY, %$D19032002.

\bibitem{PICKE1995-} A. Pickering, \emph{Concepts --- constructing quaternions}, Chap.~4 of The Mangle of Practice (Univ. Chicago Press, Chicago, 1995) 113--156. %%BOOK, %%QUATERNION, %%HISTORY, %$D19032002.

\bibitem{SINEG1995-} L. Sin\`egre, \emph{Les quaternions et le mouvement du solide autour d'un point fixe chez Hamilton}, Revue d'Histoire des Math\'ematiques {\bf 1} (1995) 83--109.  %%QUATERNION, %%HISTORY, %$D04032006. 

\bibitem{WILKI2005-} D.R. Wilkins, \emph{William Rowan Hamilton: Mathematical genius}, Physics World (August 2005) 33--36. %%QUATERNION, %%HISTORY, %$D30112005.

\end{enumerate}


\subsection{BIBLIOGRAPHY}
%------------------------
\label{BIBLIOGRAPHY}

\begin{enumerate}


\bibitem{MACFA1904-} A. Macfarlane, \emph{Bibliography of quaternions and allied systems of mathematics}, in: A. Macfarlane. ed., Bull. of the Intern. Assoc. for promoting the study of quaternions and allied systems of mathematics (Dublin University Press, Dublin, 1904) 3--86. %%BOOK, %%QUATERNION, %%BIBLIOGRAPHY, %$D20032002.

\bibitem{MACFA1910A} A. Macfarlane, \emph{Supplementary bibliography}, in: A. Macfarlane, ed., Bull. of the Inter. Assoc. for promoting the study of quaternions and allied systems of mathematics (New Era Printing Company, Lancaster PA, 1910) 10--36.  %%BOOK, %%QUATERNION, %%BIBLIOGRAPHY, %$D20032002.

\bibitem{MCCON1953-} J.R. McConnell, ed., Selected Papers of Arthur William Conway  (Dublin Institute for Advanced Studies, 1953) 222~pp.  %%BOOK, %%QUATERNION, %%BIBLIOGRAPHY, %%CONWAY, %$D20032002, %$C Obituary and reprint of 21 papers.

\bibitem{GSPON1993A} A. Gsponer and J.-P. Hurni,  \emph{Quaternion bibliography: 1843--1993}. Report ISRI-93-13 (17 June 1993) 35~pp. This is the first version of the present bibliography, with 228 entries. %%QUATERNION, %%BIBLIOGRAPHY, %$D20032002.

\bibitem{ELL---2005-} T.A. Ell, \emph{Bibliography} (Last changed 1/11/05) about 40 pp.; available in HTML at \underline{ http://home.att.net/~t.a.ell/QuatRef.htm }. %%QUATERNION, %%BIBLIOGRAPHY, %C Mixed bibliography: many references concern applications in computation, avionics, modeling, etc. %$D20032002.

\end{enumerate}

\newpage

\section{\Huge Conventions used in the bibliography}
%===================================================
\label{Conventions}

\subsection{Conventions for the \TeX ~citation-label}
%------------------------------------------------
\label{Label}


The \TeX ~{\bf citation-label} has always 10 characters, i.e.:

~~~ ~~~ {\bf LASTN1234-} ~~~ or ~~~ {\bf LASTN1234X} ~,

\noindent were LASTN are the 5 first characters of the last name, ``1234'' the four-digit year, and ``X'' is A, B, C, ..., Z if there are more than one reference for the author in that year. 

If the author last name has less that 5 characters, the missing characters are replaced by ``-''s.

{\bf N.B.:} In the case where different authors have the same last name the 10th character (i.e., ``-'') will be used to remove possible ambiguities that may arise in a given year.
 

\subsection{Conventions for the references's "style"}
%----------------------------------------------------
\label{Style}


~~~~~~Authors: Only first character of first-names. Ex: P.A.M.\ Dirac.

Composite last-names are concatenated: De Broglie  $\rightarrow$ DeBroglie.

Two authors: Use ``and'' between their names.

Three or more authors: Use ``, and'' before the last author in the list.

Book titles:    Capitalize First Letter of Each Word. Normal characters.

Article titles: Capitalize only first letter of first word. \emph{Italic} characters.

Editor(s): Use  ``ed.'' or  ``eds.''

Latin: Use only unambiguous combinations such as:  ``et al.'', ``ibid.'', `in:''...

Volume numbers: {\bf Bold} characters.

No other bold symbol should appear anywhere in the bibliography file.

Number of pages in a book:  ``123~pp.'' (Non-break space: ``$\sim$'' in~ \TeX.)

First and last pages of article: ``123--456" (Medium dash: ``-\,-'' in~ \TeX.)

Examples of typical references are given in Sec.~\ref{Format}.


\subsection{Internet URLs and {\tt arXiv.org} links}
%---------------------------------------------------
\label{internet}

In order that they are properly processed by the {\tt arXiv.org} \TeX-compiler, a white space ``\underline{ }'' is put in the front and in the back of all internet URLs and {\tt arXiv.org} links.  Moreover, they are \underline{ underlined }.  These conventions insure that the URLs and links are properly recognized, and that they are not hyphenated.  

For instance, clicking the {\tt arXiv.org} link \underline{ arXiv:name/number } automatically generates the internet URL \underline{ http://www.arXiv.org/abs/name/number }.

\subsection{Non-\TeX ~conventions: \emph{directives} and \emph{keywords}}
%------------------------------------------------------------------------
\label{Non-TeX}

In the file special \TeX ~comments are used to introduce the keywords, and to introduce directives to the bibliography update/search/display computer programs:

\begin{itemize}
\item Standard comments start with ``\%''
\item Bibliography KEYWORDS   start with ``\%\%''
\item Bibliography directives start with ``\%\$''
\end{itemize}


\subsection{Directives}
%----------------------
\label{Directives}

\begin{itemize}
\item \%\$D : date of entry in DDMMYYYY format
\item \%\$C : a commentary/appreciation of the document follows
\item \%\$L : this reference is on loan to the person whose name follows
\item \%\$M : this reference is missing or lost
\item \%\$N : this reference is not yet filed in its envelope/folder
\item \%\$O : this reference is on order from some library
\item etc. 
\end{itemize}

The date of entry is particularly important for making automatic updates.  It is normally the last item in a bibliography entry, unless there is a commentary which is appended at the very end of it.


\subsection{Format of typical references}
%----------------------------------------
\label{Format}


What follows is a list of typical references in which the label, keywords, and directives are made visible.

\begin{itemize}

\item {\bf Book:}

{\bf LANCZ1949-} C. Lanczos, The Variational Principles of Mechanics (Dover, New-York, 1949, 1986) 418~pp.  Quaternions pages 303--314. \%\%BOOK, \%\%QUATERNION, \%\%SPECIAL-RELATIVITY, \%\$D06022002.

\item {\bf Conference proceedings, festschrift or contributed volume:}

{\bf SPROS1996-} W. Spr\"ossig and K. G\"urlebeck, eds., Proc. of the Symp. ``Analytical and Numerical Methods in Quaternionic and Clifford Analysis,'' Seiffen, June 5--7, 1996 (TU Bergakademie Freiberg, 1996) 228~pp.  \%\%BOOK, \%\%QUATERNION, \%\%ANALYTICITY, \%\%CLIFFORD, \%\$D20032002.


\item {\bf Paper in a conference proceedings, festschrift, or contributed volume:}

{\bf PENRO1990-} R. Penrose, \emph{Twistors, particles, strings and links}, in: D.G. Quillen et al., eds., The Interface of Mathematics and Particle Physics (Clarendon Press, Oxford, 1990) 49--58. \%\%QUATERNION, \%\%TWISTOR, \%\$D05022002.

\item {\bf Paper in a journal:}

{\bf WEISS1941-} P. Weiss, \emph{On some applications of quaternions to restricted relativity and classical radiation theory}, Proc. Roy. Irish Acad. {\bf A 46} (1941) 129--168. \%\%QUATERNION, \%\%MAXWELL, \%\%SPINOR, \%\%LORENTZ-DIRAC, \%\$D09022002.

\item {\bf Two authors:}

{\bf EINST1932-} A. Einstein and W. Mayer, \emph{Semi-Vektoren und Spinoren}, Sitzber. Preuss. Akad. Wiss. Physik.-Math. Kl. (1932) 522--550. \%\%QUATERNION, \%\%SEMIVECTOR, \%\$D05022002.

\item {\bf Report or preprint:}

{\bf VELTM1997-} M. Veltman, \emph{Reflexions on the Higgs system}, Report 97-05 (CERN, 1997) 63~pp. \%\%BOOK, \%\%QUATERNION, \%\%LEPTODYNAMICS, \%\%GAUGE-THEORY,  \%\$D14052001, \%\$C Veltman uses quaternions in the form of 2x2 matrices.


\item {\bf Ph.D. thesis:}

{\bf GURSE1950A} F. G\"ursey, \emph{Applications of quaternions to field equations}, Ph.D.\ thesis (University of London, 1950) 204~pp.  \%\%BOOK, \%\%QUATERNION, ~~\%\%DIRAC, ~~\%\%GENERAL-RELATIVITY, ~~\%\%LANCZOS, \%\%PROCA, \%\$D20032002. %$C G\"ursey generalizes Lanczos's fundamental equation to curved space-time.\\


\item {\bf Unpublished document}:

{\bf GSPON1993A} A. Gsponer and J.-P. Hurni,  \emph{Quaternion bibliography: 1843--1993}. Report ISRI-93-13 (17 June 1993) 35~pp. This is the first version of the present bibliography, with 228 entries. \%\%QUATERNION, \%\%BIBLIOGRAPHY, \%\$D20032002.


\item {\bf ``arXived'' e-print:}


{\bf GSPON2002D} A. Gsponer, \emph{Explicit closed-form parametrization of $SU(3)$ and $SU(4)$ in terms of complex quaternions and elementary functions}, Report ISRI-02-05 (22 November 2002) 17~pp.;  e-print \underline{ arXiv:math-ph/0211056 }.  \%\%QUATERNION, \%\%ALGEBRA, \%\%HADRODYNAMICS, \%\$10092005.

\item {\bf Internet page:}

{\bf NOBIL2006-} R. Nobili, \emph{Fourteen steps into quantum mechanics}, HTML document (Posted in 2006) about 13~pp.;  available at\\ \underline{ http://www.pd.infn.it/~rnobili/qm/14steps/14steps.htm }. \%\%QUATERNION, \%\%QUANTUM-PHYSICS, \%\%QQM, \%\$D25.



\end{itemize}



\section*{ACKNOWLEDGMENTS}
%------------------------
\label{ACKNOWLEDGMENTS}

This bibliography would not exist without the help, dedication, and professionalism of Mrs.\ Claire-Lise Held and Mrs.\ Jocelyne Favre at the University of Geneva Physics department library, and many other librarians at other universities.


\end{document}

\bye
\bye
%
 
